% 本文件由 LaTeX 排版 Agent 根据有效 translation.json 自动生成。
% author=Hippolytus; work_id=0501; title=驳斥一切异端

\chapter{\WorkTitle{驳斥一切异端}}

% structure: p0001 source=h1 target=omitted reason=page-title-already-rendered-by-parent

\Emph{第一卷}\newline{}\Emph{第二卷} \Emph{缺失}\newline{}\Emph{第三卷} \Emph{缺失}\newline{}\Emph{第四卷}\newline{}\Emph{第五卷}\newline{}\Emph{第六卷}\newline{}\Emph{第七卷}\newline{}\Emph{第八卷}\newline{}\Emph{第九卷}\newline{}\Emph{第十卷}\par

\SourceRecord{Translated by J.H. MacMahon.；Ante-Nicene Fathers；Vol. 5.；Edited by Alexander Roberts, James Donaldson, and A. Cleveland Coxe.；Buffalo, NY: Christian Literature Publishing Co.,；1886.；<http://www.newadvent.org/fathers/0501.htm>.}

% structure: p0001 source=h1 target=section reason=page-title

\section{驳斥一切异端（第一卷）}

% structure: p0002 source=h2 target=subsection reason=relative-heading-rank; base=h2

\subsection{目录。}

以下是\WorkTitle{驳斥一切异端}第一卷的内容。\par

我们拟提供关于自然哲学家的学说以及这些哲学家是谁的说明，以及关于道德哲学家的学说以及这些哲学家是谁的说明，第三，关于逻辑学家的学说以及这些逻辑学家是谁的说明。\par

属于自然哲学家的有：\Person{泰勒斯}、\Person{毕达哥拉斯}、\Person{恩培多克勒}、\Person{赫拉克利特}、\Person{阿那克西曼德}、\Person{阿那克西美尼}、\Person{阿那克萨戈拉}、\Person{阿凯劳斯}、\Person{巴门尼德}、\Person{留基伯}、\Person{德谟克利特}、\Person{色诺芬尼}、\Person{厄克方托斯}、\Person{希波}。\par

属于道德哲学家的有：\Person{苏格拉底}（物理学家\Person{阿凯劳斯}的学生）、\Person{柏拉图}（\Person{苏格拉底}的学生）。这位思索者融合了三种哲学体系。\par

属于逻辑学家的有：\Person{亚里士多德}（\Person{柏拉图}的学生）。他将辩证法系统化。在斯多葛学派的逻辑学家中有\Person{克吕西波}、\Person{芝诺}。然而，\Person{伊壁鸠鲁}提出了一个几乎与所有哲学家相反的意见。\Person{皮浪}是一位学园派；这位思索者教导一切事物的不可知性。印度的婆罗门、凯尔特人中的德鲁伊，以及\Person{赫西俄德}（都致力于哲学追求）。\par

% structure: p0008 source=h2 target=subsection reason=relative-heading-rank; base=h2

\subsection{序言——进行驳斥的动机；古代秘仪的揭露；著作计划；驳斥的完整性；这部论著对后世的价值。}

我们不可忽视希腊人所谓的哲学家所杜撰的任何虚构。因为就连他们不连贯的学说也必须被认为是值得信赖的，这是由于异端者极度的疯狂；这些人因遵守沉默并隐藏自己不可言说的奥秘，而被许多人视为敬拜天主的人。我们先前也曾简要阐述过这些人的教义，并未详细说明，而是以粗略的方式加以驳斥；认为没有必要将他们的秘密教义公之于众，以便当我们用谜语解释他们的学说时，他们会感到羞愧，免得我们因揭露他们的奥秘而证明他们无神论，从而可能促使他们稍微放弃自己不合理的意见和亵渎的企图。但由于我看到他们并未因我们的容忍而感到羞愧，也不在意天主如何恒久忍耐（尽管被他们亵渎），以便他们或因羞耻而悔改，或若执迷不悟而受公义审判，我被迫继续执行我的计划，揭露他们那些秘密的奥秘；这些奥秘，他们以极具说服力的方式传授给入门者，而这些人不惯于首先向任何人披露，直到通过在一段（必要预备）时期内使他们悬而不决，并使其对真天主产生亵渎之心，从而完全控制他们，并察觉他们急切渴望所应许的启示。然后，当他们考验此人被罪奴役时，便给他入门，使他拥有万恶的圆满。然而，事先他们以誓言约束他，既不得泄露（奥秘），也不得与任何人往来，除非他先经历同样的顺服，尽管当教义被单纯地传授（给任何人）时，便不再需要誓言。因为一个甘心接受必要的净化并如此接受这些人的圆满奥秘的人，仅凭此行为以及其良心的感受，就会感到自己有足够义务不向他人泄露；因为一旦他向任何人揭露此种邪恶，他便不再被算作人类，也不配见到光明，因为就连非理性的动物也不会尝试如此大恶，我们将在论及此类主题时加以说明。\par

然而，由于理性迫使我们深入叙述的深处，我们认为不应保持沉默，而要详细阐述各学派的学说，毫无保留。现在似乎适宜，即使付出更详尽探究的代价，也不应畏缩劳苦；因为我们将为人类生活留下对抗错误复发的非小辅助，当所有人都能在明显的光照下看见这些人的隐秘仪式和秘密狂欢，这些仪式在他们的管理下，只传授给入门者。但除非是遗留给教会的圣神，才能驳斥这些；圣神是宗徒们首先领受，然后传给那些正确信仰的人的。而我们，作为他们的继承者，并分沾这恩宠、大司祭职和教导的职务，且被视为教会的守护者，决不可缺乏警惕，或倾向于压制正确的教义。然而，即便竭尽身心之力，我们也不厌倦于试图充分回报我们的神圣恩主；然而我们也不应适当地回报祂，除非我们不懈怠地履行托付给我们的信职，而是谨慎地完成我们特定时机的度量，并毫无吝啬地将圣神所供给的一切分施给众人，不仅通过我们的驳斥将（与我们主题）无关的事引人注目，而且也将真理藉圣父的恩宠所领受并服务于人的一切事项。这些，我们也将通过论证加以阐明，并以文字创造见证，毫不羞愧地宣讲。\par

因此，为了正如我们已说过的，证明他们在意见、处理问题的方式和事实上都是无神论者，并（为了表明）他们试图建立的理论如何累积到他们身上，以及他们如何竭力建立自己的教义，既不从圣经中取用任何内容——也不是因为保存任何圣人的传承而匆忙陷入这些意见——而是他们的教义源于希腊人的智慧，来自那些建立哲学体系者的结论，来自所谓的奥秘以及占星家的妄想——那么，看来首先是通过解释希腊哲学家提出的意见，使我们的读者相信这些意见比这些异端更古老，在关于神的观点上更值得尊敬；其次，将每个异端与每个思辨者的体系进行比较，以表明异端的首位倡导者利用这些试图建立的理论，通过吸收他们的原则而将其转化为自己的优势，并由此陷入更糟糕的境地，构建了自己的教义。这项任务无疑充满艰辛，需要广泛的研究。然而，我们不会缺乏努力；因为之后它将成为喜乐的源泉，如同一位运动员经过艰苦努力获得桂冠，如同一位商人在巨浪滔天后获得收益，如同一位农夫在额上流汗后享受果实，如同一位先知在遭受责备和侮辱后看到预言成真。因此，在开始时，我们将首先宣布在希腊人中谁指出了自然哲学的原理。因为这些人尤其偷偷地采纳了那些首先提出这些异端者的观点，正如我们日后在逐个比较时将证明的那样。将每位哲学领袖的独特教义分配给他们，我们将公开展示这些异端首领的赤裸与不雅。\par

% structure: p0012 source=h2 target=subsection reason=relative-heading-rank; base=h2

\subsection{第一章 \Person{泰勒斯}：他的物理学与神学；希腊天文学的奠基人。}

据说，七贤之一的\Place{米利都}的\Person{泰勒斯}首先尝试构建自然哲学体系。此人说，水之类的东西是宇宙的生成原则及其终结——因为万物由水凝固和再溶解而构成，并且万物都依赖水；由此产生地震、风的变化和大气运动，并且万物都按照生成的首要作者的本质而产生和流动——而神没有起始也没有终结。此人专注于对星辰的假说和研究，成为希腊人中这类学问的最早作者。他仰望天空，声称仔细观察天上的事物，却掉进了井里；一个名叫\Person{特拉塔}的侍女嘲笑他说，他专心注视天上之物，却不知脚下之物。他生活在\Person{克罗伊斯}时期。\par

% structure: p0014 source=h2 target=subsection reason=relative-heading-rank; base=h2

\subsection{第二章 \Person{毕达哥拉斯}：他的宇宙生成论；他教派的规则；相术的发现者；他的数论哲学；他的灵魂转世体系；\Person{查拉塔斯}论魔鬼；为什么\Person{毕达哥拉斯}禁止吃豆子；他的门徒的生活方式。}

但就在这些时期不久，还有另一种哲学，由\Person{毕达哥拉斯}创立（有人说他是\Place{萨摩斯}人），他们称之为意大利哲学，因为\Person{毕达哥拉斯}逃离\Place{萨摩斯}的\Person{波利克拉特斯}，在\Place{意大利}的一座城市定居，并在那里度过了余生。那些继承他教义的人与他的观点没有太大差异。此人研究自然现象，将天文学、几何学和音乐结合在一起。因此他宣告神是单子；他仔细了解数的本质，断言世界歌唱，其系统和谐一致，并首次将七星的运行分解为节奏与旋律。他对整个构造的管理感到惊奇，要求他的门徒首先保持静默，仿佛进入世界的人被引入了宇宙的秘密；其次，当他们似乎充分熟悉了他的教义传授方式，并能有力地哲学讨论星辰和自然时，他认为他们纯洁了，才允许他们说话。此人将他的学生分为两个等级，称为内学和外学。对于前者，他传授更高深的教义；对于后者，则给予较适度的教导。\par

他也触及了魔术——如他们所说——并且自己发现了一种面相术（physiogony），以某些数字和尺度为基础，声称它们以这种方式通过合成构成了算术哲学的原理。第一个数字成为起源原理，即一，不可定义、不可理解，在其自身内包含所有可以无限（\Emph{ad infinitum}）延伸的数字。但原初的\OriginalTerm{单一元}{monad}按实体成为数字的原理——它是男性的单一元，像父亲一样生育所有其他数字。其次，\OriginalTerm{二元}{duad}是女性数字，算术家称之为偶数。第三，\OriginalTerm{三元}{triad}是男性数字，算术家也将其归类为奇数。此外，还有\OriginalTerm{四元}{tetrad}，女性数字，也因是女性而被称为偶数。因此，从该属类衍生的所有数字是四个；但数字是无限的属类，根据他们，完美的数字即十，由之构成。因为一、二、三、四成为十，如果每个数字的本质名称得以保持。\Person{毕达哥拉斯}宣称这是一个神圣的四元，永恒自然的源头，在其自身内拥有根基；并且所有数字都从这个数字获得其起源原理。因为十一、十二等，都从十分享存在起源。这个完美的数字十，有四个部分——即数字、\OriginalTerm{单一元}{monad}、平方、立方。这些的结合与混合按本性进行，以生成增长，产生生育性数字。因为当平方自身乘以自身时，得到四方幂（biquadratic）。当平方乘以立方时，得到平方与立方的乘积；当立方乘以立方时，得到两个立方的乘积。因此，所有产生存在（数字）的数字是七个——即数字、\OriginalTerm{单一元}{monad}、平方、立方、四方幂、平方-立方、立方-立方。\par

这位哲学家还说灵魂是不朽的，并且它连续存在于身体中。因此他断言，在特洛伊时代之前他是\Person{埃塔利得斯}，在特洛伊时代是\Person{欧福耳布斯}，此后是\Place{萨摩斯}的\Person{赫耳摩提谟斯}，再后是\Place{德洛斯}的\Person{皮洛斯}；第五是\Person{毕达哥拉斯}。而\Place{厄瑞特里亚}的\Person{狄奥多洛斯}和音乐家\Person{阿里斯托克塞诺斯}断言，\Person{毕达哥拉斯}曾拜访\Place{迦勒底}的\Person{匝拉塔斯}，后者向他解释事物的两个原始起因，即父亲和母亲，父亲是光，母亲是黑暗；光的各部分分有热、干、不重、轻、快；黑暗的则分有冷、湿、重、慢；所有这一切，由女性与男性，构成了世界。但他说，世界是一种音乐的和谐；因此太阳也按和谐运行。关于从地产生的事物和宇宙体系，他们声称\Person{匝拉塔斯}作了如下陈述：有两个魔鬼，一个属天，一个属地；属地的从地上升起一种产物，即水；属天的是一种火，分有气的性质，热与冷。因此他断言，这些没有一个能毁灭或玷污灵魂，因为它们构成万物的实体。据传他命令他的追随者不要吃豆子，因为\Person{匝拉塔斯}说，在万物的起源与凝结时，当地球仍在固化过程中，腐败作用开始时，豆子产生了。他提到以下迹象：如果有人咀嚼过一颗去壳的豆子，在太阳下放置一段时间——因为这会立即有助于结果——它会散发出人精液的气味。他还提到另一个更明显的例子：当豆子开花时，我们取豆子和花，放入罐中，封好，埋在地里，几天后挖出，我们将看到它先呈现女性\OriginalTerm{私处}{pudendum}的样子，然后仔细检查，会看到儿童的头与其一起生长。此人在\Place{意大利}的\Place{克罗同}城与他的门徒一起被烧死。他有一个习惯：每当有人来找他要求成为追随者，（候选门徒必须）变卖财产，将钱密封交给\Person{毕达哥拉斯}，然后保持沉默接受教导，有时三年，有时五年。然后，若是被释放，就被允许与其他门徒交往，作为门徒留下并与之同餐；若非如此，则收回他的财产，被拒绝。这些人被称为\OriginalTerm{内部毕达哥拉斯派}{Esoteric Pythagoreans}，而其余的人称为\OriginalTerm{毕达哥拉斯追随者}{Pythagoristae}。\par

然而，在他逃脱火灾的追随者中有\Person{吕西斯}、\Person{阿耳基波斯}和\Person{毕达哥拉斯}的仆人\Person{匝摩尔克西斯}，据说后者教授\OriginalTerm{凯尔特德鲁伊德}{Celtic Druids}培育\Person{毕达哥拉斯}的哲学。他们断言\Person{毕达哥拉斯}从\OriginalTerm{埃及人}{Egyptians}学习了数字和度量的体系；我被祭司们似乎有理、奇幻且不易揭示的智慧所打动，他自己也模仿他们，命令保持沉默，并让他的门徒在地下小圣堂中过独居生活。\par

% structure: p0019 source=h2 target=subsection reason=relative-heading-rank; base=h2

\subsection{第三章 \Person{恩培多克勒}；他的双重原因；转生学说。}

但后来的\Person{恩培多克勒}也提出了许多关于魔鬼本性的说法，即他们数量众多，管理地上事务。此人断言宇宙的起源原理是不和与友爱，并且\OriginalTerm{单一元}{monad}的理智之火就是神性，万物由火构成，并将被分解为火；\OriginalTerm{斯多葛派}{Stoics}也几乎同意此观点，期待一场大火。但他最同意灵魂从身体到身体转移的学说，表达如下：——\par

\Quote{因为无疑我曾是青年和少女，\newline{}灌木、飞鸟，以及从海洋游来的鱼。}\par

这位（哲学家）主张一切灵魂都能转生为任何种类的动物。因为\Person{毕达哥拉斯}，这些智者的导师，声称他本人曾经是\Person{欧福尔布斯}，曾参与对抗\Place{伊利昂}的远征，并说他认出了自己的盾牌。以上是\Person{恩培多克勒}的学说。\par

% structure: p0023 source=h2 target=subsection reason=relative-heading-rank; base=h2

\subsection{第四章 \Person{赫拉克利特}；他的普遍独断论；他的流变论；其他体系。}

但是\Place{以弗所}的自然哲学家\Person{赫拉克利特}沉溺于普遍的悲伤，谴责全部生命及所有人的无知；不，他甚至怜悯凡人的存在，因为他宣称他自己知道一切，而其余的人类一无所知。但他也提出了几乎与\Person{恩培多克勒}一致的论断，说万物的本原是冲突与友爱，神是赋有理智的火，万物彼此承载，永不止息；并且正如\Person{恩培多克勒}，他断言我们周围的整个区域充满邪恶之物，这些邪恶之物从大地周围的区域延伸直至月球，但不更进一步，因为月球之上的整个空间更为纯净。\Person{赫拉克利特}似乎也这样认为。\par

在这些之后又兴起其他自然哲学家，我们认为不必阐述他们的意见，因为它们与已述者并无差异。然而，既然总的来说，有一个相当可观的学派由此产生，并且许多自然哲学家随后从他们中兴起，各自提出了关于宇宙本性的不同说法，我们也认为明智的做法是，在解释由\Person{毕达哥拉斯}传承下来的哲学时，我们应回溯\Person{泰勒斯}之后时代的人们所持有的意见，并且提供了这些叙述后，我们应进而考虑\Person{苏格拉底}和\Person{亚里士多德}所创始的伦理哲学和逻辑哲学，前者是伦理的，后者是逻辑的。\par

% structure: p0026 source=h2 target=subsection reason=relative-heading-rank; base=h2

\subsection{第五章 \Person{阿那克西曼德}；他的无限论；他的天文观点；他的物理学。}

\Person{阿那克西曼德}是\Person{泰勒斯}的学生。\Person{阿那克西曼德}是\Person{普拉克西亚达斯}之子，\Place{米利都}人。这个人说存在物的本原是某种\OriginalTerm{无限}{Infinite}的构造，天和其中的世界由此生成；并且这个本原是永恒不灭的，包含所有世界。他把时间说成是有限生成、存续和毁灭的东西。他宣称\OriginalTerm{无限}{Infinite}是存在物的本原和元素，是第一个使用这种对本原的称谓的人。此外，他断言有一种永恒运动，由于这种运动，天得以生成；但地球悬在空中，无所支撑，因与所有（天体）等距而保持（如此）；它的形状是弯曲、圆形的，类似于石柱。其中一表面我们踩踏，另一表面则相反。众星是火圈，与世界附近之火分离，被空气围绕。在星光照耀之处，产生某种大气蒸发；因此，当这些蒸发受阻时，便发生食。月亮有时呈现满月，有时亏缺，根据其（轨道）路径之阻塞或开启而定。但太阳的圆圈比月亮大二十七倍，太阳位于（天空）最高处，而恒星之球体则在最低处。动物由太阳蒸发（在湿气中）产生。人最初类似于另一种动物，即鱼。风由极稀薄的大气蒸发分离并凝结后的运动引起。雨由大地将从（云）中（于太阳下）接收的水汽归还而产生。当下降的风劈开云时，就会产生闪电。此人生于第四十二届奥林匹克年的第三年。\par

% structure: p0028 source=h2 target=subsection reason=relative-heading-rank; base=h2

\subsection{第六章 \Person{阿那克西美尼}；他的\Quote{无限的空气}体系；他的天文与自然现象观点。}

但是，\Person{阿那克西美尼}本人也是\Place{米利都}人，\Person{欧律斯特拉托斯}之子，断言本原是无限的气，由它生成现在存在、曾经存在以及将要存在的事物，同样也生成神祇和神圣（存在），其余事物则出自这气的产物。此气在最为均匀时有一种形式，目不能见，但能通过冷热、湿气和运动显现，并且它不断运动。因为凡经历变化之物，若无运动便不会改变。因为它根据凝聚和稀薄程度呈现不同外观：当它消散为更稀薄之物时，产生火；当它适度凝聚回气时，由于收缩而从气形成云；当凝聚更多时，形成水；当凝聚更进一步时，形成土；当凝聚到最高程度时，形成石头。因此，生成的主要原理是对立物，即热与冷。并且广阔的大地浮于气上，同样太阳、月亮和其他星辰亦是如此；因为一切具有火性之物，都借气在太空中漂浮。星辰是由大地产生的，因从这\Emph{地}升起的水汽；当水汽稀薄时，产生火，星辰由向上飘升的火构成。并且也有地上的性质在星辰区域与它们一同运行。他说星辰并非如某些人设想那样在底下移动，而是绕地而行，就像一顶帽子绕着我们的头转动；太阳隐没，不是因为它到了地下，而是被地球较高的部分遮蔽，并且因为它离我们更远之故。星辰不散发热量是因为距离遥远；风产生于凝聚的气变得稀薄而被推动时；当气进一步聚集变厚时，产生云，从而变成水。冰雹产生于从云中降下的水凝结之时；雪产生于这些云本身较为湿润而凝结之时；闪电由风的力量劈开云层引起；因为云层被撕裂时会产生明亮而炽热的闪光。彩虹由太阳光线射在聚集的空气上而产生。地震发生是因大地受冷热影响而变成更大的体量。这些就是\Person{阿那克西美尼}的见解。这位（哲学家）活跃于第五十八届奥林匹亚德的第一年左右。\par

% structure: p0030 source=h2 target=subsection reason=relative-heading-rank; base=h2

\subsection{第七章 \Person{阿那克萨哥拉}；他的心灵理论；承认一个动力因；他的宇宙起源学和天文学。}

继此之后的是\Person{阿那克萨戈拉}，他是\Person{赫格西布勒}之子，生于\Place{克拉佐美奈}。此人主张宇宙的本原是心灵与物质：心灵为动力因，物质则为被形成者。因为万物同时生成，心灵介入后引入了秩序。他说物质本原是无限的，就连其中最小的也是无限的。万物因被心灵推动而分有运动，相似物相互结合。天体由圆环运动安排而成。因此，厚、湿、暗、冷及一切沉重之物聚集于中心，它们的凝固形成了大地；而与这些相反的——即热、明亮、干、轻——则急速冲向大气的外层。大地形状平坦；它之所以悬浮高空，是由于其巨大、没有虚空，以及最强有力的空气托着飘浮的大地。地上湿润之物中有海及其中之水；当这些蒸发（受太阳）或沉降后，海洋便如此形成，同时也由不时流入的河流形成。河流也借雨水和地中实际的水维持；因为大地中空，洞穴中含水。\Place{尼罗河}夏季泛滥，是因为北方（纬度）的雪水流入其中。太阳、月亮及所有星辰都是炽热的石块，由大气的旋转带动旋转。星辰之下有太阳、月亮及某些与我们同行而不可见的天体；我们感觉不到星辰的热，既因为它们极远，也由于它们与地球的距离；再者，由于它们处于更冷的位置，其热度不及太阳。月亮低于太阳，离我们更近。太阳比\Place{伯罗奔尼撒}半岛还大。月亮自身不发光，而来自太阳。星辰的旋转在地球之下进行。月食发生于地球居中时，有时也发生在月亮之下的那些（星辰）上。日食发生于月初月亮居中时。冬至和夏至由日、月被空气排斥所致。月亮常因不能抵挡寒冷而转向。此人是第一个界定食和光照的人。他宣称月亮是土质的，其中有平原和沟壑。银河是恒星（不藉太阳发光）之光的反射；流星般飞驰的星辰由极地运动产生。风由太阳及在极地之下前进并由此被携带的燃烧天体稀薄大气而成。雷电由热落到云上引起。地震由上方空气落到地下空气所致；当此空气运动时，被它托起的大地也随之震动。动物最初产生于湿润中，而后彼此生成；当从右侧分泌的精子附着于子宫右侧时，产生男性，反之产生女性。这位哲学家活跃于第八十八届奥林匹克运动会第一年，据说\Person{柏拉图}也于此时诞生。人们声称\Person{阿那克萨戈拉}也具有预知能力。\par

% structure: p0032 source=h2 target=subsection reason=relative-heading-rank; base=h2

\subsection{第八章 \Person{阿凯劳斯}；体系与\Person{阿那克萨戈拉}相似；他的地球和动物起源；其他体系。}

\Person{阿凯劳斯}是\Place{雅典}人，\Person{阿波罗多洛}之子。他与\Person{阿那克萨戈拉}类似，主张物质的混合，并以同样方式阐述第一原理。然而此哲学家认为心灵立即内在地具有某种混合；运动的起源原则是热与冷的相互分离，热运动，冷静止。水被溶解后流向中心，在那里产生烧灼的空气和土，前者向上，后者留在下方。大地静止，并因此生成；它位于中心，可以说是宇宙中未被焚毁的一部分；由此，首先在点火状态下形成星辰的本质，其中最大的是太阳，其次是月亮；其余的有大有小。他说天倾斜成一角度，因此太阳将光散布于大地，使大气透明，地面干燥；因为起初它是海洋，由于它在地平线处高而中部凹陷。他引证凹陷的迹象：太阳并非同时为所有人升起和落下，若大地平坦则应如此。关于动物，他断言大地最初下部是火，热冷混合之处，其余动物出现，众多而不相似，都以泥为食物；它们寿命短暂，但后来彼此之间也产生繁衍；人类与其他动物分离，设立了统治者、法律、技艺、城市等。他断言心灵同样内在于所有动物；每个动物按体质差异使用（心灵），有时较慢，有时较快。\par

自然哲学于是从泰勒斯一直延续到阿凯劳斯。苏格拉底是这位（后者哲学家的）弟子。然而，还有非常多其他人，对神性和宇宙的本性提出了各种意见；倘若我们打算一一引述所有这些意见，就必须编写数量庞大的书籍。但我们要提醒读者那些我们尤其应当——那些因名声而值得提及，且可谓是后来建立哲学体系之人的引领者，并为他们提供此类事业起点的——让我们赶紧将研究转向尚待考察的内容。\par

% structure: p0035 source=h2 target=subsection reason=relative-heading-rank; base=h2

\subsection{第九章 巴门尼德；他的\Quote{合一}理论；他的末世论。}

因为巴门尼德同样认为宇宙是单一的、永恒且非受生的，且具有球形形态。他也没有逃脱大多数（思辨者）的意见，主张火和土是宇宙的本原——土为质料，火为原因，乃至动力因。他断言世界将会毁灭，但未说明方式。然而，同一（哲学家）认为宇宙是永恒的、非受生的，具有球形且同质，但自身没有形象，且不动、有限。\par

% structure: p0037 source=h2 target=subsection reason=relative-heading-rank; base=h2

\subsection{第十章 留基伯；他的原子论。}

但留基伯，芝诺的同伴，并不持相同意见，而是主张事物无限、永恒运动，生成与变化持续存在。他肯定充实和虚空为元素。他断言当许多物体从周围空间聚集并流向共同点时，世界得以产生，以至于通过相互接触，形状相同、形式相似的实体相互结合；如此交织之后，发生向其他物体的转化，被造物因必然性而生长和衰亡。但（巴门尼德）并未界定必然性的本质为何。\par

% structure: p0039 source=h2 target=subsection reason=relative-heading-rank; base=h2

\subsection{第十一章 德谟克利特；他的二原则说；他的宇宙生成论。}

德谟克利特是留基伯的相识。德谟克利特，达马西波斯之子，阿布德拉人，曾与印度许多裸体智者、埃及祭司以及巴比伦的占星家和术士交谈（而提出其体系）。他在元素方面与留基伯作出类似陈述，即充实和虚空，称充实为存在，虚空为非存在；他断言如此，因为存在者在虚空中不断运动。他主张世界无限、大小各异；有些世界既无太阳也无月亮，而另一些则比我们的更大，还有一些则有更多太阳和月亮。世界之间的间隔不等；在空间某一方向，世界较多，在另一方向较少；有些在增长，有些达到全盛，有些则缩小；在一方向生成，在另一方向衰亡；它们因相互碰撞而毁灭。有些世界缺乏动物、植物和任何水分。我们这个世界的土地先于星辰形成；月亮位于下方；其次是太阳；然后是恒星。（行星和）这些恒星并不拥有相同的高度。世界繁荣，直到无法再从外部接收任何东西。此（哲学家）将一切化为笑谈，仿佛人类的一切事务都值得嘲笑。\par

% structure: p0041 source=h2 target=subsection reason=relative-heading-rank; base=h2

\subsection{第十二章 色诺芬尼；他的怀疑论；他的天主与自然观；相信大洪水。}

但色诺芬尼，科洛丰人，是俄耳托墨涅斯之子。此人活到居鲁士时代。此（哲学家）首先断言不可能理解任何事物，他如此表述：\par

\Quote{因为即使人可能谈论大部分完美，\newline{}\newline{}他自身却不认识它，而在一切中达到猜测。}\par

他断言没有东西被生成、毁灭或运动；宇宙是单一的，超越变化。但他说天主是永恒的、单一的、完全同质且有限，具有球形，且在一切部分具有知觉。太阳每日由小火花聚集而成；大地是无限的，既不被大气也不被天穹包围。存在无限的太阳和月亮，万物从大地生出。此人断言海水是咸的，因为许多混合物流入其中。然而，墨特罗多罗斯则因其透过大地过滤，断言这正是海水变咸的原因。色诺芬尼认为大地曾与海洋混合，并随时间从湿气中分离出来，他声称能提出如下证据：在大地中部和山上发现贝壳；在叙拉古，他说在采石场发现鱼和海豹的印记；在帕罗斯，一块石头底部有月桂树的形象；在梅利塔，有各种海洋动物的部分。他说这些是在万物最初埋藏于泥中时生成的，其印记在泥中干燥；但当大地沉入海中变成泥时，所有人类都灭亡了；然后，又产生了新的生成，这种倾覆发生在所有世界。\par

% structure: p0045 source=h2 target=subsection reason=relative-heading-rank; base=h2

\subsection{第十三章 厄克梵图斯；他的怀疑论；无限性学说。}

\Person{厄克方托斯}，\Place{锡拉库萨}人，断言无法获得对事物的真知。然而，他按照自己的看法，规定原初的物体是不可分的，并且这些物体有三种变化，即体积、形状、容量，由之产生可感对象。但认为这些物体有一个可确定的数目，并且这个数目是无限的。他还认为物体既不是由重量也不是由冲击所推动，而是由一种神圣能力所推动，他称之为心智和灵魂；世界是这神圣能力的表现；因此世界也由神圣能力造成球形。而且，宇宙系统中央的地球绕其自身中心向东运动。\par

% structure: p0047 source=h2 target=subsection reason=relative-heading-rank; base=h2

\subsection{第十四章 \Person{希波}；他的二元本原论；他的心理学}

\Person{希波}，\Place{雷焦}人，主张本原为冷（如水）与热（如火）。火由水产生时，压倒了产生者的能力，并形成了世界。他说，灵魂有时是脑，有时是水；因为种子也是从湿气中产生的，据他说，灵魂由之产生。\par

至此，我们认为已充分引述了这些人的意见；既然已足够地回顾了自然哲学家们的学说，看来剩下的是现在转向\Person{苏格拉底}和\Person{柏拉图}，他们特别偏重道德哲学。\par

% structure: p0050 source=h2 target=subsection reason=relative-heading-rank; base=h2

\subsection{第十五章 \Person{苏格拉底}；他的哲学由\Person{柏拉图}传承}

\Person{苏格拉底}是自然哲学家\Person{阿基劳斯}的学生；他尊崇\Quote{认识你自己}这一准则，召集了一个大学派，其中有\Person{柏拉图}，他远胜于所有其他学生。\Person{苏格拉底}本人身后未留下任何著作。然而\Person{柏拉图}记录了他所有的智慧讲演，创立了一个学派，结合了自然哲学、伦理学和逻辑学。\Person{柏拉图}所确定的要点如下。\par

% structure: p0052 source=h2 target=subsection reason=relative-heading-rank; base=h2

\subsection{第十六章 \Person{柏拉图}；本原的三重分类；他的天主观念；关于其神学与心理学的不同意见；其末世论与转世体系；其伦理学说；关于自由意志问题的概念}

\Person{柏拉图}主张宇宙有三个本原：天主、物质和范型。天主是宇宙的创造者和调节者，并对其施行眷顾；物质则是万有的基底，他称之为承受者和养护者，由其安排产生了构成世界的四元素，即火、气、土、水，由之形成了其余一切称为具体实体的东西，以及动物和植物。范型，他亦称之为理型，是神性的理智，神性参照它，如同灵魂中的图像，造作了万物。他说，天主既是无形体又是无形状的，唯有智者才能理解；而物质则是潜在地是物体，但潜能尚未进入实现，因为它本身无形式、无性质，通过获得形式和性质而成为物体。因此，物质是一个本原，与神性同古，就此而言世界是不受造的。因为\Person{柏拉图}肯定世界是由物质造成的。而\Quote{不朽}必然属于（字面意思为\Quote{随从}）不受造之物。然而，就物体被认为是由许多性质和理型组合而成而言，它既是受造的，也是可朽的。但有些\Person{柏拉图}的门徒将这两者混合起来，用如下例子：正如一辆车，虽然不时部分修理，甚至各部分依次朽坏，但本身始终完好无损；同样，世界虽然部分朽坏，但被替换的部分得到修复，并以等价物补充，它便保持永恒。\par

有些人主张\Person{柏拉图}宣称神性是一、不受造且不朽，正如他在\WorkTitle{法篇}中所说：\Quote{天主，按古代传说，拥有万物的开始、终结和中间。}这样他表明天主是独一的，因为祂渗透万有。然而其他人主张\Person{柏拉图}断言存在无限多的神，当他使用这些话时：\Quote{神们的神，我是其创造者和圣父。}但另有人说他在以下段落中提到确定数目的神：\Quote{因此强大的\Person{朱庇特}，在天空中驾着他快速的战车；}以及当他列举天地子女的后代时。但还有一些人断言\Person{柏拉图}认为神是可生的；并且由于他们是被产生的，他们完全受腐朽的必然性所支配，但因天主的旨意他们是不朽的，在已引用的段落中，对\Quote{神们的神，我是其创造者和圣父}这句话，他补充说：\Quote{因我的旨意而不可分解}；所以如果天主愿意他们分解，他们将容易被分解。\par

他也承认魔鬼有本性，说其中一些是善的，另一些是恶的。有人认为他宣称灵魂是非受造且不朽的，当他说道：\Quote{每一个灵魂都是不朽的，因为那永远运动的东西是不朽的；}并且当他证明灵魂是自我运动、能引起运动时。然而，另一些人（说柏拉图断言灵魂是）受造的，但因天主的旨意而成为不朽的。但有些人（认为他认为灵魂是）复合的、可生可灭的；因为他也假定灵魂有一个容器，并且它有一个发光的身体，但一切受造之物都包含朽坏的必然性。然而，那些主张灵魂不朽的人，尤其因他在著作中的那些段落而坚定了意见，其中他说死后有审判，阴府中有公义的法庭，有德行的（灵魂）获得善报，而邪恶的（灵魂）获得相应的惩罚。尽管如此，有些人声称他也承认灵魂从一个身体转移到另一个身体，不同的灵魂——那些被指定为此目的的灵魂——按照各自的功德进入不同的身体，并在某些确定的时期之后被送回这个世界，以再次提供他们选择的证据。然而，另一些人（不承认这是他的教导，而认为柏拉图肯定说，灵魂）按照各自的功德获得一个位置；他们引用他的说法作为证据：有些善人与朱庇特同在，另一些与其他神祇漫游（在天上）；而许多在世时行了邪恶不义之事的人则陷入永罚。\par

人们断言柏拉图说，有些东西没有中项，有些东西有中项，有些东西本身就是中项。（例如，）清醒和睡眠以及诸如此类是没有中间状态的情况；但有些东西有中项，例如德行与恶习；还有一些是中项（介于两极之间），例如灰色介于白色与黑色之间，或其他某种颜色。他们说，他肯定与灵魂有关的东西是绝对唯一善的，但与身体有关的和外在的（东西）就不再是绝对善的，只是被认为是有福的。并且他经常也称这些为中项，因为可以善用或恶用它们。因此，他说有些德行就其内在价值而言是极端，但就其本质而言却是中项，因为没有比德行更可贵的了。但任何超越或不及这些德行的事物都归于恶习。例如，他说有四种德行——智德、节制、义德、刚毅——而每一种德行伴随着两个恶习，一为过度，一为不及：例如，智德伴随着不及的轻率和过度的奸诈；节制伴随着不及的放纵和过度的愚钝；义德伴随着不及的放弃主张和过度的坚持主张；刚毅伴随着不及的懦弱和过度的鲁莽。这些德行，当在人内时，使他完善，并带给他幸福。而幸福，他说，是尽可能地肖似于天主；当一个人将圣洁和义德与智德结合时，就实现了肖似天主。因为他认为这是至上智慧和德行的目的。他肯定德行彼此相继，是一致的，绝不相互对抗；而恶习则是多样的，有时相互继随，有时彼此对抗。他断言命运存在；当然，不是一切事物都按命运产生，而是甚至有些东西在我们能力之内，正如他在以下段落中所说：\Quote{选择者的过错在己，天主无可指责；}以及\Quote{阿德拉斯忒亚的法则。}因此，有些人（主张他支持）命运系统，而另一些人则主张自由意志。然而，他断言罪过是不自愿的。因为在我们能力之内最荣耀的东西——即灵魂——没有人会（故意）接纳邪恶的东西，即过犯；而是由于无知和对德行的错误认识，以为自己在成就某种光荣的事，而陷入恶习。他在这点上的教义在\WorkTitle{理想国}中最清楚，他说：\Quote{但是，你又断言恶习是可耻的、为天主所憎恶；那么，我问，谁会选择这样一种邪恶的东西呢？你回答，被快乐征服的人会这样做。因此，这也是不自愿的，如果得胜是自愿的；所以，无论在哪个角度看，理性证明做可耻的事是不自愿的。}然而，有人反对这一点，提出相反的主张：\Quote{那么，如果人们不自愿犯罪，为何要受惩罚？}但他回答说，他自己也尽可能早日脱离恶习并接受惩罚。因为接受惩罚不是坏事，而是好事，如果它可能成为对邪恶的净化；并且其余的人听到后就不会违犯，而是警惕这样的错误。（然而，柏拉图主张）恶的本性既非由天主创造，也非自存，而是因与善的对立以及附属而生，或由于过度与不及，正如我们先前关于德行所说的。显然，正如我们已经说过的，柏拉图将普遍哲学的三大领域汇集起来，以这种方式形成了他的思辨体系。\par

% structure: p0057 source=h2 target=subsection reason=relative-heading-rank; base=h2

\subsection{第十七章 亚里士多德；双重原理；他的范畴；他的心理学；他的伦理学；\Quote{逍遥学派}称号的起源}

亚里斯多德是\Person{柏拉图}的弟子，他将哲学化为一门技艺，并以在逻辑科学上的造诣著称。他假定一切事物的元素为\OriginalTerm{本体}{substance}与\OriginalTerm{属性}{accident}；有一个本体作为万物的基础，以及九个属性——即\OriginalTerm{数量}{quantity}、\OriginalTerm{性质}{quality}、\OriginalTerm{关系}{relation}、\OriginalTerm{地点}{where}、\OriginalTerm{时间}{when}、\OriginalTerm{所有}{possession}、\OriginalTerm{姿态}{posture}、\OriginalTerm{主动}{action}、\OriginalTerm{被动}{passion}。本体是指诸如天主、人以及任何可归入类似名称的存在。关于属性：性质见于例如白、黑；数量如二腕尺、三腕尺；关系如父、子；地点如在雅典、墨伽拉；时间如在第十届奥林匹亚赛会期间；所有如获得；主动如书写，以及一般地展示任何实践能力；姿态如躺下；被动如被击打。他还假定有些事物有中项，有些没有，正如我们关于柏拉图也说明的那样。他在多数观点上与柏拉图一致，除了关于\OriginalTerm{灵魂}{soul}的意见。柏拉图认为灵魂不朽，而亚里斯多德认为灵魂具有持续性；此后，灵魂也在\OriginalTerm{第五元素}{fifth body}中消灭，这第五元素是他与其他四种元素（即火、土、水、气）一起假定的，是一种更精微的东西，具有\OriginalTerm{精神}{spirit}的性质。因此柏拉图说，唯一真正善的东西是属于灵魂的，并且它们足以带来幸福；而亚里斯多德提出三种分类的善，断言智慧人除非同时拥有身体的和外在的善，否则并不完美。前者是美丽、力量、感官的敏捷、健康；外在的（身体的）善是财富、高贵、荣耀、权力、和平、友谊。他将灵魂的内在品质，按照柏拉图的意见，分为\OriginalTerm{智德}{prudence}、\OriginalTerm{节德}{temperance}、\OriginalTerm{义德}{justice}、\OriginalTerm{勇德}{fortitude}。这位哲学家也肯定，恶是由于善的对立而产生，它们存在于月球以下的区域，不超出月球；整个世界的灵魂是不朽的，世界本身是永恒的，而个体的灵魂，如我们先前所说，在第五元素中消灭。这位思想家在\Place{吕刻昂}（Lyceum）进行讨论，不时建立他的哲学体系；而\Person{芝诺}则在称为\Emph{彩廊}（Poecilé）的门廊下讲学。\Person{芝诺}的追随者得名于地点——即来自\Emph{廊柱}（Stoa），称为斯多亚学派；而\Person{亚里斯多德}的追随者得名于他们教学时的活动方式。由于他们习惯于在\Place{吕刻昂}漫步进行探究，因此被称为逍遥学派（Peripatetics）。这些便是亚里斯多德的学说。\par

% structure: p0059 source=h2 target=subsection reason=relative-heading-rank; base=h2

\subsection{第十八章　斯多亚学派；他们在逻辑学上的优越性；宿命论者；关于大火的学说。}

斯多亚学派自身也推动了哲学的发展，特别在推论术方面更加完善，并且几乎将所有事物都纳入定义，在这方面\Person{克吕西波}与\Person{芝诺}意见一致。他们也假定天主是万物的唯一起源原理，是一种最精微的物体，其\OriginalTerm{天主眷顾}{providential care}遍及一切；这些思想家确信命运无处不在，并用如下例子说明：就像一只狗被拴在车上，如果它愿意跟随，那么它既被牵引，又自愿跟随，在必然性即命运的结合下也行使自由能力；但如果它不愿意跟随，它终将被强制跟随。当然，同样的情况适用于人。因为即使不愿跟随，他们终将被强迫进入为他们注定的命运。（斯多亚学派）断言灵魂在死后持续存在，但灵魂是物体，是由周围大气的冷却而形成的；因此它被称为\Emph{psyche}（灵魂）。他们也承认灵魂从一身体转移到另一身体，即对于那些被预定要迁移的灵魂而言。他们接受将有大火的教义，即这个世界的净化，有些人说整个世界的净化，另一些人说部分的净化，并且世界本身正在局部毁灭；这种近乎腐败以及从其中产生另一个世界的过程，他们称之为\OriginalTerm{净化}{purgation}。他们假定所有物体存在，且物体不能穿过物体，而是发生折射，一切事物充满实体，没有\OriginalTerm{虚空}{vacuum}。以上便是斯多亚学派的意见。\par

% structure: p0061 source=h2 target=subsection reason=relative-heading-rank; base=h2

\subsection{第十九章　伊比鸠鲁；采用德谟克利特的原子论；否认天主眷顾；其伦理体系的原则。}

\Person{伊壁鸠鲁}的见解几乎与所有其他人都相反。他假定万物的本原是原子和虚空。他认为虚空是容纳将存在之物的地方，原子则是万物得以形成的质料；从原子的结合中，\OriginalTerm{神}{Deity}、所有元素、其中固有的一切，以及动物和其他受造物都得以存在；因此，除非来自原子，否则没有任何东西被产生或存在。他声称这些原子由极小的微粒构成，其中不可能存在点、符号或任何分割；因此他也称它们为原子。他承认\OriginalTerm{神}{Deity}永恒不朽，却说\OriginalTerm{神}{God}毫不眷顾万物，根本不存在眷顾或命运，一切皆由偶然发生。因为\OriginalTerm{神}{Deity}居住在\Quote{诸世界之间的空间}（这是他如此称呼的）；因为他确定在世界之外有一个\OriginalTerm{神}{God}的居所，称为\Quote{诸世界之间的空间}，并且\OriginalTerm{神}{Deity}委身于快乐，在至高的幸福中安息；既无事务缠身，也不加以关注。基于这些见解，他也提出了关于智者的理论，断言智慧的目的就是快乐。然而，不同的人对\Quote{快乐}一词有不同的理解；因为有些人（外邦人中）理解为情欲，另一些人则理解为来自德行的满足。他断定人的灵魂与身体一同消散，正如它们与身体一同产生，因为它们是血；当血流出或改变时，整个人便消亡；与此一致，\Person{伊壁鸠鲁}主张在阴间既无审判，也无公义法庭；因此，任何人在今生所犯的罪，只要能逃脱察觉，就完全不必担心来世受审。\Person{伊壁鸠鲁}就这样形成了他的见解。\par

% structure: p0063 source=h2 target=subsection reason=relative-heading-rank; base=h2

\subsection{第二十章 学院派；他们之中的意见分歧}

另一派哲学见解被称为学院派，因其在学院中辩论而得名。其创始人\Person{皮浪}（由此他们被称为皮浪派哲学家）首先引入了万物不可理解的概念，以至于对任何问题都试图从正反两方面论证，而不肯定任何确定之事；因为无论可理解之物还是可感之物，都没有什么是真的，只是对人显得如此；一切实体都处于流变之中，从不保持同一状态。因此，一些学院派追随者主张不应就任何事物的原理发表意见，而仅仅尝试放弃它；另一些人则补充了\Quote{并不更}（此物胜于彼物）的公式，说火并不更是火而非其他东西。但他们没有宣称这是什麼，而是说这是怎样的。\par

% structure: p0065 source=h2 target=subsection reason=relative-heading-rank; base=h2

\subsection{第二十一章 婆罗门；他们的生活方式；关于神性的观念；不同的种类；他们的伦理观念}

在印度人中，也有一派由婆罗门中行哲学思考者组成的教派。他们过着知足的生活，戒绝活物和一切熟食，以水果为满足；并非从树上采摘，而是取那些掉落在地上的。他们以此为生，饮用\Place{塔扎贝纳河}的水。他们赤身生活，声称身体是\OriginalTerm{神}{Deity}为灵魂所造的遮蔽。他们断言\OriginalTerm{神}{God}是光，并非人们所见的那种光，也非太阳和火那样的光；对他们而言，\OriginalTerm{神}{Deity}乃是言说，不是发出可闻声音的言说，而是那种知识的言说，智者通过它领悟自然的神秘奥秘。他们说这种光，即他们所说的言说，就是他们的神，只有婆罗门知道，因为他们唯独摒弃了一切意见的虚妄，这虚妄是灵魂最终的遮蔽。他们蔑视死亡，总是用我们先前提到的那名称以自己的特殊语言呼求\OriginalTerm{神}{God}，并向他献上赞美诗。在他们中间也没有妇女，也不生育子女。但那些向往类似生活的人，在渡河到达对岸之地后，便留居在那里，不再返回；这些人也被称为婆罗门。然而他们的生活方式不同，因为那地方也有妇女，住在那裡的人由她们所生，又生育子女。他们声称这种他们称为神的言说是具形的，并包裹在一个自身之外的躯体中，如同有人披着羊皮，但一旦脱去躯体，就会清晰可见。婆罗门说，围绕他们的身体中存在一种冲突（他们视身体为充满冲突的）；他们如同列阵对敌一样与之斗争，正如我们已经说明的。他们说所有人都被自身与生俱来的挣扎所俘虏，即感官欲、不贞、贪食、忿怒、喜乐、忧愁、贪恋等类似之物。唯有在这些方面建立功勋的人才能到神那里去；因此婆罗门将\Person{丹达米斯}奉为神，马其顿人\Person{亚历山大}曾拜访过他，认为他是在身体冲突中得胜的人。他们却责难\Person{卡拉努斯}，认为他亵渎地背离了他们的哲学。婆罗门脱去身体，如同鱼跃出水面进入纯净空气，观看太阳。\par

% structure: p0067 source=h2 target=subsection reason=relative-heading-rank; base=h2

\subsection{第二十二章 德鲁伊特；他们体系的起源}

凯尔特德鲁伊士在\Person{Zamolxis}之后，对毕达哥拉斯哲学进行了最高层次的研究。\Person{Zamolxis}是色雷斯人，\Person{毕达哥拉斯}的仆人，成为他们这门学科的创始人。\Person{毕达哥拉斯}死后，\Person{Zamolxis}前往那里，成为他们这门哲学的创始人。凯尔特人尊崇这些人为先知和先见，因为他们通过毕达哥拉斯技艺的计算和数字预言某些（事件）；关于这技艺的方法，我们也不会保持沉默，因为有些人甚至从中假借来引入异端；但德鲁伊士也使用魔法仪式。\par

% structure: p0069 source=h2 target=subsection reason=relative-heading-rank; base=h2

\subsection{第23章 赫西俄德；九位缪斯；赫西俄德的宇宙起源论；古代思辨者，唯物主义者；异端从外教哲学派生的特征。}

但诗人\Person{赫西俄德}本人也断言他如此从缪斯那里听到关于自然的道理，并且缪斯是\Person{朱庇特}的女儿。因为当连续九个日夜，\Person{朱庇特}因过度情欲与\Person{谟涅摩绪涅}同床共枕时，\Person{谟涅摩绪涅}在一个子宫里孕育了那九位缪斯，每个夜晚怀上一个。随后，他从\Place{皮埃里亚}，即\Place{奥林匹斯山}，召来了九位缪斯，鼓励她们接受教诲：—\par

首先诸神和大地如何被造，\newline{}以及河流、无底的深渊和海洋的波涛，\newline{}和闪亮的星辰以及广阔的上天；\newline{}他们如何取得冠冕并分享荣耀，\newline{}以及最初他们如何占据多谷的奥林匹斯。\newline{}这些（真理），你们缪斯，告诉我吧，他说，\newline{}从最初开始，然后哪一个最先出现。\newline{}混沌无疑最先出现；但随后\newline{}广阔的大地，永远是一切\newline{}不死者的稳固宝座，他们占据白色奥林匹斯之巅；\newline{}和广阔大地深处的和风塔耳塔罗斯；\newline{}以及爱神，他是不死之神中最美的，\newline{}驱散诸神和众人心中的忧虑，\newline{}平息胸中的心智和审慎的计谋。\newline{}但暗界从混沌和阴沉的黑夜产生；\newline{}然后，从黑夜产生了空气和白昼；\newline{}而原始大地，真理上与自身相等，生出了\newline{}暴风雨的天空，将其从四周笼罩，\newline{}永远作为幸福之神的稳固宝座。\newline{}她又孕育了高耸的山峦，那是\newline{}居住在整个多林高地的宁芙们的愉快居所。\newline{}她还不生育了荒芜的大海，生出波涛汹涌\newline{}的洪水，与甜蜜的爱神无关；但随后\newline{}拥抱天空，她生出了有深漩涡的海洋，\newline{}以及科俄斯、克里俄斯、许珀里翁、伊阿珀托斯、\newline{}忒亚、瑞亚、忒弥斯、谟涅摩绪涅、\newline{}金冠福柏和可爱的特提斯。\newline{}但在这之后，最后出生的是狡猾的克洛诺斯，\newline{}最凶狠的儿子；但他憎恶他盛开的父亲；\newline{}然后又生出了独目巨人，他们拥有野蛮的心胸。\par

所有其余从克洛诺斯所出的巨人，\Person{赫西俄德}也都列出，而后在某处说\Person{朱庇特}由\Person{瑞亚}所生。因此，所有这些人在他们的学说中对宇宙的本性和生成都做出了上述声明。但所有人都沉沦于神圣之下，忙碌于现存事物的实体，惊异于受造物的宏大，并认为它就是神性，每位思辨者偏爱选择世界的不同部分；然而，他们未能辨认出这些事物的天主和创造者。\par

因此，我认为我们已经充分阐述了那些在希腊人中试图建立哲学体系之人的意见；而异端分子则利用这些机会，试图确立那些不久后将声明的教义。然而，看来有必要先解释神秘仪式以及某些人关于星宿或巨量费力编造的幻想学说，再予以阐明；因为异端分子也同样从它们中借取，被群众视为讲述奇迹。其次，我们将阐明他们所提出的薄弱意见。\par

\Emph{第2、3卷缺失。}\par

\SourceRecord{Translated by J.H. MacMahon.；Ante-Nicene Fathers；Vol. 5.；Edited by Alexander Roberts, James Donaldson, and A. Cleveland Coxe.；Buffalo, NY: Christian Literature Publishing Co.,；1886.；<http://www.newadvent.org/fathers/050101.htm>.}

% structure: p0001 source=h1 target=section reason=page-title

\section{\WorkTitle{驳斥一切异端（卷四）}}

% structure: p0002 source=h2 target=subsection reason=relative-heading-rank; base=h2

\subsection{第一章 占星术士的系统；星辰影响；星宿配置。}

但在每个黄道星宫，他们称之为星辰的界限，其中每一颗星辰，从某一方位到另一方位，能施加最大的影响；关于这一点，根据他们的著作，他们之间并非偶然的歧见。但他们说，当星辰处于其他星辰中间，与黄道星宫连续时，它们如同有卫星伴随；例如，当某一特定星辰占据了同一黄道星宫的最初部分，而另一颗占据了最后部分，另一颗占据了中间部分时，中间的那颗被认为由两端部分的那两颗守护。它们也被说成彼此观望，彼此结合，如同呈现三角形或四边形。因此，它们假设三角形，彼此观望，其间隔距离延伸三个黄道星宫；它们假设四边形，其间隔距离延伸两个星宫。但正如下部部分与头部共鸣，头部与下部部分共鸣，同样，地上的事物也与超月的事物共鸣。但这些事情之间存在某种差异和缺乏共鸣，以至于它们不涉及同一连接点。\par

% structure: p0004 source=h2 target=subsection reason=relative-heading-rank; base=h2

\subsection{第二章 有关永世的教义；迦勒底占星术；由此产生的异端。}

运用这些（类比），\Person{幼发拉底斯的珀拉提科}、\Person{卡律斯托斯的阿肯贝斯}，以及其余这群（玄想家），为真理的教义另立名目，谈论永世的反叛，善势力对恶势力的背叛，以及善与恶（永世）的和好，称他们为\OriginalTerm{地方官}{Toparchai}和\OriginalTerm{郊区者}{Proastioi}，以及许多其他名号。但整个这种异端，正如他们所尝试的，我将在论及这些（永世）的题目时加以解释和驳斥。但现在，以免有人以为迦勒底人关于占星学说提出的意见是可靠而稳固的，我们将毫不犹豫地提供对此的简要驳斥，确立这种虚妄的技艺旨在欺骗并使沉迷于虚妄期盼的灵魂盲目，而非使其受益。我们并非根据任何这门技艺的经验，而是基于实践原则的知识来论证我们的观点。那些修习这门技艺的人，成为迦勒底人的门徒，并以仿佛奇事惊人向人传授奥秘，仅更改了名号，便从这源头炮制出他们的异端。但由于他们视占星术为强大技艺，并利用其拥护者提出的见证，希望为其自身的结论赢取信任，我们现在，依所宜，证明占星术站不住脚，因为我们接下来的意图是也要驳斥珀拉提科体系，作为从不稳定根上长出的枝干。\par

% structure: p0006 source=h2 target=subsection reason=relative-heading-rank; base=h2

\subsection{第三章 命宫是占星术的基础；命宫的不可发现性；因此迦勒底法术的虚妄。}

整个技艺的起源，以及可说基础，在于确立命宫。因为由此派生出其余基本点，以及倾斜和上升、三角形和四边形，以及星辰依此的配置；而从这一切得出预言。因此，如果命宫被移除，必然既无法在子午线、地平线或天顶对极点识别任何天体；但若这些不被理解，迦勒底人的整个体系便随之消失。然而，命宫之星标对他们而言是不可发现的，这我们可以通过多种论证表明。因为为了找到它，首先，受审视者出生的时间必须确定无疑；其次，应是无误地标示此的命宫；第三，黄道星宫的上升应被精确观测。因为从出生时刻起，应密切注视天上上升的黄道星宫的升度，因为迦勒底人由此确定命宫，并根据升度配置星辰；他们称此为配置，据此设计预言。但既不可能取得受审视者的出生（如我将说明），命宫也非无误，且上升的黄道星宫也不能被精确把握。\par

那么，加色丁人的体系如何不稳固，现在让我们来说明。他们预先划定研究的范围，然后毫无疑义地从种子播下、受孕或分娩来推算所考察之人的出生。如果有人试图从受孕来推算（星宫），其精确计算是不可理解的，因为时间流逝迅速，且自然如此。因为我们无法说出受孕是否发生在种子转移的那一刻。因为这事可以发生得像思想一样快，正如酵母放入热罐中，立即变成黏稠状态。但受孕也可以在经过一段时间后发生。因为从子宫口到基底有一段间隔，医生认为受孕发生于此，所播下的种子完全需要一些时间才能穿越这段间隔。因此，加色丁人既然不能精确知道时间的长度，就永远不能把握受孕的时刻；种子有时径直注入，落在子宫中适宜受孕的部位上，有时散落进去，由子宫的力量聚集到一处。既然这些事未知——即何时发生第一种情况，何时发生第二种，以及在那种受孕中花费多少时间，在这种中花费多少——我再说，既然在这些点上对这些（占星家）来说是无知的，那么精确理解受孕就无法实现。如果像一些自然哲学家所主张的，种子先静止不动，在子宫中发生变化，然后进入张开的血管，如同地里的种子渗入土中；由此推知，那些不知道变化所需时间量的人，也不会知道受孕的确切时刻。此外，既然女人在身体其他部分彼此不同，无论是在活力上还是其他方面，同样（可以合理假定）在子宫的活力上也彼此不同，有些受孕较快，有些较慢。这并不奇怪，因为女人自己与自己相比时，有时观察到有强烈的受孕倾向，有时却没有。既然这样，就不可能准确说出所播下的种子何时结合，以便加色丁人从此时定出出生的星宫。\par

% structure: p0009 source=h2 target=subsection reason=relative-heading-rank; base=h2

\subsection{第四章　不可能确定星宫；在出生时期尝试这样做也失败。}

因此，从受孕（时期）不可能确定星宫。但从出生（时期）也不能做到。因为首先存在这样一个困难：何时可以宣告出生？是当胎儿开始倾向产道口时，还是当它稍微突出时，还是当它落到地上时？在这些情况下，没有一种能够把握分娩的精确时刻，或确定时间。因为由于灵魂的气质、身体的适宜、部位的选择、接生婆的经验以及其他无数原因，胎儿倾向产道口、胎膜破裂、稍微突出或落地的时刻并不相同；在不同的人身上时期不同。当加色丁人不能确定而精确地计算这一点时，他们就会理所当然地无法确定出生的时刻。\par

那么，加色丁人声称知道出生时期的星宫，但实际上并不知道，从这些考虑来看是明显的。但他们的星宫也并非无误，这很容易推断。因为当他们声称，在分娩时坐在产妇旁边的人通过敲击金属盘给加色丁人一个信号，而加色丁人从高处观察星象，他仰望天空，记下升起的黄道星座；首先，我们要向他们证明，当分娩的发生是不确定的（正如我们之前稍示），甚至不容易通过敲击金属盘来传达这个（出生）。然而，即使出生是可理解的，也不可能在精确时刻传达；因为金属盘的声音可以被更长时间和延长的分割，就感知而言，声音传到高处需要（成比例的延迟）。从远处伐木的人可以看到以下证明：在斧头落下后相当长的时间后，才能听到敲击声，因此传到听者需要更长时间。因此，加色丁人不能精确地记录升起星座的时刻，因此也就不能真实地确定星宫。不仅分娩后，坐在产妇旁的人敲击金属盘，然后声音传到听者即上到高处的人，似乎有更多时间过去；而且当他环顾四周，查看月亮在哪个星座、其他星各自出现在哪里时，必然导致星的位置不同，因为极点的运动以无法计量的速度推动它们，在仔细调整天上所见与出生时刻之前。\par

% structure: p0012 source=h2 target=subsection reason=relative-heading-rank; base=h2

\subsection{第五章　另一种在出生时确定星宫的方法；同样无效；水钟在占星术中的使用；加色丁人的预言未被证实。}

这样，加色丁人所实践的技艺将被证明是不稳定的。然而，如果有人声称，通过向询问加色丁人的人提出问题，可以确定出生时间，即使通过这种方法也不可能精确确定。因为，假设他们对自己的技艺有如此关注并留有记录，这些人也无法达到精确——正如我们所证明的——那么，请问，一个未经训练的人怎么能准确理解分娩的时间，以便加色丁人从这个人那里获得必要信息后正确设置星宫图呢？但即使从地平线的外观来看，升起的星星也不会在任何地方看起来相同；而是在一个地方，其偏角被认为是星宫图，在另一个地方，升度被认为是星宫图，这取决于这些地方是较低还是较高而呈现。因此，从这个方面，精确的预测也不会出现，因为全世界许多人在同一时刻出生，每个人从不同的方向观察星星。\par

但通过水钟来理解分娩时间的设想也是徒劳的。因为当陶罐满时和半空时，里面的水流出所需的时间不同；然而，按照他们自己的说法，极轴本身以单一脉冲匀速转动。但是，如果回避论证，他们断言自己不精确地取时间，而是按照特定纬度的情况，那么他们几乎会被星辰影响本身所反驳。因为那些在同一时间出生的人并不过着同样的生活，例如，有些人成了国王，另一些人则在锁链中老去。从来没有出生过任何能与亚历山大大帝相媲美的人，尽管全世界有许多人与他同时出生；也没有能与哲学家柏拉图相媲美的人。因此，加色丁人检查特定纬度的出生时间，无法准确说出此时出生的人是否会成功。我认为，许多在这个时候出生的人都是不幸的，因此根据性情预测的相似性是徒劳的。\par

然后，我们用不同的理由和各种方法驳斥了加色丁人无效的考察方式，也不应省略这一点，即证明他们的预测将导致无法解决的困难。因为，如果按照占星家的断言，生于射手座箭头之下的人必遭横死，那么，在马拉松或萨拉米斯与希腊人作战的如此众多的蛮族，怎能同时被屠杀？毫无疑问，至少在他们所有人中，不可能有相同的星宫图。此外，据说生于宝瓶座瓮下的人会遭遇船难：那么，许多从特洛伊返回的希腊人，怎么会在优卑亚海岸的凹陷处沉入深海呢？因为令人难以置信的是，所有相隔很长时间的人都生于宝瓶座瓮下。因为说一个人注定要在海上毁灭，那么同船的所有人都应当与他一起灭亡，这是不合理的。因为为什么这个人的命运会征服所有人的命运？相反，为什么为了一个注定要在陆地上死亡的人，所有人不都应该得救呢？\par

% structure: p0016 source=h2 target=subsection reason=relative-heading-rank; base=h2

\subsection{第六章 黄道带影响；星辰名称的起源}

但是，既然他们也编造了关于黄道带星座作用的说法，并说生育的受造物与其相似，我们也不应忽略这一点：例如，生于狮子座的人勇敢；生于室女座的人头发长而直，肤色白皙，无子女，端庄。然而，这些说法以及其他类似的说法，更应引人发笑，而非严肃考虑。因为，按照他们的说法，不可能有埃塞俄比亚人生于室女座；否则他们会承认这样的人是白人，长发直而直，还有其他特征。但我更倾向于认为，古人为了便于认识，给某些特定的星星赋予了熟悉的动物名称，并非出于本性的相似；因为七颗彼此远离的星星与熊有什么共同之处？五颗星星与龙头又有什么共同之处？关于这一点，阿拉托斯说：\par

\Quote{但有两个他的太阳穴，两个他的眼睛，一个在下，\newline{} 到达巨大怪物的法则的尽头。}\par

% structure: p0019 source=h2 target=subsection reason=relative-heading-rank; base=h2

\subsection{第七章 加色丁技艺的实际荒谬；该技艺的发展}

同样，对于那些喜欢正确思考、不关注加色丁人夸张言辞的人来说，这些论点不值得如此费心，显然也是如此；加色丁人使君主陷入完全的隐晦，以助长他们的懦弱，并鼓动私人敢于大胆的壮举。但是，如果有人屈服于邪恶而行恶，那么被他这样欺骗的人并不会成为所有那些加色丁人试图用他们的错误误导的人的老师。（远非如此）；（这些占星家）使（他们的受骗者的）心灵陷入无尽的混乱，（当他们）断言，同一星座的配置只有通过大年的更新才能回到相似的位置，大年跨越七千七百七十七年。那么，请问，人的观察如何能够与如此多的年代相协调，而且不是一次，（而是多次，因为如有些人所说，世界毁灭会中断这个大年的进程；或者一次地震，即使是局部的，也会完全打破历史传统的连续性）？加色丁技艺必然要被更多论据所驳斥，尽管我们因其他情况而非特别因该技艺本身而提醒（读者）注意它。\par

然而，由于我们决定不略过外邦哲学家所提出的任何见解，以免异端者的恶名昭彰，让我们也看看那些试图就度量提出学说的人——他们观察到大多数（思辨者）徒劳无功，各自以不同方式编造谎言并赢得声誉——于是敢于提出更重大的论断，以便那些极力夸耀其可鄙谎言的人能高度颂扬他们。他们认为存在圆、度量、三角形和正方形，既有二重也有三重排列。他们关于此事的论证虽然广泛，但对于我们已着手讨论的主题而言并无必要。\par

% structure: p0022 source=h2 target=subsection reason=relative-heading-rank; base=h2

\subsection{占星家的奇谈；天文学家的体系；迦勒底人的圆周论；天体的距离。}

我认为，仅陈述这些人详述的奇谈便已足够。因此，我将用简明的叙述来阐明他们的主张，然后转向其他（尚待考虑的）要点。他们现在作如下陈述。造物主将卓越的能力赐予同类（圆）的轨道运动，因祂容许这旋转是单一且不可分的；但将此内部划分为六部分后，（形成）七个不等距的圆周，依照二倍和三倍间隔的每个距离，祂命令——因各有三个——这些圆周应彼此反向运行，其中三个（从七个总数中）以相等速度旋转，四个则以彼此不同且与其余三个不同的速度旋转，但（全部）依照确定的原则。因为他断言，同类（圆）的轨道运动被赋予主宰权，不仅因为它包含其他运动，即游星的运动，而且因为它拥有如此强大的主宰力，即如此巨大的力量，以致它甚至以其独特的力量带动那些天体——即游星——它们从西向东、同样从东向西反向旋转。\par

他主张这运动被允许是单一且不可分的，首先，因为所有恒星的旋转是在相等时间周期内完成的，并不以较大或较小时间段来区分。其次，它们都呈现出与最外层运动相同的相位；而游星则被分配了较大且变化的时间周期来完成其运动，以及到地球的不等距离。他还主张，另一（圆）的六部分运动被合理地分配为七个圆周。因为每个（圆）的部分有多少节段——我指的是部分的单子——就产生多少片段；例如，若按一节划分，则有两个片段；若按两节划分，则有三个片段；同样，若某物被分为六部分，则有七个片段。他还说，这些距离按二倍和三倍顺序交替排列，各有三个——他试图证明这一原则也适用于灵魂的构成，因其依赖于七个数字。因为其中从\OriginalTerm{单子}{monad}有三个二倍数，即2、4、8，和三个三倍数，即3、9、27。地球直径为80,108斯塔狄；地球周长为250,543斯塔狄；从地球表面到月球圆周的距离，\Person{萨摩斯的阿利斯塔克}计算为8,000,178斯塔狄，\Person{阿波罗尼奥斯}计算为5,000,000斯塔狄，而\Person{阿基米德}计算为5,544,1300斯塔狄。从月球到太阳圆周，（根据最后一位权威）为50,262,065斯塔狄；从太阳到金星圆周为20,272,065斯塔狄；从金星到水星圆周为50,817,165斯塔狄；从水星到火星圆周为40,541,108斯塔狄；从火星到木星圆周为20,275,065斯塔狄；从木星到土星圆周为40,372,065斯塔狄；从土星到黄道带和最外沿为20,082,005斯塔狄。\par

% structure: p0025 source=h2 target=subsection reason=relative-heading-rank; base=h2

\subsection{进一步的天文计算。}

诸圆周和球面之间的相互距离以及深度，由\Person{阿基米德}给出。他取黄道带周长为447,310,000斯塔狄；因此，从地球中心到最外表面的直线是前述数字的六分之一，但从我们行走的地球表面到黄道带的直线是前述数字的六分之一减去四万斯塔狄，即从地球中心到其表面的距离。他还说，从土星圆周到地球的距离是2,226,912,711斯塔狄；从木星圆周到地球是502,770,646斯塔狄；从火星圆周到地球是132,418,581斯塔狄；从太阳到地球是121,604,454斯塔狄；从水星到地球是526,882,259斯塔狄；从金星到地球是50,815,160斯塔狄。\par

% structure: p0027 source=h2 target=subsection reason=relative-heading-rank; base=h2

\subsection{依照和谐论的星体运动与距离理论。}

关于月亮，前面已经有过论述。\Person{阿基米德}对诸球体的距离和深度作了如此的叙述；但\Person{喜帕恰斯}对此提出了不同的说法，而数学家\Person{阿波罗尼奥斯}又提出了不同的说法。然而，对我们来说，遵循\Person{柏拉图}的观点，假定游星彼此间具有二倍和三倍的距离，这就足够了；因为这样便保存了宇宙由和谐构成、按照与这些距离一致的和声原理的学说。但是，\Person{阿基米德}所提出的数字以及其他人对距离的叙述，如果不符合和声原理——即\Person{柏拉图}所说的双倍和三倍距离——而是独立于和谐关系发现的，那么它们就不能保存宇宙按照和谐形成的学说。因为除非可能只有月亮，由于月亏和地球的阴影，才可以在其距离——即月球行星距地球的距离——上相信\Person{阿基米德}，否则这些距离既违背合理计划又独立于和谐与比例原则，是不可信且不可能的。然而，对于那些按照\Person{柏拉图}教义本身采用这个距离的人来说，很容易通过数值计算按照\Person{柏拉图}所要求的双倍和三倍来理解其余的距离。那么，如果按照\Person{阿基米德}，月亮距离地球表面5,544,130斯塔迪，将这些数字加倍和三倍，也很容易求出其余星体的距离，就像从月亮距地球的斯塔迪数中减去一部分一样。\par

但是，因为其余的数字——\Person{阿基米德}所声称的关于游星距离的数字——并非基于和声原理，所以很容易理解——即对于那些注意此事的人——这些数字如何相互关联以及它们依赖什么原理。然而，它们不应该和谐与和声——我指的是那些按照和谐构成的世界之部分——这是不可能的。因此，既然月亮距地球的第一个数字是5,544,130，而太阳距月亮的第二个数字是50,272,065，它大于九倍。但相对于此的更高数字，即20,272,065，大于一半。然而，高于此的数字，即50,817,165，大于一半。但高于此的数字，即40,541,108，小于五分之二。但高于此的数字，即20,275,065，大于一半。然而，最后的数字，即40,372,065，小于两倍。\par

% structure: p0030 source=h2 target=subsection reason=relative-heading-rank; base=h2

\subsection{第十一章　根据数字和谐的天体大小理论}

因此，这些（数字）关系——大于九倍、小于一半、大于两倍、小于五分之二、大于一半、小于两倍——超出了所有和声，从这些和声中不能产生任何比例或和谐的系统。但整个宇宙及其各部分在所有方面都同样按照比例与和谐构成。然而，比例与和谐的关系——如我们先前所述——是通过双倍和三倍间隔得以保存的。因此，如果我们认为\Person{阿基米德}只在第一个距离——即月亮到地球的距离——上可靠，那么通过乘以双倍和三倍也容易求得其余（间隔）。那么，按照\Person{阿基米德}，从地球到月亮的距离设为5,544,130斯塔迪；因此，太阳距月亮的斯塔迪数将是这个数字的两倍，即11,088,260。而太阳距地球为16,632,390斯塔迪；同样，金星距太阳为16,632,390斯塔迪，但距地球为33,264,780斯塔迪；水星距金星为22,176,520斯塔迪，但距地球为55,441,300斯塔迪；火星距水星为49,897,170斯塔迪，距地球为105,338,470斯塔迪；木星距火星为44,353,040斯塔迪，但距地球为149,691,510斯塔迪；土星距木星为149,691,510斯塔迪，但距地球为299,383,020斯塔迪。\par

% structure: p0032 source=h2 target=subsection reason=relative-heading-rank; base=h2

\subsection{第十二章　占星家系统中智力的虚耗}

谁不会对花费如此多劳苦的深思熟虑感到惊讶？然而，这位\Person{托勒密}——这些事情的仔细研究者——在我看来并非无用；但只有这一点令人遗憾（：他生于近代，不能为巨人的子孙们服务，他们不知这些尺度，以为天穹的高度很近，妄图建造一座塔楼。因此，如果那时他在场向他们解释这些尺度，他们就不会徒劳地做出这大胆的尝试了。但如果有人声称不相信这位（天文学家的）计算，就让他通过测量被说服（其准确性）；因为对于那些怀疑此点的人，没有比这更明显的证据了。啊，徒劳劳作的灵魂的骄傲，以及令人难以置信的信念，竟认为\Person{托勒密}在那些致力于同样智慧的人中被视为卓越超群！\par

% structure: p0034 source=h2 target=subsection reason=relative-heading-rank; base=h2

\subsection{第十三章　提及异端者\Person{科拉巴苏斯}；异端与\Person{毕达哥拉斯}哲学的同盟}

某些人，部分依附于这些（学说），仿佛提出了伟大的结论，并假设了合理的事物，却构想了巨大而无尽的异端；其中一个是\Person{科拉尔巴苏斯}，他试图用度量与数字来解释宗教。另有一些人以类似方式行事，他们的教义我们将在开始谈论那些关注\OriginalTerm{毕达哥拉斯计算法}{Pythagorean calculation}为可能的人时加以解释；他们发出虚妄的预言，轻率地假定\OriginalTerm{数字与元素哲学}{philosophy by numbers and elements}是可靠的。现在某些（思辨者）借用类似的推理，欺骗天真的人，声称自己具有\OriginalTerm{预知}{foresight}能力；有时在发出许多预言后，碰巧实现一个，并不因多次失败而羞愧，反而以此自夸。我也不会放过这些人的愚蠢哲学；但在解释之后，我将证明那些试图从上述元素中建立宗教体系的人，是一个软弱且充满欺诈的学派的门徒。\par

% structure: p0036 source=h2 target=subsection reason=relative-heading-rank; base=h2

\subsection{第十四章 算术家的体系；通过计算的预言；数字根；这些学说到字母的迁移；具体名称的实例；不同的计算方法；借此可能的预知}

那些自以为通过计算、数字、元素和名称来预言的人，他们创立的体系的起源如下。他们断言每个数字都有一个根；对于千位数，有多少千就有多少单子：例如，六千的根是六个单子；七千的根是七个单子；八千的根是八个单子；其余依此类推。对于百位数，有多少百就有多少单子作为其根：例如，七百有七个百；其根是七个单子；六百有六个百；其根是六个单子。对于十位数也是如此：八十的根是八个单子；六十的根是六个单子；四十的根是四个单子；十的根是一个单子。对于个位数，个位数本身就是一个根：例如，九的根是九；八的根是八；七的根是七。因此，我们也应当以同样的方式处理（词语的）元素，因为每个字母都按照一定的数字排列：例如，字母\Emph{n}对应五十个单子；但五十个单子的根是五，因此字母n的根是五。假设从某个名称中我们取其一些根。例如，名称\Person{阿伽门农}：字母a的根是一个单子；字母g的根是三个单子；另一个字母a的根是一个单子；字母m的根是四个单子；字母e的根是五个单子；字母m的根是四个单子；字母n的根是五个单子；字母o（长音）的根是八个单子；字母n的根是五个单子；这些数字排成一列，便是1,3,1,4,5,4,5,8,5；这些加起来总共是三十六个单子。然后，他们再取这些数的根，三十的根是三，而六的根是六。三和六相加得九；九的根是九：因此名称\Person{阿伽门农}最终以根九结束。\par

让我们用另一个名字\Person{赫克托尔}做同样的事。名字赫克托尔有五个字母——e、k、t、o、r。这些字母的根是5、2、3、8、1；这些加起来是十九个单子。再次，十的根是一个；九的根是九；它们加起来是十；十的根是一个单子。因此，名字赫克托尔经过计算后形成一个根，即一个单子。然而，更容易的计算方法是：将自字母得出的根——如现在赫克托尔的名字我们找到十九个单子——除以九，余数视为根。例如，如果我把19除以9，余数是1，因为9乘以2是18，余下一个单子；所以名字赫克托尔的根将是一个单子。再次，名字\Person{帕特罗克洛斯}的根是这些数字：8、1、3、1、7、2、3、7、2；加起来是三十四个单子。其中的余数是七个单子：三十的余数是三，四的余数是四。因此，名字\Person{帕特罗克洛斯}的根是七个单子。\par

那些按照九数规则进行计算的人，取根的总和的第九部分，将剩余部分定义为根的和。另一方面，按照七数规则进行计算的人，取根的第七部分；例如，就\Person{帕特罗克洛斯}这个名字而言，根的总和为34个单子。除以七得四，四乘七得二十八，余六个单子；因此，根据七数规则，六个单子就是\Person{帕特罗克洛斯}这个名字的根。然而，如果是43，他说六乘七得四十二，余一。因此，根据七数规则，一个\Emph{单子}就是数字43的根。但应当注意，如果所取的数字除尽无余，例如，从某个名字中，将根相加后，我找到36个单子。36除以九正好得4个\OriginalTerm{九}{\GreekText{εννεαδς}}；因为九乘四得三十六，没有余数。显然，实际的根就是九。同样，除以四十五，我们得到九且无余——因为九乘五得四十五，没有剩余；因此在这种情况下他们断言根本身就是九。至于七数，情况类似：例如，28除以7，没有余数；因为七乘四得二十八，没有剩余；因此他们说七是根。但是当计算名字时，发现同一字母出现两次，则只计算一次；例如，\Person{帕特罗克洛斯}这个名字中\Emph{pa}出现两次，\Emph{o}出现两次：因此他们只计算\Emph{a}一次，\Emph{o}一次。据此，根将是8、1、3、1、7、2、3、2，相加得27个单子；根据九数规则，这个名字的根是九本身，但根据七数规则，则是六。\par

同样，当计算\Person{萨耳珀冬}这个名字时，根据九数规则，它的根是2个单子。而\Person{帕特罗克洛斯}则产生9个单子；\Person{帕特罗克洛斯}获胜。因为当一个数奇而另一个偶时，较大的奇数取胜。但是，当有一个偶数8和一个奇数5时，8取胜，因为它更大。然而，如果有两个数，例如都是偶数或都是奇数，则较小的取胜。但是\Person{萨耳珀冬}这个名字，根据九数规则，怎么会是2个单子，因为长\Emph{o}被省略了？因为当一个名字中有长\Emph{o}和长\Emph{e}时，他们省略长\Emph{o}，只用一个字母，因为他们说两者等值；且同一字母不得重复计算，如前所述。同样，\Person{埃阿斯}这个名字是4个单子；而\Person{赫克托耳}，根据第九数规则，是1个单子。四元是偶数，而单子是奇数。在这种情况下，我们说较大的取胜——\Person{埃阿斯}获胜。又如\Person{亚历山大}和\Person{墨涅拉俄斯}。\Person{亚历山大}有一个本名（\Person{帕里斯}）。但是\Person{帕里斯}，根据九数规则，是4个单子；而\Person{墨涅拉俄斯}，根据九数规则，是9个单子。然而9胜过4：因为如前所述，当一奇一偶时，较大的取胜；但当两者皆偶或皆奇时，较小的取胜。又如\Person{阿密科斯}和\Person{波吕丢刻斯}。\Person{阿密科斯}，根据九数规则，是2个单子，而\Person{波吕丢刻斯}是7个：\Person{波吕丢刻斯}获胜。\Person{埃阿斯}和\Person{乌利西斯}在葬礼赛会上竞争。\Person{埃阿斯}，根据九数规则，是4个单子；\Person{乌利西斯}，根据九数规则，是8个单子。那么，\Person{乌利西斯}难道没有附加名字吗？他取得了胜利。根据数字，无疑\Person{埃阿斯}取胜，但历史记载\Person{乌利西斯}为胜者。又如\Person{阿喀琉斯}和\Person{赫克托耳}。\Person{阿喀琉斯}，根据九数规则，是4个单子；\Person{赫克托耳}是1个：\Person{阿喀琉斯}获胜。又如\Person{阿喀琉斯}和\Person{阿斯忒洛派俄斯}。\Person{阿喀琉斯}是4个单子，\Person{阿斯忒洛派俄斯}是3个：\Person{阿喀琉斯}征服。又如\Person{墨涅拉俄斯}和\Person{欧福耳玻斯}。\Person{墨涅拉俄斯}有9个单子，\Person{欧福耳玻斯}有8个：\Person{墨涅拉俄斯}获胜。\par

然而，有些人按照七数规则，只使用元音；但另一些人则区分元音、半元音和哑音三类；分成三类后，分别取元音的根、半元音的根和哑音的根，然后各自比较。另一些人甚至不使用这些常规数字，而使用不同的数字：例如，他们不愿意承认字母\Emph{p}的根是8个单子，而是5个，字母\Emph{x}（si）的根是4个单子；他们四处改变，发现不了任何稳妥之处。然而，当他们争论第二个字母时，他们从每个名字中去掉第一个字母；当他们争论第三个字母时，他们去掉每个名字的两个字母，然后计算剩余部分并加以比较。\par

% structure: p0042 source=h2 target=subsection reason=relative-heading-rank; base=h2

\subsection{第十五章．数字理论家的诡辩；面相术的艺术；此术与占星术的关联；白羊座下出生者的类型。}

我想，算术家的心思已被清楚地阐明，他们藉着数字和名字，以为能解释生命。如今我看出这些人，既得闲暇，又受算数训练，渴望藉着自幼习得的技艺，获得名声，被称为先知。他们上下测量字母，终归琐碎。因为倘若他们失败，便借口困难，说：\Quote{也许这名不是家传，而是后来附加的}，正如他们在讨论\Person{尤利西斯}和\Person{埃阿斯}之名时所辩称的。有谁藉此惊人的哲学，渴望被称为\Quote{异端魁首}，而不得受赞扬呢？\par

然而，在希腊人那些全智的思辨者中，还有另一门更深的技艺——异端之徒夸耀自己依附这些人为门徒，因为他们利用这些（古代哲人）的意见，来支持他们自己试图（建立的）教义，这一点稍后将证明；但这乃是一种藉检验额头进行占卜的技艺，或更确切地说，是疯狂；然而我们对此（体系）亦不缄默。有些人将塑造人的观念与性情的形状归因于星宿，将这理由归于（在特定星座下的）出生；他们如此表述：生于白羊座的人将属以下类型：长头、红发、收拢的眉毛、尖额、灰而活泼的眼睛、凹陷的面颊、长鼻、张大的鼻孔、薄唇、尖下巴、阔嘴。这些人，他说，将具有以下性情：谨慎、狡猾、敏锐、慎重、纵容、温和、过分焦虑、密谋之人、适合一切事业、凭谨慎而非力量取胜、暂时的嘲笑者、学者、可信赖、好争论、斗殴中的争吵者、贪欲的、燃起非自然的情欲、深思的、离家漂泊、事事令人不满、控告者、如醉酒的狂人、轻蔑者、年复一年因善良而在友谊中失去些许有用之物；他们大多客死异乡。\par

% structure: p0045 source=h2 target=subsection reason=relative-heading-rank; base=h2

\subsection{第十六章　生于金牛座之人的类型}

生于金牛座的人则具以下描述：圆头、厚发、宽额、方眼、粗大的黑眉；在白人身上，青筋显见、血色丰润、长眼睑、粗大的耳朵、圆嘴、厚鼻、圆鼻孔、厚唇、上身结实、双腿挺直。这些人天性讨人喜欢、善于思考、性情良善、虔诚、公正、粗野、随和、十二岁起劳碌、好争吵、迟钝。他们的胃小，易饱，多谋，谨慎，对自己吝啬，对他人慷慨，仁慈，动作迟缓；他们部分忧伤，对友谊漫不经心，因心智而有用，不幸。\par

% structure: p0047 source=h2 target=subsection reason=relative-heading-rank; base=h2

\subsection{第十七章　生于双子座之人的类型}

生于双子座的人将具以下描述：红润面色，身材不大，四肢匀称，黑眼如涂油，面颊下垂，大嘴，收拢的眉毛；他们征服一切，凡所得财物皆保留，极富，吝啬，对自己之物悭吝，纵情女色，公正，善音乐，说谎。这些人天性博学、善于思考、好奇、自作主张、贪欲的、对自己之物节省、慷慨、安静、谨慎、狡猾、多谋、算计者、控告者、纠缠不休、不顺利、受女性喜爱、商人；就友谊而言，不甚有用。\par

% structure: p0049 source=h2 target=subsection reason=relative-heading-rank; base=h2

\subsection{第十八章　生于巨蟹座之人的类型}

生于巨蟹座的人具以下描述：身材不大，发如犬毛，略带红色，小嘴，圆头，尖额，灰眼，足够美丽，四肢稍有不匀。这些人天性邪恶、狡猾、精于计谋、贪婪、吝啬、不优雅、不慷慨、无用、健忘；他们既不归还他人之物，也不索回自己之物；就友谊而言，有用。\par

% structure: p0051 source=h2 target=subsection reason=relative-heading-rank; base=h2

\subsection{第十九章　生于狮子座之人的类型}

生于狮子座的人具以下描述：圆头，红发，粗大起皱的前额，粗耳，颈部粗大，部分秃顶，红肤，灰眼，大颚，粗嘴，上身粗壮，巨胸，下肢渐细。这些人天性自专、自娱、易怒、热情、轻蔑、固执、不做计划、不多言、懒惰、滥用闲暇、亲昵、完全沉溺于女色、奸淫、不贞、不信实、纠缠不休、大胆、吝啬、掠夺者、显著；就交往而言，有用；就友谊而言，无用。\par

% structure: p0053 source=h2 target=subsection reason=relative-heading-rank; base=h2

\subsection{第二十章　生于处女座之人的类型}

生于处女座的人具以下描述：外貌清秀，眼睛不大，迷人，深色，眉毛紧凑，喜游泳；然而他们身材纤细，容貌美丽，头发梳得漂亮，前额宽阔，鼻子突出。这些人天性温驯、适中、聪慧、爱嬉戏、理性、寡言、多谋；就恩惠而言，纠缠不休；乐于观察一切；且为良善弟子，凡学习无不精通；适中，轻蔑者，非自然情欲的受害者，好交际，灵魂高尚，轻蔑，对实际事务粗心，注重教导，在关乎他人的事上比关乎自己的事上更为尊贵；就友谊而言，有用。\par

% structure: p0055 source=h2 target=subsection reason=relative-heading-rank; base=h2

\subsection{第二十一章　生于天秤座之人的类型}

凡出生于天秤座者，其相貌如下：头发细疏、下垂、微红而略长；额头尖而有皱纹；眉毛均匀紧凑；眼睛美丽，瞳孔深黑；耳朵细长；头倾斜；嘴宽大。其秉性聪明、敬畏天主、善于彼此交流、为商贾、勤劳、不积财、说谎、不和蔼；在商业或原则上真诚、直言、仁慈、无学识、欺诈、友善、粗心（对这样的人，行任何不义之事皆无益）；他们是嘲笑者、讥讽者、尖刻者、著名者、听众，而诸事不顺于此类人；就友谊而言，有益。\par

% structure: p0057 source=h2 target=subsection reason=relative-heading-rank; base=h2

\subsection{第二十二章 天蝎座出生者之类型}

凡出生于天蝎座者，其相貌如下：少女般的面容，俊美，尖刻，头发微黑，眼睛形状漂亮，额头不宽，鼻孔锐利，耳朵小而收缩，额头有皱纹，眉毛狭窄，面颊凹陷。其秉性狡猾、勤勉、说谎，不将自己的特别计划告知任何人，怀有欺瞒之心，邪恶，嘲弄者，为奸淫所害，发育良好，温顺；就友谊而言，无用。\par

% structure: p0059 source=h2 target=subsection reason=relative-heading-rank; base=h2

\subsection{第二十三章 人马座出生者之类型}

凡出生于人马座者，其相貌如下：身材高大，额头方正，眉毛浓密，显示出力量，头发整齐地突出，肤色微红。其秉性优雅，如同受过教育者，单纯，仁慈；倾向不自然的欲望，好交际，劳苦，恋爱者，被爱者，饮酒时欢乐，清洁，热情，粗心，邪恶；就友谊而言，无用；嘲弄者，有高尚灵魂，傲慢，狡猾；就伙伴关系而言，有益。\par

% structure: p0061 source=h2 target=subsection reason=relative-heading-rank; base=h2

\subsection{第二十四章 摩羯座出生者之类型}

凡出生于摩羯座者，其相貌如下：身体微红，头发灰白突出，嘴圆，眼如鹰眼，眉毛收缩，额头开阔，略秃，上半身具有更多力量。其秉性好哲学，是嘲弄者、讥讽现存事物者，热情，能让步者，可敬，仁慈，热爱音乐练习，饮酒时热情，欢乐，亲切，健谈，倾向不自然的欲望，和蔼，可爱，好争吵的恋人，就伙伴关系而言，良好。\par

% structure: p0063 source=h2 target=subsection reason=relative-heading-rank; base=h2

\subsection{第二十五章 水瓶座出生者之类型}

凡出生于水瓶座者，其相貌如下：身材方正，体型矮小；眼睛锐利、小、凶恶；专横，不和蔼，严厉，易于获得，对于友谊和伙伴关系良好；此外，他们从事海上冒险，航行并丧命。其秉性沉默寡言，谦虚，好社交，犯奸淫，吝啬，精于商业，喧闹，纯洁，友善，可敬，眉毛粗大；他们常生于琐事之中，但（在日后）追求不同的事务；尽管他们可能对某人表示过好意，却无人回报感谢。\par

% structure: p0065 source=h2 target=subsection reason=relative-heading-rank; base=h2

\subsection{第二十六章 双鱼座出生者之类型}

凡出生于双鱼座者，其相貌如下：中等身材，额头像鱼一样尖，头发蓬乱，常很快变灰白。其秉性有崇高灵魂，单纯，热情，吝啬，健谈；在生命初期，他们将昏昏欲睡；他们渴望自己管理事务，名声好，冒险，好胜，控告者，改变居住地，恋爱者，跳舞者；就友谊而言，有益。\par

% structure: p0067 source=h2 target=subsection reason=relative-heading-rank; base=h2

\subsection{第二十七章 此星辰影响理论之虚妄}

既然我们已解释了这些人的惊人智慧，并未隐瞒他们通过观星而得的过分精致的占卜技艺，我也不会对那些受骗者愚行所涉及的事业保持沉默。因为，将人的形状与性情与恒星的名字相比较，他们的体系是多么无力！因为我们知道，最初从事这类研究的人，根据意义的适当性和便于将来辨认，给恒星命名。这些（天体）与动物的形象有何相似之处？或在行为和活动上，（在这两者之间）有何本性上的共通之处？以至于有人断言，出生于狮子座的人易怒，出生于处女座的人温和，或出生于巨蟹座的人邪恶，但那些出生于……\par

% structure: p0069 source=h2 target=subsection reason=relative-heading-rank; base=h2

\subsection{第二十八章 术士之体系；恶魔之咒语；秘密魔法仪式}

……（术士）取一张纸，吩咐询问者用水写下他想要向魔鬼求问的任何问题。然后，他将纸折起，交给侍从，让他拿去投入火中，让上升的烟雾将字迹传送到魔鬼那里。然而，当侍从执行此命令时，（术士）先取下纸张的等份部分，又在其中一些部分上假装魔鬼用希伯来文写字。然后，他焚烧一种称为\OriginalTerm{基菲}{Cyphi}的埃及术士的香料，将（这些纸片）取走，放在香料旁。但询问者碰巧写了字的那张纸，他放在炭火上烧掉了。接着，（术士）看似被神灵附体，匆匆躲进（屋子的）一个角落，发出又大又刺耳、无人能懂的喊叫，……并命令所有在场的人进入，同时喊叫，并呼求\Person{弗林}或其他某个魔鬼。但进入屋子后，当在场的人并排站立时，术士将侍从扔到床上，对他说了几句话，一部分是希腊语，一部分看似是希伯来语，（包含了）术士惯用的咒语。然而，（侍从）便前去询问。而在屋内，术士将绿矾混合物注入一个装满水的容器中，熔化药物，用它喷洒那张据说字迹已被抹去的纸，迫使隐藏的字母再次显现；由此他得知询问者写了什么。如果一个人也用绿矾混合物写字，然后研磨一颗五倍子，用其蒸汽薰蒸，隐藏的字迹就会显现。如果用牛奶写字，然后将纸烤焦，刮下粉末，撒在并摩擦用牛奶写的字迹上，这些字也会显现。尿液、盐水酱、大戟汁和无花果汁也会产生类似效果。但（术士）以此方式得知问题后，便安排如何回答。接着，他命令在场的人进入，手持月桂树枝并摇晃，发出喊叫，并呼求魔鬼\Person{弗林}。因为这些人也必须呼求他；他们理应向魔鬼提出这个请求，却不愿自行提出，因为他们已丧失理智。然而，混乱和喧嚣使他们无法注意到（术士）秘密所做的那些事。但这些都是什么，现在正是我们宣告的良机。\par

因此，陷入极深的黑暗中。因为（术士）声称，凡人的本性无法看见神性之事，认为能与（这些奥秘）交谈就已足够。他让侍从头朝下躺下，并在每一边放置两块小牌，上面用看似希伯来文写上了魔鬼的名字，他说（一个魔鬼）会把其余的话传入他们耳中。但这一做法是必要的，以便在侍从耳边放置某种器具，通过它，他能够传达他选择的一切。首先，他发出一个声音，使（侍从）青年恐惧；其次，他发出嗡嗡声；然后，第三，他通过器具说出他希望青年说的话，并等待事情的结果；接着，他让在场的人保持安静，并吩咐（侍从）说出他从魔鬼那里听到的。但放在他耳边的器具是一种天然器具，即长颈鹤、鹳或天鹅的气管。如果没有这些，也会使用一些不同的人造器具；即一种由十根铜管组成的、相互套接、末端渐细的管子，适合（此目的），通过这些管子，（术士）将任何他想要的话传入耳中。青年惊恐地听到这些（话），以为是魔鬼所说，便奉命说出。如果有人将湿皮革裹在一根棍子上，干燥后拉紧，封住口，然后抽出棍子，将皮革塑造成管状，他也能达到类似目的。若以上任何一种都不在手边，他便取一本书，从里面打开，将其展开到他认为必要的长度，（从而）达到相同效果。\par

但若他事先知道有人将前来提问，他便能更好地应对一切可能情况。若他还能预先知晓问题，他便用药水写下答案，因已有准备，便被认为更显技艺，因他能清楚写出所问之事。若他对问题一无所知，他便做出推测，并说出一些意义模糊、可作多种解释的话，使神谕原本不可理解，却能适用于多种情况，且当事件发生后，预言便被认为与实际相符。接着，他将一个容器装满水，把纸片放入水中，仿佛未写一字，同时掺入铜绿混合物。这样，写有文字的纸片便会浮上水面，承载答案。于是，助手常会感到可怕的幻觉，因为他也会大力击打听众。他将乳香投入火中，又按以下方式操作：用伊特鲁里亚蜡包裹一块所谓的\Quote{化石盐}，将乳香块分成两半，投入一粒盐，再合上，放在燃烧的炭上。当乳香燃尽时，盐粒向上崩起，造成仿佛异象显现的错觉。沉积在乳香中的深蓝色染料会冒出血红火焰，如我们先前所述。但术士用蜡混合紫草根制成猩红色液体，并将蜡放入乳香。他还在炭下放置明矾粉末，使炭移动；当明矾溶解并起泡时，炭便移动。\par

% structure: p0073 source=h2 target=subsection reason=relative-heading-rank; base=h2

\subsection{第二十九章 展示不同的蛋}

他们以下列方式展示不同的蛋。在顶部两端钻孔，取出蛋白，再重新浸入，放入一些朱砂和墨水。然后用蛋壳的细屑封住开口，涂上无花果汁。\par

% structure: p0075 source=h2 target=subsection reason=relative-heading-rank; base=h2

\subsection{第三十章 羊的自杀}

那些使羊自己割头的人采用以下方法：暗中在羊的喉咙涂上腐蚀性药物，将剑放在近旁。羊因发痒而冲向刀刃，在摩擦中被宰杀，头部几乎与躯干分离。药物的配方是：泻根、盐和海葱，等量混合。为使带药者不被察觉，他携带一个用角制成的双格盒子，可见的一格盛放乳香，秘密的一格盛放上述药物。他还在即将死亡的羊耳中注入水银，但这是有毒的药物。\par

% structure: p0077 source=h2 target=subsection reason=relative-heading-rank; base=h2

\subsection{第三十一章 毒杀山羊的方法}

若有人在山羊耳上涂抹蜡膏，据说它们会因呼吸受阻而在不久后死去。据称，这便是它们呼吸时吸气的方式。他们断言，若将公羊的脖子向太阳方向折回，它便会死亡。他们通过涂抹一种称为\Quote{指贝}的鱼的汁液来焚烧房屋。这种效果因海水而起，十分有用。同样，将海洋泡沫与一些甜味材料一起在陶罐中煮沸；若将点燃的蜡烛靠近这沸腾之物，它便会着火并被烧尽；但若将混合物浇在头上，却完全不会烧伤。若再涂上加热的树脂，它会更彻底地烧尽。若再添加一些硫磺，效果更佳。\par

% structure: p0079 source=h2 target=subsection reason=relative-heading-rank; base=h2

\subsection{第三十二章 雷声的模仿及其他幻术}

雷声的制造有多种方式：将大量异常大的石块沿木板滚下，落在铜板上，产生类似雷声的响声。另外，在梳毛工用于加厚布料的薄板周围缠绕一根细绳；然后猛地一拉，使薄板旋转，旋转中发出雷声。这些把戏就是这样被耍弄的。\par

还有另外一些做法，我将加以说明，那些执行这些可笑表演的人将其视为伟大技艺。将盛满沥青的锅放在燃烧的炭上，当沸腾时，他们将手按在上面却不受灼伤；甚至赤脚在炭火上行走也不被烧伤。此外，术士在炉子上放置一座石制金字塔，使其着火，从口中喷出大量烟雾，且呈火红色。然后，将一块亚麻布放在一锅水上，同时投入大量燃烧的炭，却能保持亚麻布不被烧毁。又在屋内制造黑暗，声称能召来神祇或魔鬼；若有人要求他显现阿斯克勒庇俄斯，他便使用以下咒语：\par

\Quote{曾被杀死、又因\Person{福玻斯}而成不死之子的孩童，\newline{}我召唤你前来，相助我的祭献之礼；\newline{}你曾横渡死亡之河的无数逝者，\newline{}在永远哀伤的广大\Place{塔尔塔罗斯}之宫，\newline{}冲破致命巨浪，越过墨色洪流，\newline{}那里是凡人均须漂浮之处，\newline{}在湖畔被撕裂，无尽悲伤与痛苦，\newline{}你自身从阴郁的\Person{普罗塞尔皮娜}手中被夺走。\newline{}无论你居于圣洁的\Place{色雷斯}之境，或可爱的\newline{}\Place{帕加马}，抑或伊奥尼亚的\Place{埃皮达鲁斯}旁，\newline{}众先知的领袖，快乐的神，邀你至此。}\par

% structure: p0083 source=h2 target=subsection reason=relative-heading-rank; base=h2

\subsection{第三十三章 焚烧的\Person{阿斯克勒庇俄斯}；火术}

但在他停止说出这些笑谈之后，一个火红的\Person{阿斯克勒庇俄斯}出现在地板上。然后，他将一个装满水的锅放在中间，呼求所有神祇，他们就出现了。因为任何旁观者向锅里一看，都会看见他们全体，以及带领着吠叫猎犬的\Person{狄安娜}。然而，我们不会畏缩于叙述这些人的（伎俩的）过程，他们是如何试图（实行他们的幻术）的。因为（那术士）将手放在似乎沸腾的沥青锅上，同时加入醋、硝石和湿沥青，在锅下点火。醋与硝石混合，稍受热便搅动沥青，使气泡升到表面，只呈现出沸腾的假象。然而，（那术士）事先频繁用盐水洗手，结果锅中的东西虽然实际上沸腾，却一点也不会烫伤他。但是，如果他事先用桃金娘、硝石、没药与醋的混合液涂抹双手，再频繁用盐水洗手，他就不会被灼伤；而且，如果他双脚涂上鱼胶和火蜥蜴的混合液，他的脚也不会被烧伤。\par

至于那金字塔形物虽然由石头制成却像蜡烛一样燃烧的原因如下：白垩土被塑成金字塔形，颜色像乳白色石头，按以下方式制备：将大量油涂在泥块上，放在炭火上烘烤，再重新涂油并再次烤干，如此重复两三次并多次灼烧，（术士）便设法使它即使浸入水中也能燃烧，因为它自身含有大量油。而炉床自动点燃，当术士奠酒时，下面有热的灰烬而非冷灰，还有精炼乳香和大量\Emph{麻絮}，一捆涂油的蜡烛和内部空心并藏有火的五倍子。稍待片刻，（术士）通过同时将火放入五倍子、用麻絮缠绕它并向口中吹气，使（金字塔形物）从口中冒烟。然而，那放在锅周围并被他放上炭的亚麻布，由于下面有盐水，不会被烧着；此外，它本身也曾在盐水中洗过，然后涂上蛋白和明矾。同样地，若在其中加入景天汁和醋，并事先长时间用此药水涂抹，那么洗过这种药水后，它就能完全防火。\par

% structure: p0086 source=h2 target=subsection reason=relative-heading-rank; base=h2

\subsection{第三十四章 封缄信件的幻觉；详述这些戏法的目的}

在我们简要解释了这些（术士）所秘密行使的技艺的力量，并展示了他们获取知识的简易方法之后，我们也不打算对以下必要之点保持沉默——即他们如何解开封印，并将封缄的信件连同其印章本身恢复原状。他们将等量的沥青、树脂、硫磺以及石沥青熔化，并将此膏剂塑成形状，保存备用。然而，当需要解开小书板时，（术士们）用油涂抹舌头，然后用舌头涂抹封印，再以文火加热药膏，将其置于封印上；他们将其留在那里直至完全硬化，然后以此状态作为印章使用。但他们也说，蜡本身与冷杉胶具有类似的效力，两份乳香与一份干沥青的混合物也有同样的效力。而硫磺本身也能较好地达到目的，以及用水和树胶调和的石膏粉。这最后一种混合物对于封印熔化的铅也极其有效。而由托斯卡纳蜡、树脂渣滓、沥青、石沥青、乳香和粉末状晶石制成的混合物，全部等量煮沸，比我提到的其他药膏更优越；而用树胶制成的混合物也不逊色。因此，他们就这样试图解开封印，竭力辨认写在里面的信件。\par

然而，我犹豫是否要在此书中叙述这些伎俩，因为意识到一种危险：恐怕某个狡诈之人会（根据我的叙述）试图（实践这些戏法）。但是，对许多年轻人的关怀——他们可以免受这些行为的侵害——说服我去教导和公开（上述陈述）以便保障安全。因为即便有一人可能利用这些来学习邪恶，另一个人却可能因得到（这些实践）的教导而免受其害。而那些生命的败坏者——术士们——在行其法术时会感到羞耻。通过学习这些我们先前已阐明的内容，他们或许会被约束而不再愚妄。但为了使这封印不被打破，让我用猪油和混有蜡的毛发来封印它。\par

% structure: p0089 source=h2 target=subsection reason=relative-heading-rank; base=h2

\subsection{第三十五章 用釜占卜；火魔的幻象；魔法咒语的实例}

但是我也不能对这些（巫师）的诡计保持沉默，这诡计在于通过大锅进行占卜。因为他们建造一个封闭的房间，用\OriginalTerm{铜蓝}{cyanus}涂抹天花板以备使用，他们引入一些铜蓝的容器，并将它们朝上拉伸。然而，装满水的大锅放在地上的中央；铜蓝的反射落在大锅上，呈现出天空的模样。但地板也有一个隐蔽的开口，大锅放在上面，大锅事先装有水晶底，而本身是石头制成的。然而，在下方，观众没有注意到，有一个隔间，同谋者聚集在那里，装扮成巫师想要展示的神祇和魔怪的形象。受骗者看到这些，对巫师的诡计感到惊讶，随后相信他可能说的一切。但（巫师）制造一个燃烧的魔怪，通过在墙上描绘他想要的任何形象，然后秘密地涂抹一种按照这种方式混合的药物，即拉科尼亚和扎金索斯的沥青——接着，仿佛受到预言狂乱的影响，他将灯移向墙壁。然而，药物以相当大的光辉燃烧。而为了使一个火热的赫卡忒看起来在空中飞驰，他以下列方式策划。将某个同谋者藏在他想要的地方，将受骗者带到一边，他说服他们（相信他自己），声称他将展示一个火焰般的魔怪在空中骑行。现在他告诫他们立即注视，直到他们看到空中的火焰，然后（那时）蒙住脸，俯伏在地，直到他自己呼唤他们；并在给出这些指示后，他在一个无月的夜晚，以诗句这样说：——\par

\Quote{地下、地上和天上的邦波，来吧！\newline{}街巷之神，光辉者，夜间徘徊者；\newline{}光芒之敌，却是黑暗之友与伴侣；\newline{}在犬嚎中喜悦，在猩红血泊中，\newline{}穿越无生命尘埃的坟墓，涉足尸骸，\newline{}渴求鲜血；以恐惧使人痉挛。\newline{}戈耳戈，莫耳摩，露娜，以及多形者，\newline{}来吧，垂顾我们的祭礼！}\par

% structure: p0092 source=h2 target=subsection reason=relative-heading-rank; base=h2

\subsection{第三十六章 处理显现的方法}

当他说这些话时，火被看见在空中；但（观众）被这奇怪的显现吓住，（而且）蒙住眼睛，哑然扑倒在地。但这诡计的成功被以下装置增强。我所说的被藏匿的同谋者，当他听到咒语停止时，拿着一个包裹着麻絮的鸢或鹰，点燃它并释放它。然而，鸟被火焰惊吓，升到高处，并做出相应的更快飞行，这些受骗的人看到后，藏匿自己，仿佛看到了神圣之物。然而，那有翼的生物被火焰旋转，被带到任何可能的地方，现在烧毁房屋，现在烧毁庭院。这就是巫师们的占卜。\par

% structure: p0094 source=h2 target=subsection reason=relative-heading-rank; base=h2

\subsection{第三十七章 月亮的幻象显现}

他们以这种方式让月亮和星星出现在天花板上。在天花板的中央部分，固定一面镜子，在地板的中央部分放置一个装满水的盘子（与镜子等距），并在中央位置同样放置一根蜡烛，从比盘子更高的位置发出微弱的光——这样，通过反射，（巫师）使月亮通过镜子显现。但经常，他们也在天花板上方远处悬挂一个鼓，但这鼓被一件衣服覆盖，由同谋者藏匿，以便（天体）在适当时间之前不会显现。之后，将一根蜡烛（放在鼓内），当巫师向同谋者发出信号时，他移去足够多的覆盖物，以便模仿代表当时月亮形状的形象。然而，他用\OriginalTerm{辰砂}{cinnabar}和树胶涂抹鼓的发光部分；并事先准备好一个玻璃瓶，削去瓶颈和瓶底周围，将蜡烛放入其中，并在其周围放置一些必要的装置以使形象发光，这些装置由某个同谋者藏在高处；收到信号后，他从上方扔下这些装置，使月亮看起来从天而降。\par

同样的结果在一个树林地方通过一个罐子实现。因为正是通过一个罐子，屋内的把戏得以进行。因为设立一个祭坛，随后（在上面放置）罐子，罐中有一盏点燃的灯；然而，当有更多灯时，这种景象不会显示。然后，魔法师召唤月亮，命令所有灯都熄灭，只留一盏微弱地燃烧；然后，从罐子流出的光反射到天花板上，给在场者提供月亮的形象；罐口被遮盖一段时间，似乎是必要的，以便在天花板上展示满月的形象。\par

% structure: p0097 source=h2 target=subsection reason=relative-heading-rank; base=h2

\subsection{第三十八章 星星的幻象显现}

但鱼鳞——例如，海马的鳞片——使星星显现；将鳞片浸泡在水和树胶的混合物中，并间隔地固定在天花板上。\par

% structure: p0099 source=h2 target=subsection reason=relative-heading-rank; base=h2

\subsection{第三十九章 地震的模仿}

他们以这样的方式引起地震的感觉，使得一切似乎都在动；鼬鼠的粪便与磁石一同在炭火上燃烧（有此效果）。\par

% structure: p0101 source=h2 target=subsection reason=relative-heading-rank; base=h2

\subsection{第四十章 用肝脏的把戏}

他们以这种方式展示一个似乎带有铭文的肝脏。他用左手写下他想要的文字，附加在问题上，文字用胆汁和强醋描绘。然后拿起肝脏，用左手握住，他稍作延迟，然后它吸走印记，于是被认为上面有文字。\par

% structure: p0103 source=h2 target=subsection reason=relative-heading-rank; base=h2

\subsection{第四十一章 使头骨说话}

但将一颗头颅放在地上，他们便让它以这种方式说话。这颗头颅本身是用牛的网膜制成的；当用伊特鲁里亚蜡和准备好的胶塑成所需的形状后，将这一薄膜放置在周围，它便呈现出头颅的样子，当装置运作时，看起来所有人都能听到它说话；正如我们在（随从的）青年案例中所解释的，他们获取鹤或某种类似长颈动物的气管，秘密地连接到头颅上，同谋者说出他想要说的话。当他想让（头颅）变得隐形时，他装作烧香，在周围放置许多煤炭；当蜡受热融化时，头颅便被认为隐去了。\par

% structure: p0105 source=h2 target=subsection reason=relative-heading-rank; base=h2

\subsection{第四十二章 上述行为的欺诈；它们与异端的关联}

这些是魔术师的作为，还有无数其他此类（伎俩），它们用巧妙的言辞和貌似可信的举动来欺骗轻信者。异端首领们惊叹于这些（术士）的技巧，便模仿他们，一方面在隐秘和黑暗中传授他们的教义，另一方面将这些（主张）冒充为自己的。为此，我们渴望警告众人，便更加费心，以免遗漏魔术师所用的任何手段，供那些可能受骗的人参考。不过，我们并非不合理地详细叙述了一些术士的秘密（奥秘），这些虽然对于手头的主题并非十分必要，但对于防范魔术师那邪恶且混乱的技艺，或许被认为是有用的。因此，既然我们尽可能描述了所有（思辨者）的意见，并特别关注阐明异端首领们作为新论提出的那些（观点）；这些（观点）就虔诚而言是虚妄和伪造的，甚至在他们自己中间或许也不被认为值得认真考虑。（进行了这样的探究之后），似乎适宜用一篇简要的论述来唤起（读者）对先前所述内容的记忆。\par

% structure: p0107 source=h2 target=subsection reason=relative-heading-rank; base=h2

\subsection{第四十三章 神学与宇宙论综述；波斯人的体系；巴比伦人的体系；埃及人的天主观；基于数论的神学；他们的宇宙生成体系}

在所有遍及大地、作为哲学家和神学家进行研究的人当中，关于天主——祂的本质或本性——存在着意见分歧。因为有些人宣称祂是火，有些是灵，有些是水，而另一些人说祂是土。每个元素都有某些缺陷，一个被另一个击败。然而，对于世上的智者，发生了显而易见的事，即拥有智力的人可以看出：（我指的是）他们目睹了受造物的惊人工程，对现有事物的本质感到困惑，认为这些事物过于宏大，无法从另一者获得生成，同时又（断言）宇宙本身不是天主。但就神学而言，他们却宣称一个唯一的原因来解释视觉所及的事物，每个人都认为自己所判断的原因最为合理；因此，当他们凝视天主所造之物，以及那些与祂至高的威严相比最为微不足道之物时，由于不能将思想延伸到天主本身的实际伟大，他们便神化了这些（外部世界的作品）。\par

然而，波斯人自以为更深地进入了真理的边界，便宣称天主是光，是包含在空气中的光。巴比伦人则断言天主是暗的，这意见似乎也是前者的结果；因为白昼跟随黑夜，黑夜跟随白昼。埃及人，他们自认为比所有民族都更古老，难道不谈论天主的能力吗？（他们通过）计算这些（黄道）部分的间隔来估算（这种能力；并且，仿佛出于一种最神圣的灵感，他们断言天主是一个不可分的\OriginalTerm{单一}{monad}，既是自生的，又由祂产生了万物。因为，他们说，这个未被生的\OriginalTerm{单一}{monad}产生后续的数字；例如，单一加于自身产生二\OriginalTerm{二元}{duad}；同样，当加于（二、三等）时，产生三\OriginalTerm{三元}{triad}和四\OriginalTerm{四元}{tetrad}，直至十\OriginalTerm{十元}{decade}，这是数字的开始和终结。因此，第一个和第十个单一被产生，因为十元是等值的，并被算作一个单一，并且（因为）这个（十元）乘十次将变成百，而又成为一个单一，而百乘十次将产生千，而这将是一个单一。同样，千乘十次构成万的总和；同样它将是一个单一。但通过不可分量的比较，单一的同类数字包括三、五、七、九。\par

此外，根据六日轨道的安排（即二数的位置与偶数的划分），另一种数与\OriginalTerm{元}{monad}具有更自然的关联。但亲缘数是4和8。这些数从元的数中取得德性的概念，进展到四元素——当然，我是指灵、火、水、土。天主用这些创造了世界，把它造成\OriginalTerm{两性体}{ermaphrodite}，并为上半球分配了两个元素，即灵和火；这被称为元的半球，是仁慈的、上升的、阳性的。因为由微小粒子组成，元升到大气中最稀薄、最纯净的部分；而另外两个元素，土和水，更为粗重，祂将这些分配给\OriginalTerm{二数}{duad}；这被称为下降的半球，既是阴性的也是恶的。同样，上层的元素本身，当彼此比较时，彼此包含男性和女性，以便整个受造界繁衍增多。火是阳性的，灵是阴性的。水是阳性的，土是阴性的。因此，从一开始火与灵结合，水与土结合。因为灵的力量是火，土的力量是水；……元素本身，当通过减去\OriginalTerm{九数}{ennead}来计算和分解时，适当地终止，有些终止于阳性数，有些终止于阴性数。同样，减去九数的原因在于，整个（圆）的三百六十部分由九数组成，因此世界的四个区域被九十完全部分所限定。光被归于元，黑暗归于二数，生命归于光，死亡归于二数。正义归于生命，不义归于死亡。因此，在阳性数中生成的一切都是仁慈的，而在阴性数中生成的则是恶的。例如，他们这样计算：元——我们以此开始——变成361，通过减去九数终止于元。同样，这样计算：二数变成605；减去九数，终止于二数，每个数都回到自己特有的（功能）。\par

% structure: p0111 source=h2 target=subsection reason=relative-heading-rank; base=h2

\subsection{第四十四章：埃及关于本性的理论；他们的护符。}

因此，对于元，因其仁慈，他们断言存在上升的、仁慈的、阳性的、并被仔细观察的名字，终止于奇数；而终止于偶数的名字被认为是下降的、阴性的和恶的。因为他们肯定本性由对立面构成，即恶与善，如右与左、光与暗、夜与昼、生与死。而且他们还宣称，他们计算了\Quote{神}这个词（发现减去九数后它变为一个\OriginalTerm{五数}{pentad}）。这个名字是偶数，当他们将它写在某种材料上时，就把它系在身上，并进行治愈。同样，某种草药终止于这个数，类似地系在身体周围，通过相似的数计算起作用。甚至医生也通过相似的计算医治病人。如果计算相反，则不易治愈。注意这些数的人按照这个原则计算同族的数；有些仅按元音计算，另一些则按整个数字计算。这便是埃及人的智慧，他们夸耀说，通过这种智慧，他们认为自己认识了神圣的本性。\par

% structure: p0113 source=h2 target=subsection reason=relative-heading-rank; base=h2

\subsection{第四十五章：前述讨论的用处。}

看来，这些思辨也已由我们充分地解释了。但既然我认为没有遗漏任何在这地上的、低劣的智慧中所发现的意见，我觉察到我们在这些主题上所付出的关心并非无用。因为我们观察到，我们的论述不仅对驳斥异端有用，而且对那些持这些意见的人也有帮助。现在这些人，当他们遇到我们表现出的极大关怀时，甚至会惊叹于我们的热诚，不会轻视我们的努力，当他们看到他们自己愚昧地相信的意见时，也不会谴责基督徒为愚人。此外，那些渴求学习、献身于真理的人，将因我们的论述而得到帮助，当他们得知异端的基本原则时，变得更加聪明，不仅能轻易驳斥那些试图欺骗他们的人，而且当他们得知智者们公开承认的意见，并熟悉它们之后，他们既不会像无知者那样被它们混淆，也不会成为某些仿佛来自权威的个体的受骗者；不仅如此，他们还会警惕那些让自己成为这些幻想的受害者的人。\par

% structure: p0115 source=h2 target=subsection reason=relative-heading-rank; base=h2

\subsection{第四十六章：星辰神智学者；阿拉托斯被异端首领模仿；他的星辰配置体系。}

以上见解既已充分说明，现在让我们转而思考所讨论的主题，为的是通过证明我们关于异端所确定的论点，并迫使异端的拥护者将他们各自的教义归还给各自的思想家，从而表明异端首领缺乏体系；同时，通过宣告那些被这些异端教义说服之人的愚昧，我们就能劝导他们返回真理的宁静港湾。然而，为了使随后的陈述对读者更为清晰，也有必要宣布阿剌图关于天上星辰之排列所提出的见解。之所以必要，是因为有些人将这些学说与圣经所揭示的学说相比拟，将圣经转化为寓言，并试图诱骗那些留心他们教义之人的心智，以花言巧语引导他们接受他们希望的任何意见，展示一种奇异的奇观，仿佛他们的主张固定在星辰之中。然而，这些人专注地凝视这极其非凡的奇观，作为琐事爱好者，他们像一种称为猫头鹰的鸟一样被迷住，这个例子值得提及，因为随后的陈述需要。这种动物在大小或形状上与鹰并无太大差异，被捕获的方式如下：这种鸟的猎人看到一群鸟落在某处时，在远处挥舞双手，假装跳舞，逐渐靠近它们。但这些鸟被这奇异的景象惊呆了，对周围发生的一切视而不见。而其他已经为此目的装备好来到乡间的猎人，从后面接近这些鸟，在它们凝视着跳舞者时轻易地抓住了它们。\par

因此，我愿无人因那些解释天象之人的类似奇迹而惊奇，像猫头鹰一样被俘获。因为这种投机者所行的诡计可被视为跳舞和愚蠢，而非真理。阿剌图因此这样表达：\par

正如它们众多，此处彼处旋转\newline{}日复一日行过天上，无尽，永恒，（即每一颗星）\newline{}但这却丝毫不倾斜；而是如此精确\newline{}永远与轴心固定，在每一部分平衡\newline{}持守大地居中，并引导天穹自身。\par

% structure: p0119 source=h2 target=subsection reason=relative-heading-rank; base=h2

\subsection{第四十七章 异端从阿剌图借用的意见}

阿剌图说，天上有旋转的星，即回转的星，因为它们从东到西、从西到东永远运行，呈圆形。他还说，有一条巨大的、奇异的怪物，即巨蛇，像一条河流一样绕着\Quote{熊座}本身旋转；这正是魔鬼在约伯传中对天主所说的话，当时它用了这些词：\BibleQuote{我在地下周游，走遍了全世界}，\BibleRef{约 1:7}即我曾旋转，从而得以观察诸世界。因为他们认为，巨龙即蛇位于北极附近，从最高的极点俯视一切事物，观看所有受造物的工作，以免任何正在形成的事物逃过它的注意。因为虽然天上的星辰都落下，但这颗星的极轴从不落下，而是高悬于地平线上，观察并注视一切事物，他说，任何受造物的工作都无法逃过它的注意。\par

正当其中\newline{}日落与日出互相融合。\par

（此处阿剌图说）这个星座的头就放在那里。因为龙的头位于两个半球的西面和东面之间，他说，以便在这同一四分圆内，无论是西面还是东面的事物都不会逃过它的注意，而野兽能同时知晓一切。在龙头附近，有一个人的形象，由星辰显明，阿剌图称之为疲乏的像，像受劳苦的人，被称为\OriginalTerm{昂戈纳西斯}{Engonasis}。阿剌图于是断言他不知道这辛劳是什么，也不知道这个在天上旋转的奇物是什么。然而，异端们希望借助这种星的描述来建立自己的教义，并超乎寻常地热切致力于这些天文学体系，他们声称昂戈纳西斯是亚当，按照天主的诫命（如梅瑟所记载），把守着龙的头，而龙则把守着他的脚后跟。因为阿剌图如此表达：\par

\Quote{拥有凶猛之龙的右脚踪迹。}\par

% structure: p0124 source=h2 target=subsection reason=relative-heading-rank; base=h2

\subsection{第四十八章 里拉琴的发明；星辰形象与位置的寓意化；腓尼基人的起源；阿剌图将逻各斯与大犬座等同；大犬座对生育与一般生命的影响}

（阿剌图）说，（星座）里拉琴和冠冕安放在他两边——我现在是指昂戈纳西斯——但他弯腰屈膝，伸出双手，好像在承认罪恶。而里拉琴是一件乐器，由\OriginalTerm{逻各斯}{Logos}在还是婴孩时就制作而成，逻各斯就是希腊人称为墨丘利的那位。关于里拉琴的制作，阿剌图说：\par

然后，此外，在摇篮旁，\newline{}赫尔墨斯刺穿了它，说：称它为里拉琴。\par

它由七根弦构成，以这七根弦象征整个世界和谐而美妙的构造与组织。因为世界在六天内被创造，（造物主）在第七天安息。因此，若亚当承认（自己的罪过）并按照天主的诫命守护那野兽的头，他将效仿里拉琴——即服从天主的圣言，也就是遵守法律——他就将得到那靠近他的冠冕。然而，若他忽略本分，就将与下方躺卧的野兽一同被抛下，他说，他将与那野兽同分。而跪者似乎在两边伸出双手，一边触摸里拉琴，另一边触摸冠冕——这便是他的宣认；因此，通过这个（星象）构形本身就能辨认他。但冠冕却仍被谋害，并被另一条较小的龙——即被跪者用脚守护的那条的后裔——强行拖走。还有一个人稳固地站着，双手紧握，将冠冕后的巨蛇向后方空间拖拽；他不让那野兽有触及冠冕的机会，尽管它极力企图如此。阿拉托斯称他为持蛇者，因为他抑制了巨蛇冲向冠冕的冲动。但他说，圣言就是以人形阻碍野兽接近冠冕的那一位，他怜悯那同时被龙及其后裔谋害的人。\par

这些（星座）——他说——是\Quote{大熊}和\Quote{小熊}，是两个由七颗星组成的七组，两创造的图像。因为，他断言，第一个创造是按照劳作中的亚当，这就是那被看见\Quote{跪着}的人（跪者）。第二个创造却是按照基督，藉着他我们得以重生；这就是持蛇者，他同野兽搏斗，并阻碍它接近为那男子预留的冠冕。但\Quote{大熊}——他说——是赫利刻，象征一个强大的世界，希腊人向着它航行，即他们正在为之受训。并且，被生命的波浪吹送，他们跟随着前行，（心中存有）某种这样的旋转世界或训导或智慧，它将那些追随这世界的人领回。因为\Quote{赫利刻}一词似乎表示一种朝向同一点的盘旋与循环。还有另一个\Quote{小熊}（库诺苏里斯），仿佛是第二个创造——那按天主而形成的——的某种图像。因为——他说——走窄路的人很少。但他们断言库诺苏里斯是窄的，阿拉托斯说西顿人向它航行。但阿拉托斯提及西顿人（实指）腓尼基人，因为腓尼基人的卓越智慧。而希腊人断言，他们是从（红海）海岸迁移到他们现今所居之地的腓尼基人，因为这是希罗多德的看法。现在库诺苏拉——他说——就是这个（小）熊，第二个创造；有限制之尺度，那窄路，而非赫利刻。因为它不领他们回转，而是以直线引导，那些跟随他的人（构成）犬星的（尾部）。因为犬星就是圣言，一方面守护并保全那受狼谋害的羊群；另一方面又如狗，从受造中猎杀野兽并消灭它们；同时又产生万物，并就是他们以\Quote{犬}之名所表达的，即生产者。因此，据说阿拉托斯提到了犬星的升起，如此表述：\Quote{然而，当犬星升起，庄稼不再歉收。}这就是他说的话：在犬星升起之前植入土中的植物，虽然常常未扎根，却仍覆盖着浓密的叶子，给观看者以它们将会丰产的迹象，并且它们看来充满生机，（而实际上）自身从根部缺乏活力。但当犬星升起时，活物与死物就由犬星分开；因为凡未扎根的植物确实腐烂了。因此——他说——这犬星作为某种神圣的圣言，被设为审判活人死人的判官。正如犬星（的影响）在世上的植物产物中可见，同样，在属天的植物——人类——中也可见（圣言的）权能。由于类似的原因，库诺苏拉——第二个创造——被安置在苍穹作为由圣言而生的创造的图像。那龙则居中躺卧于两创造之间，阻止任何事物从大创造过渡到小创造；并守护着那些固着于（大）创造中的事物，如跪者，同时（又）观察在小创造中每一事物如何及以何种方式构成。并且（那龙）本身被持蛇者看守着头，他说。这个图像——他断言——被固定在天上，对能辨识它的人是一种智慧。然而，若这隐晦难明，他说，创造便通过另一个图像教导人哲学思考，关于此，阿拉托斯如此表述：——\par

\Quote{我们既非刻甫斯·伊阿西达斯的悲惨后代。}\par

% structure: p0130 source=h2 target=subsection reason=relative-heading-rank; base=h2

\subsection{第四十九章 受造物的象征；以及神魂；以及不同种类的动物。}

但阿拉托斯说，在这（星座）附近的是仙王座、仙后座、仙女座和英仙座，对于那些能辨认它们的人来说，它们是受造物的伟大轮廓。因为他断言，仙王座是亚当，仙后座是厄娃，仙女座是这两者的灵魂，英仙座是\OriginalTerm{圣言}{Logos}，朱庇特的有翼后代，而鲸鱼座是图谋不轨的怪物。他修复的不是这些中的任何一个，而唯独是杀死那怪兽的仙女座；圣言——即英仙座——同样从那里将曾被交付（并）锁在怪兽身上的仙女座娶为己有，他说，实现了她的解放。然而，英仙座是有翼的轴，它穿过地球中心的两极，并转动世界。那在世界中的精神，也被（象征为）天鹅座，一只鸟——靠近\Quote{熊座}的乐声动物——是圣神的类型，因为当它接近生命的终结本身时，惟有它生来适合歌唱，怀着美好的希望离开邪恶的受造物，向上主献上赞美诗。但螃蟹、公牛、狮子、公羊、山羊、小山羊，以及许多其他用于命名穹苍星辰的野兽，他说，都是形象和原型，受造物藉此获得（不同的）种类，充满这类动物，且易于变化。\par

% structure: p0132 source=h2 target=subsection reason=relative-heading-rank; base=h2

\subsection{第五十章 占星术之愚妄}

异端者利用这些说法，企图欺骗那些过分专心于占星术的人，从而努力建立一种宗教体系，与这些（思辨家）的思想大相径庭。因此，亲爱的诸位，让我们避免钦佩琐碎事物的习惯，那鸟（称为猫头鹰）就是被这种习惯所捕获的。因为这类和其他类似的思辨，犹如舞蹈，而非真理。因为星辰并不提供这些信息；而是人们自愿地为了给某些星辰命名，就这样称呼它们，以便于区分。因为散布在穹苍的星辰与熊、狮子、小山羊、水瓶座、仙王座、仙女座，或那些在阴间被赋予名称的幽灵有何相似之处呢？——我们必须记住，这些人及其称号本身，是在人类起源之后很久才出现的——（我请问，两者之间有何共同点），以致异端者惊叹于这一奇迹，竟试图通过这样的言论来加强自己的观点？\par

% structure: p0134 source=h2 target=subsection reason=relative-heading-rank; base=h2

\subsection{第五十一章 七元者；算术家的体系；被异端利用；西蒙和瓦伦提诺之例；宇宙之本性可从脑生理学推衍}

但既然几乎每一个异端（兴起者）通过算术技艺发现了\OriginalTerm{七元}{hebdomads}的尺度和某些\OriginalTerm{永世}{Aeons}的投射，各自以不同的方式曲解这技艺，而变化的差异仅限于名称；并且（因为）毕达哥拉斯成了这些人的导师，将这类数字从埃及引入希腊人中，看来似乎不宜忽略这一点，而是要在我们给予简明阐释之后，着手论证我们计划探究的那些事物。\par

算术家和几何学家兴起，毕达哥拉斯似乎是最初为他们提供原理的人。从那些可以通过乘法不断前进\Emph{ad infinitum}的数字，以及从图形，他们推导出第一原理，因为唯有理性才能辨认它们；几何学的原理，如同一个人可能察觉到的，是一个不可分的点。然而，通过技艺，从点出发，发现了从该点生成无限图形的方法。因为点被拉长成线，如此延伸后，以点为端点。而线流出成宽度，产生面，面的界限是线；但面流出成宽度，成为体。当固体这样从全然最小的点获得存在时，巨大物体的本性即被构成；这就是西蒙如此表达的：\Quote{小将成为大，如同一粒点，而大则无限。}现在这与几何学的点理论相符合。\par

但在通过组合包含哲学的算术技艺中，数字成为第一原理，它是一种不可界定和不可理解（的实体），本身包含所有可以通过聚合不断前进\Emph{ad infinitum}的数字。但第一个\OriginalTerm{单子}{monad}成为数字的原理，依据实体，它是一个男性的单子，像父亲一样生育所有其余的数字。其次，\OriginalTerm{二元}{duad}是女性的数字，算术家自己也称之为偶数。第三，\OriginalTerm{三元}{triad}是男性的数字；算术家通常也称之为奇数。除所有这些之外，\OriginalTerm{四元}{tetrad}是女性的数字；这同一个，因为是女性的，也被称为偶数。因此，所有数字按类而言是四个——然而数字，就其种类而言，是无限的——根据他们的体系，由此形成完美的数字——我指\OriginalTerm{十元}{decade}。因为一、二、三、四成为十——正如先前已证明的——如果对于每个数字，依据实体，保持适当的名称。这就是毕达哥拉斯的神圣四元组，本身包含无限本性的根源，即所有其他数字；因为十一、十二等数字的生成原理都来自十。对于这个十元——完美的数字——有四个部分被称为：数字、单子、幂、立方。它们的结合和混合发生在于生成增加的过程中，按照自然完成有生产力的数字。因为当平方乘以自身时，成为四次方；但当平方乘以立方时，成为平方与立方的乘积；但当立方乘以立方时，成为立方与立方的乘积。因此，所有数字是七个；这样，所产生之物的生成可以来自七元——即数字、单子、幂、立方、四次方、平方与立方的乘积、立方与立方的乘积。\par

关于这个\OriginalTerm{七元组}{hebdomad}，\Person{西蒙}和\Person{华伦提努}更改了名称，仔细讲述了奇妙的故事，并从中匆忙采用了一个体系。西蒙使用的名称如下：心、智、名、\OriginalTerm{声音}{Voice}、推理、反思，以及那站着的、站着、将站者。华伦提努（列举的）是：心、真理、\OriginalTerm{圣言}{Word}、生命、人、教会，以及与这些同列的\OriginalTerm{圣父}{Father}，依据与算术哲学倡导者相同之原理。而（异端首领）赞叹这仿佛不为人知的（哲学），并追随它，编造了他们自己设计的异端教义。\par

确实，有些人同样试图从医学（技艺）中形成七元组，对大脑的解剖感到惊奇，声称宇宙的本质、生育的能力以及神性可以从大脑的构造中确定。因为大脑作为全身的主宰部分，平静不动地休息，内含精气。这样的描述并非不可信，但与这些（异端）试图从中推导的结论相去甚远。大脑被解剖时，内部有一个可称为拱形腔的结构。其两侧有薄膜，他们称为小翼。这些被精气轻轻推动，又将精气推向小脑；精气通过像芦苇一样的血管，奔向松果体。附近是小脑入口，接纳精气流，并将其分布到所谓的脊髓中。从脊髓，整个框架参与属灵能量，因为所有动脉像分枝一样从此血管附着，其末端终结于生殖血管；所有（动物）种子从脑经过腰部分泌（到精囊中）。大脑的形状像蛇头，关于此，那些假称为博学之士的人进行了冗长的讨论，正如我们将证明的。另外六条联结韧带从大脑长出，围绕头部，终止于（头部）自身，固定身体；但第七条（韧带）从小脑到身体其余下部，如我们所述。\par

关于此有广泛讨论，由此会发现\Person{西蒙}和\Person{华伦提努}都从这源头取得其观点的出发点，而且，尽管他们可能不承认，他们首先是说谎者，然后是异端者。既然我们似乎已经充分解释了这些教义，并且所有这尘世哲学之著名观点已包含在四卷书中；继续考虑这些人的门徒，毋宁说，那些窃取他们教义的人，似乎是合适的。\par

\SourceRecord{Translated by J.H. MacMahon.；Ante-Nicene Fathers；Vol. 5.；Edited by Alexander Roberts, James Donaldson, and A. Cleveland Coxe.；Buffalo, NY: Christian Literature Publishing Co.,；1886.；<http://www.newadvent.org/fathers/050104.htm>.}

% structure: p0001 source=h1 target=section reason=page-title

\section{\WorkTitle{驳斥一切异端}（第五卷）}

% structure: p0002 source=h2 target=subsection reason=relative-heading-rank; base=h2

\subsection{目录。}

以下是\WorkTitle{驳斥一切异端}第五卷的目录：——\par

自称诺斯替派的纳阿森派的主张是什么，以及他们提出的那些观点是希腊哲学家先前所倡导的，也是那些传授神秘（礼仪）的人所传下来的，（从两者中）纳阿森派趁机建构了他们的异端。\par

以及佩拉特派的教义是什么，他们的体系并非由他们从圣经中构建，而是来自占星术。\par

塞特派的学说是什么，他们从希腊智者那里窃取理论，并用来自\Person{穆赛欧斯}、\Person{利诺斯}和\Person{俄耳甫斯}的意见碎片拼凑自己的体系。\par

犹斯定的教义是什么，他的体系并非由他根据圣经构建，而是来自历史学家\Person{希罗多德}提供的奇闻细节。\par

% structure: p0008 source=h2 target=subsection reason=relative-heading-rank; base=h2

\subsection{第一章 概述；异端的特征；纳阿森派名称的起源；纳阿森派的体系。}

我想，在前四卷中，我已经非常详尽地解释了所有希腊和蛮族中的思辨者关于天主性体和宇宙创造所提出的意见，并且我也没有忽略对他们的魔法体系的考察。因此，我为了读者付出了非同寻常的努力，热切地推动多人走向学习的渴望，并在对真理的认识上坚定不移。因此，剩下的就是要赶着去驳斥异端；但我们提出我们已经作出的陈述，正是为了提供这种（驳斥）。因为异端首领从哲学家那里获取起点，（并且）像鞋匠一样按照他们自己的特殊解释拼凑古人的错误，将它们作为新奇事物呈现给那些能够受骗的人，正如我们在以下各卷中将证明的那样。在余下的（工作中），这个机会邀请我们去处理我们拟议的主题，并从那些胆敢赞美蛇的（人）开始，蛇是通过他自身（才智）的能量所设想的某些表达方式而来的错误的始作俑者。那么，这一体系的司祭和捍卫者，首先是那些被称为纳阿森派的人，他们是从希伯来语得名的，因为蛇在希伯来语中称为\OriginalTerm{纳亚斯}{naas}。后来，他们却自称诺斯替派，声称唯有他们探究了知识的深处。现在，从这些（思辨者）的体系中，许多人切取部分，构建了一个异端，虽然有几个分支，但本质上是一个，并且他们精确地解释了相同的（教义）；尽管以不同意见的面貌传达，正如随后的讨论将随着进展而证明的那样。\par

这些（纳阿森派）根据他们提出的体系，尊崇（作为万有之原）一个人和一个人子。而这个人是一个雌雄同体，在他们中被称作亚当；并且为他创作了众多各样的赞美诗。然而，这些赞美诗——简言之——在他们当中大致以这样的形式表达：\Quote{从你（而来）父亲，通过你（而来）母亲，两个不朽的名字，永世者之祖，啊，天上的居民，你，显赫的人。}但他们将他如\Person{革律翁}般分为三部分。因为，他们说，这个人的一部分是理性的，另一部分是魂性的，另一部分是地上的。并且他们假定对他的认识是能够认识天主的本原，他们这样表达：\Quote{完美的本原是对人的认识，而对天主的认识则是绝对完美。}然而，所有这些品质——理性的、魂性的和地上的——（纳阿森派）说，都已退隐并同时降入一个人——耶稣，他由\Person{玛利亚}所生。并且这三个人（纳阿森派）说，习惯于同时一起说话（通过耶稣），各自从自己的实体向那些属于自己的人说话。因为，根据这些人的说法，一切存在物有三种——天使的、魂性的、地上的；并且有三个教会——天使的、魂性的、地上的；这些的名称是被选的、被召唤的、被掳的。\par

% structure: p0011 source=h2 target=subsection reason=relative-heading-rank; base=h2

\subsection{第二章 纳阿森派通过\Person{玛利亚纳}将他们的体系归给主的兄弟\Person{雅各伯}；实则可追溯到古代神秘；他们的心理学见于\Quote{\WorkTitle{多默福音}}；亚述人的灵魂理论；纳阿森派与亚述人的体系比较；纳阿森派从弗里吉亚和埃及神秘中得到的支持；伊西斯神秘；这些神秘被纳阿森派以寓意解释。}

这些是无数论述的要点，（纳阿森派）断言主的兄弟\Person{雅各伯}将这些传给了\Person{玛利亚纳}。那么，为使这些不敬的（异端者）不再诽谤\Person{玛利亚纳}或\Person{雅各伯}，或救主本身，让我们来考察神秘礼仪（他们从中得出了他们的虚构）——（如果合适的话，对蛮族和希腊（神秘）的考察）——让我们看看这些（异端者）如何收集所有外邦人的秘密和不可言说的神秘，对基督撒谎，并使那些不熟悉这些外邦人的秘仪的人受骗。因为，他们教义的根基是人亚当，并且他们说关于他写着：\Quote{谁能述说他的世代？}\BibleRef{依 53:8} 要知道，他们部分地从外邦人那里取得那人的难以发现和多变的世代，如何虚构地将其应用于基督。\par

\Quote{大地，}希腊人说，\Quote{大地产生了人，它首先带来了一份美好的礼物，希望不是成为无感觉植物或无情野兽的母亲，而是成为温顺且蒙恩造物的母亲。} \Quote{但难以确定，}（纳阿塞人）说，\Quote{要查明\Person{阿拉耳科墨诺斯}作为第一个人类，是否在\Place{玻俄提亚}的\Place{刻菲索斯湖}畔升起；或\Place{伊达}的\Person{枯瑞忒斯}，一个神性种族；或\Place{弗里吉亚}的\Person{科律班忒斯}，太阳最先看到他们像树木生长般冒出来；或\Place{阿耳卡狄亚}产生了\Person{珀拉斯戈斯}，比月亮更古老；或\Place{厄琉西斯}（产生了）\Person{狄奥洛斯}，\Place{拉里亚}居民；或\Place{楞诺斯}生了\Person{卡比洛斯}，秘仪的美童；或\Place{帕勒涅}（产生了）\Place{佛勒格赖}的\Person{阿尔库俄纽斯}，巨人之最。但\Place{利比亚人}断言\Person{伊阿尔巴斯}，最早出生，从干旱平原出现，开始食用\Person{尤皮特}的甜橡实。但埃及的\Place{尼罗河}，}他说，\Quote{至今仍使泥土肥沃，（因此）产生动物，交出活的身体，从湿气中获得肉体。}亚述人却说，食鱼的\Person{俄安内斯}（是第一个人，并且）在他们中间产生。迦勒底人却说，这个\Person{亚当}是大地单独产生的人。他躺卧无生气、不动弹、静止如雕像；他是那位在上面者的肖像，他被颂扬为亚当那个人，由许多力量所生，对于这些力量各有详细的讨论。\par

因此，为使上面的大人最终被制服，正如他们所说，\Quote{从祂而来}，地上和天上得名的全家便形成了，他也被赐予了一个灵魂，好藉着灵魂受苦；并且使被奴役的肖像受伟大、最荣耀、完美的人的惩罚，因为他们正是这样称呼他。于是，他们又问灵魂是什么，从何处来，其本性如何，以至于它来到人身并驱使人时，竟能奴役并惩罚完美之人的肖像。然而，他们不（在此事上）从圣经探求，而是也从神秘仪式询问这个问题。他们断言灵魂极难发现，难以理解；因为它不总保持同一形状或同一形式，也不处于同一被动状态，使人能用一个标记表达它，或从本质上理解它。\par

但他们在题名为\WorkTitle{依照埃及人的福音}中记载了灵魂的这些多样变化。于是，他们如同外邦人中的其余人一样，怀疑（灵魂）究竟是从某种前存之物而来，还是从自生者而来，还是从广布混沌而来。他们首先逃向亚述人的神秘，察觉人的三分：因为亚述人首先提出灵魂有三部分，但（本质上）是一体。因为，他们说，每种本性都渴望灵魂，且各以不同方式。因为灵魂是一切被造之物的原因；纳阿塞人说道，一切被滋养并生长之物都需要灵魂。因为他说，在灵魂不在场之处，不可能获得任何滋养或生长。他甚至断言，石头也是有生命的，因为它们具有增长的能力；但增长若无滋养绝不会发生，因为增长之物是通过增加而成长，而增加是被养育之物的滋养。因此，每一种本性——天界的和（纳阿塞人说）天界、地界、阴间之物——都渴望灵魂。亚述人称这样的存在为\Person{阿多尼斯}或\Person{恩底弥翁}；当它被称为阿多尼斯时，他说，\Person{维纳斯}爱慕并渴望以这名字称呼的灵魂。但据他们说，维纳斯是生产。但每当\Person{珀耳塞福涅}或\Person{科瑞}恋上阿多尼斯时，他说，便产生了一个与维纳斯（即生育）分离的必死灵魂。但若月亮对恩底弥翁生出贪欲，并爱上他的形态，他说，更高存在的本性同样需要灵魂。但若诸神之母阉割了\Person{阿提斯}，她自己将此（人）作为爱慕对象，他说，至高永恒者的福乐本性独能将灵魂的男性力量召回自身。\par

因为（纳阿森派）说，有一个\OriginalTerm{阴阳人}{hermaphrodite}。按照他们这种说法，男女交合被证明是极其邪恶污秽的行为。因为（纳阿森派）说，\Person{阿提斯}被阉割，即他从地下世界的尘世部分过渡到了永远实体\Quote{上面}，那里，他说，既没有女性也没有男性，而是一个新受造物，一个新的人，即阴阳人。至于他们使用\Quote{上面}这一说法之处，我将在适当地方（对此事）说明。但他们断言，按他们的说法，他们证明\Person{瑞亚}并非完全孤立，而是——容我这样说——普世受造物；他们宣称这就是圣言所肯定的。\Quote{因为祂的无形事自创世以来，借着祂所造之物就被看见，祂永远的大能和神性，都叫他们无可推诿。因此，他们虽然认识天主，却不当作天主荣耀祂，也不感谢祂；他们的心反而变为虚妄。他们自称为聪明，反成了愚拙，将不能朽坏之天主的荣耀，变为朽坏的人、飞鸟、四足走兽和爬虫的样式。所以天主任凭他们放纵可耻的情欲；因为连他们的女人也把自然的用途变为反自然的。}然而，按照他们，自然的用途是什么，我们以后将说明。\Quote{同样，男人也放弃了女人的自然用途，彼此欲火中烧，男与男行可耻的事}——纳阿森派认为\Quote{可耻的事}指着那第一和蒙福的实体，无形，是一切形状的原因——\Quote{并在自己身上受他们错谬所当得的报应。}\BibleRef{罗 1:20}因为这些保罗所说的话，他们认为是包含了他们全部的秘密和隐藏的福乐奥秘。因为洗濯的应许，按照他们，无非是将受洗者引入生命之水（按照他们），并以不可言喻的油膏敷抹（将其引入）不朽福乐。\par

但他们断言，不仅可以从亚述人的奥秘中，也可以从弗吕家人关于那有福本性的奥秘中，找到有利于他们学说的证据——这有福本性是隐蔽而又同时显露的，关乎已存在、正在存在、并将要存在之物——（这有福本性）纳阿森派说，就是那当在人里面寻求的天国。\BibleRef{路 17:21}关于这（本性），他们传下一个明确的段落，出现在称为\BibleBook{多默}{福音}的经卷中，他们这样说：\Quote{寻找我的人，会在七岁孩童里找到我；因为我隐藏在那里，将在第十四岁显现出来。}然而，这不是基督的（教训），而是\Person{希波克拉底}的，他用了这些话：\Quote{七岁孩童是父亲的一半。}所以这些（异端者）将宇宙的本源本性置于生殖种子中，并确认了\Person{希波克拉底}这一格言，即七岁孩童是父亲的一半，便说，按照多默，他在十四岁时显现。在他们看来，这就是不可言说且神秘的圣言。他们继而断言，埃及人——已被证实比弗吕家人更古老，比全人类更古老，并且公认为首先向其余人宣告所有诸神的礼仪和秘仪，以及（万物的）种类和能量——拥有神圣威严、对未受启蒙者不可言传的\Person{伊西斯}奥秘。然而，这些无非就是那位穿七件衣袍、披黑色长袍者所寻求并夺走的东西，即\Person{奥西里斯}的阴部。他们说\Person{奥西里斯}是水。而那七重衣袍的本性，被以太质料的七重幔子环绕并装饰——因为他们如此称呼行星，将其寓意化并称之为以太衣袍——仿佛就是可变受造物，被那不可言说、不可描绘、不可思议、无形者所转化。纳阿森派说，这就是圣经所宣告的：\BibleQuote{义人七次跌倒，仍必兴起。}\BibleRef{箴 24:16}因为他说，这些跌倒就是星辰的变动，由那推动万物者所推动。\par

他们关于那作为万有之因的种子的本质所主张的，是它并非这些中的任何一个，而是它产生并形成所有被造之物，他们如此表达：\Quote{我成为我所愿的，且我是我所是的；为此我说，那推动万有者本身是不动的。因为那存在者始终形成万有，而被造之物中无一被造成。} 他说唯有这（一位）是善的，而救主所说的正是关于这一位：\Quote{你为什么称我是善的？只有一位是善的，就是我在天的父，祂使祂的太阳升起在义人和不义之人身上，降雨给圣人和罪人。} \BibleRef{玛 5:45} 至于祂降雨给谁是圣人，又降雨给谁是罪人，这一点我们将在后面连同其余一并说明。这就是宇宙伟大、秘密且未知的奥秘，在埃及人中被隐藏又被揭示。因为\OriginalTerm{纳阿塞人}{Naassene}说，\OriginalTerm{奥西里斯}{Osiris}在\OriginalTerm{伊西斯}{Isis}面前的庙宇中；他的\OriginalTerm{生殖器}{pudendum}暴露着，向下注视，并以它自己一切被造之物的果实为冠冕。他断言，这样的（形象）不仅立于最神圣的偶像庙宇中，而且为了众人的认识，它仿佛一盏灯，不是放在斗底下，而是放在灯台上，在屋顶上、在所有旁道、所有街道和住宅附近宣告其信息，作为住所前放置的某种指定的界限和终结，且所有人都称之为善的（实体）。因为他们称这为善的生产者，却不知道自己在说什么。希腊人从埃及人那里得到这神秘的（表达），并保留至今。因为\OriginalTerm{纳阿塞人}{Naassene}说，我们看到\OriginalTerm{墨丘利}{Mercury}的雕像是如此形状在他们当中受尊敬。\par

然而，他们以特别的尊崇敬拜\OriginalTerm{库勒尼乌斯}{Cyllenius}，称他为\OriginalTerm{逻各斯之神}{Logios}。因为\OriginalTerm{墨丘利}{Mercury}是\OriginalTerm{逻各斯}{Logos}，祂是同时被造、正在产生和将要存在之物的解释者和制造者，在他们当中受尊敬，被塑造成类似男人的\OriginalTerm{生殖器}{pudendum}的形状，具有从下部向上部的冲动力量。而且\OriginalTerm{纳阿塞人}{Naassene}说，这个神——即这种描述的墨丘利——是召唤死者的人，是离世灵魂的引导者，也是灵魂的起源者；这一点也没有逃过诗人们的注意，他们如此表达：——\par

\Quote{库勒尼的赫耳墨斯也召唤\newline{}凡间追求者的灵魂。}\par

他说，不是\OriginalTerm{珀涅罗珀}{Penelope}的追求者们，啊，你们这些可怜的人！而是（灵魂）被唤醒并被带回到对自己的记忆之中，\par

\Quote{从如此伟大的尊荣，从如此长久的福乐。}\par

也就是说，从上面有福的人，或原初之人或\OriginalTerm{亚当}{Adam}，如他们所认为的，\Emph{灵魂}被传送到这里进入泥土的受造物中，好让他们服侍这位受造物的\OriginalTerm{德穆革}{Demiurge}，即\OriginalTerm{伊尔达巴奥特}{Ialdabaoth}，一位火神，第四个数；因为他们如此称这位有形世界的德穆革和父亲：——\par

\Quote{他手中持着一根可爱的\newline{}金杖，能迷惑人眼，\newline{}他愿意的人，以及再次唤醒了那些沉睡的人。}\par

他说，这就是那唯一掌握生死权柄者。关于这一点，他说，经上记载：\Quote{你必用铁杖管辖他们。} 然而，他说，诗人渴望装饰逻各斯有福本性那不可理解的能力，便给了他一根金杖而非铁杖。如他所说，他迷惑死人的眼睛，并再次举起那些沉睡的人，即被从睡眠中唤醒并成为追求者之后的。关于这些人，他说，圣经说：\Quote{你这睡眠者，醒来，从死者中起来，基督必要光照你。} \BibleRef{弗 5:14}\par

他说，这就是那基督，在所有被生者中，是从不可描绘的\OriginalTerm{逻各斯}{Logos}而出的被描绘的人子。他说，这就是\OriginalTerm{厄琉息斯秘仪}{Eleusinian rites}那伟大而不可言说的奥秘：\OriginalTerm{许厄}{Hye}、\OriginalTerm{库厄}{Cye}。他断言，万有都已服从于祂，而这就是那所说的：\Quote{它们的声音传遍普世，} \BibleRef{罗 10:18} 正如与这些表达一致：\Quote{\OriginalTerm{墨丘利}{Mercury}挥动他的金杖，\Emph{引导灵魂}，但它们吱吱地跟随。} 我的意思是，脱体的灵魂持续跟随，正如诗人通过其意象所描绘的，使用这些话语：——\par

\Quote{犹如在魔幻洞穴的深处\newline{}蝙蝠嗡嗡飞翔，当一只掉下}\par

\Quote{从岩石的突脊，每一只都彼此紧贴。}\par

他说，他用\Quote{岩石}这一表达指\OriginalTerm{亚当}{Adam}。他断言，这就是亚当：成为了角石头的头角石。因为在头中，实体是形成性的脑，整个家族由此被塑造。\BibleRef{弗 3:15} 他说：\Quote{谁，我立他为\OriginalTerm{锡安}{Zion}根基上的岩石？} 他说，他是用寓意方式谈论人的受造。岩石被插入（在）牙齿之间，如荷马所说：\Quote{牙齿的围墙，} 即一道墙和堡垒，其中存在着内在这人，他从上面的原初之人亚当堕落。并且他被\Quote{无手分割而劈开，} 并被带入遗忘的形象，\Emph{带着属土和属泥的性质}。他断言，吱吱叫的灵魂跟随他，即\OriginalTerm{逻各斯}{Logos}：——\par

于是这些吱吱叫的聚在一起：然后灵魂。\par

也就是说，他引导他们；\par

\Quote{温和的赫耳墨斯带领他们穿过宽广延伸的道路。}\par

也就是说，他说，进入那与一切邪恶分离的永恒之地。因为，他说，他们从何处来：——\par

\Quote{他们越过海洋的溪流和\OriginalTerm{琉卡斯悬崖}{Leuca's cliff}，\newline{}经过太阳之门和梦境之地。}\par

他说，这就是海洋，\Quote{诸神之诞生与人类之诞生}，被水的旋涡不断旋转，一时向上，一时向下。但他说，当海洋向下流时，便有人类的诞生；而当它向上流向路卡斯的壁垒、堡垒和悬崖时，便有诸神的诞生。他断言，这就是所记载的：\BibleQuote{我说：你们是神，都是至高者的子女；}\Quote{如果你们急忙飞出埃及，穿过红海进入旷野，}也就是说，从地上的交往进入天上的耶路撒冷，它是永生的母亲；\BibleRef{迦 4:26}\Quote{此外，如果你们再回到埃及，}即回到地上的交往，\Quote{你们就必像人一样死亡。}因为他说，凡在下的受生都是必死的，而在上受生的则是不朽的，因为它仅由水和圣神所生，是属神的，不是属血肉的。但下面所生的（是属血肉的），也就是说，他所写的：\BibleQuote{那由圣神生的就是神，那由神生的就是神。}\BibleRef{若 3:6}按照他们的说法，这就是属神的诞生。他说，这就是伟大的约但河\BibleRef{苏 3:7}，它在这（下）面流淌，阻止以色列子民离开埃及——我是指地上的交往，因为埃及对他们而言就是身体——耶稣使它后退，并使之向上流。\par

% structure: p0036 source=h2 target=subsection reason=relative-heading-rank; base=h2

\subsection{第三章　对纳阿森派异端的进一步阐释；宣称追随荷马；承认三一原则；该三一原则的技术名称；用希腊诗人的权威支持这些名称；将救主的奇迹寓言化；萨莫色雷斯人的奥秘；主为何拣选十二门徒；色雷斯人和弗里吉亚人使用的\OriginalTerm{科里巴斯}{Corybas}之名的解释；纳阿森派声称在圣经中找到他们的体系；他们对雅各伯异象的解释；他们对\Quote{完美的人}的观念；弗里吉亚人称\Quote{完美的人}为\Quote{爸爸}；纳阿森派和弗里吉亚人论复活；圣保禄的狂喜；基督所暗示的宗教奥秘；对撒种比喻的解释；应许之地的寓言；弗里吉亚人体系与圣经陈述的比较；对上下厄琉息斯奥秘意义的阐释；纳阿森派看来在这里发现的降生成人。}

采纳了这些及类似（观点），这些最神奇的诺斯替主义者，发明了一种新奇的文法艺术，将荷马尊为先知——（据他们说）他按照奥秘的方式宣告了这些真理；他们嘲笑那些未受圣经训导的人，使其陷入这些观念。然而，他们提出以下断言：说万物源自一体的，是错误的；但说万物源自三体的，则拥有真理，并将提供宇宙（现象）的解答。因为（纳阿森派）说，有一位有福者的有福本性，即在上者（即亚当）；有一位必死的本性，即在下的；还有一位无王的受生，是在上受生的，他说，那里有被寻觅的玛利亚姆、大智的约托尔、凝视的塞福拉，以及其受生不在埃及的梅瑟，因为他的子女是在米德杨生的；他说，连这一点也没有逃过诗人们的注意。\par

我们的分份有三，各得\newline{}其应得的尊荣。\par

因为他说，这些伟大者必须被宣告，并且必须这样被所有人到处宣告，\BibleQuote{使他们听是听见，却不明白；看是看见，却不晓得。}\BibleRef{玛 13:13}因为他说，如果这些伟大者不被宣告，世界就不可能存在。这些是（这些异端者）的三个膨胀表达：\OriginalTerm{考劳考}{Caulacau}、\OriginalTerm{扫劳素}{Saulasu}、\OriginalTerm{扫劳扫}{Saulasau}。扫劳素即至上的亚当；扫劳扫即在下必死的；齐萨尔即向上流的约但河。他说，这就是存在于万物中的阴阳人。但那些不认识他的人，称他为三体革律翁——革律翁，仿佛出自\Quote{从地上流溢}之意——但希腊人一致（称）为\Quote{月亮的角}，因为他将万物混合并融合于万物中。\BibleQuote{因为万物是藉着他造的，凡被造的，没有一样不是藉着他造的，在他内有生命。}\BibleRef{若 1:3}他说，这就是生命，即完美之人的不可言喻的受生，这在先前的世代是未被认识的。而\Quote{没有一样不是藉着他造的}这句话，指的是有形的世界，因为它是没有藉着他的工具，由（上述四元组的）第三和第四（原则）所造。因为他说，这就是那杯\BibleQuote{康迪，君王饮取之时，借此占卜。}\BibleRef{创 44:2}他说，这已发现在本雅明的美丽种子里隐藏着。同样，希腊人，他说，也以以下方式谈论它：——\par

水到怒吼的口中；奴隶，拿酒来；\newline{}醉倒我，使我陷入麻木。\newline{}我的酒杯告诉我\newline{}我必将成为何种。\par

他说，这单凭这一点就足以使人们理解；（我是指）阿那克里翁的杯，它宣告了（虽无声）一个不可言喻的奥秘。因为他说，阿那克里翁的杯是哑的；然而阿那克里翁断言它用无声的语言对他说话，告知他必将成为何种——即属神的，而非属血肉的——如果他沉默地倾听那隐藏的奥秘。这就是耶稣在那美好婚筵中将水变成的酒。他说，这就是耶稣在加里肋亚的加纳所行的伟大而真实的最初神迹，并（由此）彰显了天国。他说，这就是那如同宝藏、如同藏于三斗面中的酵母、安息在我们内的天国。\par

这是，他说，萨莫色雷斯人的伟大而不可言喻的奥秘，只有我们这些受过启蒙的人才能知道。因为萨莫色雷斯人在他们举行的奥秘仪式中明确传授了（那同一个）\Person{亚当}作为\OriginalTerm{原始人}{primal man}。在萨莫色雷斯人的神庙里通常立着两尊裸体男子的像，双手朝天高高举起，\Emph{生殖器勃起，如同}\Place{库勒涅山}上的\Person{墨丘利}雕像。上述像就是原始人和那重生之\OriginalTerm{属神的人}{spiritual one}的形像，与那人完全相同。这，他说，就是救主所说的话：\Quote{如果你不喝我的血，不吃我的肉，你就不能进入天国；但是即使，他说，你喝了我所喝的杯，我往哪里去，你也不能进去。}因为他说他知道每个门徒是什么样的本性，并且每一个都必须达到他自己特有的本性。因为他说他从十二支派中拣选了十二门徒，并通过他们向各支派说话。为此，他说，十二门徒的宣讲并不是所有人都听见，即使听见了，也不能接受。因为不合乎本性的事，在他们看来是违背本性的。\par

这，他说，居住在哈伊莫斯附近的色雷斯人，以及与其相似的弗里吉亚人，称之为\OriginalTerm{科里巴斯}{Corybas}，因为（虽然）他从上面的首脑和不可描绘的脑髓开始下降，并渗透现存事物的所有原则，（然而）我们却察觉不到他如何以及以何种方式下降。这，他说，就是所说的：\Quote{我们确实听见了他的声音，但没有看见他的形状。}\BibleRef{若 5:37}因为那个被分离和描绘者的声音可被听见；但他的形状——从上面从不可描绘者下降——是什么样，无人知晓。它住在一个属地的模子里，却无人认识它。这，他说，就是\WorkTitle{圣咏}中所说的\Quote{居住在大水中的神}，\Quote{从众水中发言呐喊}。\Quote{众水}，他说，就是必死之人的多样化世代，从这（世代）他向不可描绘的人呼喊并大声说：\Quote{从狮子口中保全我的独生子。}作为对他的回应，他说，已宣告：\Quote{以色列，你是我的孩子：不要害怕；即使你经过河流，它们也不会淹没你；即使你经过火，它也不会烧灼你。}\BibleRef{依 49:15}他说，河流意指生殖的湿润物质，火则指冲动原则和生养欲望。\Quote{你是我的；不要害怕。}他又说：\Quote{如果母亲忘记她的孩子，不顾惜他们，也不给他们食物，我也要忘记你。}他说，\Person{亚当}对他的族人说：\Quote{即使妇人忘记这些，我也不会忘记你。我已将你刻在我的手掌上。}然而，关于他的上升，即他的重生，为使他成为属神的，而非属血肉的，他说，\WorkTitle{圣经}（如此）说：\Quote{你们这些统治者，打开城门；永久的门户，你们要被举起！那光荣的君王要进来。}这是奇妙中的奇妙。\Quote{因为，他说，这光荣的君王是谁？是虫，不是人；是人的耻辱，是人民的弃物；他自己就是光荣的君王，是战场上的勇士。}\par

他说，\Quote{战争}是指身体内的战争，因为其结构是由敌对元素构成的；正如经上所记，他说：\Quote{当记得身体内的争斗。} 他说，\Person{雅各伯}在往\Place{美索不达米亚}的旅途中看见了这道入口和这门，也就是说，当他从一个孩子正在长成青年和成人时；也就是说，当他行经美索不达米亚时，那（入口和门）就被启示给了他。但\Place{美索不达米亚}，他说，是从\OriginalTerm{完美的人}{Perfect Man}中间涌出的大洋之流；他对那天上的门感到惊奇，喊道：\BibleQuote{这地方何等可畏！这若不是天主的殿，必是天堂的门。} 因此，他说，\Person{耶稣}说：\BibleQuote{我就是那真门。} \BibleRef{若 10:9} 说这些话的那一位，他说，就是那按上方不可描绘的那位所塑成的完美的人。因此，完美的人不能得救，除非通过这门进入而重生。但\Person{弗里吉亚人}，他说，也称这同一位为\OriginalTerm{帕帕}{Papa}，因为他在显现之前平静了一切混乱不和地运动着的事物。因为他说，\OriginalTerm{帕帕}{Papa}这名同时属于所有受造物——天上的、地上的和地下的——它们喊说：\Quote{止息，止息世界的纷争}，并\BibleQuote{给远方的人带来和平}，即给物质和属地的存在；且\BibleQuote{给近处的人带来和平} \BibleRef{弗 2:17}，即给那些属神的和具有理性的完美的人。但\Person{弗里吉亚人}也称这同一位为\Quote{尸体}——如同埋葬在身体里，像在陵墓和坟墓中。他说，这就是所宣告的：\BibleQuote{你们是粉饰的坟墓，里面却装满了死人的骨骸}，\BibleRef{玛 23:27} 因为你们内没有那活人。他又喊道：\BibleQuote{死人要从坟墓里出来}，\BibleRef{玛 27:52} 即从地上的身体里出来，重生为属神的，而非属血肉的。因为他说，这就是通过天国之门发生的复活；通过这门的，那些不进入的就仍留在死中。然而他说，这些同样的\Person{弗里吉亚人}又断言，这同一位（人）因着变化而成为一位神。因为他说，当他从死者中复活，通过这样的一扇门进入天堂时，他就成为神。他说，\Person{保禄}宗徒知道这门，在奥秘中部分地开启它，说：\BibleQuote{他被一位天使提到第二层和第三层天，直到乐园里；他看见了异象，听到了不可言说的话语，是人不可能说的。} \BibleRef{格后 12:2}\par

他说，这些就是被众人称为隐秘的奥秘的，\BibleQuote{我们也讲这些，不是用人的智慧所教的言语，而是用圣神所教的言语，用属神的事解释属神的事。然而属血气的人不接受圣神的事，因为在他看是愚妄的。} \BibleRef{格前 2:13} 他说，这些就是圣神不可言喻的奥秘，只有我们知道。关于这些，他说，救主曾宣告：\BibleQuote{除非我天上的父吸引某人，谁也不能到我这里来。} \BibleRef{若 6:44} 因为他说，要接受并接纳这伟大而不可言喻的奥秘是非常困难的。他又说，救主曾宣告：\BibleQuote{不是凡向我说\InnerQuote{主啊，主啊}的人，就能进入天国，而是那承行我在天之父旨意的人。} \BibleRef{玛 7:21} 而实行这（旨意）的人，不仅听，还必须进入天国。他又说，救主曾宣告：\BibleQuote{税吏和娼妓要在你们之前进入天国。} \BibleRef{玛 21:31} 因为他说，\Quote{税吏}是收取一切事物税捐的人；但我们，他说，就是税吏，\Quote{时代的终结已临到我们身上。} 因为他说，\Quote{终结}是从不可描绘的那位散布在世界上的种子，整个宇宙体系借着这些种子得以完成；因为也是借着它们开始存在的。他说，这就是所宣告的：\BibleQuote{撒种者出去撒种。有的落在路旁，被践踏了；有的落在石头地上，发芽了，}他说，\BibleQuote{因为没有泥土的深度，就枯萎而死；有的，}他说，\BibleQuote{落在好地里，结出果实，有一百倍的、六十倍的、三十倍的。有耳可听的，}他说，\BibleQuote{就听吧！} 他说，这意思如下：除了完美的\OriginalTerm{诺斯替信徒}{Gnostics}，没有一个人能成为这些奥秘的聆听者。他说，这就是\Person{梅瑟}所说的好地：\BibleQuote{我要领你们进入那美好宽阔的流奶流蜜的地方。} 他说，这就是奶和蜜，完美的人尝了它们就成为无王的，并分享\OriginalTerm{圆满}{Pleroma}。他说，这就是圆满，借着它，一切已产生的存在之物都从\OriginalTerm{非受生者}{ingenerable one}被产生并完成。\par

这位同样的一位也被弗里吉亚人称为\Quote{不结果实的}。因为当他属于血肉时便不结果实，并引起肉体的欲望。他说，这就是所说的：\Quote{凡不结好果子的树，就砍下来丢在火里。}因为他说，这些果子仅是理性的活人，他们从第三道门进入。他们确实说：\Quote{你吞吃死者，却使生者存活；（但）你若吃活的，你将做什么？}然而，他们声称这活者\Quote{是理性的官能和心智，以及人——那不可描绘者投在下面受造物面前的珍珠。}他说，这就是（耶稣）所断言的：\Quote{不要把圣物给狗，也不要把你们的珍珠丢在猪前。}现在，他们声称猪和狗的工作就是女人与男人的交合。他说，弗里吉亚人将此同一位置称为\Quote{牧羊人}（\OriginalTerm{牧羊人}{Aipolis}），不是因为，他说，他习惯喂养母山羊和公山羊，如世俗之人所用之名，而是因为，他说，他是\Quote{\OriginalTerm{牧羊人}{Aipolis}}——即，永远巡游——他藉其旋转运动使整个宇宙体系旋转并携带。因为\Quote{\OriginalTerm{转动}{Polein}}一词意指转动和改变事物；因此，他说，他们都将天的两极称为\Quote{\OriginalTerm{极}{Poloi}}。而诗人说：—\par

那海生的无罪智者来到这里，\newline{}不死的埃及普罗透斯？\par

他未被毁灭，他说，而是仿佛旋转并围绕自身周转。此外，我们居住的城市，因为我们在其内转向并走动，也被称为\Quote{\OriginalTerm{城市}{Poleis}}。照此方式，他说，弗里吉亚人称这位为\Quote{\OriginalTerm{牧羊人}{Aipolis}}，因为他处处不息地转动万物，并将它们变为各自特有的（功能）。而弗里吉亚人，他说，也称呼他为\Quote{\OriginalTerm{极其多产}{very fruitful}}，\Quote{因为，}他说，\Quote{不怀孕、不生养的女子，她的儿女比有丈夫的儿女更多；}即，重生的事物成为不朽并永远大量存留，即便被产生的事物可能很少；而属血肉的事物，他说，全是可朽坏的，即使很多（此类事物）被产生。因此，他说，\Quote{辣黑耳痛哭她的儿女，不肯受安慰，}（先知）说，\Quote{她为他们悲伤，因为她知道，}他说，\Quote{他们不在了。}但耶肋米亚也为下面的耶路撒冷哀哭，不是腓尼基的城市，而是下面可朽坏的世代。因为耶肋米亚，他说，也认识那完美之人——那重生之人——不是血肉的，而是从水和圣神生的。至少耶肋米亚本人曾说过：\Quote{他是人，谁能认识他？}照此方式，（那塞尼派）说，对完美之人的知识是极其深奥、难以理解的。因为，他说，成全的开端是对人的认识，而对天主的认识则是绝对的成全。\par

然而，他说，弗里吉亚人还声称他是\Quote{被收割的青麦穗}。在弗里吉亚人之后，雅典人在给人们举行厄琉息斯仪式时，也向那些被接纳进入这些奥迹最高等级的人，显明那适合一个被引入最高奥理之人的伟大、奇妙且最完美的秘密：（我指的是）一穗被默默收割的麦子。但这麦穗在雅典人看来也被认为构成了从不可描绘者而来的完美巨大的光照，就如掌秘祭师亲自（声明）一样；并非像阿提斯那样被阉割，而是因毒芹而成为太监，并且藐视一切血肉的生育。（现在）夜间在厄琉息斯，在一个巨大的火把之下，（司礼者）举行伟大而隐秘的奥迹，高声呼喊并叫道：\Quote{庄严的布里摩生了一个神圣的儿子布里摩斯；}即，一个强有力的（母亲）生下了一个强有力的孩子。但神圣的，他说，是那属灵的、属天的、从上而来的生育，而如此出生的那位是强有力的。因为这奥迹被称为\Quote{\OriginalTerm{厄琉息斯}{Eleusin}}和\Quote{\OriginalTerm{安那克托里翁}{Anactorium}}。\Quote{\OriginalTerm{厄琉息斯}{Eleusin}}，他说，是因为我们这些属灵的人是从上面的亚当流下来的；因为\Quote{\OriginalTerm{来临}{eleusesthai}}一词，他说，与\Quote{来}之表达同义。但\Quote{\OriginalTerm{安那克托里翁}{Anactorium}}与\Quote{上升}之表达同义。他说，这就是那些在厄琉息斯人的奥迹中被启明者所肯定的。然而，法律规定，那些已被接纳进入小奥迹的人应当再次被启明进入大奥迹。因为更大的命运获得更大的份。而较低级的奥迹，他说，是属于下面普洛塞耳皮娜的；关于这些奥迹和那通往那里的道路，那路是宽阔的，引向毁灭之人到普洛塞耳皮娜，诗人也说：—\par

但在她下面一条可怕的道路延伸，\newline{}空洞泥泞，却是最佳的引导\newline{}通往备受尊崇的阿佛洛狄忒可爱的树林。\par

这些，他说，就是较低的奥迹，那些属于血肉生育的。现在，那些被引入这些较低（奥迹）的人应当停留，然后（再）被引入伟大的（和）属天的（奥迹）。因为，他说，那些获得（此奥迹）份的人，领受了更大的份。因为，他说，这是天堂的门；这是天主的殿，唯有善神居住其中。而进入这（门），他说，凡不洁的人不得进入，也没有属自然或属血肉的人；它只为属灵的人存留。那些来到此处的人应当脱去衣服，全都成为新郎，藉着童贞之灵而被阉割。因为这就是那位童贞女，她腹中怀孕、成孕并生下一个儿子，不是动物的，不是属物质的，而是永远蒙福的。关于这些，据说救主曾明确宣称：\Quote{引往生命的门是窄的，路是狭的，找到它的人很少；但引往丧亡的门是宽的，路是阔的，走那条路的人很多。}\par

% structure: p0052 source=h2 target=subsection reason=relative-heading-rank; base=h2

\subsection{第四章 对弗里吉亚人体系的进一步利用；举行奥秘仪式的方式；\Quote{大母}的奥秘；这些奥秘与\OriginalTerm{纳塞尼人}{Naasseni}共同崇拜的对象；纳塞尼人寓言化地解释圣经中伊甸园的记载；寓言应用于耶稣的生平。}

然而，弗里吉亚人进一步断言，宇宙之父是\Quote{Amygdalus}，他说，不是一棵树，而是先前存在的\Quote{Amygdalus}；他自身拥有完美的果实，仿佛在深处悸动和移动，撕裂了自己的胸脯，生出了他那如今不可见、无名字、不可言喻的孩子。关于他，我们将要讲述。因为\Quote{Amyxai}这个词的意思，他说，就如同爆裂和撕开，正如在发炎的身体上且有肿块时发生的情况；当医生切开这个时，他们称之为\Quote{Amychai}。他说，以这种方式，弗里吉亚人称他为\Quote{Amygdalus}，由此那不可见者出来并被生出，\BibleQuote{万物是藉着祂造的；凡受造的，没有一样不是藉着祂造的。} 弗里吉亚人说，由此产生的是\OriginalTerm{Syrictas（吹笛者）}{Syrictas}，因为所生的圣神是和谐的。\Quote{因为天主是神，}他说，\Quote{所以真正朝拜的人，既不在这山上，也不在耶路撒冷，而是在心神内朝拜。因为对成全者的朝拜，}他说，\Quote{是属神的，不是属血肉的。}然而，他说，圣神就在那里，那里也有圣父被称呼，圣子从这圣父那里被生出。他说，这就是那多名的、千眼的\OriginalTerm{不可理解者}{Incomprehensible One}，每个本性——虽然各不相同——都渴慕祂。他说，这就是天主的圣言，他说，这是对\OriginalTerm{大能者}{Great Power}启示的圣言。因此，它将被密封、隐藏和遮蔽，躺卧在寓所中，那里有宇宙根基的基础，即：\OriginalTerm{永世}{Aeons}、\OriginalTerm{权能}{Powers}、\OriginalTerm{理智}{Intelligences}、\OriginalTerm{神}{Gods}、\OriginalTerm{天使}{Angels}、\OriginalTerm{委派的神灵}{delegated Spirits}、\OriginalTerm{存在者}{Entities}、\OriginalTerm{非存在者}{Nonentities}、\OriginalTerm{可生成者}{Generables}、\OriginalTerm{不可生成者}{Ingenerables}、\OriginalTerm{不可理解者}{Incomprehensibles}、\OriginalTerm{可理解者}{Comprehensibles}、年、月、日、时，（以及）\OriginalTerm{不可见的点}{Invisible Point}，从这一点，那最小的开始逐渐增大。他说，那无有的、由无组成的，由于是不可分的——（我是指）一个点——将凭借自身的反思能力成为某种不可理解的大小。他说，这就是天国、芥菜子、身体内不可分的点；他说，除了属神的人，没有谁知道这一点。他说，这就是所说的：\BibleQuote{没有言语，没有辞句，也听不到它们的声音。}\par

他们以这种方式轻率地假定，所有人所说所做的一切，都可以与各自的心智观点相协调，声称一切事物都成为属神的。因此他们断言，那些在戏院中表演的人——甚至这些人——也不是毫无预谋地说或做任何事。所以，他说，当人们聚集在戏院时，如果有人穿着引人注目的袍子进入，带着竖琴，弹奏曲调，并配合着伟大奥秘的歌曲歌唱，他便说出以下的话，不知道自己在说什么：\Quote{无论你是萨图尔努斯的种族，还是快乐的朱庇特，或是强大的瑞亚，欢呼，阿提斯，瑞亚的阴郁阉割者。亚述人称你为三次渴望的阿多尼斯，整个埃及称你为奥西里斯，月亮的角；希腊人称你为智慧；萨莫色雷斯人称你为可敬的亚当；哈埃蒙人称你为科里巴斯；弗里吉亚人有时称你为帕帕，有时为尸体，或神，或不结果实者，或艾波洛斯，或绿色收获的谷穗，或极丰产的Amygdalus所生——一个人，一位音乐家。} 他说，这就是千变万化的阿提斯，当他们用赞美诗庆祝时，他们吟唱这些话：\Quote{我将歌颂阿提斯，瑞亚之子，不是用号角的嗡鸣声，或伊达山吹笛人的声音，它们与库里特斯的声音相和；而是我要用阿波罗的竖琴音乐来调和我的（歌），\OriginalTerm{evoe}{evoe}，\OriginalTerm{evan}{evan}，因为你是潘，你是巴克斯，你是璀璨星辰的牧者。}\par

由于这些和类似的原因，这些人经常参加所谓的\Quote{大母}的奥秘，尤其认为通过那里举行的仪式，他们看到了整个奥秘。因为除了他们没有受阉割之外，他们所拥有的并不比那里举行的仪式多：他们只是完成了被阉割者的工作。因为极其严厉和警惕地，他们命令（其信徒）如同被阉割者一样，与女人断绝来往。然而，其余的程序（在这些奥秘中观察到的），正如我们已经较详细地阐述的，（他们遵循）就像（他们是）被阉割者一样。他们并不崇拜任何别的对象，只崇拜\OriginalTerm{纳亚斯}{Naas}，（由此）被称为\OriginalTerm{纳塞尼人}{Naasseni}。但纳亚斯就是蛇，他说，从该词纳亚斯，天下所有被称为\OriginalTerm{圣殿}{Naos}的事物都得名。并且（他声称）唯独归于祂——即纳亚斯——每一座圣所、每一个入教仪式和每一个奥秘都是奉献的；而且，总的来说，在天底下找不到一种宗教仪轨，其中没有圣殿；而在圣殿本身之中就是纳亚斯，圣殿由此得名。这些人断言，蛇是一种潮湿的实体，正如米利都的泰勒斯（也将水视为本原）一样，并且现存的一切，不朽的或有死的，有灵的或无灵的，没有祂根本无法存在。万物都隶属于祂，祂是善的，祂在自身内拥有一切，如同独角兽的角中一样；以致祂按各自的本性和特性赋予一切存在者美丽和繁茂，仿佛穿过一切，正如\BibleQuote{从伊甸流出的河水，分作四道。}\par

然而，他们断言，厄登好比是大脑，被环绕的衣袍捆绑并紧紧束缚，犹如在天上。但他们认为，人只就头部而言是乐园，因此那\Quote{从厄登流出的河}，即从大脑流出，\Quote{分为四支：第一支名叫丕雄，环绕哈威拉全地，那里有黄金，那地的黄金是好的，还有芳香树脂和玛瑙石。}他说，这就是眼睛，它以其在其余身体器官中的尊贵及色彩，为所说的话提供了见证。\Quote{第二支河名叫基红，环绕雇士全地。}他说，这就是听觉，因为基红是弯曲的河，类似某种迷宫。\Quote{第三支河名叫底格里斯，流在亚述之东。}他说，这就是嗅觉，以河流极为湍急的水流作比喻。但它流在亚述之东，因为在每一次呼吸后，从外部大气吸入的气息以更快的运动和更大的力量进入。他说，这就是呼吸的本性。\Quote{第四支河是幼发拉底。}他们断言，这就是口，通过它实现祈祷的外出和营养的进入。口使人欢喜，滋养并塑造属灵的完美之人。他说，这就是\Quote{穹苍以上的水}，关于这水，救主曾宣称：\Quote{你若知道那向你求水的是谁，你必早已求祂，祂也必早给了你活水，涌流的水。}他说，每一种本性都进入这水，选择自己的实体；而这水的特质来到每一本性，就像铁被磁石吸引，金子被海隼的脊骨吸引，谷壳被琥珀吸引一样。\par

但他说，若有人生来瞎眼，从未见过\Quote{那照亮每一样来到世界的人的真光}，就让我们使他重见光明，并看见，如同经过一座栽种各样树木、果实丰盛的乐园，水流穿行于所有树木和果实之间；他会看到，从同一水中，橄榄为自己选取并汲取油，葡萄汲取酒；其余植物也各按其类。然而，他说，那人在世上无声誉，在天上却显赫，被无知者出卖给不认识他的人，被视为桶中滴落的一滴。但我们，他说，是属灵的，从流经巴比伦中心的幼发拉底的生命之水中，当我们经过真门——即蒙福的耶稣时，选择自己的特质。在所有世人中，唯独我们基督徒在第三道门举行奥迹，并用角中的不可言喻的油傅油，正如达味受傅，而非用瓦器，他说，正如撒乌耳，他与肉体欲望的邪魔交谈。\par

% structure: p0058 source=h2 target=subsection reason=relative-heading-rank; base=h2

\subsection{第五章．取自拿先森人一首赞美诗对其体系的解释}

以上所述，虽仅从许多中选取少数，我们认为适宜提出。因为愚妄的企图数不胜数。但既然我们已尽力解释了那未识的诺斯，似乎同样宜于提出以下一点。他们谱写了这首赞美诗，在其中以如下措辞的颂歌赞美他们谬误的所有奥迹：——\par

世界产生的法则是原初心灵，\newline{}接着是首生者倾泻的混沌；\newline{}第三，灵魂接受了辛劳的律法：\newline{}因此，被水状形式环绕，\newline{}被忧虑压倒，它屈服于死亡。\newline{}现在掌握权柄，它凝视光明，\newline{}现在它哭泣，被抛入苦难；\newline{}现在它哀恸，现在它战栗喜悦；\newline{}现在它悲号，现在它听见自己的厄运；\newline{}现在它听见自己的厄运，现在它死去，\newline{}现在它离开我们，永不回返。\newline{}它，不幸迷途，踏遍灾殃的迷宫。\newline{}但耶稣说：父啊，看，\newline{}一股灾殃的冲突横跨大地\newline{}从你的气息（忿怒）中游荡；\newline{}但那人寻求逃避苦涩的混沌，\newline{}却不知如何穿越它。\newline{}为此，父啊，派遣我；\newline{}带着印玺，我将降临；\newline{}我将扫荡所有世代，\newline{}我将解开所有奥秘，\newline{}我将展示诸神之形；\newline{}并将圣洁道路的秘密，\newline{}名为\Quote{诺斯,}传授。\par

% structure: p0061 source=h2 target=subsection reason=relative-heading-rank; base=h2

\subsection{第六章．奥斐特人——异端之大源头}

这些教义，拿先森人力图建立，自称为诺斯替派。但既然谬误多头而多样，正如历史上所记载的九头蛇；当我们用驳斥之杖一击砍下这幻象的头颅，以真理之杖，我们将彻底消灭这怪物。因为其余的异端与此并无多大差别，由同一谬误之灵彼此相连。但既然他们更改蛇的言词和名称，希望蛇有许多头，即使如此，我们将如他们所愿，彻底驳斥他们。\par

% structure: p0063 source=h2 target=subsection reason=relative-heading-rank; base=h2

\subsection{第七章．佩拉特人的体系；他们的三神论；降生成人的解释}

无疑还有一种（许德拉的头，即异端）\OriginalTerm{培拉特派}{Peratae}，他们对基督的亵渎多年来未曾被发现。现在正是揭露这些（异端者）秘密奥秘的适当时机。这些异端者声称世界是一元的，分为三部分。按照他们的三分法，一部分是一个单一的本原原则，就像一个巨大的泉源，可以在理智上分割为无限多的片段。第一个片段，按照他们的说法，是（片段中）优先的，是一个三元组，被称为完美的善、父的宏大。第二个部分，则像是从自身产生的无限潜能之众，第三个部分是形式的。第一个，即善，是非受生的；第二个是自我产生的善；第三个是被造的。因此他们明确宣称有三位神、三个\OriginalTerm{逻各斯}{Logoi}、三个心智、三个人。因为对于世界的每一部分，经过分割之后，他们指派给神、逻各斯、心智、人以及其余的一切；但是从非受生和世界的第一片段，当后来世界达到圆满时，在黑落德时期，有一个被称为基督的人带着三重本性、三重身体、三重能力从天而降，原因我们以后会说明，他在自身内拥有来自世界三个部分的全部具体化和潜能；而这就是（培拉特派者）所说的：\BibleQuote{乐意让一切的圆满居住在他内}，并且整个神性居住在他内，如同这样划分的三元组。因为他声称，从两个较高的世界——即从（三元组中）非受生的部分和从自我产生的部分——所有种类的潜能之种子被传递到我们所在的这个世界。然而，下降的方式是什么，我们以后会说明。\par

（培拉特派者）接着说，基督从上面、从非受生下降，以便通过他的下降，一切三分的事物都能得救。因为有些事物从上面被带下，将藉着他上升；而那些攻击从上面带下之事物的，将被抛弃，并被置于惩罚状态而被弃绝。他说，这就是所说的：\BibleQuote{因为人子来到世界，不是要毁灭世界，而是要藉着他使世界得救。}他说，他称世界为位于上面的两部分，即（三元组的）非受生部分和自我产生的部分。当圣经说\BibleQuote{免得我们和世界一同被定罪}时，是指第三部分，即形式的世界。因为第三部分，他称之为我们所在的这个世界，必须灭亡；而位于上面的两部分，必须脱离败坏。\par

% structure: p0066 source=h2 target=subsection reason=relative-heading-rank; base=h2

\subsection{培拉特派从占星学家那里获取他们的体系；这一点由黄道带的占星学理论的说明证明；因此产生了培拉特派异端的术语。}

因此，我们首先来了解（培拉特派信徒）如何从占星学家那里获得这一教义，从而恶意对待基督，对跟随他们这种错误的人造成毁灭。因为占星学家声称世界是一元的，将其分为黄道十二宫的十二个固定部分，称固定黄道带的世界为一个不动的世界；而另一个世界，他们断言是运行（星体）的世界，无论在其能力、位置还是数目上，都延伸到月亮。他们规定，一个世界从另一个世界获得某种能力和相互分享，处在下面的从上面获得此种分享。然而，为了使所断言的内容清晰明了，我将逐一采用占星学家本身的那种表达方式；我这样做只是提醒读者我们在解释占星学家整个技艺的作品部分中先前所作的陈述。那么，这些（思辨者）所持的意见如下——\par

（他们的教义是）从星宿的流溢中，完成了下属（部分）的生成。因为，当他们渴望地仰望天空时，加色丁人断言（七曜星）包含一种原因，即所有发生在我们身上的事件的效果原因，并且固定黄道星座的部分也合作（于这种影响）。他们将（黄道圈）分为十二（部分），每个黄道星座分为三十份，每份再分为六十个微小部分；因为他们如此称呼最微小且不可分的部分。在黄道星座中，他们称一些为男性，另一些为女性；一些是双体的，另一些则不是；一些是回归的，而另一些是固定的。男性星座是（或雌性或雄性，它们具有协作本性以产生男性，或者自身产生女性）。因为白羊座是男性黄道星座，而金牛座是女性；其余（星座）依此类推，一些为男性，另一些为女性。我想毕达哥拉斯派受此（考虑）影响，称单子为男性，双元为女性；并且，三元为男性，以及其余偶数和奇数依此类推。然而，有些人将每个黄道星座分为十二部分，采用几乎相同的方法。例如，在白羊座中，他们称十二部分中的第一部分为白羊座且为男性，第二部分为金牛座且为女性，第三部分为双子座且为男性；对剩余的（部分）也采用同样的方法。他们断言有些星座是双体的，即双子座和与之相对的星座，即射手座、室女座和双鱼座，其余则不是双体的。并且（他们说）有些也是回归的，当太阳处于这些星座时，会引起周围（星座）的巨大转折。白羊座就是这样的星座，以及与之相对的星座，如天秤座、摩羯座和巨蟹座。因为白羊座是春天的转折点，摩羯座是冬天的，巨蟹座是夏天的，天秤座是秋天的。\par

然而，关于这一体系的细节，我们已在前一本书中详细说明；任何想获得（这方面）教导的人，都可以从中了解佩拉特异端的创始人，即佩拉特的欧弗拉特和卡里斯托的塞尔贝斯，如何（在将这些观点转化为他们自己的体系时）仅仅更改了名称，而实际上提出了类似的教义。（不仅如此），他们还以过分的热情，亲自致力于（占星家的）技艺。因为占星家们也谈论星辰的界限，他们断言在这些界限中，主导的星辰具有更大的影响力；例如，在一些星辰上它们有害地作用，而在另一些星辰上则有好处。在这些星辰中，他们称一些为恶意的，一些为有益的。据说星辰彼此相望，并相互协调，以至于它们呈现为三角形或正方形的形状。彼此相望的星辰按照三角形的形状排列，间隔三个黄道星座的距离；而（间隔）两个黄道星座的则按照正方形的形状排列。（他们的教义是），正如在人身体中，下属部分与头部相互感应，头部也与下属部分相互感应一样，所有地上的事物都与月上的事物相互感应。但是（占星家更进一步）；因为（根据他们）这些事物之间存在某种差异和不兼容，以至于它们不包含同一结合。星辰的这种结合与分歧，是加色丁人的（教义），已被我们之前提到的那些人据为己有。\par

如今，这些人歪曲真理之名，将永世的叛乱和善向恶势力的叛变宣扬为基督的教义；他们谈论善势力与恶势力的联盟。因此，他们称这些为\OriginalTerm{托帕尔科}{Toparchai}和\OriginalTerm{普罗斯提奥}{Proastioi}，并且（虽然）为自己编造了许多其他并非（来自其他来源）暗示的名号，却仍不熟练地将占星家关于星辰的全部虚构教义体系化。由于他们提出的假说蕴含着巨大的错误，他们将通过我们绝妙的安排被驳斥。因为我将用一些佩拉特派的论文来对照前面提到的加色丁人的占星术，从中通过比较，将有机会看出佩拉特派的教义确实是占星家的教义，而不是基督的教义。\par

% structure: p0071 source=h2 target=subsection reason=relative-heading-rank; base=h2

\subsection{第九章　从佩拉特人自己的书中解释其体系}

看来，有必要展列他们当中一本受推崇的书籍，其中包含以下段落：\Quote{我是从黑夜世代之沉睡中唤醒的声音。从今往后，我开始剥夺那出自混沌的力量。那力量是淤泥最深处之物，它撑起那不朽（且）潮湿的空间广袤的污泥。它是那动荡的全部力量，永远运动，呈现水的颜色，旋转静止之物，抑制颤抖之物，释放前行之物，减轻滞留之物，移除增长之物，是微风轨迹的忠实管家，享受从法律之十二只眼中吐出的东西，并向那与自己一同分配向下不可见之水、被称为塔拉萨的力量显明一个印记。无知常称这力量为克洛诺斯，用锁链看守，因他紧紧束缚了浓密、多雾、昏暗、阴沉的塔尔塔罗斯的褶皱。按照这形象产生了刻甫斯、普罗米修斯、伊阿珀托斯。被托付塔拉萨的力量是雌雄同体的。它将出自十二张嘴的嘶嘶声固定到十二根管子，并倾泻而出。这力量本身是精细的，它移除（海的）控制性的、狂暴的、向上的运动，并封印其路径的轨迹，免得（任何敌对力量）发动战争或引入任何改变。此力量的风暴之女是各种水的忠实保护者。她的名字是霍尔扎尔。无知常称这（力量）为尼普顿，按照其形象产生了格劳科斯、墨利刻忒斯、伊诺、涅布洛厄。那被十二天使的金字塔环绕，并以各种颜色使金字塔之门变暗，并以夜的黑色完成全部者：无知称这为克洛诺斯。他的臣仆有五位——第一乌，第二阿奥艾，第三乌奥，第四乌奥阿布，第五……其他忠实的掌管者（存在）于他的夜与昼的领域中，安息于自身的力量中。无知称这些为流浪之星，腐败的世代依赖于它们。星的升起的掌管者是卡尔法卡塞梅奥希尔，（而）埃卡巴卡拉（是同一者）。无知常称这些为风之首的库里特斯；第三顺位的是阿里埃尔，按照其形象产生了埃俄罗斯、布里阿瑞俄斯。十二时辰夜（力量）的首领是索克朗，无知常称他为奥西里斯；（而）按照这形象诞生了阿德墨托斯、美狄亚、海伦、埃图萨。十二时辰昼力量的首领是欧诺。这是星之升起者普罗托卡马鲁斯和以太（区域）的掌管者，但无知称他为伊西斯。此者的标志是天狼星，按照其形象诞生了阿尔西诺埃之子托勒密、狄杜玛、克利奥帕特拉、奥林匹亚斯。天主的右手力量是无知称之为瑞亚的力量，按照其形象产生了阿提斯、密格冬、（和）奥诺涅。左手力量掌管滋养，无知常称这为刻瑞斯，（而）她的名字是贝娜；按照这形象诞生了刻琉斯、特里普托勒摩斯、密叙尔、普拉克西狄刻。右手力量掌管果实。无知称这为墨娜，按照其形象诞生了布墨伽斯、奥斯塔奈斯、赫尔墨斯·特里斯墨吉斯托斯、库里特斯、佩托西里斯、佐达里翁、贝罗苏斯、阿斯特兰普苏科斯、（和）琐罗亚斯德。左手力量（掌管）火，无知称这为武尔坎努斯，按照其形象诞生了埃里克托尼俄斯、阿喀琉斯、卡帕纽斯、法厄同、墨勒阿革洛斯、提丢斯、恩克拉多斯、拉斐尔、苏列尔、（和）翁法勒。有三个中间力量悬于空中，是产生者。无知常称这些为命运女神；按照这些形象产生了普里阿摩斯家族、拉伊俄斯家族、伊诺、奥托诺厄、阿高厄、阿塔玛斯、普洛克涅、达那伊得斯、珀利阿得斯。有一个雌雄同体的力量，永远处于婴幼期，永不衰老，是美、快乐、成熟、欲望、贪欲的原因；无知常称这为厄洛斯，按照其形象诞生了帕里斯、那喀索斯、伽倪墨得斯、恩底弥翁、提托诺斯、伊卡里俄斯、勒达、阿弥莫涅、忒提斯、赫斯珀里得斯、伊阿宋、勒安德耳、（和）赫洛。}这些是\OriginalTerm{普罗阿斯提奥伊}{Proastioi}直至以太，因为他也以此标题签署那本书。\par

% structure: p0073 source=h2 target=subsection reason=relative-heading-rank; base=h2

\subsection{第十章 培拉特异端名义上与占星术不同，但实为同一体系之寓意化}

对于所有人来说，显而易见，\OriginalTerm{佩拉特派}{Peratae}的异端只是改变了占星术的名称。这些（异端者）的其他著作也包含同样的方法，如果有人愿意通读它们的话。因为，如我所说，他们认为万有生成的原因是无生而超然者，而我们的世界是按流溢的方式产生的，他们称这世界为有形的世界。他们（还主张）天上所有的星辰共同造成了这世界的生成。然而，他们更改了这些星辰的名称，正如人们可以通过比较（两种体系）从\OriginalTerm{普洛阿斯提奥}{Proastioi}中察觉到的。其次，按照与世界由上层流溢而产生的相同方法，他们断言这里的万物也是从星辰的流溢中获得生成、败坏与次序。既然占星家熟悉上升点、中天、下降点以及中天对点；由于这些星辰因宇宙永恒的旋转而在不同时间占据不同的位置，因此不同时期就有不同的下降（向中心）和上升（至中心）。（佩拉特异端者）给占星家的这一布局附会上寓意，将中心描绘为神、元一和统管宇宙生成的主，而下降（被视为）左边的能力，上升是右边的能力。因此，当任何人阅读这些（异端者）的著作时，若在其中发现左右能力的提及，他应当（从迦勒底智者的）中心、下降和上升中追溯，并会清楚地看到这些（佩拉特派）的整个体系就建立在占星学说的基础上。\par

% structure: p0075 source=h2 target=subsection reason=relative-heading-rank; base=h2

\subsection{第十一章   他们为何自称佩拉特派；他们的生成理论借助古代的论证；他们对以色列出埃及的解释；他们从圣经推导出的\Quote{蛇}的体系；这正是占星家学说的真正含义。}

他们自称\OriginalTerm{佩拉特派}{Peratae}，认为没有哪个由生成产生的东西能逃脱为那些源自生成的事物所注定的命运。因为（佩拉特派）说，如果某物完全是生成的，它也会灭亡，正如\Person{西比拉}的意见。但我们，他说，只有我们这些通晓生成的必然性、以及人进入世界的路径，并受到精确教导（这些事）的人，才有能力穿越并超越毁灭。而水，他说，就是毁灭；世界，他说，没有任何东西比水更快地毁灭它。但水，他说，是那在\OriginalTerm{普洛阿斯提奥}{Proastioi}之间流转的，（他们）断言（它就是）\OriginalTerm{克洛诺斯}{Cronus}。因为这种能力，他说，是水色的；而这种能力——即克洛诺斯——他说，没有哪个由生成产生的东西能够逃脱。因为克洛诺斯是每一个生成在屈服于毁灭方面的原因，并且不可能存在没有克洛诺斯干涉的生成。他说，这正是诗人们所肯定的，甚至令诸神恐惧：\par

因为，他说，要知道这大地和广阔的上天，\newline{}以及\OriginalTerm{斯提克斯}{Styx}的洪水之水，那正是\newline{}最大且最受幸福诸神敬畏的誓言。\par

他说，不仅诗人们这样说，而且希腊人中最有智慧的人也已如此主张。\Person{赫拉克利特}甚至就是其中之一，使用了以下的话：\Quote{对于灵魂，水成为死亡。}（佩拉特派）说，这死亡在红海中抓住了埃及人和他们的战车。他说，所有对此无知的人都是埃及人。他们断言，这就是出离埃及，即离开身体。因为他们认为小埃及就是身体，它渡过红海——即败坏的之水，也就是\OriginalTerm{克洛诺斯}{Cronus}——并到达红海之外的地方，即生成；它进入旷野，即达到独立于生成的状态，在那里所有毁灭之神和救恩之神混杂在一起。\par

他说道，众星是毁灭之神，它们将可变之生成的必然性强加于存在之物。这些，他声称，被梅瑟称为旷野的蛇\OriginalTerm{旷野的蛇}{serpents of the wilderness}，它们咬噬并彻底毁灭那些自以为已经渡过红海的人。他说道，对于那些在旷野中被咬的以色列子民，梅瑟展示了真实而完美的蛇\OriginalTerm{完美的蛇}{perfect serpent}；凡信这蛇的人，在旷野中就没有被咬，也就是说，没有被邪恶的权势侵袭。因此，他说，没有人能够拯救和释放那些离开埃及的人，即离开肉身和这个世界的人，除非那唯一完美且充满圆满的蛇。他说道，凡寄望于此蛇者，就不会被旷野的蛇，即生成之众神所毁灭。这事，他声称，记载在梅瑟的一卷书中。这条蛇，他说道，就是伴随梅瑟的力量，即那变为蛇的棍杖。然而，术士的蛇——即毁灭之神——在埃及对抗梅瑟的力量，但梅瑟的棍杖使它们全都臣服并杀死了它们。这宇宙之蛇，他说道，是厄娃的智慧言论\OriginalTerm{厄娃的智慧言论}{wise discourse of Eve}。他说道，这就是厄登的奥秘\OriginalTerm{厄登的奥秘}{mystery of Edem}，这就是厄登的河流；这就是加音身上的记号，免得遇见他的人击杀他。他说道，这就是加音，此世之神不接纳他的祭品。然而，他认可了亚伯尔的血祭；因为此世的统治者喜爱血祭。他说道，这就是那在末期、在黑落德时代、以人形显现的那一位，仿照若瑟的样式而生，被弟兄的手出卖，唯独他拥有彩色长衣。他说道，这就是那仿照厄撒乌样式的那一位，他的衣裳——他本人不在场——蒙受了祝福；他并未接受那眼目昏花者所发出的祝福。然而，他从别处获得了财富，没有从那眼睛昏花者那里得到任何东西；雅各伯看见了他的面容，如同人看见天主的面容。关于这一点，他说道，经上记载：\BibleQuote{尼默洛得在上主面前是个英勇的猎人}。他说道，有许多人紧密模仿他；他们如同那些在旷野中被以色列子民看见的咬人之蛇一样众多，梅瑟树立的那条完美蛇将那些被咬的人拯救出来。他说道，这就是那所宣告的：\BibleQuote{梅瑟怎样在旷野里高举了蛇，人子也应照样被举起来}。按此样式，在旷野中铸造了铜蛇，梅瑟将它举起。他说道，唯独这蛇的形象在天上，永远在光明中显现。\par

他说道，这就是经上所说的伟大的开始\OriginalTerm{伟大的开始}{the great beginning}。关于这，他说道，经上记载：\BibleQuote{在起初已有圣言，圣言与天主同在，圣言就是天主。这圣言在起初就与天主同在。万物是藉着他而造成的；没有他，什么也没有造成。凡在他内形成的，就是生命。} 在他内，他说道，形成了厄娃；厄娃是生命\OriginalTerm{厄娃是生命}{Eve is life}。然而，他说道，这就是厄娃，众生之母——即众神、天使、不朽者、有死者、无理性者、有理性者的共同本性。因为，他说道，\Quote{所有}一词是他对所有存在者说的。他说道，如果任何人的眼睛蒙福，仰望穹苍，就会在那高耸的苍穹顶上看见蛇的美妙形象\OriginalTerm{蛇的美妙形象}{beauteous image of the serpent}，它转动自身，成为一切被生成之物各种运动的起源。他将知道，没有它，天上、地上或地下的任何事物都不能存在。黑夜、月亮、果实、生成、财富、生计，以及一切存在之物，没有一件是在它的统御之外的。他说道，关于这一点，就是那穹苍中目睹的伟大奇观\OriginalTerm{伟大奇观}{great wonder}，是那些有能力观察者所见到的。因为，他说道，在它的头顶上——这对无知者而言是比一切更难以置信的事实——\Quote{日落与日出相互混合}。这就是无知者惯常所声称的：在天上\par

\Quote{天龙旋转，可怖怪兽之奇妙奇迹。}\newline{} 在他两侧安放着北冕座\OriginalTerm{北冕座}{Corona}和天琴座\OriginalTerm{天琴座}{Lyra}；在上方，靠近头顶处，可见可怜之人\Quote{跪者座\OriginalTerm{跪者座}{Engonasis},}\newline{} \Quote{紧握凶猛天龙右脚的末端。}\par

在跪者座背后，有一条不完美的蛇\OriginalTerm{不完美的蛇}{imperfect serpent}，双手被持蛇者\OriginalTerm{持蛇者}{Anguitenens}紧紧缚住，不能触及完美蛇旁边的北冕座。\par

% structure: p0083 source=h2 target=subsection reason=relative-heading-rank; base=h2

\subsection{第十二章 培拉特派学说概述}

这就是培拉特派\OriginalTerm{培拉特派}{Peratae}的多样化智慧，它难以完整陈述，因为它似乎包含占星术，如此复杂。那么，只要可能，我们将简要解释这异端的全部要义。为了通过简要陈述阐明这些人的全部教义，似乎适宜补充以下观察。根据他们，宇宙是圣父\OriginalTerm{圣父}{Father}、圣子\OriginalTerm{圣子}{Son}、物质\OriginalTerm{物质}{Matter}；这三者每一个自身都有无穷的能力。于是，在物质与圣父之间，坐着圣子、圣言\OriginalTerm{圣言}{Word}、蛇\OriginalTerm{蛇}{Serpent}，总是动态地朝向不动的圣父和动态的物质本身。有时他转向圣父，将能力接入自身；有时他拿起这些能力，转向物质。物质虽然无属性且未成形，却从圣子那里接受样式，而这些样式是圣子从圣父塑造的。\par

但聖子以不可言喻、不可描述、且不變的方式——即是梅瑟所說，從放在飲水槽中的木杖流到受孕牲畜身上的顏色那樣——從聖父那裏獲得形狀。同樣地，能力也從聖子那裏流入物質，正如那來自木杖、作用於受孕牲畜的受孕能力。他說，顏色之差異與不同，即從木杖通過水而流到羊群上的那種差異，就是可朽與不朽受生的差異。然而，正如一位按自然作畫的畫家，雖然他沒有從動物身上取走任何東西，卻用畫筆將所有形狀轉移到畫布上；照樣，聖子以隸屬於自己的能力，將聖父的印記從聖父轉移到物質中。所有聖父的印記都在此，並無更多。他說，因為在此世中的任何一人若有能力察覺自己是從上頭轉移而來的聖父印記，並且是降生成人的——就像因木杖受孕而產生的白色之物——那麼他就與天上的聖父全然同體，並返回那裏。然而，若他未遇到這學說，也就不會理解受生的必然性，就像夜裏出生的流產兒將在夜裏滅亡。因此，他說，當救主說：\BibleQuote{你們的天父}時，祂指的是那一位，聖子從祂那裏取得特徵並轉移到此處。然而，當（耶穌）說：\BibleQuote{你們的父從起初就是殺人者}時，祂指的是物質的統治者和造物主，他將從聖子傳下來的印記據為己有，在此處產生了那從起初就是殺人者，因為他的工作導致腐敗和死亡。\par

因此，他說，沒有一個人能得救或返回（天上）而不藉著聖子，而聖子就是那蛇。因為正如祂從上面帶下聖父的印記，照樣，祂又從那裏將這些印記（從沉睡狀態喚醒，使之成為聖父的特性，從非實體之存有變為實體的特性）帶上來，並將它們從那裏轉移到此處。他說，這就是所說的：\BibleQuote{我是門}。祂將這些（印記）轉移給那些閉上眼瞼的人，正如石腦油將火從各方向引向自身；不，更確切地說，正如磁石（吸引）鐵而不吸引其他東西，或者正如海鷹的脊骨（吸引）金子而不吸引其他東西，或者正如琥珀引導碎屑。照樣，他說，那被描繪的、完美的、同體的族類被蛇從世界再次吸引出來；祂不（吸引）任何其他東西，正如它被祂派下來一樣。為此，他們引用大腦的解剖結構，將大腦本身因其不動性比作聖父，將小腦因其運動和形狀（如蛇頭）比作聖子。他們聲稱，這個（小腦）以一種不可言喻、難以測度的方式，透過松果腺吸引從拱形腔室（大腦所在）發出的屬靈和賦予生命的實質。在接受此實質後，小腦以不可言喻的方式將觀念傳遞給物質，正如聖子所作的一樣；或者，換句話說，按肉體所產生之物的種子和種類沿脊髓流下。利用這個例子，（這些異端者）似乎巧妙地引入了他們秘密的奧蹟，這些奧蹟是在靜默中傳授的。我們若宣告這些奧蹟，是不敬的；但根據已作出的許多陳述，不難對它們有所認識。\par

% structure: p0087 source=h2 target=subsection reason=relative-heading-rank; base=h2

\subsection{第十三章　培拉特異端並不廣為人知。}

但由於我認為我已經清楚地說明了培拉特異端，並用許多（論證）揭示了（一個迄今）一直隱藏、完全由虛構構成、且掩蓋其特有毒素的系統，看來，除了已經提出的陳述之外，不宜再作任何進一步的說明；因為（異端者）自己所提出的意見已足夠判決他們自己有罪。\par

% structure: p0089 source=h2 target=subsection reason=relative-heading-rank; base=h2

\subsection{第十四章　塞特派的體系；他們的無限三原則；他們的異端解說；他們對降生成人的解釋。}

现在让我们看看西提安派的说法。他们认为宇宙有三个明确的本原，且每个本原都拥有无限的能力。当说到能力时，听者须注意他们所说的这话：凡你藉由理智活动所察觉的，或由于不被理解而被忽略的，其本性都适合于成为每个本原的一部分，如同在人的灵魂中，任何一门被教导的技艺。例如，他说：这个孩子将成为音乐家——等待了必要的时间（以掌握竖琴）；或几何学家——（事先经过必要的学习以获得几何学知识）；或文法家——（充分研习文法之后）；或工匠——（获得手工艺的实际经验）；而接触其他技艺的人也会发生类似的情况。他说，本原的实质是光与黑暗；在这两者之间，有一纯净的中间之灵。这灵被安置在下面的黑暗与上面的光之间。它不是如同气流或可感觉的微风，而仿佛是某种由复合物制成的油膏或香料的气味。它是一种微妙的力，藉由芬芳中某种不可名状且超越言辞的冲动性质而渗透。由于光在上、黑暗在下，而灵如所述居于其间；又因光的性质如同太阳的光线，从上方照耀下面的黑暗；且灵的芬芳处于中间位置，弥漫并向各方扩散——如同置于火上的香品，我们能察觉向各方飘散的香气——当我说，这属于被分为三类的本原的能力具有这种描述时，灵与光的能力同时存在于它们下方的黑暗之中。黑暗是可怕的水，光被吸收并转化为与灵相同性质的存在。然而黑暗并非没有理智，而是完全反思的，它意识到：当光从黑暗中被抽离后，黑暗便孤立、不可见、昏暗、无力、无作为、虚弱。因此，它被迫藉其一切反思和理解，将光的光辉和闪烁连同灵的芬芳收集到自己内。我们可以在人的面容上看到这类性质的图像：例如，眼睛的瞳孔，因下面的体液而呈黑色，却被灵照亮。因此，正如黑暗寻求光辉以束缚火花并拥有感知力，同样，光和灵也寻求属于它们自身的能力，并努力向上托举，彼此将相互混合的能力提升至下面黑暗而可怕的水中。\par

但三个本原的所有能力，在数目上无限，各自按其自身的实质是反思且智慧的，数目不可计量。既然那些反思且智慧者数目无限，它们自身静息时便皆安宁。然而，若能力接近能力，并置的差异会产生某种运动与动力，这运动与动力由并置中合力的汇聚所产生。因为能力的汇聚，如同一个印章的印记，经由与印模相应的汇聚而印在被带上来（与之接触）的实质之上。因此，既然三个本原的能力数目无限，且从无限能力产生无限汇聚，那么必然产生无限印章的图像。这些图像便是各类动物的形态。于是，从三个本原的第一次大汇聚产生了一个伟大的形态，即天地之印。天与地具有类似子宫的形状，中间有肚脐；他说，若有人希望用视觉器官来观察这一形态，便应仔细审视任何他所愿动物的怀孕子宫，他将发现天地以及一切位于万物中间不可改变之物的图像。\par

（三个本原的首次巨大汇聚）产生了天和地的形状，类似于首次交合后的子宫。然而，在天与地之间，又产生了无数力量的汇聚。每一次汇聚仅仅成就和塑造了一个类似子宫的天与地的印记。然而，在地面上，从这无数的印记中产生了各种动物的无数群体。但在天上之下所有不同动物的这无限多样性中，来自上界的灵的芬芳与光明一同散布并分配。因此，从水中产生了第一个受生的本原，即风，它猛烈而狂暴，是一切生成的原因。风在水中产生某种发酵作用，从水中掀起波浪；波浪的运动，正如某种怀孕的冲动力量成为人或思想的产生之起源，是在灵的冲动力量的激发下向前推进的。然而，当这波浪被风从水中掀起，并在其本质中怀孕时，它在自身内获得了女性所拥有的生殖能力，它凝聚了从上界散落的光明以及灵的芬芳——即在不同物种中塑造的思想。这（光明）是一位完美的天主，他从上界那非受生的光辉中，并借由灵，在自然冲动力量和风的运动下，降临到人性中，如同进入一座圣殿。它从水中产生，与身体混合交融，仿佛是有存在之物的盐和黑暗的光明。它努力从身体中释放出来，却无法为自己找到解放和出口。因为一个非常微小的火花，一个从上界剥离的碎片，如同星辰的光芒，被混杂在许多（存有）的复合水中，正如\Person{达味}在圣咏中所言。因此，上界光明的每一个思想和挂虑都是：如何以及以何种方式，通过那败坏而黑暗的身体的死亡，使思想从下面的父——即那喧嚣骚动地掀起波浪的风，并生成了自己儿子完美思想的父——中解放出来；然而，这儿子并非本质上的特有（产物）。因为他是从上界那完美光明中降下的一道光，被黑暗、可怕、苦涩、污秽的水所压制；他是降临在水面上的光明之灵。因此，当从水中掀起的波浪在自身内接受了女性所拥有的生殖能力时，它们像某种子宫一样，在不同物种中容纳了注入的光辉，以便在所有动物身上都能看见。而那同时猛烈而可怕、旋转吹动的风，就其嘶嘶声而言，如同蛇。\par

首先，从风——即从蛇——产生了如所述方式的生成本原，万物同时接受了生成的本原。然后，当光明和灵被接受进入那污秽、有害、混乱的子宫后，他说，蛇——黑暗之风，水中首生者——进入并产生人，而那不洁的子宫既不喜爱也不认识任何其他形式。因此，上界光明的完美圣言，与兽（即蛇）相似，进入了那玷污的子宫，借着与那兽本身的相似而欺骗了（子宫），以便圣言能解开那包围完美思想的锁链，这完美思想是在子宫的污秽中由水的首生者（蛇、风、兽）所生成的。他说，这就是奴仆的形象，而这正是天主圣言降孕于童贞女子宫的必要性。但他认为这并不足够：那完美的人，圣言，进入童贞女的子宫并解开了那黑暗中的产痛。不，还有更大的需要；因为在他进入子宫的污秽奥秘之后，他被洗涤，并饮用了生命涌流之水的杯。那将要脱去奴仆形象、穿上天上衣裳的，必须饮这杯。\par

% structure: p0094 source=h2 target=subsection reason=relative-heading-rank; base=h2

\subsection{第二十章 \Person{塞特}派以\WorkTitle{圣经}寓意解释支持其学说；其体系实源自自然哲学家与\Person{俄耳甫斯}礼仪；采纳\Person{荷马}宇宙起源论。}

这些是\Person{塞特}教义的拥护者所能用少数话语阐述的内容。然而，他们的体系由自然（哲学家）的命题和涉及其他不同主题的表述组成；并将这些（含义）转用于永恒圣言，按我们所述解释它们。但他们也声称\Person{梅瑟}证实了他们的学说，因为他说：\BibleQuote{黑暗、密云和暴风}。这些，（\Person{塞特}）派说，就是（我们体系的）三个本原；或者当他提到有三个生在乐园中——\Person{亚当}、\Person{厄娃}、蛇；或者当他谈到三个（人）——\Person{加音}、\Person{亚伯尔}、\Person{舍特}；以及另外三个（人）——\Person{闪}、\Person{含}、\Person{雅弗}；或者当他提到三位族长——\Person{亚巴郎}、\Person{依撒格}、\Person{雅各伯}；或者当他提到在太阳和月亮之前存在三天；或者当他提到三条法律——禁止性、允许性和惩罚裁定性。现在，禁止性法律如下：\BibleQuote{乐园中各树上的果子，你都可吃，只有知善恶树上的果子你不可吃。}但在\BibleQuote{离开你的故乡、你的家族，到我指给你的地方去}这段中，他说这是允许性的法律；因为愿意的人可以离去，不愿意的人可以留下。而惩罚裁定性的法律是发表以下声明的：\BibleQuote{不可奸淫，不可杀人，不可偷盗}；因为对这些邪恶行为中的每一种都规定了惩罚。\par

然而，他们全部学说的体系，皆衍生自古代神学家\Person{穆赛欧斯}、\Person{利诺斯}和\Person{俄耳甫斯}；俄耳甫斯尤其阐明入教仪式以及奥秘本身。因为他们关于子宫的教义也是俄耳甫斯的信条；而脐（的意象），即和谐，在俄耳甫斯的酒神秘仪中也具有同样的象征意义。但在\Person{刻琉斯}、\Person{特里普托勒摩斯}、\Person{刻瑞斯}、\Person{珀耳塞福涅}和\Person{巴克斯}于\Place{厄琉息斯}举行神秘仪式之前，这些秘仪已在\Place{阿提卡}的\Place{弗利乌斯}举行并传授于人。因为在厄琉息斯秘仪之前，在弗利乌斯已举行了被称为\Quote{大母神}的秘仪。而该地有一门廊，门廊上刻有至今仍可见的（所有）仪式中念诵之词的图画。那么，门廊上所刻的许多词，正是\Person{普鲁塔克}在其反对\Person{恩培多克勒}的十卷著作中所探讨的内容。这些著作的大部分卷帙中还绘有一位老人的形象：灰发、有翼、\Emph{\OriginalTerm{阳具勃起}{pudendum erectum}}，追逐一位后退的蔚蓝色女子。老人上方刻有题字\Quote{\OriginalTerm{光流}{phaos ruentes}}，女子上方刻有题字\Quote{\OriginalTerm{彼方}{peree}}。\par

\SourceRecord{Translated by J.H. MacMahon.；Ante-Nicene Fathers；Vol. 5.；Edited by Alexander Roberts, James Donaldson, and A. Cleveland Coxe.；Buffalo, NY: Christian Literature Publishing Co.,；1886.；<http://www.newadvent.org/fathers/050105.htm>.}

% structure: p0001 source=h1 target=section reason=page-title

\section{\WorkTitle{驳斥一切异端（第六卷）}}

% structure: p0002 source=h2 target=subsection reason=relative-heading-rank; base=h2

\subsection{目录}

以下是\WorkTitle{驳斥一切异端}第六卷的内容：——\par

西满试图（确立）的见解是什么，以及他的教义从术士和诗人的（著作）中获得力量。\par

瓦伦廷所提出的见解是什么，以及他的体系不是出于圣经，而是出于柏拉图与毕达哥拉斯的学说。\par

塞孔都斯、托勒密和赫拉克利翁的见解是什么，这些人自己也宣扬与希腊哲学家相同的教义，只是以不同的措辞表达。\par

马尔库斯与科拉巴西乌斯提出的假设是什么，以及他们中有些人致力于巫术和毕达哥拉斯的数字。\par

% structure: p0008 source=h2 target=subsection reason=relative-heading-rank; base=h2

\subsection{第一章——奥菲特派：后世异端的始祖}

凡是从蛇那里获取（教义）首要原理、并逐渐刻意将其主张公之于众的那些人所持的见解，我们已在前一本即\WorkTitle{驳斥异端}第五卷中予以说明。但如今我也不对追随这些（奥菲特派）的异端首领的见解保持沉默；不，我将不放过任何一个（猜测），只要可能记住所有（他们的主张）以及这些异端者的秘仪——人们完全可以称之为秘仪——因为传播如此大胆主张的人离（天主的）义怒不远——以便我能借助词源学的帮助。\par

% structure: p0010 source=h2 target=subsection reason=relative-heading-rank; base=h2

\subsection{第二章——西满术士}

看来现在同样适宜说明撒玛利亚村庄吉塔人西满的见解；我们也将证明，他的后继者从他出发，试图以换名的方式建立类似的见解。此西满精通巫术，嘲弄许多人，部分按特拉叙墨得斯的技艺，如我们以上所说明的方式，部分也藉鬼魔的帮助行恶，试图神化自己。（但）这人（不过是）骗子，充满愚妄，宗徒们曾在\BibleBook{宗徒}{大事录}中斥责他。\BibleRef{宗 8:9}利比亚人阿普塞图斯比西满以多得多的智慧和节制，被类似的渴望所激动，竭力使自己在利比亚被视为神明。既然他那传说体系与那愚昧西满的过度欲望并无太大差异，叙述此事似乎是适宜的，因它值得此人所作的努力。\par

% structure: p0012 source=h2 target=subsection reason=relative-heading-rank; base=h2

\subsection{第三章——利比亚人阿普塞图斯的故事}

利比亚人阿普塞图斯过度渴望成为神；但当他多次密谋后完全未能实现其愿望时，他仍希望显得自己已成了\Emph{一位神}；随着时间的推移，他确实显得真成了神。因为愚昧的利比亚人惯于向他献祭，如同向某种神性，以为他们听信了从天上传下来的声音。因为他将大量的鸟——鹦鹉——收集到同一只笼子里，关了起来。利比亚境内有很多鹦鹉，它们非常清晰地模仿人声。这人喂养这些鸟一段时间后，习惯教它们说：\Quote{阿普塞图斯是神。}当鸟练习这语句很长时间，并惯于说出他所想的、说出来会让人以为阿普塞图斯是神的话之后，他就打开（鸟的）住所，放出鹦鹉，每只飞向不同方向。然而，当鸟振翅飞翔时，它们的声音传遍了整个利比亚，它们的言语一直传到希腊地区。于是利比亚人惊讶于鸟的声音，没有察觉阿普塞图斯所行的诡计，就视阿普塞图斯为神。但有一位希腊人通过仔细探究，识破了这假神的伎俩，藉同样的鹦鹉不仅驳斥，而且彻底毁灭了那自夸可厌的家伙。希腊人关了许多鹦鹉，重新教它们说：\Quote{阿普塞图斯关了我们，强迫我们说：阿普塞图斯是神。}利比亚人听到鹦鹉的撤回声明，聚集起来，全体一致决定烧死阿普塞图斯。\par

% structure: p0014 source=h2 target=subsection reason=relative-heading-rank; base=h2

\subsection{第四章——西满对圣经的强解；剽窃赫拉克利特与亚里士多德；西满关于可感与可知存有的体系}

我们当这样思想术士西满，使我们得以将他比作那个利比亚人，远胜过比作那一位虽成为人、却是真实天主者。然而，若这相似的说法本身是准确的，术士也受了与阿普塞图斯相似的激情，就让我们重新教导西满的鹦鹉：基督，祂曾站立、站立、将要站立（即昔在、今在、永在的），并不是西满。（耶稣）却是人，女人的后裔，由血气和肉欲所生，与其余（人类）一样。至于这些事如此，我们将在随后的讨论中容易地证明。\par

现在，\Person{西满}既愚蠢又狡诈地曲解\Person{梅瑟}的法律，作出以下陈述：因为当\Person{梅瑟}断言说：\BibleQuote{天主是燃烧的烈火和吞噬的火，}\BibleRef{申 4:24}——\Person{西满}未按正确意义理解\Person{梅瑟}所说的话，便断定火是宇宙的起源原则。但他并未考虑所陈述的内容，即：天主不是火，而是燃烧的烈火和吞噬的火，这不仅是对\Person{梅瑟}法律本身的强解，甚至是从晦涩的\Person{赫拉克利图斯}那里剽窃来的。\Person{西满}将宇宙的起源原则称为一种无限的能力，他这样表达自己：\Quote{这是藉着对伟大无限能力的理智洞察而对声音和名称的启示的论著。因此它将被密封、保密、隐藏，并安息在那居于其根基为万有之根的住所中。}并且他断言，这个由血而生的人就是前述住所，在他内住着一种无限的能力，他声称这就是宇宙的根。\par

按照\Person{西满}的说法，作为火的无限能力并非任何一种单一本质（与那些主张四元素是单纯的人之意见一致，他们也因此同样以为作为四元素之一的火是单纯的）。但实际情况远非如此：因为火具有某种双重本性；他将这双重本性的一部分称为隐藏之物，另一部分称为显现之物，并且隐藏之物隐藏在火的显现部分中，而火的显现部分从其隐藏部分获得存在。然而，这正是\Person{亚里士多德}用\Quote{潜能}和\Quote{现实}所称呼的，或者\Person{柏拉图}称之为\Quote{可理知的}和\Quote{可感觉的}。火的显现部分包含一切事物在其自身内，无论是任何人可能察觉的，乃至他可能忽略的可见受造界的任何对象。但整个隐藏部分（即火的隐藏部分）若有人能察觉，则是由理智认知的，逃脱感官的能力；或者，由于缺乏那种特定感知的能力，人便无法观察到它。但总的来说，既然一切存在之物都归属于范畴，即可感觉的对象和可理知的对象，而对于这些的命名，\Person{西满}使用了\Quote{隐藏}和\Quote{显现}这两个术语；那么可以断言，火——我指的是超天界的火——是一个宝藏，如同一棵大树，就像\Person{拿步高}在梦中所见的那棵，一切血肉都由它而得滋养。他将火的显现部分视为树干、树枝、树叶和覆盖其上的外皮。他说，这棵大树的所有这些附属物一旦被点燃，就会因吞灭一切的火焰而消失。然而，当树的果实完全成熟并接受了它自己的形状时，便被存放在谷仓里，而不是投入火中。因为他说，果实被产生是为了收入仓库，而糠秕被交付给火。糠秕就是树干，它的产生不是为了自身，而是为了果实。\par

% structure: p0018 source=h2 target=subsection reason=relative-heading-rank; base=h2

\subsection{第五章 \Person{西满}援引圣经以支持其体系}

他说，这就是经上所记载的：\BibleQuote{因为万军上主的葡萄园是以色列家，犹大人是他所喜爱的植物。}然而，如果犹大人是所喜爱的植物，他说，那就证明除了那人之外没有别的树。但是关于这棵树的分泌和消解，他说，圣经已经充分论述了。至于对那些按他的形象被塑造之人的教导，经上所说的那句话就够了：\BibleQuote{凡有血肉的尽都如草，他的一切荣美都像草上的花；草必枯干，花必凋谢；但主的话却永远长存。}他说，主的话就是那在口中产生的话语，是一个\OriginalTerm{逻各斯}{Logos}，此外没有任何其他地方有生成之处。\par

% structure: p0020 source=h2 target=subsection reason=relative-heading-rank; base=h2

\subsection{第六章 \Person{西满}的体系在\WorkTitle{伟大宣言}中详述；追随\Person{恩培多克勒}}

现在，简而言之，既然按照\Person{西满}，火具有这样的描述，并且既然所有事物都是可见的和不可见的，同样地鸣响的和不鸣响的，可数的和不可数的；他在\WorkTitle{伟大宣言}中将一种完美的可理知实体如此命名，以至于那些无限存在、可以被无限理解的事物中的每一个，都以\Person{恩培多克勒}所说的方式说话、理解和行动：——\par

因为大地确实由大地我们看见，水由水，\newline{}神圣的空气由空气，凶猛的火焰由火，\newline{}爱由爱，以及恨由阴暗的恨。\par

% structure: p0023 source=h2 target=subsection reason=relative-heading-rank; base=h2

\subsection{第七章 \Person{西满}关于成对的三重流溢体系}

因为，他说，他习惯于认为，这些火的各部分，无论是可见的还是不可见的，都具有知觉和一份智能。因此，世界，即受造的世界，是从非受生的火产生的。然而，它开始存在，他说，按照以下方式。那从该火的本原所生者取用了六条根，即生成原始本原的初始根。而且，他说，这些根由火成对地产生，他称这些根为\Quote{\OriginalTerm{心智}{Mind}}和\Quote{\OriginalTerm{理智}{Intelligence}}，\Quote{\OriginalTerm{声音}{Voice}}和\Quote{\OriginalTerm{名}{Name}}，\Quote{\OriginalTerm{推理}{Ratiocination}}和\Quote{\OriginalTerm{反思}{Reflection}}。并且在这六条根中，同时潜在地存有全部无限的能力，但实际上并不存有。他说，这个无限的能力就是那曾站立、站立、将要站立者。因此，无论何时，当他被塑造成一个形象时，既然他存在于这六种能力中，他将在实质、潜势、量度、完全上存在（并且他将是一种能力），与那非受生的、无限的能力相同，并不比那非受生的、不变的和无限的能力更欠缺。然而，如果他仅仅潜在地存在于这六种能力中，而没有被塑造成一个形象，他便会消失，他说，并会像人的灵魂中的文法或几何能力那样被毁灭。因为当能力接纳了一门技艺时，就产生了现存之物的光；但若能力没有接纳（一门技艺），结果便是无能和无知；正如当（能力）不存在时，它随着人断气而消亡。\par

% structure: p0025 source=h2 target=subsection reason=relative-heading-rank; base=h2

\subsection{第八章。这三重流溢的进一步进展；与双三一体共存之第七存在。}

在这六种能力以及与之共存的第七种能力中，第一对，即心智和理智，他称之为天和地。其中一方，为男性，从上俯瞰并照管其伴侣；而地则在下接受从天而降的理性果实，相似于地。为此，他说，\OriginalTerm{圣言}{Logos}常常注视那从心智和理智，即从天和地，所生成之物，并呼喊说：\BibleQuote{诸天，请听！大地，请侧耳！因为主发言：我养育了子女，使他们长大，他们却背叛了我！}现在，说出这些话的，他说，是第七种能力——那曾站立、站立、将要站立者；因为他本人就是摩西所称赞的那些美丽造物的原因，摩西说它们非常好。而声音和名（三对中的第二对）是太阳和月亮；推理和反思（三对中的第三对）是空气和水。在这一切之中，如我所言，那伟大的、无限的、自在的能力已混合并交织其中。\par

% structure: p0027 source=h2 target=subsection reason=relative-heading-rank; base=h2

\subsection{第九章。西蒙对摩西六日创世的解释；他对乐园的寓意描绘。}

因此，当摩西说到\BibleQuote{六天之内，天主造了天和地，并在第七天歇了祂所做的一切工作}\BibleRef{创 2:2}时，西蒙，以已述的方式，给予（这些和其他圣经段落）一种不同的应用（与圣作者的原意不同），而自比为神。因此，当（西蒙的追随者）声称在太阳和月亮之前有三日被生时，他们是以隐晦的方式谈论心智和理智，即天和地，以及第七种能力，（即）无限者。因为这三重能力是在其他一切之前产生的。但当他们说\BibleQuote{祂在万世之前生了我}\BibleRef{箴 8:22}时，他说，这类说法被认为是适用于第七种能力。现在，这第七种能力，是存在于无限能力中的一种能力，在万世之前就产生了，这就是，他说，第七种能力，关于它，摩西说了以下的话：\BibleQuote{天主的神在水面上运行}；即，西蒙派说，那包含万有于自身之中、且是无限能力之形象的神——\Quote{一个来自不朽形式之形象，惟有它使万有归于秩序}。因为那运行于水面上的能力，他说，是从一个不朽的形式所生，使万有归于秩序。因此，根据这些（异端者），当世界有了某些这样的安排和（一个）类似的（安排）时，他说，天主就着手造人，用地上的土塑造人。\BibleQuote{祂造人不是单一的，而是双重的，按照（祂自己的）肖像和模样。}\BibleRef{创 2:7}而肖像就是运行于水面上的神；凡未成形为这肖像者，将与世界一同灭亡，因为他仅潜在地存有，实际上并不存在。他说，这就是所说的，\BibleQuote{免得我们和世界一同被定罪}。\BibleRef{格前 11:32}但是，如果一个人被塑造成（那神）的形象，并且从一个不可分的点所生，正如\WorkTitle{宣告}中所写的，（这样的人虽）小，却将成为大的。而那大的将延续至无限不变的存在，因为它不再受制于受造物的条件。\par

那么，他说，天主怎样并如何形成人呢？在乐园里；因为他似乎这样认为。他说，请允许把乐园当作子宫；而这是一个真实的（假定），圣经会教导我们，当它说出这话时：\Quote{是我在你母胎中形成了你。}\BibleRef{耶 1:5} 因为他希望这也是如此写下的。他说，梅瑟采用寓意手法，已经宣告乐园就是子宫，如果我们可以信赖他的陈述。然而，如果天主在母胎中形成人——即在乐园中——如我所肯定的那样，那么就让乐园是子宫，厄登是胎衣，\Quote{有一条河从厄登流出，灌溉乐园，}\BibleRef{创 2:10}（意指）脐。他说，这脐分为四个本原；因为脐的两侧各有一条动脉，是气的通道，以及两条静脉，是血的通道。但是，他说，当脐带从厄登（即胎衣）发出时，包裹胎儿的胎膜长入正在上腹部附近形成的（胎儿）中——（现在）所有人都称之为脐——这两条静脉，血液通过它们流动并从胎衣厄登输送到所谓的肝门；（我说，这些静脉）滋养胎儿。但我们所说的作为气之通道的动脉，在骨盆周围包裹着膀胱两侧，并将其与脊柱附近的大动脉（称为主动脉）相连。这样，气通过心室进入心脏，产生胎儿的运动。因为在乐园中成形的婴孩既不通过口接受营养，也不通过鼻孔呼吸：因为它处于水分之中，若试图呼吸，死亡就在它脚下；因为它会（这样）被抽离水分而（相应）消亡。但（还可以更进一步）；因为整个（胎儿）被一种称为胎膜的覆盖物紧紧包裹，由脐带滋养，并通过（主动脉）在脊柱附近，如我所说，接收气的实质。\par

% structure: p0030 source=h2 target=subsection reason=relative-heading-rank; base=h2

\subsection{第十章 西满对梅瑟前两卷书的解释。}

因此，他说，那从厄登流出的河分为四个本原，四条支流——即分成四种感觉，属于那即将诞生的受造物：视觉、嗅觉、味觉和触觉；因为在乐园中成形的孩子只有这些感觉。他说，这就是梅瑟所立的法律；而正是针对这法律，他的每一卷书都是这样写成的，正如题署所表明的。第一卷书是创世纪。他说，该书的题署足以认识宇宙。因为这就是（等同于）生成，（即）视觉，河流的一个分支进入其中。因为世界是通过视觉的力量被看见的。同样，第二卷书的题署是出谷纪。因为那被产生的，经过红海，必须进入旷野——他们说红（海）是血——并尝苦水。他说，因为那（喝）在红海之后的（水）是苦的；这（水）是一条必须踏上的道路，引导（我们）在（此）生中认识我们劳苦和苦涩的命运。然而，被梅瑟——即\OriginalTerm{圣言}{Logos}——改变后，那苦的（水）变为甜的。这一点，我们可以普遍地从所有按照（诗人）情感表达自己的人那里听到：—\par

\Quote{根暗，如乳，花朵，\newline{}神祇称之为\InnerQuote{摩吕}，凡人难\newline{}挖掘，但神能无限。}\par

% structure: p0033 source=h2 target=subsection reason=relative-heading-rank; base=h2

\subsection{第十一章 西满对梅瑟最后三卷书的解释。}

外邦人所说的，对于那些有（能）听之耳的人来说，足以认识宇宙。他说，因为凡尝过这果实的，不仅被基尔刻变成野兽；而且，运用这果实的效力，他再次形成和重塑，并诱使那些已变成野兽的回到他们最初特有的性格。他说，一个忠实的人，被那女巫所爱，是通过那如乳般的神圣果实而被发现的。同样，第三卷书是肋未纪，即嗅觉或呼吸。因为那整卷书是（关于）祭献和供物的。然而，在祭献之处，通过焚香，从祭献中升起某种芳香气味；而针对这气味，（感官）嗅觉是一种检验。第四卷书户籍纪，表示味觉，在那里论述是起作用的。因为从它讲述一切事物这一事实，它被数字编排命名。他说，但申命纪是针对正在成形孩子的（感官）触觉而写的。因为正如触觉通过抓住其他感官所见之物来总结并确认它们，测试粗糙、温暖、粘湿（或寒冷）之物；同样，法律的第五卷书构成了前四卷书的概要。\par

因此，他说，一切事物，当非受生时，在我们中是潜在地，而非现实地，如文法或几何（技艺）。如果有人接受了适当的指导和教导，（因而）那苦涩的将被变为甘甜的——即\BibleQuote{将矛打成镰刀，将刀打成犁头}\BibleRef{依 2:4}——就不会有糠秕和木头为火而生，而是成熟的、完全成形的果实，如我所说，与那非受生和无限的力量相等且相似。然而，如果一棵树单独存在，不产出完全成形的果实，它就被彻底毁灭。因为他说，斧子已放在树根上。他说，\Quote{凡不结好果子的树，就砍下来，丢在火里。}\par

% structure: p0036 source=h2 target=subsection reason=relative-heading-rank; base=h2

\subsection{第十二章 火是元始本原，按西满之说。}

\Person{西蒙}，因此，那个有福且不朽的以潜在状态存在于每个人中——即，潜在而非现实；这就是那曾站、站、且将站者。他曾站在上方，在\OriginalTerm{非受生之力}{unbegotten power}中。他站在下方，当他在水流中按相似受生时。他将在上方站在有福的无限之力旁，若他被塑造成一个形象。因为，他说，有三位曾站者；若没有三位\OriginalTerm{永世}{Aeons}曾站，那非受生者便不得装饰。（非受生者）按照他们，飘浮在水上，并按（永恒本性的）相似被重造为完美的属天存在，在（任何）智能方面都不逊于非受生之力：这就是他们所说的——\Quote{我与你，一；你在我之前；我，在你之后。}他说，这是一个分为上下的能力，自我生成，自我成长，寻找自己，找到自己，是自己的母亲、父亲、姐妹、配偶、女儿、儿子、母亲、父亲、一个单元，是整个存在之圆的根源。\par

他说，受生之物的生成的本原来自火，他通过以下方式辨认。万物（即，凡有生成之物）的生成欲望之始来自火。因此，欲求可变生成的欲望被称为\Quote{被点燃}。因为火是一，却容许两种转变。因为，他说，人中的血既温暖又黄色，如同有形状的火焰转变为种子；但女人中同样的血转变为乳汁。男性的转变成为生成，女性的转变成为对胎儿的滋养。他说，这就是\BibleQuote{那火焰的剑，转动把守生命树的道路}。\BibleRef{创 3:24}因为血转变为种子和乳汁，而这一能力成为母亲和父亲——生成中之物的父亲，以及被滋养之物的增长；（这能力）无所缺乏，自足。他说，如我们所说，生命树被挥舞的火焰之剑守护。它是第七能力，从自身产生，包含所有（能力），安息于六能力中。因为若火焰之剑不被挥舞，那好树将被毁坏而消失。然而，若这些转变为种子和乳汁，那潜在存在于其中、并拥有适当位置的原则，其中发展出灵魂的原则，（这样的原则）如同从极小的火花开始，将被完全放大，增长并成为无限而不变的能力，（相等且相似于）不变的永恒，不再进入无限永恒。\par

% structure: p0039 source=h2 target=subsection reason=relative-heading-rank; base=h2

\subsection{十三、他的流溢说进一步展开}

因此，按照这一推理，\Person{西蒙}对他愚蠢的追随者来说成了公认的神，如那位\Place{利比亚}人，即\Person{阿普塞图斯}——无疑受生而可受情欲，当他潜在地存在时，却无倾向。（这也，尽管从有倾向者出生，并非受造却从受生者出生），当他被塑造成一个形象，并成为完美时，可能从两个原初能力，即天和地出来。因为\Person{西蒙}在\WorkTitle{启示录}中明确如此说：\Quote{那么，我对你们说我所讲的事，并（对你们）写我所写的。这书写如下：从所有\OriginalTerm{永世}{Aeons}有两个分枝，既无始也无终，来自一个根。这是一个能力，即\OriginalTerm{希格}{Sige}，（她是）不可见（且）不可理解的。其中一个（分枝）从上面显现，构成了一个伟大的能力，（创造性）宇宙的\OriginalTerm{理智}{Mind}，管理一切，（是）雄性。另一个（分枝），却从下面，（构成了）一个伟大的\OriginalTerm{智能}{Intelligence}，是雌性，生产一切。从此，成对相对排列，它们经历结合，显现出一个中间间隔，即不可理解之气，无始无终。但在其中有一位父亲，支撑一切，滋养有始有终之物。这就是那曾站、站、且将站者，是按先存的无限能力的\OriginalTerm{两性之力}{hermaphrodite power}，无始无终。现在这（能力）独存。因为\OriginalTerm{理智}{Intelligence}，（存在于）统一中，从这（能力）出来，（并）成为二。那位（父亲）是一，因为在他内拥有此（能力），他是孤立的，但，他并非原初，虽是先存的；但由于从自身使自己显现，他进入二重状态。但在此之前，他并未被称为父亲，直到这（能力）称他为父亲。因此，正如他自身，通过自身引出自已，向自己显示了他自己的特殊\OriginalTerm{理智}{Intelligence}，同样，\OriginalTerm{智能}{Intelligence}在显现时，并没有行使创造功能。但看见他，她把父亲隐藏在自己内，即那能力；这是一个\OriginalTerm{两性之力}{hermaphrodite power}和一个\OriginalTerm{智能}{Intelligence}。因此，它们成对排列，一个相对另一个；因为能力与\OriginalTerm{智能}{Intelligence}绝无不同，因为它们是 一。因为从上面之物中发现了能力；从下面之物中发现了\OriginalTerm{智能}{Intelligence}。所以，同样，从它们显现的，是一，被发现为二，一个自身包含雌性的两性体。这（因此）是（存在于\OriginalTerm{智能}{Intelligence}中的）\OriginalTerm{理智}{Mind}；这些彼此可分离，（虽然合在一起）是一，（并）被发现为二重状态。}\par

% structure: p0041 source=h2 target=subsection reason=relative-heading-rank; base=h2

\subsection{十四、西蒙用特洛伊的海伦的神话解释他的体系；关于他自己与特洛伊女英雄的叙述；其追随者的不道德；西蒙关于基督的观点；西蒙派为其恶习的辩护}

西满发明了这些（道理）之后，不仅用恶计随意解释梅瑟的著作，甚至（也解释）诗人的（作品）。因为他也把（关于）木马和持火炬的海伦，以及许多其他（故事）赋予寓意，并将其转移到与他自身以及智有关的内容，（并）对它们作出虚构的解释。然而，他说这个（海伦）就是亡羊。她常住在妇女中间，因她的超凡美貌而混乱了世界上的权能。因此，特洛伊战争也因她而起。因为当时出生的海伦身上住着这个智；因此，当所有权能都想（把她）占为己有时，就发生了纷争和战争，在此期间（这个主要权能）显示给了万民。由此，无疑我们可以相信，通过（某些）诗句辱骂过她的斯特西科罗斯被剥夺了视力；而当他悔改并创作了翻案诗，在其中歌颂（海伦的）赞美时，他又恢复了视觉。但是天使和下界的权能——他说这些权能创造了世界——导致了（海伦灵魂）从一个身体到另一个身体的转移；后来她站在腓尼基城提洛的一座房屋的屋顶上，当（西满）下去那里时，（他声称）找到了她。因为他宣称，他来到（提洛）主要是为了寻找这个（女人），以便把她从束缚中解救出来。在这样赎回她之后，他习惯带着她同行，声称这个（女孩）就是亡羊，并肯定他自己是超越万有的权能。但这个污秽的家伙迷恋这个名叫海伦的可怜女人，买下她（为奴），并享用了她。他，（然而，）也对门徒感到羞愧，就编造了这个虚构之事。\par

但是，再次，那些跟随这个骗子——我指的是巫师西满——的门徒也从事类似的行为，并非理性地断言混杂交合的必要性。他们表达自己如下：\Quote{所有的土地都是土地，无论在哪里播种，只要播种，就没有区别。}然而，他们甚至为这种混杂交合而自我庆幸，声称这是完美的爱，并使用\Quote{至圣所}和\Quote{彼此圣化}这些说法。因为（他们想让我们相信）他们没有被所谓的恶习征服，因为他们已被赎回。\Quote{（耶稣）通过这种方式赎回海伦，}（西满说，）通过他自己的独特智慧赐予了人类救恩。因为既然天使因其贪图优越而错误地管理了世界，（耶稣基督）被改变，并变得与统治者和权能和天使相似，来（复兴万物）。因此（耶稣）以人的形象出现，而实际上他并不是人。并且（同样地）他受苦——尽管实际上没有遭受苦难，只是向犹太人看似如此——在犹太地作为\Quote{子}，在撒玛黎雅作为\Quote{父}，在其余外邦人中作为\Quote{圣神}。（西满声称）耶稣容忍人们用这三者中任意一个名字来称呼他。\Quote{先知们从创造世界的天使那里获得灵感，发出了（关于他的）预言。}因此，（西满说，）那些相信西满和海伦的人至今对这些（先知）毫不关心，他们作为自由人，随心所欲；因为他们声称是靠恩宠得救。因为没有一个理由需要惩罚，即使有人行为邪恶；因为这样的人不是本性邪恶，而是因律法而邪恶。\Quote{因为创造世界的天使们，}他说，\Quote{制定了他们喜欢的任何律法，}以为通过这样（立法）的言辞就能奴役那些听从他们的人。但是，再次，他们谈到世界的毁灭，为了他特别追随者的救赎。\par

% structure: p0044 source=h2 target=subsection reason=relative-heading-rank; base=h2

\subsection{第十五章 西满的门徒接纳奥秘；西满在罗马遇见圣伯多禄；西满晚年纪事。}

这术士的门徒举行魔法的仪式，并诉诸咒语。他们传授爱情咒语和符咒，以及被称为梦的使者的魔鬼，以迷惑他们愿意的任何人。他们也使用那些称为\OriginalTerm{帕雷德罗}{Paredroi}的灵体。\Quote{他们有一个用\Person{朱庇特}形象塑造的\Person{西满}像，还有一个\Person{弥涅耳瓦}形象的\Person{海伦}像，并朝拜它们。}但他们称一位为主，另一位为夫人。如果有人看见\Person{西满}或\Person{海伦}的像，直呼其名，就会被驱逐，因为他不了解奥秘。这个\Person{西满}在撒玛黎雅用巫术欺骗许多人，被宗徒们责备，并受到诅咒，正如\BibleBook{宗徒}{大事录}所记载的。但他后来背弃了信仰，并尝试了上述行为。他一直旅行到罗马，遇到了宗徒们；\Person{伯多禄}多次反对他，因为他用巫术欺骗许多人。此人最终去到……，坐在一棵梧桐树下，继续教导他的教义。的确，最后，当定罪迫在眉睫，他若再拖延，就声称如果他被活埋，他第三天会复活。于是，他命令门徒挖一条壕沟，并指示自己被埋在那里。他们执行了命令，而他一直到今天仍在那坟墓里，因为他不是基督。这就是\Person{西满}所提出的传奇体系，而\Person{瓦伦廷}由此获得了出发点了。事实上，这教义与西满的教义相同，只是瓦伦廷用不同的名称称呼它：因为\Quote{努斯}、\Quote{阿勒忒亚}、\Quote{逻各斯}、\Quote{佐厄}、\Quote{安特罗波斯}、\Quote{厄克利西亚}和瓦伦廷的永世，是西满的六个根源，即\Quote{心灵}、\Quote{理智}、\Quote{声音}、\Quote{名字}、\Quote{推理}、\Quote{反思}。但既然我们觉得已经充分解释了西满传说的网络，让我们也看看瓦伦廷的主张。\par

% structure: p0046 source=h2 target=subsection reason=relative-heading-rank; base=h2

\subsection{第十六章 瓦伦廷的异端；源自柏拉图和毕达哥拉斯。}

\Person{瓦伦廷}的异端确实与\Person{毕达哥拉斯}和\Person{柏拉图}的学理有关。因为\Person{柏拉图}在\WorkTitle{蒂迈欧篇}中完全从\Person{毕达哥拉斯}获取其印象，因此\Person{蒂迈欧}本人就是他的毕达哥拉斯式陌生人。所以，我们似乎应当先提醒读者\Person{毕达哥拉斯}和\Person{柏拉图}学理中的几个要点，然后再阐述\Person{瓦伦廷}的意见。因为尽管在我们先前费力完成的书籍中已包含了\Person{毕达哥拉斯}和\Person{柏拉图}所提出的意见，但我现在通过概要回顾这些思想家的主要信条，也并非不合理。这种回顾将有助于我们通过更近的比较和两个体系的相似构成，来了解\Person{瓦伦廷}的教义。因为他们从埃及人那里获得了这些信条，并将其新奇的意见引入希腊。但\Person{瓦伦廷}从他们那里获取意见，因为尽管他隐瞒了他对这些希腊哲学家的亏欠，并以此试图构建一个似乎是独特的教义，但事实上，他仅仅在名称和数字上改变了那些思想家的教义，并采用了独特的术语。\Person{瓦伦廷}按度量制定了他的定义，以便建立一个希腊化的异端，无疑是多样的，但不稳定，且与基督无关。\par

% structure: p0048 source=h2 target=subsection reason=relative-heading-rank; base=h2

\subsection{第十七章 希腊哲学的起源。}

那么，\Person{柏拉图}在\WorkTitle{蒂迈欧篇}中获取其理论的起源，是埃及人的智慧。因为根据某种古老而预言的传说，\Person{梭伦}从这一源泉将他的整个关于世界生成和毁灭的体系教给希腊人——正如\Person{柏拉图}所说——希腊人在知识上如同幼童，不熟悉更古老的神学教义。因此，为了准确追溯\Person{瓦伦廷}建立其信条的论证，我现在将解释萨摩斯的\Person{毕达哥拉斯}的哲学原理——一种与希腊人如此颂扬的\OriginalTerm{沉默}{Silence}相连的哲学。然后以此方式，我将阐明\Person{瓦伦廷}从\Person{毕达哥拉斯}和\Person{柏拉图}那里获取的那些意见，但他却以极其庄重的言辞将其归诸基督，在基督之前则归诸宇宙的圣父以及与圣父相连的沉默。\par

% structure: p0050 source=h2 target=subsection reason=relative-heading-rank; base=h2

\subsection{第十八章 毕达哥拉斯的数字系统。}

\Person{毕达哥拉斯}宣称宇宙的起源原则是非受生的\OriginalTerm{单子}{monad}、受生的\OriginalTerm{二元}{duad}以及其余的数字。他说单子是二元之父，二元是万有之母——受生者乃受生之物的母。\Person{毕达哥拉斯}的学生\Person{匝拉塔斯}习惯于称一为父，二为母。因为根据\Person{毕达哥拉斯}，二元由单子生成；单子是男性且首要的，二元是女性且次要的。又从二元，如\Person{毕达哥拉斯}所言，生成了\OriginalTerm{三元}{triad}及直到十的后续数字。因为\Person{毕达哥拉斯}知道只有这个完美数字——我指的是\OriginalTerm{十}{decade}——因为十一和十二是十的增加和重复；然而，所加的并不构成另一个数字的生成。他还从无形本质生成所有固体。因为他断言，无形与有形实体的元素和原则是不可分的点。他说，从点生成线，从线生成面；面流入高度，他说，就成为固体。因此，毕达哥拉斯学派有一种起誓对象，即四元素的和谐。他们用以下话语宣誓：——\par

由那赋予我们之首者以\OriginalTerm{四元组}{quaternion}者，\newline{}一个具有永恒本性之根的泉源。\par

如今四元组是自然固体物体的根源原则，如同一元组是理智物体的根源原则。而他也说四元组产生完全数，正如在理智物中一元组产生十进制数，他们这样教导。若有人开始数数，说一，然后加二，同样再加三，这些（总和）将是六，再给这些加四，总和同样将是十。因为一、二、三、四成为十，即完全数。如此，他说，四元组在各方面模仿了能产生完全数的一元组。\par

% structure: p0054 source=h2 target=subsection reason=relative-heading-rank; base=h2

\subsection{第十九章　毕达哥拉斯的双重实体；他的\Quote{范畴}。}

那么，据毕达哥拉斯，有两个世界：一个理智世界，以一元组为根源原则；另一个是可感世界。而后者有四元组，含有约塔那一点，\BibleRef{玛 5:18}一个完全数。而且，据毕达哥拉斯派，约塔，即那一点，是首要且最具统治力的，它使我们能把握那些可通过理智和感觉理解之理智实体的实质。（在此实体中）内蕴着九个不能脱离实体而存在的非形体偶性，即：\Quote{性质}、\Quote{数量}、\Quote{关系}、\Quote{地点}、\Quote{时间}、\Quote{位置}、\Quote{拥有}、\Quote{动作}和\Quote{承受}。这些就是（内蕴于）实体的九个偶性，与这些（实体）合计时，便构成完全数，即约塔。因此，宇宙被划分为理智世界和可感世界，正如我们所说，我们也有来自理智（世界）的理性，以便通过理性，我们得以注视那些通过理智认知、无形且神圣之事物的实质。但他又说，我们有五种感官——嗅觉、视觉、听觉、味觉和触觉。通过这些，我们获得对可感知事物的知识；所以，他说，可感世界从理智世界划分出来。而且，我们为每种都有获取知识的工具，这从以下考虑可知。他说，没有任何理智之物能通过感官为我们所知。因为眼睛未曾看见，耳朵未曾听见，任何其他感官也未曾知晓那（由心智认知之物）。同样，也无法通过理性达到对任何可感官事物的知识。但人必须看到某物是白色的，尝到它是甜的，通过听觉知道它是悦耳的或不和谐的。任何气味是芬芳的还是难闻的，是嗅觉的功能，而非理性。触摸对象也是如此；因为粗糙、柔软、温暖或寒冷，不可能通过听觉得知，而是由触觉判断这些（感觉）。如此构成之后，已造和正在造之物的安排被观察到符合数字（组合）。因为正如从一元组开始，通过增加一元组或三元组，以及后续数字的集合，我们得出某个非常庞大的数字复合整体；然后，又从如此累积的数字，通过某种减法和重新计算，我们完成总集合数的分解；同样，他断言，世界被某种算术和音乐的锁链束缚，通过其张弛、加减，永远且持续地保持不朽。\par

% structure: p0056 source=h2 target=subsection reason=relative-heading-rank; base=h2

\subsection{第二十章　毕达哥拉斯的宇宙生成论；类似恩培多克勒。}

毕达哥拉斯派因此以类似这样的方式宣告他们对世界延续的观点：——\par

因为过去它已是，将来也仍是；我料想，\newline{}这两者永不会使永恒时代空虚。\par

\Quote{这些；}——但它们是哪两者？纷争与爱。如今在他们的体系中，爱使世界成为不朽的、永恒的，正如他们所设想。因为实体与世界是一回事。然而，纷争则分离、拆散，并通过细分而多次尝试形成世界。正如有人将万军切成小块，并算术地分为千、百、十；将银币分为小钱和铜币。这样，他说，纷争将世界的实体分为动物、植物、金属以及类似之物。而万有所生的创造者，据他们说是纷争；反之，爱则管理并供应宇宙，使之享有恒久。它将宇宙中分散的部分聚拢到统一中，并带领它们脱离（各自）存在方式，将之结合并添加到宇宙中，使之享有恒久；由此构成一个体系。因此，它们不会停止——不论是纷争分裂世界，还是爱将分裂部分附着于世界。依此看来，毕达哥拉斯关于世界的分布似乎就是这样。但毕达哥拉斯说星辰是太阳的碎片，动物的灵魂从星辰传来；当它们在身体内时，如同被埋葬在坟墓中一样是有死的；而当我们脱离身体时，它们（从这世界）升起并成为不朽的。因此，有人问柏拉图\Quote{什么是哲学？}时，他回答说：\Quote{哲学是灵魂与身体的分离。}\par

% structure: p0060 source=h2 target=subsection reason=relative-heading-rank; base=h2

\subsection{第二十一章　毕达哥拉斯的其他意见。}

毕达哥拉斯同样成为这些学说的学生，他以谜语及诸如以下的言辞说话：\Quote{当你离开自己的（帐幕）时，不要返回；但是，如果你不（这样做），作为正义辅助者的复仇女神将追上你，}——称身体为自己的（帐幕），称其情欲为复仇女神。因此，他说，当你离开时，即当你离开身体时，不要恳切渴求此事；但若你急切渴求（离开），情欲将再次把你限制在身体内。因为他们认为灵魂从一个身体转移到另一个身体，正如恩培多克勒采纳毕达哥拉斯的原则所肯定的那样。因为，他说，正如柏拉图所述，那些喜爱快乐的灵魂，如果当它们处于人类所经历的痛苦状态时没有形成哲学理论，就必须通过所有动物和植物（再）回到一个人类身体。而当（灵魂）在同一个身体中三次形成思辨体系时，（他主张）它上升到某种同源星辰的本性。然而，如果（灵魂）不从事哲学，（它必须）再次经历同一（系列变化）。那么，他断言，灵魂有时甚至可能成为可朽的，若它被复仇女神——即（身体的）情欲——所战胜；并且成为不朽的，若它成功地逃脱了复仇女神，即那些情欲。\par

% structure: p0062 source=h2 target=subsection reason=relative-heading-rank; base=h2

\subsection{第二十二章 毕达哥拉斯的\Quote{格言}}

但既然我们也选择提及毕达哥拉斯通过象征模糊地对其门徒所说的格言，似乎同样适宜提醒（读者）其余（的学说。我们触及此话题）也是因为异端首领们，他们试图用某种类似的方法通过象征来交谈；而这些并非他们自己的，而是他们（在提出时）利用了毕达哥拉斯派所用的表达。于是毕达哥拉斯教导他的门徒，对他们这样说：\Quote{捆起盛放被褥的袋子。}（现在，）既然那些打算旅行的人将衣服捆成行囊，以便上路；因此，（同样地，）他希望他的门徒预备好，因为死亡随时可能突然临到他们。（通过这种方式毕达哥拉斯寻求）使（其追随者）在门徒所需素质方面没有缺乏。因此，他习惯在拂晓时教导毕达哥拉斯派互相鼓励捆起盛放被褥的袋子，即为死亡预备好。\Quote{不要用刀拨火；}（意思是，）不要通过对他说话与一个愤怒的人争吵；因为发怒的人像火，而刀是说出的话语。\Quote{不要践踏扫帚；}（意思是，）不要轻视小事。\Quote{不要在屋里种棕榈树；}（意思是，）不要在家庭中煽动不和，因为棕榈树是战斗与屠杀的象征。\Quote{不要从凳子上吃饭；}（意思是，）不要从事卑贱的技艺，以免你成为身体——它是可朽的——的奴隶，而要通过文学谋生。因为滋养身体并使灵魂更美好在你能力范围之内。\Quote{不要从没切开的整条面包上咬一口；}（意思是，）不要减少你的财产，而靠其利润生活，并保护你的财物如同整条面包。不要吃豆子；（意思是，）不要接受城市的统治，因为当时他们惯常用豆子为长官投票。\par

% structure: p0064 source=h2 target=subsection reason=relative-heading-rank; base=h2

\subsection{第二十三章 毕达哥拉斯的天文系统}

毕达哥拉斯派提出这些以及诸如此类的断言；异端者模仿这些，被一些人视为道出重要真理。然而，毕达哥拉斯体系主张，所有所谓存在物的创造者是伟大的几何学家和计算者——一个太阳；并且这一个已被安置在整个世界中，就像在身体中的灵魂一样，按柏拉图的说法。因为太阳（具有火的本质）类似于灵魂，而地球（类似于）身体。而且，若与火分离，便没有可见之物，也没有可触之物而没有某种坚固的东西；但没有任何坚固的物体没有地球。因此，神性将空气置于中间，用火和地球塑造了宇宙的身体。并且，他说，太阳以下列方式计算和几何地度量世界：世界是一个可被感官认知的统一体；关于这个（世界）我们现在作这些陈述。但一个精通数字科学且为几何学家者已将其分为十二部分。这些部分的名称为：白羊、金牛、双子、巨蟹、狮子、室女、天秤、天蝎、射手、摩羯、宝瓶、双鱼。他又将十二部分中的每一部分分为三十部分，这些是月份的日子。他又将三十部分中的每一部分分为六十小份，并将这些小（份）再细分为更小的份，（这些再分为）更小的份。并且总是这样做，不间断，而是从这些被分的份中收集（总量），并构成一年；又再次分解与复合，（太阳）完全完成了伟大而永恒的世界。\par

% structure: p0066 source=h2 target=subsection reason=relative-heading-rank; base=h2

\subsection{第二十四章 瓦伦提尼被证实剽窃柏拉图与毕达哥拉斯哲学；瓦伦提尼派相偶发散的学说}

我曾精确研究他们的体系并试图扼要陈述的，就是\Person{毕达哥拉斯}和\Person{柏拉图}的意见。我们已经证明，\Person{瓦伦提努斯}正是从这个体系（而不是从福音）收集了异端的素材——即他自己的异端——故可公正地被算作毕达哥拉斯派和柏拉图派，而非基督徒。因此，\Person{瓦伦提努斯}、\Person{赫拉克利翁}、\Person{托勒密}以及这些（异端者）的整个学派，作为\Person{毕达哥拉斯}和\Person{柏拉图}的弟子，追随这些导师，将算术体系定为其教义的基本原则。同样，按照这些（瓦伦提努斯派），宇宙的起源因由是一个\OriginalTerm{元一}{Monad}，非受生、不朽坏、不可理解、不可思虑、能生产，并是一切存在物生成的原因。上述元一被他们称为父。然而，在他们中间可以发现有相当大的意见分歧。有些人为了使\Person{瓦伦提努斯}的毕达哥拉斯学说完全不受（其他教条）混杂，假设父是无女性、未婚配、独居的。但另一些人，认为从单一男性根本不可能产生任何存在物的生成，所以必然在宇宙之父旁边加上\OriginalTerm{沉默}{Sige}作为配偶，以便他成为父。至于沉默是否在任何时候与（父）结合，这一点我们留给他们自己去争论。我们目前，遵守毕达哥拉斯的原则——即单一、未婚配、无女性、毫无所缺——将继续按照他们自己所教诲的来叙述他们的教义。\Person{瓦伦提努斯}说：没有任何东西是受生的，只有父是非受生的，不受地方条件限制，不受时间条件限制，没有谋士，也不是按照通常感知方式所能实现的任何其他实体。然而，（父）是独居的，如同他们所说，存在于一种静止状态，他自己在他内隐退安息。但是，当他变得能生产时，他认为在某个时候生成并带出他自己内所拥有的最美丽最完美的（存在种子）是合适的，因为（父）并不喜欢孤独。因为，他说，他是全部的爱，但爱若不是有被爱的对象，就不是爱。父自己既然是独居的，就投影并产生\OriginalTerm{理智}{Nous}和\OriginalTerm{真理}{Aletheia}，即一个\OriginalTerm{二元}{Duad}，这二元成为所有永世的女主人、起源和母亲，这些永世被他们算作存在于\OriginalTerm{圆满}{Pleroma}之内。理智和真理从父投影出来，一个是能继续生殖的，从能生产的存有获得存在；理智自己也模仿父，投影出\OriginalTerm{逻各斯}{Logos}和\OriginalTerm{生命}{Zoe}；而逻各斯和生命投影出\OriginalTerm{人}{Anthropos}和\OriginalTerm{教会}{Ecclesia}。但当理智和真理看到他们自己的后代已是能生产的，就向宇宙之父献上感恩，并献上他一个完美的数目，即十个\OriginalTerm{永世}{Aeons}。因为，他说，理智和真理不能献给父比这个数目更完美的（东西）。因为父是完美的，理当由完美的数目来称颂，而十是完美的数目，因为十是由多数形成的第一个数目，因此完美。然而，父更完美，因为唯独他是非受生的，借着理智和真理那最初的结合，找到了投影一切存有之根的方法。\par

% structure: p0068 source=h2 target=subsection reason=relative-heading-rank; base=h2

\subsection{第二十五章　二元论作为\Person{瓦伦提努斯}永世流溢体系的基础}

\OriginalTerm{圣言}{Logos}自己与\OriginalTerm{生命}{Zoe}，看见\OriginalTerm{心灵}{Nous}和\OriginalTerm{真理}{Aletheia}曾以一个完美的数字庆祝了宇宙之父；\OriginalTerm{圣言}{Logos}自己也与\OriginalTerm{生命}{Zoe}想要尊崇他们自己的父母，即\OriginalTerm{心灵}{Nous}和\OriginalTerm{真理}{Aletheia}。然而，由于\OriginalTerm{心灵}{Nous}和\OriginalTerm{真理}{Aletheia}是受生的，并不拥有父亲的（与）完美的非受生性，\OriginalTerm{圣言}{Logos}和\OriginalTerm{生命}{Zoe}并不以完美的数字尊崇他们的父亲\OriginalTerm{心灵}{Nous}，而是远非如此，以不完美的数字。因为\OriginalTerm{圣言}{Logos}和\OriginalTerm{生命}{Zoe}向\OriginalTerm{心灵}{Nous}和\OriginalTerm{真理}{Aletheia}奉献十二个\OriginalTerm{永世}{Aeon}。因为，照\Person{瓦伦廷}所说，即\OriginalTerm{心灵}{Nous}和\OriginalTerm{真理}{Aletheia}、\OriginalTerm{圣言}{Logos}和\OriginalTerm{生命}{Zoe}、\OriginalTerm{人}{Anthropos}和\OriginalTerm{教会}{Ecclesia}，这些乃是\OriginalTerm{永世}{Aeon}的首要根源。但从\OriginalTerm{心灵}{Nous}和\OriginalTerm{真理}{Aletheia}出来十个\OriginalTerm{永世}{Aeon}，从\OriginalTerm{圣言}{Logos}和\OriginalTerm{生命}{Zoe}出来十二个，总共二十八个。他们给这些（十个）起以下名称：\OriginalTerm{深渊}{Bythus}与\OriginalTerm{混合}{Mixis}、\OriginalTerm{不老}{Ageratus}与\OriginalTerm{合一}{Henosis}、\OriginalTerm{自生}{Autophyes}与\OriginalTerm{快乐}{Hedone}、\OriginalTerm{不动}{Acinetus}与\OriginalTerm{混合}{Syncrasis}、\OriginalTerm{独生}{Monogenes}与\OriginalTerm{福乐}{Macaria}。这些是十个\OriginalTerm{永世}{Aeon}，有些人说是由\OriginalTerm{心灵}{Nous}和\OriginalTerm{真理}{Aletheia}所发出，但有些人说是由\OriginalTerm{圣言}{Logos}和\OriginalTerm{生命}{Zoe}。然而，另有些人断言（这些）十二个（\OriginalTerm{永世}{Aeon}）是由\OriginalTerm{人}{Anthropos}和\OriginalTerm{教会}{Ecclesia}所发出，还有些人说是由\OriginalTerm{圣言}{Logos}和\OriginalTerm{生命}{Zoe}。他们给这些起以下名字：\OriginalTerm{护慰者}{Paracletus}与\OriginalTerm{信德}{Pistis}、\OriginalTerm{父性}{Patricus}与\OriginalTerm{望德}{Elpis}、\OriginalTerm{母性}{Metricus}与\OriginalTerm{爱德}{Agape}、\OriginalTerm{永恒}{Aeinous}与\OriginalTerm{理解}{Synesis}、\OriginalTerm{教会}{Ecclesiasticus}与\OriginalTerm{福乐}{Macariotes}、\OriginalTerm{意志}{Theletus}与\OriginalTerm{智慧}{Sophia}。但在十二个当中，所有二十八个\OriginalTerm{永世}{Aeon}中第十二个也是最年幼的，是一个女性，名叫\OriginalTerm{智慧}{Sophia}，她观察到周围\OriginalTerm{永世}{Aeon}的数量和能力，就匆忙退回到圣父的深渊中。她察觉所有其余的\OriginalTerm{永世}{Aeon}，作为受生的，都通过配偶交合而生育；另一方面，圣父独自一圣，不经交合，就产生了（后裔）。她想要效法圣父，没有婚姻伴侣而凭自己产生（后裔），以便成就一项绝不次于圣父的工程。然而，（\OriginalTerm{智慧}{Sophia}）不知道，\OriginalTerm{非受生者}{Unbegotten One}作为宇宙的起源原则，同时也是根、深处和深渊，唯独拥有自我生成的能力。但\OriginalTerm{智慧}{Sophia}是受生的，在许多（\OriginalTerm{永世}{Aeon}）之后才出生，不能获得\OriginalTerm{非受生者}{Unbegotten One}所固有的能力。因为在\OriginalTerm{非受生者}{Unbegotten One}中，他说，万有同时存在，但在受生的（\OriginalTerm{永世}{Aeon}）中，女性投射实体，男性形成女性所投射的实体。因此，\OriginalTerm{智慧}{Sophia}准备好只投射她所能（投射）的，即一个无形式且未成形的实体。而这，他说，就是\Person{梅瑟}所断言的：\BibleQuote{大地还是混沌空虚。}\BibleRef{创 1:2}这（实体），他说，就是那善的（与）天上的耶路撒冷，天主应许要领以色列子孙进入其中，说：\BibleQuote{我要领你们进入流奶流蜜的地方。}\BibleRef{出 3:8}\par

% structure: p0070 source=h2 target=subsection reason=relative-heading-rank; base=h2

\subsection{第二十六章 瓦伦廷关于基督与圣神存在的解释。}

在\OriginalTerm{圆满}{Pleroma}内因\OriginalTerm{智慧}{Sophia}而产生无知，因\OriginalTerm{智慧}{Sophia}的后裔而产生无形状态，\OriginalTerm{圆满}{Pleroma}内便产生了混乱。（因为所有）受生的\OriginalTerm{永世}{Aeon}（变得充满忧虑，想象）同样无形不完全的后裔会被产生出来；并且某种毁灭，不久之后，最终会抓住\OriginalTerm{永世}{Aeon}。于是，所有\OriginalTerm{永世}{Aeon}都向圣父祈求，求他安抚悲伤的\OriginalTerm{智慧}{Sophia}；因为她继续哭泣哀伤，因她所产生的流产——因为他们如此称呼。于是圣父怜悯\OriginalTerm{智慧}{Sophia}的眼泪，且接纳了\OriginalTerm{永世}{Aeon}的祈求，下令进一步发出。因为他（瓦伦廷）说，圣父并非亲自发出，而是\OriginalTerm{心灵}{Nous}和\OriginalTerm{真理}{Aletheia}（发出）\OriginalTerm{基督}{Christ}和\OriginalTerm{圣神}{Holy Spirit}，为恢复形式、消灭流产，并为安慰和停止\OriginalTerm{智慧}{Sophia}的呻吟。于是连同\OriginalTerm{基督}{Christ}和\OriginalTerm{圣神}{Holy Spirit}，三十个\OriginalTerm{永世}{Aeon}出现了。有些（这些瓦伦廷派）希望这应是\OriginalTerm{永世}{Aeon}的一个\OriginalTerm{三十元体}{triacontad}，而另一些则希望\OriginalTerm{沉默}{Sige}应与圣父一起存在，且\OriginalTerm{永世}{Aeon}应连同它们一同被计数。\par

因此，\OriginalTerm{基督}{Christ}和\OriginalTerm{圣神}{Holy Spirit}由\OriginalTerm{心灵}{Nous}和\OriginalTerm{真理}{Aletheia}发出之后，这\OriginalTerm{智慧}{Sophia}的流产立即——无形式的、仅由她自己所生、未经交合而产出的——从所有\OriginalTerm{永世}{Aeon}中分离出来，免得完美的\OriginalTerm{永世}{Aeon}看见这（流产）因它的无形而受到干扰。于是，为使流产的无形完全不显现在完美的\OriginalTerm{永世}{Aeon}面前，圣父又进一步额外发出一个\OriginalTerm{永世}{Aeon}，即\OriginalTerm{十字架}{Staurus}。他由一位伟大而完美的圣父所生，且为大，为保卫和守护\OriginalTerm{永世}{Aeon}而被发出，成为\OriginalTerm{圆满}{Pleroma}的界限，其中包含所有三十个\OriginalTerm{永世}{Aeon}，因为这些就是已被发出的。现在这个（\OriginalTerm{永世}{Aeon}）被称为\OriginalTerm{界限}{Horos}，因为它从\OriginalTerm{圆满}{Pleroma}中分离出外面的\OriginalTerm{欠缺}{Hysterema}。它又被称为\OriginalTerm{分享者}{Metocheus}，因为它也分有那\OriginalTerm{欠缺}{Hysterema}。它还被命名为\OriginalTerm{十字架}{Staurus}，因为它被固定得坚定不可动摇，使得\OriginalTerm{欠缺}{Hysterema}的任何部分都不能靠近\OriginalTerm{圆满}{Pleroma}内的\OriginalTerm{永世}{Aeon}。因此，在\OriginalTerm{界限}{Horos}（或\OriginalTerm{分享者}{Metocheus}，或\OriginalTerm{十字架}{Staurus}）之外，是所谓的\OriginalTerm{八元体}{Ogdoad}，照他们所说，就是那在\OriginalTerm{圆满}{Pleroma}之外的\OriginalTerm{智慧}{Sophia}，基督——由\OriginalTerm{心灵}{Nous}和\OriginalTerm{真理}{Aletheia}额外发出——塑造她并使之成为一个完美的\OriginalTerm{永世}{Aeon}，使她在能力上绝不亚于\OriginalTerm{圆满}{Pleroma}内的任何一个\OriginalTerm{永世}{Aeon}。然而，因为\OriginalTerm{智慧}{Sophia}在外面被塑造，且从\OriginalTerm{心灵}{Nous}和\OriginalTerm{真理}{Aletheia}发出的\OriginalTerm{基督}{Christ}和\OriginalTerm{圣神}{Holy Spirit}不可能也不适宜留在\OriginalTerm{圆满}{Pleroma}之外，基督和\OriginalTerm{圣神}{Holy Spirit}急忙离开那已获得形状的\OriginalTerm{智慧}{Sophia}，回到界限内的\OriginalTerm{心灵}{Nous}和\OriginalTerm{真理}{Aletheia}那里，以便与其余的\OriginalTerm{永世}{Aeon}一同光荣圣父。\par

% structure: p0073 source=h2 target=subsection reason=relative-heading-rank; base=h2

\subsection{第二十七章 瓦伦廷关于耶稣存在的解释；耶稣对人类的能力。}

于是，在普累若麻内所有移涌之间有了某种和平与和谐的协议之后，他们觉得不仅通过婚姻结合尊崇了圣子，而且也应该以成熟果实的奉献来光荣圣父。于是所有三十位移涌同意投射一位移涌，即普累若麻的共果，作为他们团结、一致与和平的保证。他是唯一由全体移涌为纪念圣父而投射的。在他们中间，这位被称为\Quote{普累若麻的共果}。这些事就这样发生在普累若麻之内。那\Quote{普累若麻的共果}被投射出来，即耶稣——因为这是他的名字——伟大的大司祭。然而，在普累若麻之外的索菲亚，寻找那赋予她形式的基督和圣神时，陷入了极大的恐惧，恐怕那赋予她形式和存有者会离开她。她陷入悲伤，并处于相当困惑的状态，思考谁赋予了她形式，圣神是什么，祂往何处去了，是谁阻止了它们同在，是谁嫉妒了那荣耀而蒙福的景象。陷入如此苦难时，她转而向那遗弃她者祈祷和恳求。在她祈求期间，普累若麻内的基督怜悯了她，所有其余的移涌也同样怜悯；他们派\Quote{普累若麻的共果}超出普累若麻，作为在外面的索菲亚的配偶，并作为她寻找基督时所受苦难的纠正者。\par

那\Quote{果实}到达普累若麻之外，发现索菲亚处于那四种原始情感之中，即恐惧、悲伤、困惑和恳求，他纠正了她的情感。然而，在纠正时，他观察到既不宜毁灭这些情感，因为它们本质上是永恒的，且为索菲亚所特有；也不宜让索菲亚处于如此情感之中，即恐惧、悲伤、恳求和困惑。因此，作为如此伟大的移涌和整个普累若麻的后裔，他使情感离开她，并使它们成为实体性的本质。他把恐惧改变为动物性欲望，把悲伤改变为物质，把困惑改变为魔鬼的情感。而悔改、恳求和祈求，他确立为通往悔改的道路和对动物性本质的权能，这动物性本质被称为右。创造者来自恐惧；这正是，他说，圣经所肯定的：\Quote{敬畏主是智慧的开端。}因为这是索菲亚情感的开始，她先被恐惧抓住，接着是悲伤，然后是困惑，于是她投靠恳求和祈求。他说，动物性本质是火性的，也被他们称为超天穹、七重天和\Quote{万古者}。他们所提出的关于这位移涌的任何其他此类陈述，他们都声称适用于那动物性本质者，他们断言他是世界的创造者。他是火的形象。梅瑟，他说，也这样表述：\Quote{你的天主是燃烧的烈火。}因为他也希望如此记载。然而，他说，火有一种双重能力；因为火是吞噬一切的，不能熄灭。因此，根据这种区分，存在着一种有死的灵魂，它作为某种中介，因为它是七重天和安息。\BibleRef{创 2:2}因为在八重天（索菲亚所在）之下，但在物质（即创造者）之上，形成了一天，以及\Quote{普累若麻的共果}。如果灵魂被塑造成上面者（即八重天）的形象，它就变成不朽的，并归向八重天，他说，那是天上的耶路撒冷。如果它被塑造成物质的形象，即肉体的情感，那么灵魂是可朽坏的，并因而被毁灭。\par

% structure: p0076 source=h2 target=subsection reason=relative-heading-rank; base=h2

\subsection{第二十八章 瓦伦廷论创造的起源}

因此，正如动物性本质的第一和最大的能力成为存在，是（独生子的）肖像；同样，魔鬼——这世界的统治者——构成物质本质的能力，正如贝尔则步是出于焦虑的魔鬼本质的能力。因此，上面的索菲亚从八重天向七重天发出她的能力。因为德穆革，他们说，一无所知，按照他们，是缺乏悟性的、愚笨的，并且不知道自己在做什么或工作。但当他处于这种无知状态，甚至不知道自己正在创造时，索菲亚在他身上施加了各种能力，注入了力量。虽然索菲亚实际上是运行的原因，但他自己想象是从自己中演化出世界的创造；因此他开始说：\Quote{我是天主，在我以外没有别的神。}\par

% structure: p0078 source=h2 target=subsection reason=relative-heading-rank; base=h2

\subsection{第二十九章 其他瓦伦廷流溢符合毕达哥拉斯数字体系}

因此，瓦伦廷所倡导的四元组是\Quote{具有根源的永恒本性之源}；索菲亚是（能力），动物性和物质性创造从中获得了其现有状态。索菲亚被称为\Quote{灵}，德穆革被称为\Quote{灵魂}，魔鬼被称为\Quote{这世界的统治者}，贝尔则步被称为\Quote{魔鬼的（统治者）}。这是他们提出的陈述。但除此之外，正如我先前所提，将他们的整个教义体系等同于算术技艺，他们规定，普累若麻内的三十位移涌除了这些之外，又按照（毕达哥拉斯派采用的）数字比例投射了其他移涌，以便普累若麻得以按照一个完全数形成一个整体。因为毕达哥拉斯派如何将（天球）分为十二、三十和六十部分，以及他们如何有微小部分的细分，已经显明。\par

如此，这些（瓦伦廷的追随者）细分了普累若麻内部的部分。同样，八元体中的部分也被细分，并产生了索菲亚——按照他们的说法，她是所有生物的母，以及\Quote{普累若麻的联合果实}（即）逻各斯，（以及其他永世，）他们是天上的天使，他们的籍贯在上面的耶路撒冷，就是天上的耶路撒冷。因为这里的耶路撒冷就是索菲亚，她（在普累若麻）之外，她的配偶是\Quote{普累若麻的联合果实}。德缪哥产生了灵魂；因为（这索菲亚）是灵魂的本质。按照他们的说法，这（德缪哥）就是亚巴郎，而（这些灵魂）是亚巴郎的子女。德缪哥从物质的和魔鬼的本质为灵魂塑造了身体。这正是经上所记载的：\Quote{天主用地上的尘土造了人，把生命的气息吹在他的脸上，人就成了一个活灵。} \BibleRef{创 2:7} 按照他们的说法，这就是内在的人，属魂的人，居于物质的身体中：而物质的人是可朽坏的、不完全的，由魔鬼的本质形成。这就是物质的人，按照他们的说法，就像一个客栈或住所，有时只住灵魂，有时住灵魂和魔鬼，有时住灵魂和逻各斯。这些逻各斯是从上边，从\Quote{普累若麻的联合果实}和对世界的索菲亚那里分散下来的。他们带着灵魂居住在属地的身体里，当魔鬼不与那灵魂同住时。他说，这就是经上所写的：\Quote{因此，我向我们的主耶稣基督的天主和父屈膝，恳求天主赐给你们基督居住在内在的人里} \BibleRef{弗 3:14} ——即属魂的人，不是属肉体的人—— \Quote{使你们能够理解什么是深奥}，就是宇宙之父，\Quote{什么是广阔}，就是斯塔乌罗斯，普累若麻的界线，\Quote{什么是长度}，就是永世的普累若麻。因此，他说，属魂的人不接受天主圣神的事，因为这对他是愚妄；\BibleRef{格前 2:14} 但愚妄，他说，是德缪哥的能力，因为他愚昧无知，以为自己是在创造世界。他却不知道，索菲亚，母亲，八元体，才是他一切作为的真正原因，而他对自己在世界创造中的角色毫无意识。\par

% structure: p0081 source=h2 target=subsection reason=relative-heading-rank; base=h2

\subsection{第三十章。瓦伦廷对耶稣诞生的解释；关于耶稣身体本质的两种学说；意大利学派的意见，即赫拉克里昂和托勒密；东方学派的意见，即阿克西奥尼库和巴尔德萨内斯。}

因此，所有的先知和法律都是借着德缪哥说话的——他说，他是一个愚昧的神，而（先知们）也是愚昧的人，什么都不知道。为此，他说，救主说：\Quote{凡在我以前来的，都是贼和强盗。} \BibleRef{若 10:8} 而宗徒（说）：\Quote{这奥秘在以前的世代没有让人知道。} \BibleRef{哥 1:26} 因为，他说，众先知没有说过任何关于我们所谈论的事；因为（先知）不可能不知道这一切，因为这些事肯定只是由德缪哥说出的。因此，当创造完成，并且在那之后应当有天主子女的启示——即德缪哥，这启示一直隐藏着，在这隐藏中属魂的人被遮盖，心上蒙着一层帕子——当时机到来，帕子应当被除去，这些奥秘应当被看见，耶稣由童贞玛利亚诞生，正如经上所言：\Quote{圣神要临到你身上}——索菲亚就是圣神——\Quote{至高者的能力要庇荫你}——至高者就是德缪哥——\Quote{因此，那要诞生的圣者，将称为圣者。} \BibleRef{路 1:35} 因为他不是独自从至高者而生，如同那些在（像）亚当里受造的人那样独自从至高者而生——即来自索菲亚和德缪哥。然而耶稣，新人，（是由）圣神——即索菲亚和德缪哥——而生的，以便德缪哥可以完成他身体的构型和组织，圣神可以供给他的本质，并且一个天上的逻各斯可以从八元体出来，由玛利亚所生。\par

关于这（逻各斯），他们中间有一个大问题——既是分裂也是争端的起因。因此他们的学说分裂了：按照他们的说法，一种学说称为东方学派，另一种称为意大利学派。来自意大利的人，其中有赫拉克里昂和托勒密，说耶稣的身体是属魂的。因此，他们主张，在耶稣受洗时，圣神如鸽子降临——即来自上面母亲（我指索菲亚）的逻各斯——成为（声音）对属魂的人，并使他从死人中复活。他说，这正是经上所记：\Quote{那使基督从死者中复活的，也必使你们有死的和属魂的身体活过来。} \BibleRef{罗 8:11} 因为尘土受了诅咒；\Quote{因为，}他说，\Quote{你是尘土，你还要归于尘土。} \BibleRef{创 3:19} 另一方面，东方学派，其中有阿克西奥尼库和巴尔德萨内斯，声称救主的身体是属神的；因为圣神——即索菲亚和至高者的能力——临到了玛利亚身上。这是创造的艺术，被赐予是为了使由圣神赐予玛利亚的东西得以成形。\par

% structure: p0084 source=h2 target=subsection reason=relative-heading-rank; base=h2

\subsection{第三十一章。瓦伦廷关于永世的进一步学说；道成肉身的原因。}

让，那么，那些（异端者）在他们当中继续这些探究吧，（也让其他人如此），如果有谁愿意探究（这些要点）。然而，\Person{瓦伦廷}补充了以下陈述：属于圆满内\OriginalTerm{永世}{Aeons}的过程已经得到纠正；同样，属于圆满外的\OriginalTerm{奥格多}{Ogdoad}（即\Person{索菲亚}）的过程也已经得到纠正；还有属于\OriginalTerm{赫布多美}{Hebdomad}（的过程）也得到了纠正。因为\Person{德穆革}已经从\Person{索菲亚}那里学到，他自己并非如他所想象的那样是唯一的天主，而在他自己之外并没有另一个（神祇）。但当他被\Person{索菲亚}教导后，他得以认识那更高的（神祇）。因为他受她指导，并被引入和传授了\Person{圣父}和永世的伟大奥秘，并对此保持缄默。正如他所说，这就是（天主）对\Person{梅瑟}所说的：\BibleQuote{我是\Person{亚巴郎}的天主，\Person{依撒格}的天主，\Person{雅各伯}的天主；我没有向他们宣布我的名字；}也就是说，我没有宣告那奥秘，也没有解释我是谁的天主，而是从我由\Person{索菲亚}那里听到的奥秘秘密地保留在我自己之内。那么，当上面那些的过程被纠正后，根据同样的结果，这里的（过犯）同样需要得到纠正。为此，\Person{耶稣}\Person{救主}由\Person{玛利亚}所生，为要纠正这里的（过犯）；正如那位由\Person{努斯}和\OriginalTerm{真理}{Aletheia}从上面额外流出的\Person{基督}，已经纠正了\Person{索菲亚}的情欲——即圆满外的流产（者）。同样，由\Person{玛利亚}所生的\Person{救主}前来纠正灵魂的情欲。因此，根据这些（异端者），共有三位\Person{基督}：第一位是由\Person{努斯}和真理与\Person{圣神}一起额外流出的；第二位是\Quote{圆满的联合果实}，是那在圆满外的\Person{索菲亚}的配偶。而她本身也被称为\Person{圣神}，但次于第一位（流出者）。第三位（\Person{基督}）是为了修复我们这个世界的由\Person{玛利亚}所生的那位。\par

% structure: p0086 source=h2 target=subsection reason=relative-heading-rank; base=h2

\subsection{第三十二章 \Person{瓦伦廷}被控剽窃\Person{柏拉图}。}

我认为，\Person{瓦伦廷}的异端源于\Person{毕达哥拉斯}（主义），已经得到充分、甚至过于充分的描述。因此，在解释了他的观点之后，停止（对其体系的）进一步反驳似乎也是合适的。然后，\Person{柏拉图}在阐述有关宇宙的奥秘时，给\Person{狄奥尼修斯}写信，以类似如下的方式表达自己：\Quote{我必须用谜语对你说话，以便这封信如果在海上或陆地上遭遇不测，读信的人也不能理解。因为事情就是这样。万物都是关于万有之王的，并且因他而存在，他是所有荣耀（受造物）的原因。第二是关于第二的，第三是关于第三的。但关于王，我所说的这些没有一个与他有关。但在这之后，灵魂热切地想要学习这些是什么，看着那些与自己相似的事物，而没有一个（事物）是足够的。这是\Person{狄奥尼修斯}和\Person{多莉丝}的儿子啊，你的问题是一切恶的原因。不，毋宁说对此的焦虑是灵魂内在的；若有人不除去它，他永远不会真正获得真理。但关于此事令人惊异之处，请听。因为有些人已经听到了这些事——（这些人）具备学习的能力，具备记忆的能力，并且完全以各种方式禀赋一种以推理为目的的探究才能。（这些人）是年迈的思辨者。他们断言，一度可信的意见现在变得不可信，而一度不可信的意见现在变得相反。因此，当你将审视的目光转向这些（探究）时，要谨慎，免得你有时会因为那些事以有损你尊严的方式发生而后悔。为此，我没有写过任何关于这些（论点）的东西；也没有\Person{柏拉图}的任何论文（关于它们），也永远不会有了。然而，现在所作的这些观察是\Person{苏格拉底}的，他在年轻时就已经以美德著称。}\par

\Person{瓦伦廷}附和了这些（言论），在他的体系中设定了一个基本原则，即\Person{柏拉图}提到的\Quote{万有之王}，这个异端者称之为\Person{圣父}、\OriginalTerm{深渊}{Bythos}、\OriginalTerm{超越万有之始}{Proarche}，位居其他永世之上。当\Person{柏拉图}使用\Quote{关于第二事物的第二}这些词时，\Person{瓦伦廷}认为所有在界限（内的永世）以及（界限本身）都是第二。当\Person{柏拉图}使用\Quote{关于第三事物的第三}这些词时，他（将）整个在界限和圆满之外的安排（视为）第三。\Person{瓦伦廷}非常简洁地在一首赞美诗中阐明了这一点，这首诗从下面开始（不像\Person{柏拉图}从上面开始），如此表达自己：\Quote{我看见万物被灵悬在空中，我感知万物被灵飘动；肉体（我看见）悬于灵魂，但灵魂从空中照耀，空气源于以太，果实源于\OriginalTerm{深渊}{Bythus}，胎儿从子宫生出。}就这样，\Person{瓦伦廷}形成了关于这类（论点）的意见。根据这些（异端者），肉体是悬于\Person{德穆革}灵魂的物质。灵魂从空中照耀；也就是说，\Person{德穆革}从灵（即圆满外的灵）中出现。但空气源于以太；也就是说，在界限内的圆满和整个圆满（一般地）流出了那在圆满外的\Person{索菲亚}。果实源于深渊；也就是说，永世的整个流出都来自\Person{圣父}。因此，\Person{瓦伦廷}提出的意见已经充分阐述。剩下的是要解释那些出自他学派的人的观点，尽管每个追随者（信奉\Person{瓦伦廷}者）都持有不同的意见。\par

% structure: p0089 source=h2 target=subsection reason=relative-heading-rank; base=h2

\subsection{第三十三章 \Person{塞孔都斯}的永世体系；\Person{厄丕法乃}；\Person{托勒密}。}

有一位（异端者）\Person{塞孔都斯}，与\Person{托勒密}大约同时代，这样表述：（他说）存在着一个右四元体和一个左四元体——即光明与黑暗。并且他断言，那退隐并在匮乏中劳苦的能力，并非出自三十个永世，而是出自这些永世的果实。然而另一位（异端者）——他们中的导师\Person{埃皮法内斯}——这样表述：\Quote{最初的本原是无法构想、不可言说、无名的；}并称此为\OriginalTerm{独一性}{Monotes}。并且（他主张）与此（本原）并存有一种能力，他称之为\OriginalTerm{单一性}{Henotes}。这个单一性和这个独一性，并非通过（自身）流溢，而发出了一种（掌管）所有可理智者的本原；（这既是）非受生的，也是不可见的，他称之为\OriginalTerm{单子}{Monad}。\Quote{与此能力并存一个同本质的能力，我称之为\OriginalTerm{合一}{Unity}。这四种能力发出了永世其余的所有流溢。}但另一些人又以如下名称称呼首要和起源的（且）第四（个）不可见的\OriginalTerm{八元体}{Ogdoad}：第一，\OriginalTerm{太初}{Proarche}；第二，\OriginalTerm{不可思想的}{Anennoetus}；第三，\OriginalTerm{不可言说的}{Arrhetus}；第四，\OriginalTerm{不可见的}{Aoratus}。并且，从第一（即太初），由第一和第五位置流出了\OriginalTerm{元始}{Arche}；从不可思想的，由第二和第六位置流出了\OriginalTerm{不可把握的}{Acataleptus}；从不可言说的，由第三和第七位置流出了\OriginalTerm{无名的}{Anonomastus}；从不可见的，流出了\OriginalTerm{非受生的}{Agennetus}，这是第一个八元体的完成。他们希望这些能力存在于\OriginalTerm{深渊}{Bythus}和\OriginalTerm{沉默}{Sige}之前。然而，关于\OriginalTerm{深渊}{Bythus}本身，有许多不同的意见。有些人断言他是无配偶的，既非男性也非女性；但另一些人（主张）\OriginalTerm{沉默}{Sige}，一个女性，与他同在，并且这构成了第一个婚配结合。\par

但\Person{托勒密}的追随者断言，（深渊）有两个配偶，他们也称之为性情，即\OriginalTerm{意念}{Ennoia}和\OriginalTerm{意愿}{Thelesis}（意念和意志）。因为首先构想出流溢某物的意念；接着，如他们所说，产生了如此做的意愿。因此，这两种性情和能力——即意念和意愿——仿佛彼此混合，便通过婚配结合产生了\OriginalTerm{独生者}{Monogenes}和\OriginalTerm{真理}{Aletheia}。结果是，从不可见的（永世）中，即从\OriginalTerm{意志}{Thelema}产生了\OriginalTerm{心灵}{Nous}，从\OriginalTerm{意念}{Ennoia}产生了\OriginalTerm{真理}{Aletheia}，显出了父的那两种性情的可见肖像和形象。因此，后来产生的意志的肖像是男性；而非受生意念的（肖像）是女性，因为意愿仿佛是意念的一种能力。因为意念总是怀有流溢的想法，但本身至少不能流溢自身，而是怀有（流溢的）想法。然而，当意愿的能力（出现）时，它便流溢了那已被构想出的想法。\par

% structure: p0092 source=h2 target=subsection reason=relative-heading-rank; base=h2

\subsection{第三十四章 马尔谷体系；一个纯粹的骗子；他在圣体圣事杯上的恶毒把戏。}

他们中的另一位导师\Person{马尔谷}，一个精通巫术的人，部分通过手法技巧，部分通过魔鬼，不时欺骗了许多人。这个（异端者）声称，来自不可见、不可名之处的最强大能力居住在他里面。他经常拿起杯，仿佛在献上圣体圣事的祈祷，并将呼求之词延长到比通常更长的长度，使杯中出现紫色、有时是红色的混合物，以致他的受骗者以为某种\OriginalTerm{恩宠}{Grace}降下，并将血红色的能力传递给了饮料。然而，那个恶棍当时成功地逃脱了许多人的察觉；但现在，因（骗局）被揭露，他将被迫停止。因为，他秘密地向混合物中注入某种具有赋予那种颜色的能力的药物，（正如上面所提到的），长时间念诵无意义的词句，他习惯等待，直到（药物）获得足够的水分而溶解，并与饮料混合，赋予其颜色。然而，我们在《\WorkTitle{论魔术师之书}》中已经提到了具有产生这种效果的特性的药物。在此我们借机解释他们如何欺骗许多人，并彻底毁灭他们。如果他们愿意更仔细地关注我们所说的陈述，他们就会意识到\Person{马尔谷}的欺骗。\par

% structure: p0094 source=h2 target=subsection reason=relative-heading-rank; base=h2

\subsection{第三十五章 马尔谷进一步的戏法。}

这位（马谷斯）将（前述的）混合物注入一个较小的爵中，习惯上将它递给一位妇女，由她献上感恩祈祷，而他本人则站在一旁，手中拿着另一个较大的空（爵）。当他的女受骗者念完祝圣词后，他（从她手中）接过（爵），然后将（其中的内容）注入较大的（爵）中，并频繁地将它们从一个爵倒入另一个爵，同时习惯念出以下呼求：\Quote{愿那存在于宇宙之前、不可言喻、不可言说的恩宠，充满你的内在之人，并使你对这（恩宠）的认识丰盈，如同她将芥菜籽撒在好土地上一样。}在念出类似这样的语句的同时，他使他的女受骗者和在场的人惊讶，被视为行奇迹的人；而较大的爵从较小的爵中被充满，以至于（内容物）过满溢出。这（戏法）的诡计我们也已在上述（第四）书中解释过，我们在那里证明了许多药物，当以这种方式与液体物质混合时，具有产生膨胀的特性，尤其是在酒中稀释时。现在，当（这些骗子之一）事先偷偷地用任何一种这类药物涂抹一个空杯，并（向观众）展示它好像什么也没装时，通过将（内容物）从另一个杯注入其中，并再次倒回，药物因其膨胀性质，在与湿物质混合时溶解。其效果是混合物过满，并如此膨胀，以至于注入的物质开始运动，这就是药物的性质。如果有人将（杯）藏起来，当它被充满后，不久便会恢复其自然体积，因为药物的效力因湿度的持续而消失。因此他习惯于急忙将杯递给在场的人喝；他们既恐惧又渴望（品尝杯中的内容），于是开始喝（那混合物），仿佛它是某种神圣的、由天主设计的东西。\par

% structure: p0096 source=h2 target=subsection reason=relative-heading-rank; base=h2

\subsection{第36章 马西特派在圣洗圣事上的异端实践}

这样的和其他（伎俩）是这位骗子试图施行的。因此他被他的受骗者尊崇，有时被认为能发出预言。但有时他试图使别人（发预言），部分是通过运行这些操作的魔鬼，部分是通过施展手法，正如我们先前所说的。因此，他蒙蔽了众多人，引导许多这类成为他门徒的受骗者（进入恶行），教导他们，他们无疑倾向于犯罪，但处于危险之外，因为他们属于完美的力量，并参与了不可言喻的权能。在第一次圣洗之后，他们还向他们应许另一种圣洗，他们称之为救赎。藉着这（另一种圣洗），他们邪恶地颠覆那些留在他们那里期待救赎的人，好像人在受圣洗一次之后，能够再次获得赦免。正是通过这种诡计，他们似乎留住了听众。当他们认为这些人已经受过考验，能够保守（交托给他们的秘密）时，他们就接纳他们到这种（圣洗）中。然而，他们并不满足于仅仅这一点，还应许（他们的信徒）某种其他（恩惠）以坚定他们的希望，以便使他们成为（他们教派的）不可分离的（追随者）。因为当他们按手在领受救赎者身上时，他们用一种不可言喻的（声调）说出一些话。他们声称，除非一个人受过高度考验，或者处于死亡的时刻，（当）主教来并低声（向临终者）耳语时，否则很难（向他人）宣告所说的内容。这诡计（是为了）确保马谷斯的门徒们持续地追随主教，因为他们急切渴望知道最后所说的是什么，藉着（知道）它，学习者将被提升到那些被接纳进入更高奥迹者的行列。关于这些，我保持了沉默，唯恐有人会认为我犯了贬低这些（异端者）的过错。因为这并不属于我们当前工作的范围，除非有助于证明（异端者）是从什么源头获得了他们的立足点，并由此引入他们所提出的见解。\par

% structure: p0098 source=h2 target=subsection reason=relative-heading-rank; base=h2

\subsection{第37章 依勒内对马谷斯体系的解释；马谷斯的异象；华伦提努启示其体系的异象}

因为可敬的长老依勒内，以一种更自由的精神着手进行驳斥的题目，已经解释了这样的洗濯和救赎，并以一种粗略纲要的方式陈述了他们的做法。（似乎有些马谷斯派的人，）在遇到（依勒内的著作）时，否认他们如此领受了（刚才提到的秘密话语），但他们学会了总是应该否认。因此，我们更关切地要更准确地调查，并详细发现他们在第一次圣洗（他们用某种名称称呼它）的情况下，以及在他们称为救赎的第二次的情况下所传授的（教导）。但这个秘密也没有逃过（我们的审视）。然而，对于这些意见，我们同意宽恕华伦提努和他的学派。\par

但是马尔谷，模仿他的老师，也自己假装一个神视，妄想以此方式他被高举。因为瓦伦提诺同样声称他曾看见一个刚出生的婴孩；并询问（这个孩子），他继而探询这可能是谁。（这儿童）回答说，他本人就是圣言，然后添加了一种悲剧式的传说；而瓦伦提诺希望由此构成他所尝试的异端。马尔谷作出与此（异端者）相似的尝试，断言那四元体以一个女人的形像来到他面前——因为世界不能承受，他说，这四元体的男性（形像）——而她启示了她是谁，并向这位（马尔谷）单独解释了宇宙的产生，这是她从未向任何神或人启示过的，以下述方式表达自己：当最初自在的圣父，那位不可揣测且无实体的，祂既非男亦非女，愿意祂自己的不可言说性在某种言说中得以实现，且祂的不可见性在形式中得以实现时，祂开口，并发出一个类似于祂自身的圣言。而这（圣言）站在祂旁边，并向祂启示了祂是谁，即：他自身已显明为那不可见者之（实现）形式。而名字的发音如下所述。他习惯于说出名字本身的第一个词，即\Quote{元始}，而这个词的音节由四个字母组成。然后他加上第二个（音节），这也由四个字母组成。接着他发出第三个（音节），由十个字母组成；他发出第四个（音节），由十二个字母组成。然后整个名字的发音由三十个字母组成，但四个音节。而每个元素有它自己独特的字母、自己独特的形式、自己独特的发音，以及形像和肖像。没有一个是能看到那（字母）的形式的，而这（字母）本身是它的一个元素。当然，他们中没有一个能知道下一个（字母）的发音，而只能按照他自己发音的方式（并且以这种方式），以致在发出整个（词）时，他以为他是在发出整个（名）。因为这些（元素）中的每一个，作为整个（名）的一部分，他按其自身独特的声音称呼，仿佛整个（词）一样。他不停歇地发音，直到他到达最后一个元素的最后一个字母，并以一个单一的发音发出它。然后他说，当所有（元素）降入一个字母，发出一个相同的声音时，整体就恢复了，而且他以为当我们同时说出\Emph{阿们}这个词时，就有发音的肖像。并且这些声音就是那些给予那无实体且非受生的永世以形式的，而那些形式就是主所宣告为天使的——那些不断注视圣父面容的（形式）。\par

% structure: p0101 source=h2 target=subsection reason=relative-heading-rank; base=h2

\subsection{马尔谷的字母体系}

但他称之为元素的普遍和表达的名称是永世、圣言、根源、种子、圆满和果实。（他主张）这些中的每一个，以及每一个所特有的，都被认为包含在\Quote{教会}这个名字中。而最后一个元素的最后一个字母发出了它自己的独特发音。而这个（字母）的声音出来，并按照元素的形像产生它自己独特的元素。从这些，他说，现在的事物得到了装饰，而这些之前的事物也被产生出来。这个字母本身，其声音与下面的声音相伴，他说，被它自己的音节收回到整个（名）的圆满中；但那声音，好像被扔在外面，留在了下面。而那个元素本身，字母连同它自己的发音从它下降到下面，他说，由三十个字母组成，并且这三十个字母中的每一个本身包含其他字母，通过它们，该字母的名称被称呼。并且，其他的（字母）由不同的字母命名，其余由不同的（字母）命名。这样，逐一写出字母，数字将无穷无尽。这样，人们可以更清楚地理解所说的内容。元素德尔塔，（他说，）本身有五个字母，（即，）德尔塔、厄普西隆、兰布达、陶和阿尔法；而这些字母本身是由其他字母（写成的）。因此，如果德尔塔的整个实体趋向无穷，并且不同的字母总是产生不同的字母，并前后相继，那么比那个元素更大的字母的浩瀚海洋有多大呢？如果一个字母这样无限，请看整个名字的字母深度，马尔谷的耐心劳作——不，我应说徒劳——希望这（宇宙的）原始圣父由之构成！因此（他也主张）那圣父，知道祂与自身不可分离，将这（深度）赐给了元素，他同样称之为永世。并且对每一个元素大声说出它自己独特的发音，因为没有人能发出整个（名）。\par

% structure: p0103 source=h2 target=subsection reason=relative-heading-rank; base=h2

\subsection{四元体展示\Quote{真理}}

（\Person{马尔谷}断言）四元体在解释这些事后，如此发言：\Quote{现在我也愿意向你们展示真理本身，因为我已将她从高天居所带下，好使你们能看见她赤露的形体，并认识她的美丽；不仅如此，也使你们能听见她发言，并惊叹她的智慧。那么，请观察，}四元体说，\Quote{首先，头部上方是\GreekText{α}（长）\GreekText{Ω}；颈项是\GreekText{Β}与\GreekText{Ρ}［\GreekText{σι}］；肩膀连同手臂是\GreekText{Γ}与\GreekText{Χ}［\GreekText{χι}］；胸部是\GreekText{Δ}与\GreekText{Π}［\GreekText{πι}］；横膈膜是\GreekText{ΕΥ}；腹部是\GreekText{Ζ}与\GreekText{Τ}；\Emph{私处}是\GreekText{Η}与\GreekText{Σ}；大腿是\GreekText{Θ}与\GreekText{Ρ}；膝盖是\GreekText{ΙΠ}；小腿是\GreekText{ΚΟ}；脚踝是\GreekText{ΛΞ}［\GreekText{ξι}］；双足是\GreekText{Μ}与\GreekText{Ν}。}这就是真理的身体，根据\Person{马尔谷}。这是元素的外形；这是字母的特征。他也称这个元素为人，并断言它是每个言语的源头、每个声音的起源原则、一切不可言说之物的言语实现，以及静默沉默的口。而这就是（真理）本身的身体。但你们应当高举理智的理解力量，从真理的口中聆听那自我生成且是起源者的圣言。\par

% structure: p0105 source=h2 target=subsection reason=relative-heading-rank; base=h2

\subsection{第四十章 基督耶稣的名字}

但在说出这些话之后，（\Person{马尔谷}详述道）真理注视着他，张开她的口，说出了（刚才提及的）话语。而（他告诉我们）那话语成了一个名字，而那名字就是我们知晓并呼号的，即基督耶稣；并且她一说出这（名字）就沉默了。然而，当\Person{马尔谷}正期待她还要说更多时，四元体再次走到中间，发言如下：你将你从真理之口听到的这话语视为可鄙的。然而，这个你知晓并似乎早已拥有的并不是名字；因为你仅仅有其声音，却不知道其能力。因为耶稣是一个非凡的名字，由六个字母组成，被所有属于（基督的）蒙召者所呼号；而另一个（名字，即基督）则由许多部分构成，并居于圆满之（五）永世之中。这（名字）具有另一种形态和不同的类型，并被与他同源的那些存有所认识，这些存有的伟大也与他一同恒久存在。\par

% structure: p0107 source=h2 target=subsection reason=relative-heading-rank; base=h2

\subsection{第四十一章 马尔谷对字母的神秘解释}

所以，（你）要知道，在你们那里被算作二十四个的这些字母，是来自三种能力的流溢，并且是那些甚至包含全部元素之数的（能力的）代表。因为假设有些字母是哑音——它们有九个——来自\OriginalTerm{帕特尔}{Pater}和\OriginalTerm{阿勒忒亚}{Aletheia}，因为它们是哑的——即不可言说、不可表达的。又（假设）另一些是半元音——它们有八个——来自\OriginalTerm{逻各斯}{Logos}和\OriginalTerm{佐厄}{Zoe}，因为它们是辅音和元音之间的中介，接受上面字母的流溢，但产生下面字母的回流。又（同样认定）有些是元音——它们有七个——来自\OriginalTerm{安特罗波斯}{Anthropos}和\OriginalTerm{厄克利西亚}{Ecclesia}，因为安特罗波斯的声音发出，并赋予了宇宙（万物）以形态。因为声音的声响产生了形状，并将之赋予它们。由此得出，\OriginalTerm{逻各斯}{Logos}和\OriginalTerm{佐厄}{Zoe}有八个（半元音）；\OriginalTerm{安特罗波斯}{Anthropos}和\OriginalTerm{厄克利西亚}{Ecclesia}有七个（元音）；\OriginalTerm{帕特尔}{Pater}和\OriginalTerm{阿勒忒亚}{Aletheia}有九个（哑音）。但由于逻各斯缺少一个（以成为八位一体），那位在父内的（者）被移离（他在天主右手的座位），并降下（到地上）。他被（父）派遣给那与他分离者，为纠正已犯的行为。（他的降下）是为了使圆滿中那内在于\OriginalTerm{阿伽托斯}{Agathos}的统一过程，能在所有中产生那从一切而来的单一能力。于是，那属于七个（元音）的获得了八个（半元音）的能力；并且产生了三个\Emph{\OriginalTerm{托波斯}{topoi}}，与（九、七、八这三个）数目相对应——这三个\Emph{\OriginalTerm{托波斯}{topoi}}是八位一体。这三个彼此相加，便显示了二十四个（字母）的数目。而且，（他当然）主张，那三个元素——（他本人断言它们通过婚姻联合与三种能力（相连），并且（由于这种二元状态）成为六个，并从它们流溢出二十四个元素——通过四元体不可言说的话语被四重化，产生与这些相同的数目（二十四）。并且，他说，这些属于\OriginalTerm{阿诺诺玛斯托斯}{Anonomastus}。又（他断言），这些通过六种能力被传达到与\OriginalTerm{阿奥拉托斯}{Aoratus}的相似。又（他说），这些元素中有六个双字母，是图像之图像，与二十四个字母一起计算，通过一个类比的能力，产生数目三十。\par

% structure: p0109 source=h2 target=subsection reason=relative-heading-rank; base=h2

\subsection{第四十二章 他的体系应用于解释我主的生活与死亡}

他又说，根据这一计算及比例，一位在六日之后亲自登上高山者，以肖像的相似形出现，成了第四位，并成为第六位。他还断言，祂降下并被七重体（\OriginalTerm{七重体}{Hebdomad}）所拘禁，因而成为光辉的八重体（\OriginalTerm{八重体}{Ogdoad}）。祂在自身内包含了祂所显示的万有的全部数目，正如祂前来受洗时那样。显示的表记是鸽子的降下，这鸽子是\GreekText{Ω}（欧米伽）和\GreekText{Α}（阿尔法），其数目是801。为此，他主张\Person{梅瑟}说人是在第六日受造的。他还断言，苦难的安排发生在第六日，即预备日；因此在这日，末后的人为首先之人的重生而显现。这一安排的开端与终结是第六时辰，其时祂被钉在木头上。因为他说，完美的理智（\OriginalTerm{理智}{Nous}）知道六重数目具有生产与重生的能力，便通过那显现的光辉者，将引入该数目的重生显示给光明之子。因此他又说，双字母包含那卓越的数目。因为那光辉的数目与二十四元素相混合，产生了由三十个字母组成的名字。\par

% structure: p0111 source=h2 target=subsection reason=relative-heading-rank; base=h2

\subsection{第四十三章　字母，诸天的象征}

然而，他利用了七个数目的总和为工具，以使那自我筹划的结果得以显示。他说，目前你要明白，那卓越的数目就是由那光辉者所形成的、并被分割而留在外面者。祂借着祂自己的权能与智慧，通过祂自身的射发，仿效七种权能，将生命赋予这个世界，使之由七种权能构成，并使之成为可见宇宙的灵魂。因此，这一位亲自进行了这样一切运作，如同自发而为；这些运作，由于是不可模仿之物的模仿，服务于母亲的智慧。第一重天发出\GreekText{Α}（阿尔法）音，第二重天发出\GreekText{Ε}（厄普西隆）音，第三重天发出\GreekText{Η}（厄塔）音，第四重天——即七元音中间的那个——发出\GreekText{Ι}（约塔）的权能，第五重天发出\GreekText{Ο}（奥密克戎）音，第六重天发出\GreekText{Υ}（宇普西隆）音，第七重天即中央之音后的第四重天发出\GreekText{Ω}（欧米伽）音。所有权能连合为一，发出声音，赞美那从其射发者。那赞美之声的荣耀上达至元父（\OriginalTerm{至元父}{Progenitor}）。此外，他还说，这颂赞之声传到地上，成为地上万物的创造者和生产者。他主张，其证明可从刚出生的婴儿取得，他们的灵魂在离开子宫的同时，同样发出每一个元素的这种声音。因此，他说，正如七种权能赞美圣言（\OriginalTerm{圣言}{Logos}）一样，婴儿中悲伤的灵魂也赞美祂。为此，\Person{达味}也曾宣告说：\BibleQuote{从婴儿和吃奶者的口中，你完美了赞美。}又说：\BibleQuote{诸天宣扬天主的荣耀。}然而，当灵魂遭遇困苦时，它发出的呼喊只有\GreekText{Ω}（欧米伽），因为它受苦，为的是使上面的灵魂意识到与她相似者，从而派遣一位来帮助这地上的灵魂。\par

% structure: p0113 source=h2 target=subsection reason=relative-heading-rank; base=h2

\subsection{第四十四章　论二十四字母的产生}

就此为止。然而，关于二十四元素的产生，他表述如下：单一（\OriginalTerm{单一}{Henotes}）与独一（\OriginalTerm{独一}{Monotes}）共存，由此生出两个射发，即单一者（\OriginalTerm{单一者}{Monas}）与一（\OriginalTerm{一}{Hen}），这两者相加成为四，因为二加二为四。再者，二与四的射发相加，显出了数目六；这六被四倍化，产生二十四形式。这些是第一四元（\OriginalTerm{四元}{tetrad}）的名字，它们被视为至圣所，不可言说，惟有圣子认识。它们的名字，圣父知道是什么。那些在祂那里以沉默与信心被念出的名字，乃是深不可言者（\OriginalTerm{深不可言者}{Arrhetus}）、沉默（\OriginalTerm{沉默}{Sige}）、父（\OriginalTerm{父}{Pater}）与真理（\OriginalTerm{真理}{Aletheia}）。这四元的全部数目是二十四字母。因为深不可言者有七个元素，沉默有五个，父有五个，真理有七个。同样，第二四元也是如此：圣言（\OriginalTerm{圣言}{Logos}）与生命（\OriginalTerm{生命}{Zoe}）、人（\OriginalTerm{人}{Anthropos}）与教会（\OriginalTerm{教会}{Ecclesia}），展现了相同数目的元素。他还说，救主所表达的名字——即\Person{耶稣}——由六个字母组成，但祂那不可言说的名字，按字母逐一计算，由二十四个元素组成；而\Person{基督}，圣子，则由十二个字母组成。他说，在基督内的不可言说之名由三十个字母组成，这是根据祂内里的字母，逐一计算元素而得。因为基督一名由八个元素组成：\GreekText{Χ}（希）由三个，\GreekText{Ρ}（若）由两个，\GreekText{ΕΙ}由两个，\GreekText{Ι}（约塔）由四个，\GreekText{Σ}（西格玛）由五个，\GreekText{Τ}（陶）由三个，\GreekText{ΟΥ}由两个，\GreekText{Σ}（桑）由三个。因此，他们声称，在基督内的不可言说之名由三十个字母组成。他们断言，为此祂说：\BibleQuote{我是阿耳法和敖默加}，显示那鸽子，它象征性地有此数目，即八百零一。\par

% structure: p0115 source=h2 target=subsection reason=relative-heading-rank; base=h2

\subsection{第四十五章　为何耶稣被称为阿尔法}

现今耶稣拥有这种不可言喻的诞生。因为从宇宙之母，即第一\OriginalTerm{四元体}{tetrad}，以女儿的方式，产生了第二\OriginalTerm{四元体}{tetrad}，它成为一个\OriginalTerm{八元体}{ogdoad}，从它产生了\OriginalTerm{十元体}{decade}；这样便产生了十，接着是十八。因此，\OriginalTerm{十元体}{decade}与\OriginalTerm{八元体}{ogdoad}结合，使其变为十倍，产生了八十；再次使八十变为十倍，生成了八百。因此，从\OriginalTerm{八元体}{ogdoad}到\OriginalTerm{十元体}{decade}产生的所有字母的总数是八百八十八，这就是耶稣；因为耶稣这个名字，按照字母的数值，是八百八十八。同样，希腊字母有八个\OriginalTerm{一元体}{monad}、八个\OriginalTerm{十元体}{decade}和八个\OriginalTerm{百元体}{hecatontad}；这些显示了八百八十八的计算之和，即耶稣，祂由所有数字组成。正因如此，祂被称为\OriginalTerm{阿耳法}{Alpha}（和\OriginalTerm{敖默加}{Omega}），表明祂的诞生来自万有。\par

% structure: p0117 source=h2 target=subsection reason=relative-heading-rank; base=h2

\subsection{第四十六章 \Person{马尔库斯}论主的诞生与生活}

但论及这位（耶稣）的受造，他如此表述：从第二\OriginalTerm{四元体}{tetrad}发出的权能塑造了在地上显现的耶稣；天使\Person{加俾额尔}占据了\OriginalTerm{逻各斯}{Logos}的位置，圣神占据了\OriginalTerm{柔厄}{Zoe}的位置，\Quote{至高者的能力}占据了\OriginalTerm{人}{Anthropos}的位置，童贞女占据了\OriginalTerm{教会}{Ecclesia}的位置。因此，在\Person{马尔库斯}的体系中，按\OriginalTerm{救恩安排}{dispensation}出现的那人经由\Person{玛利亚}诞生。当祂来到水边时，（他说）祂如鸽子降临在那已上升并充满第十二数目者身上。在祂内居住着这些的种子，即那些与祂一同播种、与祂一同降下、与祂一同上升的。他说，那降临在祂身上的这个权能，是\OriginalTerm{圆满}{Pleroma}的种子，它自身包含圣父和圣子，以及\OriginalTerm{缄默}{Sige}的不可名状的能力，这能力透过这些和所有\OriginalTerm{永世}{Aeon}而被认识。他说，这（种子）是祂内的那神灵，它藉着圣子的口在祂内说话，宣认自己为人子，并宣认祂是彰显圣父的那一位；（并且）当这神灵降临在\Person{耶稣}身上时，祂便与耶稣合一。他说，那属于\OriginalTerm{救恩安排}{dispensation}的救主消灭了死亡，而祂（即基督耶稣）则彰显了圣父。他说，因此耶稣是那\OriginalTerm{救恩安排}{dispensation}中的人的名字，它是为那将要降临在祂身上的\OriginalTerm{人}{Anthropos}的相似和成形而设立的；当祂将祂收到自己内时，祂便保有了祂。（他）说祂是\OriginalTerm{人}{Anthropos}，（祂）是\OriginalTerm{逻各斯}{Logos}，（祂）是\OriginalTerm{父}{Pater}，以及\OriginalTerm{不可言说者}{Arrhetus}、\OriginalTerm{缄默}{Sige}、\OriginalTerm{真理}{Aletheia}、\OriginalTerm{教会}{Ecclesia}和\OriginalTerm{柔厄}{Zoe}。\par

% structure: p0119 source=h2 target=subsection reason=relative-heading-rank; base=h2

\subsection{第四十七章 从\Person{马尔库斯}追随者的著作摘录证明\Person{马尔库斯}体系是\Person{毕达哥拉斯}的体系}

因此，我相信，关于这些教义，凡心智健全者都显而易见：（这些主张）毫无论据，且远非合乎宗教的知识，而只是占星术发现和\OriginalTerm{毕达哥拉斯学派}{Pythagorean}算术技艺的部分内容。渴望学习的你们，将能确证这一断言（通过参考之前的书籍，在我们所阐明的其他意见中，我们也同样解释了这些教义）。然而，为了能证明一个更清晰的陈述，即这些（\OriginalTerm{马尔库斯派}{Marcosian}）不是基督的门徒，而是\Person{毕达哥拉斯}的门徒，我将继续解释这些（异端者）从\Person{毕达哥拉斯}有关星辰的\OriginalTerm{天象}{meteoric phenomena}中所汲取的那些意见，以撮要的形式尽可能这样做。\par

现在毕达哥拉斯学派提出以下说法：宇宙由\OriginalTerm{一元}{Monad}和\OriginalTerm{二元}{Duad}构成，从一元算到四，这样就产生了一个\OriginalTerm{十}{decade}。再者，二元一直推进到特殊的（字母）——例如二、四和六——便显出了（数字）十二。再者，如果从二元算到十，就产生了三十；其中包含了\OriginalTerm{八}{ogdoad}、十和\OriginalTerm{十二}{dodecade}。因此，由于十二具有特殊的（字母），十二便伴随有一种奇特的情感。正因如此，（他们主张）当有关第十二个数字出现错误时，那只羊就从羊群中跳出而迷失；因为背教的发生同样是从十开始。以类似方式关联到十二，他们谈论一块钱币，妇人丢失后点灯仔细寻找。对九十九只羊中丢失的一只，他们也有类似的应用；将这些数字彼此相加，他们对数目进行了神话般的解说。他们断言，十一乘以九就产生九十九；因此据说阿们这个词包含了九十九这个数。关于另一个数字，他们这样表述：字母Eta连同特殊的（字母）构成全部八，因为它位于从Alpha起的第八位。然后，计算这些元素（字母）的数目，除去特殊的（字母），从Alpha加到Eta，就得出三十。任何从Alpha到Eta的人，减去特殊的（字母），都会发现元素的数目是三十。因此，三十这个数由三种力量统一；当它自乘三次就产生九十，因为三乘三十是九十（而这三元自乘产生九）。这样，八从第一个八、十和十二中产生了九十九这个数。有时他们把这个（三元组）的数目集合成一个总和，产生一个\OriginalTerm{三十}{triacontad}；而有时他们减去十二，算作十一。同样，减去十，变成九。将这些彼此连接，乘以十倍，他们便完成了九十九这个数。然而，第十二\OriginalTerm{永世}{Aeon}离开了十一个（上面的永世）而向下离去，他们声称这也与字母相关。因为字母的形状教导我们这一点。因为L是第十一个字母，这个L就是数字三十。他们说，这是按照上界安排的形象来安放的；因为从Alpha起，不计特殊的（字母），字母本身的数目，加上到L（包括L本身），产生九十九。但是，位于（字母表）第十一位的L下来寻求与自己相似的数字，以便填补第十二个数字，当它被发现时就被填满，这从字母本身的形状就显明了。因为Lambda，当它达到对相似自身的探究时，找到并夺取了它，填补了第十二个位置，即字母M，它由两个Lambda组成。正因如此，这些（马库斯的追随者）通过他们的知识，避开九十九的位置，即\OriginalTerm{缺损}{Hysterema}，左手的象征，而追求那一个（即加上九十九）据说是被转移到他右手的。\par

% structure: p0122 source=h2 target=subsection reason=relative-heading-rank; base=h2

\subsection{第四十八章. 他们的宇宙论根据这些字母的神秘教义而构建。}

他们声称，由母亲首先创造了四个元素，即火、水、土、气；这些被投射为上面\OriginalTerm{四元}{tetrad}的影像；而计算这些元素的能量——例如热、冷、湿、干——他们断言它们准确地描绘了八。接着他们这样计算十种力量。（他们说）有七个球体，他们称之为天。还有一个包含这些的圆圈，他们称之为第八天；此外，他们肯定存在太阳和月亮。这些是十个，他们说，是从\OriginalTerm{逻各斯}{Logos}和\OriginalTerm{生命}{Zoe}发出的无形之十的影像。（他们断言）然而，十二由所谓的黄道带指示。因为十二个黄道星座，他们说，最明显地预示了\OriginalTerm{人}{Anthropos}和\OriginalTerm{教会}{Ecclesia}的女儿，即十二。并且，他说，上层天从相反方向与整个（星座）的最快旋转运动结合；而且因为（这天）在其腔内减缓，以其缓慢平衡了那些（星座）的速度，以致在三十年内完成从星座到星座的运行——他们因此断言这（天）是\OriginalTerm{界限}{Horos}的影像，他环绕这些的母亲，她有三十个名字。再者，（他们肯定）月亮，三十天内穿越天空，根据天数描绘了永世的数目。并且（他们说）太阳，完成其运行并在十二个月内精确返回其轨道起始位置，显明了十二。并且（他们说）日子本身，涉及十二小时的度量，构成了空洞之十二的预象；而实际黄道圈的圆周由三百六十度组成，每个黄道星座有三十部分。因此，即使通过圆圈，他们坚持保存了十二与三十连接的影像。此外，他们声称大地被分为十二个区域，每个区域通过从天上发送下来的力量接受一种力量，并产生与通过流溢传递（相似）的力量相似的子女；（因此）他们断言大地是上面十二的预象。\par

% structure: p0124 source=h2 target=subsection reason=relative-heading-rank; base=h2

\subsection{第四十九章 巨匠造物主的作品是会朽坏的}

除这些（论点）之外，他们（规定）超级\OriginalTerm{八神体}{Ogdoad}的\OriginalTerm{巨匠造物主}{Demiurge}，渴望模仿那无限、永恒、无界限者，以及不受时间条件限制者；而（巨匠造物主）因他是\OriginalTerm{缺体}{Hysterema}的果实，不能表现此（八神体）的稳定和永恒，为此他指定了时间、季节和数字，用许多年来衡量此（八神体）的永恒，以为通过时间的众多可以模仿其无限性。他们又说，当真理逃避他的追寻时，虚谎便紧随着他；因此，当时候满足时，他的作品便瓦解了。\par

% structure: p0126 source=h2 target=subsection reason=relative-heading-rank; base=h2

\subsection{第五十章 马尔谷与科拉巴苏被依勒内所驳斥}

因此，\OriginalTerm{瓦伦提诺}{Valentinus}学派的人就创造和宇宙提出了这些主张，在每种情况下都宣扬更加空洞的意见。他们认为，如果有人也做出更大的发现，似乎能行奇事，这就是他们体系中的生产力。而且，（他们暗示）发现圣经的每个细节都与上述数字相符，他们（试图）控告梅瑟和先知，声称他们寓言式地谈论了永世的度量。既然这些说法琐碎而不稳定，我以为将之呈现给（读者）并不妥当。（这并不必要，）因为如今可敬的司铎\Person{依勒内}已经有力而详尽地驳斥了这些（异端者）的意见。我们感谢他使我们知道了他们的发明，（并因此成功地）证明这些异端者从\OriginalTerm{毕达哥拉斯哲学}{Pythagorean philosophy}和占星家的过度理论中窃取了这些见解，并将它们归咎于基督，好像基督曾传授这些（教义）一样。但既然我认为这些人的无价值意见已经足够解释清楚，并且已经清楚地证明了\Person{马尔谷}和\Person{科拉巴苏}是谁的门徒，他们继承了瓦伦提诺学派的学说，让我们看看\OriginalTerm{巴西利得}{Basilides}同样提出了什么主张。\par

\SourceRecord{Translated by J.H. MacMahon.；Ante-Nicene Fathers；Vol. 5.；Edited by Alexander Roberts, James Donaldson, and A. Cleveland Coxe.；Buffalo, NY: Christian Literature Publishing Co.,；1886.；<http://www.newadvent.org/fathers/050106.htm>.}

% structure: p0001 source=h1 target=section reason=page-title

\section{\WorkTitle{驳斥一切异端（第七卷）}}

% structure: p0002 source=h2 target=subsection reason=relative-heading-rank; base=h2

\subsection{目录}

以下是《\Emph{\WorkTitle{驳斥一切异端}}》第七卷的内容：\par

\Person{巴西理得}的主张是什么，以及他如何受到\Person{亚里士多德}学说的影响，并据此构建了他的异端。\par

以及\Person{撒图尼鲁斯}的陈述，他与\Person{巴西理得}大约同时期活跃。\par

以及\Person{默南德}如何提出世界是由天使创造的论断。\par

\Person{马尔基翁}的愚昧是什么，以及他的观点并非新颖，也不是出自圣经，而是从\Person{恩培多克勒}得来。\par

\Person{卡尔波克拉特}如何愚昧地行事，他自己也声称万有是由天使创造的。\par

\Person{塞林图}完全不依赖圣经，他的观点并非来自圣经，而是来自\OriginalTerm{埃及人}{Egyptians}的学说。\par

\OriginalTerm{厄俾翁派}{Ebionaeans}提出了什么意见，以及他们更倾向于遵守犹太习俗。\par

\Person{德奥多图}如何堕入错误，他从\OriginalTerm{厄俾翁派}{Ebionaeans}（部分从\Person{塞林图}）吸取了其体系的要素。\par

以及\Person{塞尔东}的意见是什么，他不仅阐述了\Person{恩培多克勒}的学说，还邪恶地引诱\Person{马尔基翁}出面。\par

以及\Person{路奇安}在成为\Person{马尔基翁}的门徒后，如何抛弃了一切羞耻，不断亵渎天主。\par

\Person{阿佩莱斯}也成为这位异端者的门徒，但他并不常提出与其导师相同的意见；相反，他受自然哲学家的学说影响，将宇宙的物质作为万物的根本原理。\par

% structure: p0015 source=h2 target=subsection reason=relative-heading-rank; base=h2

\subsection{第一章 异端比拟为（1）风暴海洋，（2）塞壬的礁石；从\Person{乌利西斯}和\Person{塞壬}的道德教训。}

这些人的弟子，当他们看到异端的教义像被狂风激荡的海洋时，应当航向宁静的港湾。因为这样的海既布满野兽又难以航行，好比（我们可以说）西西里海，传说那里有独目巨人、卡律布狄斯、斯库拉和塞壬的礁石。希腊诗人称\Person{乌利西斯}巧妙地航过此通道，利用了这些奇异怪物的恐怖。因为这些怪物对航行者的野蛮残暴是显著的。而\Person{塞壬}甜美和谐地歌唱，迷惑航行者，以悦耳的声音引诱听到的人将船驶向岬角。诗人报告说\Person{乌利西斯}得知此事后，用蜡封住同伴的耳朵，将自己绑在桅杆上，清楚地听到\Person{塞壬}的歌声，却安全地驶过。我建议我的读者采取类似的方法：或因自身的软弱，用蜡封住耳朵，径直驶过异端的教义，甚至不听那些能像\Person{塞壬}的甘美歌声一样轻易引诱他们进入快乐的道；或将自己绑在\Person{基督}的十字架上，忠诚地听他的话语，不被分心，因为他已将信赖寄托于他早已坚定结合的那一位，并应坚守在这信德中。\par

% structure: p0017 source=h2 target=subsection reason=relative-heading-rank; base=h2

\subsection{第二章 \Person{巴西理得}的体系源自\Person{亚里士多德}}

因此，既然在前六卷中我们已经解释了先前的异端意见，现在似乎不宜对\Person{巴西理得}的学说保持沉默，这些学说是斯塔基拉的\Person{亚里士多德}的观点，而非\Person{基督}的。但即使以前已经阐明过\Person{亚里士多德}提出的意见，我们现在也不迟疑地以某种概要的形式预先陈述它们，以使读者通过更接近地比较两个体系，能轻易看出\Person{巴西理得}所提出的学说实际上是\Person{亚里士多德}的巧妙诡辩。\par

% structure: p0019 source=h2 target=subsection reason=relative-heading-rank; base=h2

\subsection{第三章 \Person{亚里士多德}哲学纲要}

因此，\Person{亚里士多德}将实体分为三类。一部分是某种属，另一部分是某种种，如那位哲学家所述，第三部分是某种个体。个体之所以是个体，并非由于身体的微小，而是因为它在本性上不能有任何分割。相反，属是一种集合体，由许多不同的胚芽组成。从这个属，就像从一个堆中，所有存在事物的种获得其区别。而属为所有生成物构成一个充足的原因。为了使上述陈述清晰，我将用一个例子来证明。通过这个例子，我们有可能追溯整个逍遥学派的思辨。\par

% structure: p0021 source=h2 target=subsection reason=relative-heading-rank; base=h2

\subsection{第四章 \Person{亚里士多德}的一般理念}

我们肯定绝对动物存在，而不是某个动物。这动物既不是牛，也不是马，也不是人，也不是神；它根本不表示任何这些，而是绝对动物。从这绝对动物，所有个别动物的种获得其实体。这动物性本身，即\Emph{\OriginalTerm{最高属}{summum genus}}，构成了所有在那（个别）种中产生的动物的（起源）原则，但（本身）不是任何被生成的东西之一。因为人是从那动物性获得存在原则的动物，马是从那动物性获得存在原则的动物。马、牛、狗以及其余每一种动物，都是从绝对动物获得存在原则，而动物性本身不是这些中的任何一个。\par

% structure: p0023 source=h2 target=subsection reason=relative-heading-rank; base=h2

\subsection{第五章 非存在作为原因}

然而，如果这动物性不是这些种中的任何一个，那么根据\Person{亚里士多德}，所生成之物的实体是从非存在的实体获得其实在性的。因为动物性，从那里这些个别之物被衍生出来，不是它们中的任何一个；尽管它不是它们中的任何一个，它却成为存在物的某个起源原则。但谁建立了这个实体作为随后产生之物的起源原因，我们将在适当时候讨论这类问题时宣告。\par

% structure: p0025 source=h2 target=subsection reason=relative-heading-rank; base=h2

\subsection{第六章 亚里士多德论实体；谓词。}

既然，如我所言，实体是三重性的，即属、种、（和）个体；（且）我们已设定动物性为属，而人为种，因之已与大多数动物相区别，但仍需与同种者相认同，因其尚未成为现实实体的种——（因此，）当我从属中借用一名，赋予人以形相时，我称他为苏格拉底或第欧根尼，或众多名称之一。并且既然（以此方式，我重复）我用一个名称来涵盖那从属而生的种的人，我便称此类型的实体为个体。因为属被分为种，种被分为个体。但（至于）个体，既然它被一个名称所涵盖，按其本性，它不能像我们分割前述（属和种）那样再被分割为其他任何东西。\par

亚里士多德首先地、尤其地、且首要地将此称为实体，因为它既不能述说一个主体，也不能存在于一个主体之中。然而，他述说主体，正如用属，即我所说的构成动物性者，通过一个共同的名称来述说所有具体动物，如牛、马及其他位于此属之下的。因为说\Quote{人是动物}、\Quote{马是动物}、\Quote{牛是动物}以及其余每一个，都是正确的。所谓\Quote{述说一个主体}的意思就是：因为它是一，它同样可以述说许多（殊相），即使它们种属不同。因为马或牛就其为动物而言，与人并无不同，因为动物的定义平等地适用于所有动物。动物是什么？如果我们定义它，一个普遍的定义将涵盖所有动物。因为动物是一个有生命的实体，赋有感觉。如牛、人、马以及其余（动物界）每一个。但\Quote{在一个主体中}的意思是：那内在于某物之中，但不是作为部分，不可能与它所内在于的东西分开存在。而这构成实体中的每一个偶性，被称为性质。据此，我们说某些人具有某种性质，例如，白的、灰的、黑的、公正的、不公正的、节制的，以及其他类似（特性）。但它们中任何一个都不能自存，而必须内在于其他事物之中。然而，如果既不是我用来述说所有个体动物的动物，也不是在所有非本质属性之物中可发现的偶性，能够自存，而（且如果）个体是由这些构成的，（那么结论就是）三重划分的实体，它并非由其他事物构成，乃是由非存在物组成。那么，如果那首要地、卓越地且特别地被称为实体的东西是由这些构成的，则根据亚里士多德，它从非存在物中取得存在。\par

% structure: p0028 source=h2 target=subsection reason=relative-heading-rank; base=h2

\subsection{第七章 亚里士多德的宇宙论；他的\WorkTitle{心理学}；他的\Quote{隐德莱希}；他的神学；他的伦理学；巴西里德追随亚里士多德。}

但关于实体，以上所述已经足够。然而，实体不仅被称为属、种、（和）个体，而且还被称为质料、形相和缺失。不过，即使允许这样划分，在这些当中也没有区别。既然实体是这样描述的，世界的安排就按照某种如下的方式发生了。根据亚里士多德，世界被分成众多且多样的部分。从地球到月亮的那部分世界缺乏预见、无引导，只受其自身本性的支配。而另一部分（世界）在月亮之上，延伸到天的表面，则安排于一切秩序、预见和治理之中。现在，（天体的）表面构成了某种第五实体，并且远离所有那些构成宇宙体系的自然元素。根据亚里士多德，这是一种第五实体——仿佛是一种超越世界的本质。而（这种本质）在他的体系中成为（一种逻辑必然性），以便与（逍遥学派的）世界划分相一致。并且（关于这第五本性的主题）构成了哲学中一个独特的研究。因为存在一篇题为\WorkTitle{物理（现象）讲演}的讨论，在其中他详尽地处理了从地球直到月亮（的区域）中由自然而非上智安排所引导的运作。他还有另一篇独特的论著，讨论月亮以外（区域）的万物原理，其标题如下：\WorkTitle{形而上学}。另一篇独特的论文是由他（所写），题为\WorkTitle{论第五实体}，亚里士多德在此作品中阐述了他的神学观点。\par

世上存在着某种我们刚刚试图勾画轮廓的宇宙划分，而（与之对应的是划分）的亚里士多德哲学。然而，他那部（题为）\WorkTitle{论灵魂}的著作晦涩难懂。因为在整三卷书（他处理此主题）中，不可能清楚地说出亚里士多德关于灵魂的观点。因为就其提供的灵魂定义而言，宣布这个定义是容易（足够）的，但定义所意指的东西却难以发现。因为灵魂，他说，是自然有机身体的\OriginalTerm{隐德莱希}{entelecheia}；（而要解释）这究竟为何物，则需要大量论证和冗长研究。然而，至于天主，万有中一切光辉之物的创造者，（本质）这一（第一因）——即使对于那些进行比（灵魂）更持久研究的思考者而言——也比灵魂本身更难认识。但亚里士多德所提供的关于天主的定义，我承认，不难确定，但理解其意义却不可能。因为，他说，（天主）是一个\Quote{概念的概念；}但这是完全不存在（实体）。世界则是依照亚里士多德，是不朽（且）永恒的。因为世界本身毫无瑕疵，因为它受天主眷顾和自然所引导。亚里士多德不仅奠定了关于自然、宇宙体系、天主眷顾和天主的学说，而且（不止于此）；因为还存有他关于伦理学主题的某篇论文，他题为\WorkTitle{伦理学}。但贯穿这些，他的目标是将听众的习惯从毫无价值变为卓越。因此，当发现巴西利德不仅是在精神上，而且在实际表达和名称上将亚里士多德的教义移植到我们福音和救恩的教义中时，除了通过将他从他人那里偷取的东西归还，以此向这位（异端者）的门徒证明基督对他们毫无益处，因为他们仍是异教徒似的，还能做什么呢？\par

% structure: p0031 source=h2 target=subsection reason=relative-heading-rank; base=h2

\subsection{第八章 巴西利德与依西多禄声称其体系获宗徒认可；实则跟随亚里士多德。}

因此，巴西利德与依西多禄，即巴西利德的真正儿子和门徒，说玛弟亚传授给他们秘密言论，这些言论，在我受特别教导之后，是他从救主那里听到的。那么，让我们看看巴西利德与依西多禄以及整个这些（异端者）的团伙如何不仅彻底欺骗玛弟亚，甚至欺骗救主本人。（时间）曾是，（巴西利德）说，那时无物存在。然而，那个\InnerQuote{无}甚至也不构成存在物的任何东西；但是，坦率、真诚且毫不诡辩地说，它完全是\Quote{无}。但是，他说，当我使用\Quote{曾是}这一表达时，我并不是说它曾是；而是（我这样说）为了表达我意欲阐明的东西的意思。我因而肯定，他说，它曾是\Quote{完全是无}。因为，他说，那被称为（如此）的并非绝对不可言说——尽管我们无疑称这为不可言说——而是那\Quote{非不可言说}。因为那\Quote{非不可言说}并不被称为不可言说，而是，他说，超乎一切可称呼的名之上。因为，他说，这些名称绝不适用于世界，但其划分如此繁多，以致存在（名称的）不足。而且，他说，我不承担为所有事物发现恰当名称的责任。然而，无疑，人应当在思想上，而非通过名称，以一种不可言说的方式，去领会那些被命名（事物的）特性。因为歧义的术语（当被教师使用时）已为他们的学生造成混乱，并成为关于对象错误的源头。（巴西利德派）首先抓住了这个从逍遥学派（哲人）那里借来并暗地获取的教义，利用那些与他们聚集的人的愚昧。因为亚里士多德，生于巴西利德之前许多代，最早在\WorkTitle{范畴篇}中建立了关于\OriginalTerm{同名异义词}{homonymous words}的系统。而这些异端者将这个（系统）公之于众，仿佛它完全是他们自己的，仿佛是某种新奇的（教义），以及来自玛弟亚言论的某种秘密启示。\par

% structure: p0033 source=h2 target=subsection reason=relative-heading-rank; base=h2

\subsection{第九章 巴西利德采纳亚里士多德关于\Quote{无物}的教义。}

既然，因此，\Quote{无物}存在——（我是指）没有物质，也没有实体，也没有非实体，也没有绝对者，也没有复合体，（也没有可思想的，也没有不可思想的，（也没有可感觉的，）也没有无感觉的，也没有人，也没有天使，也没有神，简言之，没有任何那些有名称、可被感官把握或理智认识的对象，但（它们）是这样（被认识的），甚至更精细——当一切事物被完全移除时——（既然，我再说，\Quote{无物}存在，）天主，\Quote{不存在的}——亚里士多德称其为\Quote{思想的思考}，但这些人（巴西利德主义者）称其为\Quote{不存在的}——不可思想地、不可感觉地、无定地、无意地、无情感地、无欲望地、无动机地，愿意创造世界。他说，我使用\Quote{愿意}这一表达，是为了表示（祂这样做是）无意的、不可思想的、不可感觉的。而通过\Quote{世界}这一表达，我并非指后来按广度和分离形成并独立存在的那个；不，（远非如此，）因为（我是指）世界的胚种。然而，世界的胚种自身包含一切事物，正如芥菜籽同时包含一切事物，将它们（收集）在一起置于最微小（的范围）内，即根、茎、枝、叶和无数从植物产生的籽粒，（作为）其他植物的种子，并常常还有其他（更多）从它们产生的。这样，不存在的天主出自非存在物创造了世界，投下并放置了某一个种子，它自身包含世界诸胚种的聚集体。但为了使我更清楚这些（异端者）所肯定的， （我将提及他们的如下比喻。）正如某个斑驳多彩的鸟的卵——例如孔雀，或其他更斑斓多彩的鸟——实际上是一个，却包含众多形态的、斑驳的、多种复合物质的形式；同样，他说，由不存在的天主所投下的世界之非存在种子，同时构成众多形式和众多物质的胚种。\par

% structure: p0035 source=h2 target=subsection reason=relative-heading-rank; base=h2

\subsection{第十章 世界之起源；巴西利德对\Quote{子性}的说明}

一切事物，因此凡可能表达者，以及凡尚未被发现而必须略过者，都应当从种子中生出而适应于将要生成的世界。而这一（种子），在适当的季节，以独特的方式，根据增添而增大体积，如同通过如此伟大且如此性质的天主的工具。 （但这种天主）受造物既不能言说也不能用感知把握。（现在，所有这些东西）都内在于种子中，被珍藏起来，正如我们后来在新生的孩子身上观察到牙齿的生长、父亲的禀赋、理智以及一切从年轻时期起逐渐增长于人的东西，尽管先前并不存在。但若说不存在的天主的任何投射成为任何非存在物，那是荒谬的（因为巴西利德完全避开并恐惧那些以投射方式生成的事物的实体；因为，他问，有何种投射的必要，或我们必须假定何种物质先前存在，以使天主构造世界，如同蜘蛛结网；或（如同）一个必死的人，为了工作而取一块黄铜或木头，或物质的某些其他部分？），——（投射，我再说，被排除），当然，巴西利德说，天主说话，事就成了。而这些（人）断言，这就是梅瑟所说的：\Quote{有光，就有了光。}他说，光从何而来？从无。因为经上并未写，他说，从何处，而只有这一点，（它）来自说话者的声音。而说话者，他说，是不存在的；那被产生的，也不存在。宇宙的种子，他说，是从非存在物产生的；（我所说的种子）就是所说的那句话：\Quote{有光。}而这一点，他说，就是福音中所记载的：\Quote{他是真光，照亮每个来到世界的人。}他从那种子中获得他的起源原则，并从同一源头得到他的光照能力。这就是那在其自身拥有整个胚种聚集的种子。亚里士多德将这一点称为\OriginalTerm{本性}{genius}，他把它分布到无限种类中；正如从非存在的动物，我们区分出牛、马、人。因此，当宇宙的种子成为（后续发展的）基础时，那些（异端者）断言，（引用巴西利德自己的话：）\Quote{我肯定，}他说，\Quote{在此之后所造的一切，不要问从何而来。因为（种子）将所有种子珍藏并安放在自身之内，正如非存在物，它们注定由不存在的天主所产生。}\par

因此，让我们看看他们所说从宇宙种子所生成的事物中，什么是第一，什么是第二，什么是第三。他说，在种子自身内存在一个\OriginalTerm{子性}{Sonship}，三重，每一方面都与不存在的天主同质，从非存在物所生。这（包含）三重划分的\OriginalTerm{子性}{Sonship}，一部分是精细的，（另一部分是粗糙的，）还有一部分需要净化。因此，那精细的部分，首先，在不存在者最早投下种子同时，立即冲出，向上急升，从下方向高处飞驰，运用某种诗意描述的速度——\par

\Quote{……如同翅膀或思想，}——\par

他说，并达到了那非存在者。因为每一个本性都渴望那（非存在者），因其超丰盈的美和光华。然而，每个（本性）以不同的方式（渴望它）。但是，那较粗糙的部分，（子性）仍留在种子中，（并且）是一个模仿性的（原则），不能快速向上。因为（这部分）远不及子性所具有的精微，那子性凭借自身快速向上，（所以较粗糙的部分）被抛在后面。因此，较粗糙的子性为自己装备了某种翅膀，正如\Person{亚里士多德}的导师\Person{柏拉图}在\WorkTitle{费德鲁斯篇}中为灵魂所配的飞翼。而\Person{巴西里德斯}称这样的（东西）不是翅膀，而是圣神；子性披戴这（圣神）施予恩惠，也接受恩惠。它施予恩惠，因为正如鸟的翅膀，若从鸟身上取下，本身不会高高向上翱翔；同样，鸟若脱离羽翼，也永远不会高高向上翱翔；（同样地，）子性涉及与圣神的某种关系，圣神也涉及与子性的关系。因为子性被圣神如翅膀般托起，也（反过来）高举它的羽翼，即圣神。它接近那精微的子性和非存在的天主，即那从无中创造世界的（天主）。然而，它不能将这（圣神）与（子性）本身一同拥有；因为（圣神）与（天主）不是同一实体，也与子性没有（任何共同）本性。正如纯净干燥的空气与（鱼的）本性相反，并毁灭鱼类；同样，与圣神的本性相反的是那同时是非存在神性与子性的地方——一个比不可言喻者更不可言喻，超越一切名号的地方。\par

因此，子性将这（圣神）留在那不能以任何言语构思或表达的蒙福之地附近。（他将圣神）并非完全遗弃或与子性分离；不，（远非如此，）正如最芬芳的香膏放入器皿中，即使（器皿）被极其小心地倒空，香膏的一些气味仍然存留，在（香膏）与器皿分离后仍被留下；器皿虽没有（香膏本身），却保留香膏的气味。同样，圣神继续没有与子性共享任何部分，与（子性）分离，自身如香膏般拥有自己的力量，即子性的香气。这就是所宣告的：\BibleQuote{正如香膏倒在头上，流到亚郎的胡须上。}这是圣神的香气从上降下，直至无形，以及我们世界附近的（空间）间隔。而圣子开始由此上升，好像被鹰的翅膀和背托起，\Person{巴西里德斯}说。因为，他说，一切（实体）从下方向向上方，从低劣者向优越者急升。因为在优越者中，没有一个愚蠢到下降至下方。然而，第三子性，即需要净化的那一部分，他说，继续留在所有种子的巨大聚合中，施予恩惠并接受恩惠。至于（第三子性）如何接受恩惠并施予恩惠，我们将在以后讨论此问题的适当场合说明。\par

% structure: p0041 source=h2 target=subsection reason=relative-heading-rank; base=h2

\subsection{第11章 \Person{巴西里德斯}的\Quote{大掌权者}}

因此，当\OriginalTerm{子性}{Sonship}的第一次和第二次上升发生，而圣神本身也按上述方式留下时，穹苍被安置在超世俗（空间）与世俗之间。因为按照巴西里德的说法，现有事物被分为两个连续的主要部分，并且据他称，一部分在某方面被称为世界，另一部分在某方面被称为超世俗（空间）。但作为世俗与超世俗（空间）分界线的圣神，既是圣的，又将子性的芬芳持守于自身。因此，当穹苍在天之上产生时，从宇宙种子和万种聚合中迸发并诞生了伟大的元首，世界之首，一种美、伟大和不可分解的能力。因为，他说，他比不可言说的实体更不可言说，比强者更强，比智者更智，超越所有你能想到的美者。这位（元首）一诞生，便升起并高飞，整个地被提升至穹苍。他在那里停住，以为穹苍是他上升和升高的终点，并认为穹苍之外再无一物。他比所有下界（实体）中任何仍留在世俗中的事物都更智慧、更强有力、更俊美、更光辉，（事实上）在美丽上超过任何你能想到的实体，唯独除却仍留在万种（聚合）中的子性。因为他不知道有一位比自己更智慧、更强大、更善美的（子性）。因此，他自以为是主、治理者和智慧的建造师，就转向创造宇宙体系中的每一个物体。首先，他认为自己不应独处，而是为了自己，从邻近（实体）中制造并生出了一个远比自己更优越、更智慧的儿子。因为非存在之神在抛下万种（聚合）时早已预定这一切。因此，他看见这儿子，便惊异、爱慕并感到惊叹。因为某种如此的美出现在大元首看来属于儿子，元首便使他坐在自己右边。按照这些（异端者）的说法，这就是所谓的八度，大元首的宝座设在那里。于是，整个天上的受造物，即以太，由大元首——伟大的智慧造物主——亲自形成。然而，从这位（元首）所生的儿子在他内行动，向他提出建议，因其自身被赋予了远超过造物主本人的智慧。\par

% structure: p0043 source=h2 target=subsection reason=relative-heading-rank; base=h2

\subsection{第十二章 巴西里德采用亚里士多德的\OriginalTerm{圆成实现}{Entelecheia}。}

这构成了亚里士多德所论的自然有机身体的\OriginalTerm{圆成实现}{entelecheia}，即一个在身体内运作的灵魂，没有它身体便无法成就任何事；我没有说任何更伟大、更辉煌、更强有力、更智慧于身体的事。因此，亚里士多德先前关于灵魂和身体所作的解释，巴西里德将其阐明为适用于大元首和他的儿子。因为按巴西里德的说法，元首生了一个儿子；而亚里士多德断言灵魂作为运作和完成，是自然有机身体的\OriginalTerm{圆成实现}{entelecheia}。因此，正如\OriginalTerm{圆成实现}{entelecheia}控制身体，按巴西里德的说法，儿子控制那位比不可言说者更不可言说的天主。因此，一切事物都是由大元首的威严所供给和管理；我说的是以太区域中直至月亮的一切物体，因为从那里起，气与以太分开了。当以太区域的一切物体都安排好后，又从万种（聚合）中升起了另一位元首，当然比所有下界（实体）更伟大，但次于被留下的\OriginalTerm{子性}{Sonship}，并远逊于第一位元首。这第二位元首他们称为瑞图斯。这个区域被称为七度，这位（元首）是下界一切事物的管理者和创造者。他也像第一位元首那样，从万种（聚合）中为自己造了一个比自己更精明、更智慧的儿子。他说，存在于这一区域（宇宙）的，正是万种的实质聚合与收集；所生成之物都是按本性产生的，正如那位预知未来者所宣告的那样，它们当在何时、当为何种、当如何。这些事物中没有一个是首领、守护者或创造者。因为对于它们而言，非存在之天主当行使创造功能时所形成的计算便是（存在）的充分（原因）。\par

% structure: p0045 source=h2 target=subsection reason=relative-heading-rank; base=h2

\subsection{第十三章 进一步解释\OriginalTerm{子性}{Sonship}。}

当因此，按照这些（异端者），整个世界和超世间实体都已完成，且（当）没有任何事物处于缺陷之下时，在所有胚种的（凝聚）中仍保留了第三\OriginalTerm{子性}{Sonship}，它被留在种子中以便施恩并受恩。而被留下的子性也必须被揭示并重新安置于上方。祂的位置应在\OriginalTerm{邻接之灵}{Conterminous Spirit}之上，靠近纯净且仿效的子性和\OriginalTerm{非存有者}{Non-Existent One}。但这是按照所写的，他说：\BibleQuote{受造之物一同叹息，一同经历产痛，等待天主子女的显现。} 现在，我们是属神的子女，他说，是被留在这里整理、塑造、修正并完善那些按其本性注定留在此宇宙区域的灵魂。\BibleQuote{于是，罪从亚当直到摩西统治了，} \BibleRef{罗 5:14} 正如所写的。因为\OriginalTerm{大元首}{Great Archon}统治，其帝国界限延伸至穹苍。祂想象自己是唯一的天主，并且之上没有别的，因为万物都被未揭示的\OriginalTerm{沉默}{Siope}所守护。他说，这是从前世代未曾显明的奥秘；但在那些日子里，大元首，即\OriginalTerm{第八层}{Ogdoad}，似乎是宇宙的王和主。但（实际上）\OriginalTerm{第七层}{Hebdomad}是此宇宙区域的王和主，第八层是\OriginalTerm{不可言说者}{Arrhetus}，而第七层是\OriginalTerm{可言说者}{Rhetus}。他说，这是第七层的元首，他曾对梅瑟说话，并说：\BibleQuote{我是亚巴郎、依撒格和雅各伯的天主，我没有将天主的名显示给他们} \BibleRef{出 6:2}（因为他们希望如此记载）——即是那不可言说者的天主，第八层的元首。因此，他说，所有在救主之前的先知都从这个（灵感）源头发预言。既然，他说，我们必须作为天主的子女被揭示，受造之物常在盼望其显现时叹息并经历产痛，福音便进入了世界，并经过了所有领导、能力和权势，以及一切被称呼的名号。\BibleRef{弗 1:21} 而（福音）确实来了，尽管没有从上面降下什么；那有福的子性也没有从那不可思议、有福、（且）非存有的天主退去。不，（远非如此；）因为如同\OriginalTerm{印度石油}{Indian naphtha}，只要从相当远的距离点燃，仍会吸引火（向其靠近），同样，从下面，从胚种凝聚的无形式状态中，权势向上上升，直至子性。因为，根据印度石油的比喻，第八层大元首之子，仿佛某种石油，领悟并攫取来自有福子性的概念，有福子性的居所位于邻接（之灵）之后。因为在圣神之中，即（邻接）之灵之中的子性的权势，与第八层大元首之子分享子性流动且涌流的思想。\par

% structure: p0047 source=h2 target=subsection reason=relative-heading-rank; base=h2

\subsection{福音从何而来；根据巴西里德的诸天数目；基督奇妙受孕的解释}

于是福音——据\Person{巴西里得}说——首先从\OriginalTerm{子性}{Sonship}，通过坐在\OriginalTerm{掌权者}{Archon}旁边的\OriginalTerm{子}{Son}，来到掌权者那里。掌权者得知他并不是宇宙的天主，而是受生的。但他（得知）在他之上有那不可言说、不可名状、非实存者和子性的藏宝，他就既悔改又充满恐惧，当他被引而明白自己陷入何种无知时。他说，这正是经上所载：\BibleQuote{敬畏主是智慧的开始。}\BibleRef{箴 1:7}因为，由坐在近处的基督口授教导，他开始获得智慧，（因为他由此）学到谁是非实存者、子性是什么、圣神是什么、宇宙的装置是什么，以及万物的结局可能是什么。这就是那在奥秘中讲述的智慧，关于这智慧，（巴西里得）说，圣经用了以下的话：\BibleQuote{不是用人的智慧所教的言语，而是用圣神所教的言语。}\BibleRef{格前 2:13}于是，掌权者受了口授教导和训诲，并（因此）充满恐惧，就开始为他妄自尊大所犯的罪作宣认。他说，这正是经上所载：\Quote{我认清了我的罪，我知道我的过犯，我为此将永远宣认。}当大掌权者受了口授教导，以及\OriginalTerm{八度}{Ogdoad}的一切受造物都受了口授教导和训诲，（并且）奥秘为天上（诸权力）知晓之后，福音也必须随后降临到\OriginalTerm{七度}{Hebdomad}，为使七度的掌权者也能同样受教导并被引入福音。大掌权者的子（因此）在他自己从上方从子性所点燃的光中，点燃了七度掌权者之子的光。七度掌权者之子得了光照，并将福音宣报给七度的掌权者。同样，根据之前的叙述，他本人也畏惧并被引而作宣认。因此，当七度的一切（存有）同样被光照、接受了福音（因为在这些宇宙区域中，根据这些异端者，有无限的受造物，即宰制者、掌权者和统治者，关于这些，在（巴西里得派）中有一部极冗长繁琐的论文，他们声称有三百六十五个天，这些天的伟大掌权者是\Person{阿布拉克萨斯}，因他的名字包含计算出的数字365，以致该名号的演算包含了所有（存有），且因此原因一年由这么多天构成）——但是，他说，当这些事情（即七度蒙光照和福音显现）如此发生后，也必须随后使我们这区域中的\OriginalTerm{无形状态}{Formlessness}得以光照，并且将奥秘启示给那留在无形状态中的子性，如同一个流产儿。\par

如今这（奥秘）并未告知以前的世代，正如他所说，经上记载：\BibleQuote{这奥秘借启示告知了我；}\BibleRef{弗 3:3}以及\BibleQuote{我听见了不可言说的话，是人不能说出的。}\BibleRef{格后 12:4}于是，（那）从上方的八度降临到七度之子身上的光，从七度降临到玛利亚之子耶稣身上，他因受照在他身上的光而得了光照。他说，这正是所记载的：\BibleQuote{圣神要临于你，}\BibleRef{路 1:35}（意思是）那从子性通过相邻之神而到达八度和七度，直到玛利亚的灵；\BibleQuote{至高者的能力要庇荫你，}（意思是）那（从上方（通过\OriginalTerm{造化者}{Demiurge}流到受造物，直到那子的）傅油的能力。他说，世界一直如此构成。直到那（子），整个子性——它被留下为使无形状态中的灵魂受益，并反过来接受益处——（这子性，我说，）当它被转化时，跟随了耶稣，向上疾驰，并出来得以净化。它变得极为精纯，以致它能如同第一个（子性）一样，通过自身而向上疾驰。因为它拥有一切能力，这能力按其本性稳固地与从上方照下来的光相联结。\par

% structure: p0050 source=h2 target=subsection reason=relative-heading-rank; base=h2

\subsection{第15章 天主与受造物的交往；\Person{巴西里得}关于（1）内在的人，（2）福音的概念；他对吾主生活与受苦的诠释。}

因此，他说，当整个\OriginalTerm{子性}{Sonship}已来临，且居于\OriginalTerm{邻接之神}{conterminous spirit}之上时，那时受造物将成为怜悯的对象。因为（受造物）至今叹息\BibleRef{罗 8:19}，并痛苦，等待着天主子女的显扬，以便所有属于子性的人都能从那里上升。当这事发生时，他说，天主将把巨大的无知带到整个世界，好使万物能按本性延续，且无一物会不当地渴望任何违反本性的事物。但（远非如此）；因为在这一区域的受造物中，所有拥有在此唯一区域中保持不朽之本性的灵魂，都继续（在其中），对任何更高或更好的（存在）一无所知。在下界领域中，不会流传任何关于居住在更高处之存有的传言或知识，以免下层灵魂因渴望不可能之事而备受折磨。（这）正像一条鱼渴望与羊群一起在山坡上觅食。（因为）他说，这样的愿望将是它们的毁灭。因此，所有停留在（这一）区域的事物都是不朽的，但如果它们倾向于离开并跨越出那些合乎本性的事物，则必朽。如此，\OriginalTerm{七重天}{Hebdomad}的\OriginalTerm{掌权者}{Archon}将对上界存有一无所知。因为巨大的无知也将抓住这位掌权者，以便悲伤、忧愁和呻吟能远离他；因为他将不渴望任何不可能之事，也不会遭受痛苦。然而，同样，同样的无知也将抓住\OriginalTerm{八重天}{Ogdoad}的大掌权者，并同样抓住所有臣属于他的受造物，以便在任何方面，没有任何事物会渴望那些违反本性的事物，从而不会（因此）被悲伤压倒。于是，将有一切事物的复原，这些事物合乎本性，从起初就在宇宙的种子中有了根基，但将在（它们自己的）适当时期复原。而每一事物，他说（巴西利德），都有其特定的时间，救主在祂说\BibleQuote{我的时辰尚未来到。}\BibleRef{若 2:4}时就是充分的（证据）。而贤士们也是如此（提供类似见证），当他们渴望地注视着（救主的）星时。因为（耶稣）自己，他说，在众星生成之时，以及在所有季节完全回归到它们起点的巨大聚集中（即所有胚种），就已经在理智中被预想了。这就是，根据这些（巴西利德派），那在自然本性中被构想为\OriginalTerm{内在属神的人}{inner spiritual man}（现在这就是那子性，它离开了那里的灵魂，不是为了让灵魂成为必朽的，而是为了它按本性停留在此处，正如第一子性在上方其适当的地方留下了圣神，即那邻接之神），——（这就是，我说，那被构想为内在属神的人，）然后被穿上它自己特有的灵魂。\par

为了不遗漏这（\Person{巴西利德}）的任何教义，我将同样解释他们所提出的关于福音的论述。因为，正如已阐明的，福音在他们眼中是对超尘世实体的知识，这是大元首所不了解的。因此，当大元首得知同样存在圣神——即边界之（神）——和子性，以及作为这一切原因的非存在之天主时，他对此传达感到喜悦，并充满欢欣。根据他们的说法，这就是福音。然而，耶稣是按照这些（异端者）的说法出生的，正如我们已经阐述过的。当先前解释过的世代发生时，我们主生活中的一切事件，按照他们的说法，都与福音书中所描述的方式相同。他说，这些事发生，是为了使耶稣成为区分混杂在一起的各类受造物的初果。因为当世界被分为一个八元体——即整个世界的头，现在大元首是整个世界的头——和一个七元体——即七元体的头，下层实体的创造者——以及我们中间盛行的这种受造物秩序，其中存在无定形时，必须通过耶稣进行的分离过程来区分那些混杂在一起的各类受造物。这分离过程是：他的肉体部分受苦，这是无定形部分，回归无定形；他的灵魂部分复活，这是七元体部分，回归七元体；他复活了属于大元首所居崇高区域的特有元素，并停留在大元首旁边；他将边界之神的特有元素提升到之上，并停留在边界之神中；通过他，第三个子性被净化，这子性曾被留下以赐予利益和接受利益，并通过耶稣上升到有福的子性，穿越所有这一切。因为这一切的全部目的就是混合所有种子的混杂，区分各类受造物，并将混杂之物恢复到其适当的组成部分。因此，耶稣成了区分各类受造物的初果，他的苦难发生的原因无非是通过苦难达成的各类混杂受造物的区分。因为照此方式，\Person{巴西利德}说，留在无定形中以赐予利益和接受利益的整个子性，按照在耶稣身上发生的本性区分方式，被分成其组成部分。这些就是\Person{巴西利德}在\Place{埃及}居留后所详述的传说；受（该国智者）教导如此伟大的智慧体系，（这位异端者）结出了这样的果实。\par

% structure: p0053 source=h2 target=subsection reason=relative-heading-rank; base=h2

\subsection{第十六章 \Person{撒图尼鲁斯}的体系。}

但一位与\Person{巴西利德}大约同时期、却待在\Place{叙利亚}的\Place{安提约基雅}城的\Person{撒图尼鲁斯}，提出了与\Person{默南德尔}所主张的（任何）教义类似的看法。他断言有一位为所有人所不知的圣父——祂创造了天使、总领天使、宰制者、（和）掌权者；世界及其中万物是由某些天使——七位——所造的。而（\Person{撒图尼鲁斯}）主张，人是天使的作品。曾从绝对主权（之存在者）那里显现出一个光辉的形象；当（天使们）无法拘留这形象时，他说，因为它迅速向上返回，他们就彼此劝勉说：\Quote{让我们按照我们的肖像和样式造人。}\BibleRef{创 1:26}当形体被造出后，他说，由于天使的无能，它不能抬起自己，而是像虫一样扭动，上面的能力怜悯它，因为它按自己的形象而生，便发出一道生命火花，将人举起，使他有了活力。（\Person{撒图尼鲁斯}）断言这道生命火花在死后迅速返回同类的存在者；而其余部分，即它们由之产生的，则分解回归到那些之中。他以为救主是无生、无体、无形状的。（\Person{撒图尼鲁斯}）却（坚持）耶稣只是外表上显现为人。他说犹太人的天主是天使之一，并且因圣父想要剥夺所有元首的统治权，基督来是为了推翻犹太人的天主，并拯救那些信靠祂的人；这些人在他们内有生命火花。因为他断言天使形成了两种人——一种邪恶，另一种良善。并且，由于魔鬼时常帮助邪恶的人，（\Person{撒图尼鲁斯}）肯定救主来是为了推翻无价值的人和魔鬼，但为了拯救良善的人。他肯定婚姻和生育来自\OriginalTerm{撒殚}{Satan}。然而，大多数属于这（异端者）学派的人也戒绝动物食物，（并且）藉此苦行伪装（使许多人成为他们的受骗者）。他们坚持说，预言部分是由创造世界的天使所发，部分是由撒殚所发，他们认为这撒殚就是与宇宙的（天使）尤其与犹太人的天主敌对的那位天使。这些就是\Person{撒图尼鲁斯}的真实教义。\par

% structure: p0055 source=h2 target=subsection reason=relative-heading-rank; base=h2

\subsection{第十七章 \Person{马西昂}；他的二元论；他的体系源自\Person{恩培多克勒}；\Person{恩培多克勒}学说概要。}

但是，本都人马尔基昂比这些（异端者）更为狂乱，他略去了大多数（思辨者）的大部分主张，进入一种更无耻的教义，假设宇宙有两个起源因，声称一个是某个善的（本原），另一个是恶的。他自己想象在引进某些新颖的（意见），便建立了一个充满愚昧的学派，跟随者都是些感官生活的人，因为他本人也是好色之徒。这个（异端者）以为众人会忘记他并非基督的门徒，而是远在他之前的恩培多克勒的门徒，便炮制和形成了同样的意见——即宇宙有两个原因：不和与友爱。因为恩培多克勒关于世界计划说了什么？即使我们先前已经（就此主题）说过，但现在为了比较这位剽窃者的异端（与其源头），我们也不会沉默。\par

这位（哲学家）断言，构成世界并使之存在的所有元素有六个：两个是物质的，（即）土和水；两个是物质对象得以安排和改变的工具，（即）火和气；两个是通过工具作用于物质并塑造它的，（即）不和与友爱。（恩培多克勒）大致这样表达：—\par

万物的四个根源，你首先听吧：\newline{}光辉的朱庇特，生命赐予的朱诺，哈迪斯，\newline{}以及涅斯提斯，她用泪水浸湿了必死的源泉。\par

朱庇特是火，生命赐予的朱诺是土，它产生果实以维持生存；哈迪斯是气，因为虽然通过他我们看见万物，但他本身我们却看不见。涅斯提斯是水，因为这是（食物）的唯一载体，从而成为所有被滋养者的sustenance的原因；（但）它本身不能给予被滋养者营养。因为如果它确实能给予营养，他说，动物生命就永远不会因饥荒毁灭，因为水在世界中总是过剩的。因此他称涅斯提斯为水，因为（虽然间接地）是营养的原因，但它本身不能给予被滋养者营养。因此，这些——大致描绘——是包含恩培多克勒整个世界理论的诸本原：（即）水和土，产生生成物；火与精神，是工具和有效（原因）；但不和与友爱，是艺术性制造（宇宙的）（本原）。友爱是一种和平、一致和爱，其全部努力是使世界成为一个完成的整体。而不和则不断分割那个（整体），并细分它，或从一中造出多。因此不和是整个受造界的一个原因，他称之为\Quote{\OriginalTerm{毁灭性的}{oulomenon}}，即破坏性的。因为这不和的关切是，在每个时代，受造界本身应继续维持其现有状态。而毁灭性的不和一直是所有生成物产生的制造者和有效原因；而友爱则是将现存事物引向、改变和恢复到一个系统中的（原因）。关于这些（原因），恩培多克勒断言它们是两个不死且无生的本原，并且还没有获得存在的原因。（恩培多克勒）在某处（这样表达）：—\par

因为如果两者曾经存在，且将会存在；我想，\newline{}永恒的时代绝不会缺少两者之一。\par

（但）这些是什么？不和与友爱；因为它们并非开始存在，而是预先存在并永远存在，因为由于它们是无生的，它们不能遭受朽坏。但火、（水、）土和气是可朽坏并复活的。因为当这些生成物（物体）因不和而停止存在时，友爱抓住它们，将它们带出来，并将它们与她自身附上并联合于宇宙。（这发生）为了使宇宙保持为一，一直以同一方式被友爱有序地治理，并具有（不间断的）一致性。\par

然而，当友爱从多中造出一，并将分离的实体与一联合时，不和又强行将它们从一中分开，使之成为多，即火、水、土、气，以及由此产生的动物和植物，以及我们观察到的世界的任何部分。关于世界的形状，由友爱安排的是什么样子，（恩培多克勒）这样表达：—\par

因为从背脊不伸出两臂，\newline{}没有脚，没有敏捷的膝，没有生殖的腹股沟，\newline{}而是一个球，且它自身相等。\par

这种运作由友爱维持，并从多中造出世界的一个最美形式。而不和，作为宇宙各部分安排的原因，强行从那一（形式）中分开并造出多。这就是恩培多克勒关于他自己生成所断言的：—\par

\Quote{在这些中，我也是来自天主的流浪放逐者。}\par

即，（恩培多克勒）称天主为那一（形式）的合一与统一，其中（世界）存在于不和分离和产生之前，那些（事物）中多数（依据不和的安排而存在）。因为恩培多克勒断言不和是狂暴、混乱且不稳定的\OriginalTerm{德穆革}{Demiurge}，（从而）称不和为世界的创造者。因为这就是灵魂的定罪和必然性，不和强行将它们从合一中分开，并塑造和操作它们，（据恩培多克勒说，）他大致这样表达：—\par

谁在罪上叠加重誓，\newline{}而精灵获得延长的生命；\par

他借精灵指长寿的灵魂，因为它们不死且活很长年代：—\par

\Quote{因为三十万年从福乐中被放逐；}\par

将那些由\OriginalTerm{友爱}{\GreekText{Φιλία}}从大多数实体中收集到统一过程（源自理知世界）的（灵魂）称为有福的。他断言那些是流放者，且\par

随着时间的流逝，各式各样的必死之人生出，\newline{}改变着烦扰的生活方式\newline{}\par

他断言烦扰的方式就是灵魂进入（相继）身体的变化与变形。这就是他所说的：—\par

\Quote{改变着烦扰的生活方式。}\par

因为灵魂\Quote{改变}，身体一个接一个地变化，被\OriginalTerm{纷争}{\GreekText{Νεῖκος}}惩罚，且不被允许停留在同一个（形体）中，而是灵魂通过\OriginalTerm{纷争}{\GreekText{Νεῖκος}}从一个身体到另一个身体而遭受各种惩罚。他说：—\par

以太之力驱赶灵魂入海，\newline{}大海又将它们喷向地面，\newline{}地面又抛向灼热太阳的光芒，太阳将（灵魂）\newline{}掷入以太深处，且每一（灵）\newline{}取自另一，全都燃着恨意。\newline{}\par

这就是\OriginalTerm{巨匠造物主}{\GreekText{Δημιουργός}}施行的惩罚，恰如某个铁匠塑造（一块）铁，连续地将其从火中浸入水中。因为火就是以太，\OriginalTerm{巨匠造物主}{\GreekText{Δημιουργός}}从那里将灵魂转入海中；而陆地就是大地：因此他用了这些词——从水中到大地，从大地到空气。这就是（恩培多克勒）所说的：—\par

而地面又抛向\newline{}灼热太阳的光芒，太阳将（灵魂）\newline{}掷入以太深处，且每一（灵）\newline{}取自另一，全都燃着恨意。\newline{}\par

因此，按照恩培多克勒，灵魂这样被憎恨、折磨并在这个世界中受惩罚，却被\OriginalTerm{友爱}{\GreekText{Φιλία}}作为某种善（力量）收集起来，且（是）怜悯这些的呻吟，以及狂怒的\OriginalTerm{纷争}{\GreekText{Νεῖκος}}混乱邪恶的计谋（的）。同样（\OriginalTerm{友爱}{\GreekText{Φιλία}}）渴望并辛劳地将灵魂一点一点从世界领出，并使它们与统一驯顺，为要使万物被自己引导，达到统一。因此由于这分裂世界的毁灭性\OriginalTerm{纷争}{\GreekText{Νεῖκος}}的安排，恩培多克勒劝诫他的门徒戒除所有种类的动物食物。因为他断言动物的身体正是那些以受惩罚的灵魂住所为食的。他教训那些听信此类学说（他的）的人，要避免与妇女交往。（他颁布此训令）是为了使（他的门徒）不要与\OriginalTerm{纷争}{\GreekText{Νεῖκος}}所制造的那些作品合作并协助，\OriginalTerm{纷争}{\GreekText{Νεῖκος}}总是在分解并强力割裂\OriginalTerm{友爱}{\GreekText{Φιλία}}的作品。恩培多克勒断言这是宇宙管理的最伟大法则，他自己这样说：—\par

有物受命运支配，古老、\newline{}无尽的神明之律，以强大起誓封印。\newline{}\par

他称命运为从统一到多元的变化，按\OriginalTerm{纷争}{\GreekText{Νεῖκος}}，以及从多元到统一的变化，按\OriginalTerm{友爱}{\GreekText{Φιλία}}。而且，如我所言，（恩培多克勒）断言有四个可朽的神（即）火、水、地、气。然而他断言有两个不死、非受生、彼此持续敌对的（神），即\OriginalTerm{纷争}{\GreekText{Νεῖκος}}和\OriginalTerm{友爱}{\GreekText{Φιλία}}。而且（他断言）\OriginalTerm{纷争}{\GreekText{Νεῖκος}}常犯不义和贪婪，强夺\OriginalTerm{友爱}{\GreekText{Φιλία}}之物并归于自己。然而（他断言）\OriginalTerm{友爱}{\GreekText{Φιλία}}，由于总是恒定不变的某种善（力量）且致力于联合，召回、带向（自身）并还原为统一那些被\OriginalTerm{巨匠造物主}{\GreekText{Δημιουργός}}在创造中强力割裂、折磨和惩罚的宇宙部分。恩培多克勒为我们提出了这样一套关于世界生成、毁灭和构成（作为由善与恶组成的）的哲学体系。他还说存在一个可被理智认识的第三力量，这可以从这些（即\OriginalTerm{纷争}{\GreekText{Νεῖκος}}和\OriginalTerm{友爱}{\GreekText{Φιλία}}）中理解，他自己这样说：—\par

因为，若在橡木心中把这些真理固定，\newline{}并以清净的默观友善地观看它们，\newline{}这些中的每一个，随着时间的流逝，将萦绕你，\newline{}并有许多其他源于这些的，降下。\newline{}因这些将长入每一种习性，如本性所促；\newline{}但若你渴望别的事物，它们无数，滞留着\newline{}在人间毫不掩饰，钝化了忧虑之锋，\newline{}年岁流转，它们将迅速离你而去，\newline{}因它们渴望达至自己亲爱的族类；\newline{}因为知道一切拥有感知和一份心灵。\newline{}\par

% structure: p0082 source=h2 target=subsection reason=relative-heading-rank; base=h2

\subsection{第十八章。马西昂主义的来源；恩培多克勒重新被确认为该异端的提议者。}

因此，当玛尔奇翁或他的某只猎犬吠叫攻击德穆革，并从善恶的比较中引出论据时，我们应当对他们说：宗徒保禄和残缺手指的马尔谷都未曾宣布过这种（主张）。因为在这些学说中，没有一条写在《马尔谷福音》里。（这体系的真正作者）乃是恩培多克勒，墨托之子，阿格里根特人。而（玛尔奇翁）剥夺了这位（哲学家的思想），并想象至今仍能未被察觉地将他的整个异端体系，在相同的措辞下，从西西里转移到福音叙事中。因为请容忍我，噢，玛尔奇翁：正如你设立了善恶的比较，我今天也要设立一个比较，追随你自己的主张——你认为它们如此。你断言世界的德穆革是邪恶的——为何不羞愧地掩面，（以此）将恩培多克勒的学说教导给教会？你说有一位善神毁灭德穆革的工作：那么你不正是向你的门徒明白宣讲恩培多克勒的\Quote{友谊}作为善神吗？你禁止婚姻、生育儿女，以及戒食天主为信者及认识真理的人所造、可以享用的食物。\BibleRef{弟前 4:3}（那么你认为，）你能逃避察觉，（同时）命令恩培多克勒的净化礼仪？因为事实上你在每一点上都追随这位（外教哲学家），同时教导你自己的门徒拒绝肉类，以免吃到任何可能是一个被德穆革惩罚的灵魂残骸的身体。你解散天主所结合的婚姻。在这里你又符合恩培多克勒的主张，以便为你，友谊的工作可以保持为一（且）不可分割。因为，按照恩培多克勒，婚姻分裂统一，并从其中产生多元，正如我们已经证明的。\par

% structure: p0084 source=h2 target=subsection reason=relative-heading-rank; base=h2

\subsection{第十九章 普雷彭的异端；追随恩培多克勒；玛尔奇翁拒绝救主的诞生。}

玛尔奇翁的主要异端，以及它最纯正（不混杂其他异端）的一种，是由善恶（之神）理论构成的体系。我们已经表明，这属于恩培多克勒。但是，既然在我们这个时代，有一位玛尔奇翁的追随者——亚述人普雷彭——试图引入某些更新奇的东西，并在题献给亚美尼亚人巴尔达桑的一部著作中说明了他的异端，我也不将对此保持沉默。普雷彭声称，正义构成第三本原，且被置于善恶之间，他显然无法避免（灌输）恩培多克勒的观点。因为恩培多克勒断言，世界由邪恶的纷争治理，而另一个（由友谊治理的）世界是理智可认知的。（他断言）这些是善恶的两个不同本原，且在这两个不同本原之间的是公正理性，按照它，被纷争分离的事物得以结合，（并）按照友谊的影响，被整合为一。那公正理性本身，即友谊之助者，恩培多克勒称之为\Quote{缪斯}。他自己也恳求她帮助他，并表达如下：—\par

\Quote{因为如果对于短暂必死的凡人，不朽的缪斯，\newline{}您的关怀是让思绪占据我们的心意，\newline{}卡利俄珀，再次垂听我现今的祈祷，\newline{}当我揭示关于幸福诸神的纯净叙述。}\par

玛尔奇翁采纳了这些观点，完全拒绝了我们救主的诞生。他认为，认为作为友谊之助者的圣言——即善神——应属于由破坏性纷争所造之受造物的（范畴）是荒谬的。然而，（他的教义是）独立于出生，祂自己在提庇留凯撒在位第十五年从上降下，并作为介于善恶神之间的中保，开始在会堂里教导。因为如果祂是中保，他说，祂已从恶神的整个本性中被释放出来。此外，他断言德穆革及其工作是邪恶的。为此原因，他断言，耶稣无生地降下，以便祂可以从一切（邪恶的）混杂中被释放。祂也已，他说，从善神的本性中被释放，以便祂可以成为中保，如保禄所说（\BibleRef{迦 3:19}），也如祂自己承认：\Quote{你为什么称我为善？只有一位是善的，}这些就是玛尔奇翁的观点，他借此使许多人受骗，采用了恩培多克勒的结论。他将那位（古代思辨者）发明的哲学转移到他自己的思想体系中，（从恩培多克勒）构建了他（自己）的邪恶异端。但我认为这已被我们充分驳斥，并且我没有忽略任何那些从希腊人那里窃取观点并对基督的门徒行恶之人的意见，好像他们已成为这些（教义）的导师给门徒。但既然我们似乎已充分解释了这位（异端者）的教义，让我们看看卡波克拉特怎么说。\par

% structure: p0088 source=h2 target=subsection reason=relative-heading-rank; base=h2

\subsection{第二十章 卡波克拉特的异端；关于耶稣基督的邪恶教义；施行魔法；采纳灵魂转世。}

\Person{卡尔波克拉特斯}肯定说，世界及其中的事物是由天使造的，天使远不及未生之圣父；而\Person{耶稣}是\Person{若瑟}所生，祂与（其他）人相似地出生，但比（其余人类）更正义。而且（\Person{卡尔波克拉特斯}断言）\Person{耶稣}的灵魂，因它变得坚强无玷，就记起了它在与未生之天主的交往中所见的事物。并且（\Person{卡尔波克拉特斯}主张）因此，由那（天主）有一德能降在（\Person{耶稣}）身上，为的是藉着它，祂能逃脱造世（天使）。而（他说）这德能经过了万有，在万有中获得了自由，又升到（天主）自己那里。并且（他声称）与（\Person{耶稣}的灵魂）同样状态的是所有那些拥抱与（刚才提到的德能）相似欲望对象的（灵魂）。他们断言，\Person{耶稣}的灵魂，虽然按法律曾在犹太风俗中受训练，（实际上）却鄙视它们。而且（他说）因此（\Person{耶稣}）领受了能力，藉以祂使为惩罚人而附着于人的情欲变为无效。而（他论证）因此，那与\Person{基督}的灵魂相似地能够鄙视造世执政者的（灵魂），同样领受能力去行相似的事。因此，也（根据\Person{卡尔波克拉特斯}，有人）竟达到如此骄傲的程度，以至声称他们中有些人与\Person{耶稣}自己平等，而另一些人甚至更有能力。但（他们还争辩说）有些人享有超过那（救主）的门徒的卓越，例如\Person{伯多禄}和\Person{保禄}，以及其余宗徒，而这些人在任何方面都不亚于\Person{耶稣}。而且（\Person{卡尔波克拉特斯}断言）这些人的灵魂起源于那至高的德能，因而他们，同样鄙视造世（天使），被认为配得同样的能力，并（有特权）升到同一（地方）。然而，若有人比那（救主）更鄙视世俗之事，（\Person{卡尔波克拉特斯}说）这样的人就能成为超越（\Person{耶稣}）的。这异端的追随者）行他们的魔术和咒文、符咒和纵欲的宴饮。并且（他们惯于召唤）属下的鬼魔和送梦者，以及（诉诸）其余（巫术）的骗术，声称他们拥有能力，现在就能获得对执政者和这世界的制造者，甚至对其中一切作为的统治权。\par

（现在这些异端者）本身是由\Person{撒殚}派遣来的，目的是在外邦人面前诽谤教会的神圣之名。（魔鬼的意图是）使人听到那些（异端者）的教义，时而是这种方式，时而是那种方式，以为我们所有人都是同一种人，就掩耳不听真理的宣讲，或者他们（不加摈弃地）看那些（异端者）的一切主张，也对我们加以诽谤。（\Person{卡尔波克拉特斯}的追随者）声称灵魂从身体转移到身体，直到它们满全所有罪的（量度）。然而，当没有留下（任何罪）时，（卡尔波克拉特派肯定说灵魂）就得了释放，并离开到那在造世天使之上的天主那里，并且这样所有的灵魂都将得救。然而，若有些（灵魂），在灵魂居于身体的一次生命中，就预先卷入了全部过犯的量度，它们（根据这些异端者）就不再经历灵魂转世。（这类灵魂）然而，一旦立即清偿所有过犯，就（卡尔波克拉特派说）从再居于身体中得释放。这些（异端者）中有些人还在自己门徒的右耳垂后面打上烙印。他们制造\Person{基督}的假像，声称这些是在（我主在世时）存在的，且是（由）\Person{比拉多}（制作的）。\par

% structure: p0091 source=h2 target=subsection reason=relative-heading-rank; base=h2

\subsection{第二十一章 \Person{克林妥}关于\Person{基督}的体系}

但有一个\Person{克林妥}，本人受教于\Place{埃及}人的学说，断言世界不是由原初神性造的，而是由某种德能造的，这德能是从那高于万有的大能分出的一枝，而（它）却不认识那高于万有的天主。他设想\Person{耶稣}不是由童贞女所生，而是\Person{若瑟}和\Person{玛利亚}的儿子，与其余人完全相同地出生，且（\Person{耶稣}）比（全人类）更正义更智慧。并且（\Person{克林妥}声称）在（我主）受洗之后，\Person{基督}以鸽子形式降在祂身上，来自那高于万有的绝对主权。然后，（根据这异端者）\Person{耶稣}开始宣讲未知的圣父，\BibleRef{宗 17:23}并为证实（祂的使命）而行奇迹。但是，（\Person{克林妥}的意见是）最终\Person{基督}离开了\Person{耶稣}，而\Person{耶稣}受难并复活；而\Person{基督}，因为是属神的，保持在不可能受苦的状态。\par

% structure: p0093 source=h2 target=subsection reason=relative-heading-rank; base=h2

\subsection{第二十二章 厄彼翁派的教义}

厄彼翁派却承认世界是由那真实的天主创造的，但他们关于\Person{基督}所提出的传说却与\Person{克林妥}和\Person{卡尔波克拉特斯}相似。他们按犹太人的习俗生活，声称他们按法律成义，并说\Person{耶稣}是因遵行法律而成义的。因此，（根据厄彼翁派，）（救主）被称为天主的\Person{基督}和\Person{耶稣}，因为其余（人类）中没有一人完全遵守了法律。因为即使另有他人遵行了法律中所含的诫命，那人也会是那\Person{基督}。并且（厄彼翁派声称）他们自己，当同样遵行（法律）时，也能成为\Person{基督}；因为他们断言我们的主自己是一个人，与所有（其余人类）意义相同。\par

% structure: p0095 source=h2 target=subsection reason=relative-heading-rank; base=h2

\subsection{第二十三章 \Person{特奥多图}的异端}

但有一位特奥多图斯，原籍拜占庭，他引入了一种新异的异端。他关于宇宙本原的教义，部分与真教会之道理一致，因为他承认万物皆由天主所造。然而，他从诺斯替派以及则林托和厄彼雍那里强行挪用关于基督的概念，声称（我主）显现的方式正如我将要描述的那样。（据此，特奥多图斯主张）耶稣是一个（纯粹的）人，由童贞女所生，依照圣父的旨意，在祂与众人共同生活、并成为卓越虔敬之人后，随后在约但河领受圣洗时，领受了来自高处、以鸽子形状降下的基督。这就是为何（根据特奥多图斯）在降下的圣神——即宣告祂为基督的那位——显现之前，祂内没有运行异能的原因。但在特奥多图斯的追随者中，有人倾向于认为，即使在圣神降临时，这个人也从未成为天主；而另一些人则主张，他在死后复活后才成为天主。\par

% structure: p0097 source=h2 target=subsection reason=relative-heading-rank; base=h2

\subsection{第二十四章 默基瑟德派；尼苛劳派}

然而，在他们中间出现了不同的争论，一位同样被称为特奥多图斯、职业为钱币兑换商的异端者试图确立如下教义：某位默基瑟德掌握着最大的权能，并且这位（默基瑟德）大于基督。他们进一步宣称，基督只是按默基瑟德的肖像而形成。他们自己，与前面提到的特奥多图斯追随者一样，断言耶稣是一个（纯粹的）人，依照相同的说法，基督降临到了祂身上。\par

然而，在诺斯替派中有多种意见分歧；但我们觉得不值得一一列举这些异端者的愚蠢教义，因为它们数量众多、毫无理性且充满亵渎。即使那些在神性方面较为严肃的异端者，且其思辨体系衍生自希腊人，也必须被定罪。尼苛劳是这些恶者广泛勾结的原因。他作为被宗徒选为执事职的七人之一。但尼苛劳偏离了正确的教义，惯于教导对生活与食物抱持漠不关心的态度。当尼苛劳的门徒继续侮辱圣神时，若望在\WorkTitle{默示录}中责备他们为行淫者和吃祭偶像之物的人。\par

% structure: p0100 source=h2 target=subsection reason=relative-heading-rank; base=h2

\subsection{第二十五章 切尔东异端}

但有一位切尔东，同样从这些（异端者）和西满那里获得契机，断言由梅瑟和先知所宣讲的天主并非耶稣基督的父。因为（他争辩说）这位（父）是已知的，而基督的父是未知的；前者是公义的，后者是良善的。玛尔基翁在其试图撰写的、称作\WorkTitle{Antitheses}的著作中，证实了这位（异端者）的教义。他惯于在这部书中，向宇宙的创造主发出任何他心中存想的诽谤。类似地，切尔东的门徒路恰诺也如此行。\par

% structure: p0102 source=h2 target=subsection reason=relative-heading-rank; base=h2

\subsection{第二十六章 阿培肋的教训；他的女先知斐禄美纳}

但阿培肋源自这些人，如此表达自己，说是有一位善神，正如玛尔基翁所假定的，而创造万物的那一位是公义的。如今，据阿培肋说，他是受造物的造物主。这位异端者还主张有第三位神，即惯常与梅瑟说话的那一位，这位是火性的；还有第四位神，是诸恶的原因。但阿培肋称这些为天使。他诽谤法律和先知，声称所写的内容属于人类且是虚假的。阿培肋从福音书或宗徒保禄的著作中选取他所喜欢的内容。但他致力于某个斐禄美纳的言论，视之为一位女先知的启示。然而，他肯定基督是从高天之上的权能，即从善神那里降下，并且他是那善神的儿子。他断言耶稣并非由童贞女所生，且在显现时并非无肉身。他主张基督从宇宙的实体中取部分构成了他的战利品，即热与冷、湿与干。他说，基督带着这样的身体领受宇宙力量，在世界上生活了一段时期。但耶稣后来被犹太人钉在十字架上，断了气，三天后复活，显现给门徒。救主向他们展示钉痕和肋旁，想要说服他们祂确实不是幻影，而是以肉身临在。阿培勒说，在向他们显示祂的肉身后，救主将肉身归还给大地，因为肉身正是从大地的实体而来。祂这样做是因为祂不贪恋属于他人的任何东西。虽然耶稣暂时可使用他人的东西，但祂在适当的时候将各元素特有的部分归还给了它们。因此，当祂再次解开祂身体的束缚时，祂将热归还给热，冷归还给冷，湿归还给湿，干归还给干。在此状态下，主去到良善的圣父那里，在世上为那些通过祂的门徒而相信祂的人留下生命的种子。\par

在我们看来，这些教义已经得到了充分的说明。然而，既然我们决定不留下任何由异端者提出的观点不予驳斥，让我们也看看幻象论派发明了什么体系。\par

\SourceRecord{Translated by J.H. MacMahon.；Ante-Nicene Fathers；Vol. 5.；Edited by Alexander Roberts, James Donaldson, and A. Cleveland Coxe.；Buffalo, NY: Christian Literature Publishing Co.,；1886.；<http://www.newadvent.org/fathers/050107.htm>.}

% structure: p0001 source=h1 target=section reason=page-title

\section{驳斥一切异端（卷八）}

% structure: p0002 source=h2 target=subsection reason=relative-heading-rank; base=h2

\subsection{目录。}

以下是\WorkTitle{驳斥一切异端}第八卷的目录：\newline{}\par

\OriginalTerm{德西特派}{Docetae}的主张是什么，以及他们如何从自然哲学中形成他们所宣称的教义。\par

\OriginalTerm{摩诺伊穆斯}{Monoïmus}如何信口开河，专注于诗人、几何学家和算术家。\par

\OriginalTerm{塔提安}{Tatian}的（体系）如何源自\OriginalTerm{瓦伦提尼}{Valentinus}和\OriginalTerm{马西昂}{Marcion}的主张，以及这位异端者（由此）如何形成他自己的教义。\OriginalTerm{赫尔摩该尼}{Hermogenes}则利用苏格拉底的学说，而非基督的学说。\par

那些主张在第十四天守逾越节的人是如何错误的。\par

\OriginalTerm{弗吕吉亚人}{Phrygians}的错误是什么，他们以为\OriginalTerm{孟他努}{Montanus}、\OriginalTerm{普黎斯奇拉}{Priscilla}和\OriginalTerm{马克西米拉}{Maximilla}是先知。\par

\OriginalTerm{恩克拉提派}{Encratites}的自负是什么，以及他们的见解并非来自圣经，而是出自他们自身，以及印度的\OriginalTerm{裸形哲学家}{Gymnosophists}。\par

% structure: p0010 source=h2 target=subsection reason=relative-heading-rank; base=h2

\subsection{第一章 先前所驳斥的异端；德西特派的主张。}

既然（异端者）的大多数人不采用主的劝告，眼中有梁木（\BibleRef{玛 7:3}），却宣称自己能看见，实则劳苦于盲目，我们认为绝不可对这些人的主张缄默不言。我们的目的是：通过我们所完成的驳斥，（异端者）因自觉羞愧，或许能被引导认识救主曾劝告（人）先取出梁木，然后才清楚看见你兄弟眼中的木屑。因此，在之前的七卷书中，我们已经充分并足够地解释了大多数（异端者）的教义，现在也不应对随后（由此而来）的（异端）意见保持沉默。我们要借此显示圣神恩宠的丰盈；并且要驳斥那些（自认为）已获得教义坚固性的人，而实际上那只是表面。这些人自称为德西特派，并宣扬以下主张：\newline{}\par

（德西特派主张）天主是原初（存在），如同一粒无花果树的种子，体积甚小，但能力无限。（据德西特派称，这种子构成）一种卑微的度量，在数量上不可计数，在生成方面毫无欠缺。（这种子是）恐怖者的避难所，赤裸者的遮盖物，羞怯的帷幔，以及所寻求的果实，他（指救主）三度前来寻找（果子），但未找到。因此，他说，他诅咒了无花果树，因为他在其上没有找到那甘甜的果实——所寻求的产出。由于按他们的说法——简言之——神性是这样且如此广大，即微小而细微，世界按他们的看法，是按以下方式形成的：当无花果树的枝条变得柔嫩时，首先发芽（如通常所见），然后依次结果。在此（果实）中，保存着无花果树无限且不可计数的种子。因此，（德西特派说）我们认为，从无花果树的种子最初产生三部分：树干（构成无花果树）、叶子和果实——即无花果本身，如前所述。同样，（德西特派）断言，产生三个\OriginalTerm{永世}{Aeon}，它们是从宇宙的原初根源而来的本原。梅瑟对此并未保持沉默，他说天主有三个话语：\Quote{黑暗、密云、暴风，且不再加添。}（\BibleRef{申 5:22}）因为（德西特派）说，天主对那三个永世没有增加任何东西；但它们在各方面已经足够为那些被生成者所用，而且仍然足够。天主自己则与自己同在，远远地与那三个永世分离。当每个永世获得了生成的本原后，就如所述，逐渐增长，放大，最终成为完美的。但他们认为完美的是以十为度。因此，当永世在数目和完美上相等时，按（德西特派）的意见，它们总共构成三十个永世，每个在十全中达成完全的完美。这三个永世相互有别，彼此之间享有同等的（尊荣），仅在位置上不同，因为其中一个为第一，另一个为第二，第三个为第三。位置赋予它们能力的差异。因为那获得最接近原初神性（如同一粒种子）位置的，拥有比其余更大的生产能力，因为他自己，即无量者，以十倍体积量度了自己。而那在位置上排第二的，由于他是不可领悟者，以六倍领悟了自己。但排第三的，因其兄弟们的扩张，被带到无限的距离。（当第三个永世）三次在思想中实现自己后，他用某种永恒的联合环围绕了自己。\par

% structure: p0013 source=h2 target=subsection reason=relative-heading-rank; base=h2

\subsection{第二章 德西特关于降生成人的观念；他们的永世教义；他们对创造的描述；他们对火的天主的观念。}

这些（异端者）以为这就是救主所说的：\BibleQuote{有一个撒种的出去撒种；那落在好地里的，结了一百倍、六十倍、三十倍的果实。}因此，幻身论者说，（救主）说了这句话：\BibleQuote{有耳可听的，就应当听！}因为这些（真理）并非全然是传闻。所有这些永世，无论是三个，还是从那无限中以不可计量的方式产生的无限多的永世，都是雌雄同体的永世。所有这些永世，在从那个原种增长、壮大并生出之后，都怀有和谐与合一的精神，并且它们全都融合成一个永世。就这样，他们由一位童贞玛利亚生出一个共同的产物，即一位中保，也就是所有处于（盟约）中保之中者的救主。这位（救主）在各方面都与无花果树的种子权能相等，唯独他是受生的。而那个原种，无花果树从它而生，却是非受生的。因此，当那三个永世以及那个独生子——因为他是由那无限永世从直接参与他诞生的三个永世所生的，三个不可测量的永世同心合意地生了他——被装饰以一切德行和一切圣洁之后，（整个可理知的本性就被造就成毫无缺欠的了。现在，所有那些可理知且永恒的本体都构成了光。然而，光并非无形、无效，或如同需要任何（其他能力）的帮助。（相反地，）光与那无限（永世）的众多数目成比例，它们照着无花果树的样式无限地（产生），在自身中拥有无数种居住在创造之地的各类动物，并且照耀到下面的混沌之上。当混沌同时被光照，并由那些来自上面的多样种类赋予形式时，它就获得了坚固性，并从那使自己成为三重的第三永世获得了所有那些上层的种类。\par

这第三永世，却看见他自己所有独特的属性被下面的黑暗所整体攫取，并非不知道黑暗的能力，同时也知道光的安全与丰盛，就没有允许他来自上面的光辉属性长久地被下面的黑暗夺走。相反地，他采取了相反的行动：他将黑暗置于永世之下。然后，他在下界之上造了穹苍，\BibleQuote{就把黑暗从光中分开，并称光为昼，称暗为夜。}正如我所说，当第三永世的所有无限种类被阻截在这最下面的黑暗中时，这永世本身的形象，正如他所描述的那样，也连同其余属性一同被印在上面。（这个形象是）一种从光而生的生命之火，大元首即由此产生。关于这位（元首），梅瑟说：\BibleQuote{在起初，天主创造了天地。}\BibleRef{创 1:1} 梅瑟提到\BibleRef{出 3:2}这位火神曾从荆棘——即黑暗的空气中说话。因为整个黑暗之下的空气，作为传递光的媒介。现在幻身论者说，梅瑟使用 \OriginalTerm{荆棘}{batos} 一词，是因为所有光的种类从上经过，借着空气作为传递媒介而下降。并且，从荆棘向我们说的上主的话，同样可以被认知；因为声音，作为意义的表达，是空气的回响，没有空气，人的言语就无法被认知。不仅从荆棘向我们说的话语——即空气——立法并与我们同作公民；而且，气味和颜色也借着空气向我们显明它们自身的特性。\par

% structure: p0016 source=h2 target=subsection reason=relative-heading-rank; base=h2

\subsection{第三章　基督解除德穆革的作为；幻身论者关于耶稣受洗与死亡的记述；为何他在世上生活了三十年。}

这个火之神，在从光变成火之后，就按照\Person{梅瑟}所描述的方式着手创造世界。然而，他自身缺乏实质，以黑暗为（他的）质料，并不断地侮辱那些来自上界的、被下界（黑暗）所捕获的光的永恒属性。因此，直到救主显现之时，由于火之光的神性（即德缪格），一种广泛的灵魂幻象盛行。因为那些形态被称为灵魂，它们是从上界的（永世）冷却而来，并停留在黑暗中。但当（灵魂）从一个身体转移到另一个身体时，它们仍受德缪格的监护。幻象说者说，这些事也可以从\Person{约伯}的话中看出，他说了以下的话：\Quote{我是一個流浪者，不斷更換地方和房屋。} 并且（根据幻象说者，我们也可以从）救主的话中（学到同样的道理）：\BibleQuote{如果你愿意接受，他就是那要来的\Person{厄里亚}。有耳可听的，就应当听。} \BibleRef{玛 11:14} 但藉着救主的工具，这种灵魂从一个身体到另一个身体的转移停止了，并且为了赦免罪恶而宣讲信德。以类似的方式，那独生子，当他注视从上界转移到昏暗身体中的上层永世的形式时，降下来，希望下降并拯救它们。然而，当（圣子）意识到那些集合的永世不能看见所有永世的圆满，而是惊慌失措，害怕自己会朽坏而遭受败坏，并且被能力的伟大和辉煌所压倒时——（当圣子，我再说，察觉了这一点，）他收缩了自己——如同在一个极小的身体中的一道极大的闪光，更确切地说，如同压缩在眼睑下的视觉光芒，并且（在这种状态下）他前进到天空和明亮的星辰。在这个受造物的区域，他再次按照自己的意愿将自己收集在视觉的眼睑之下。现在视觉之光完成同样的效果；因为它虽然无处不在，并使一切事物可见，但对我们来说却是不可感知的。我们只是看到视觉的眼睑，而眼角（的）组织宽大、曲折、极其纤维化，是角膜膜；其下是瞳孔，形状像浆果，网状且圆形。（我们还观察到）其他属于眼睛之光并包裹着它的膜。\par

因此，幻象说者说，那从上界来的独生永恒之子，以与三个永世中每个永世相对应的形式着装；当他处于三十永世之中时，他像我们描述的那样进入这个世界，不被注意、不为人知、隐蔽且不被相信。因此，幻象说者说，为了让他穿上在更远的受造界区域盛行的黑暗——（这里的黑暗他指）肉身——一位天使与他一起从上界下来，向\Person{玛利亚}报喜，正如所记载的那样。并且从她所生的（孩子）出生了，正如所记载的那样。那从上界来的穿上了那所生的；他做了一切事，正如福音中所记载的（关于他的）。他在\Place{约但河}中洗了澡，当他受洗时，他在那（另一属灵战利品的）水中接受了形象和印记，除了那由童贞玛利亚所生的身体。（其目的是）当执政官将自己特有的（肉身的）捏造品判处死刑，即十字架时，那在（童贞玛利亚所生的）身体中受滋养的灵魂可以脱去那身体，并把它钉在（受咒诅的）木架上。（这样，灵魂）藉着这（身体）胜过掌权者和权势者，并且不会赤身露体，而是代替那肉身穿上另一个身体，就是他受洗时在水中所代表的。幻象说者说，这就是救主所肯定的：\BibleQuote{除非人由水和圣神而生，不能进天主的国，因为从肉身生的就是肉身。} \BibleRef{若 3:5} 因此，从三十永世中，（圣子）取了三十种形式。为此原因，那永恒者在地上存在了三十年，因为每个永世以一种特殊的方式在（其）一年中显现。灵魂就是那被每个三十永世所捕获的所有那些形式；每个灵魂被这样构造，以至于能辨认\Person{耶稣}，他的本性（与他们类似）。（而正是这个\Person{耶稣}的本性）那独生永恒者从永恒之处所取的。然而，这些（处所）是多样的。因此，相应数量的异端，以最大的热忱，寻找\Person{耶稣}。现在所有这些异端都有自己的特殊\Person{耶稣}；但他被不同地看见，根据他所说的每个灵魂被带往和赶往的不同处所而定。（每个灵魂）认为（从其特定处所看见的\Person{耶稣}）是唯一那与自己同族同城的\Person{耶稣}。并且在第一次看见（这\Person{耶稣}）时，那灵魂认出他是自己特殊的兄弟，而其余的是私生子。因此，那些从下层处所得着本性的，不能看见在他们之上的救主的形式。然而，他说，那些从上界来的，从中间十元组和最卓越的八元组来的——幻象说者说，我们就是从那来的——他们自己认识耶稣救主，不是部分地，而是完全地。那些从上界来的，唯独是完美的，而其余的都只是部分性的。\par

% structure: p0019 source=h2 target=subsection reason=relative-heading-rank; base=h2

\subsection{第四章　幻象说者的学说源自希腊诡辩家}

这些（陈述），因此，我认为对于有适当心智的人来说，足以达到认识\OriginalTerm{幻影派}{Docetae}那复杂而不稳定的异端。但那些提出关于不可及和不可理解之物质的论证尝试的人，自称幻影派。现在，我们认为其中有些人行事愚蠢，我并不是说在表面上，而是在实际上。至少，我们已经证明，来自这种物质的一道梁木被携带在眼中，如果他们无论如何能够看到它的话。然而，如果他们看不见（它，我们的目的是）他们不应使别人瞎眼。但事实是，古代希腊的诡辩家在许多细节上早已设计了这些（幻影派）的教义，正如我的读者（费心）所能查明的。这些就是幻影派提出的见解。但同样地，关于\Person{摩诺伊穆斯}的教义，我们也不会保持沉默。\par

% structure: p0021 source=h2 target=subsection reason=relative-heading-rank; base=h2

\subsection{第五章 \Person{摩诺伊穆斯}；\Person{摩诺伊穆斯}认为人是宇宙；他的单子体系。}

\Person{阿拉伯人摩诺伊穆斯}远离了那位高调诗人的荣耀。（因为\Person{摩诺伊穆斯}）认为有某个这样的人，正如那位诗人（所称呼的）\Person{俄刻阿诺斯}，以某种方式如此表达：——\par

\Quote{俄刻阿诺斯，诸神之源和人之源。}\par

\Person{摩诺伊穆斯}说，人是宇宙。现在宇宙是万物的起因，非受生，不朽，（且）永恒。而（他）说，先前所提到的（那）人的儿子是受生的，且具有可受难的属性，（并且）他是独立于时间、无意地、且没有被预定而生成的。因为，他说，那就是那人的能力。而当他这样在能力中构成时，（\Person{摩诺伊穆斯}宣称）那儿子比思想和意愿更快地诞生。而这，他说，就是\WorkTitle{圣经}中所说的：\Quote{他曾是，并被生成。}而它的意思是：人曾存在，他的儿子被生成；正如有人可以说，火曾存在，而独立于时间、无意地、且没有被预定，光与火的存在同时生成。而这个人构成一个单一的\OriginalTerm{单子}{monad}，它是非复合而不可分的，（但同时）复合的（和）可分的。（这个单子）在一切方面友善且和平，在一切方面争吵且与自己敌对，不相似而相似。（这个单子）也如同某种音乐和谐，它自身包含万物，正如一个人所能表达或忽略而未考虑的；它显明一切，并生成一切。这是母亲，这是父亲——两个不朽的名字。然而，作为一个类比，他说，考虑那完美者的最大图像，那一个\OriginalTerm{约塔}{jot}——那一个\OriginalTerm{点}{tittle}。而这一个点是一个非复合、单纯、纯粹的单子，它根本不从任何事物获得构成。（然而这一点也是）复合的、多形的、分成许多部分、由许多部分组成的。那一个不可分的点，他说，是约塔（字母）的一点，具有许多面孔、无数眼睛、数不清的名字，而这（点）是那完美不可见之人的图像。\par

% structure: p0025 source=h2 target=subsection reason=relative-heading-rank; base=h2

\subsection{第六章 \Person{摩诺伊穆斯}的\Quote{约塔}；他关于\Quote{人子}的概念。}

因此，他说，单子（即那一点）也是一个十进制。因为通过这一点本身的动力，产生了二元、三元、四元、五元、六元、七元、八元、九元，直至十。因为这些数字，他说，能够被许多方式分割，并且它们居住在那个简单而不可分的一点（约塔）之中。而这就是所宣告的：\Quote{（天主）乐意使一切圆满有形地居住在\Person{人子}内。}因为这种从简单而不可分的一点（约塔）而来的数字组合，他说，成为有形体的实在。因此，他说，\OriginalTerm{人子}{Son of man}是从那完美之人（没有人认识他）所生成的；然而，每一个不知晓那儿子的受造物，都将祂想象为女人的后代。而这个人子的某些非常模糊的光芒，接近这个世界，控制和掌管变化（与）生成。而那个人子的美直到现在仍对所有人不可理解，正如许多受欺骗的人关于女人的后裔那样。因此，他说，在我们这区域的创造物中，没有一样是由那个人产生的，也不会有一样（从祂）将会（被生成）。然而，万物都是从那个人子的一部分、而非从整体所生成的。因为他说，人子是一个点中的一点（约塔），它从上而来，圆满而完全地充满所有（从上而下的光芒）。而它自身包含了那个人（即人子之父）所拥有的一切事物。\par

% structure: p0027 source=h2 target=subsection reason=relative-heading-rank; base=h2

\subsection{第七章 \Person{摩诺伊穆斯}论安息日；将\Person{梅瑟}的棍杖寓意化；关于十诫的看法。}

世界，如梅瑟所说，是在六天内造成的，即通过六种能力，这些能力内含于一个\OriginalTerm{约塔的笔划}{one tittle of the iota}中。第七天，即安息和安息日，是由掌管地、水、火、气的\OriginalTerm{七元组}{Hebdomad}产生的。世界由这些元素通过一个笔划形成。因为立方体、八面体、金字塔以及所有类似的图形（火、气、水、地由它们构成）来源于包含在那个简单约塔笔划中的数字。这个笔划构成了一个完美之人的完美儿子。因此，当梅瑟提到杖为了埃及的灾祸而被变化地挥舞时——这些灾祸，他说，是创造之寓意的象征——他没有为多于十个的灾祸塑造杖。这个杖构成一个约塔的笔划，是双重的和多样的。这十个灾祸的序列，他说，是尘世的创造。因为所有事物通过被打击而产生并结果实，就像葡萄树一样。人，他说，迸发而出，并通过被某种打击切割而与另一个人强行分离。这是为了使人能够被产生，并宣布梅瑟所立的律法，这律法是梅瑟从天主领受的。与那个一个笔划一致，律法构成十诫的系列，这系列寓意地表达了那些诫命的神圣奥秘。因为，他说，关于宇宙的全部知识都包含在与十个灾祸的序列和十诫系列相关的事物中。没有人认识这知识，如果他是那些被关于女人后裔的谎言所欺骗的人中的一员。然而，如果你说\WorkTitle{五经}构成全部律法，它来自于包含在一个笔划中的\OriginalTerm{五元组}{Pentad}。但对于那些尚未完全领悟奥秘的人来说，全部律法是一个新的、非过时的、合法的、永久的节日，是主天主的逾越节，要由那些能辨别这奥秘的人在我们世代遵守，始于十四日，他说，那是十进制的开始，他们从那里计算。因为一元直到十四是那个完美数字的一个笔划的总结。因为一、二、三、四成为十；这就是一个笔划。但从十四到二十一，他断言有一个七元组内在于世界的一个笔划中，并在所有这些中构成一个无酵的受造物。因为，他说，一个笔划为何需要来自外部的像酵母一样的实物来用于主逾越节，这永恒的节日，它被赐予世世代代？因为整个世界和创造的所有原因构成一个逾越节，即主的节日。因为天主因创造的转变而喜乐，这通过一个笔划的十次打击而实现。这个笔划是梅瑟的杖，由天主交在梅瑟手中。梅瑟用这杖击打埃及人，以改变物体——例如，将水变为血；其他物质也类似——例如蝗虫，它是草的象征。他以此意指元素转变为肉体；他说：\BibleQuote{凡有血肉的，如同草。}\BibleRef{依 40:6}这些人甚至以类似的方式接受全部律法；我认为他们很可能采纳了那些希腊人的观点，即肯定存在实体、性质、数量、关系、地点、时间、姿态、动作、拥有和承受。\par

% structure: p0029 source=h2 target=subsection reason=relative-heading-rank; base=h2

\subsection{第八章：\Person{莫诺伊慕斯}在致\Person{特奥夫拉斯图斯}信中解释其观点；何处寻找天主；其体系源自毕达哥拉斯。}

\Person{莫诺伊慕斯}本人，在致\Person{特奥夫拉斯图斯}的信中，明确地陈述如下：\Quote{省略寻求天主和创造以及类似事物，从你自己内寻求祂，并学习那绝对将你内一切归于己有者是谁，祂说：\InnerQuote{我的天主是我的理智、我的理解、我的灵魂、我的身体。}并学习从何处来悲伤、喜乐、爱、恨、不自愿的苏醒、不自愿的睡意、不自愿的愤怒、不自愿的情感；如果你仔细调查这些，他说，你将发现祂自己，单一与多元，在你自己内，根据那个\OriginalTerm{笔划}{tittle}，并且祂发现（神性）的出路由你自己而出。}那些异端者，因此，已做了这些陈述。但我们没有必要将这样的教义与之前希腊人沉思过的内容相比较，因为这些异端者提出的主张显然源自几何和算术技艺。然而，毕达哥拉斯的门徒们以一种更卓越的方法阐述了这门技艺，我们的读者可以通过查阅我们之前提供希腊人全部智慧阐释的那些段落来确认。但既然\Person{莫诺伊慕斯}的异端已被充分反驳，让我们看看其余异端者也虚构了哪些学说，他们渴望为自己建立一个虚名。\par

% structure: p0031 source=h2 target=subsection reason=relative-heading-rank; base=h2

\subsection{第九章：\Person{塔提安}。}

\Person{塔提安}，虽然他自己是殉道者\Person{犹斯定}的门徒，却未持有与他老师相同的意见。他试图建立一些新奇的理论，并断言存在某些不可见的永世。他像\Person{瓦伦廷}一样编造了关于它们的传说。他又像\Person{马尔基昂}一样，主张婚姻是败坏。但他声称\Person{亚当}不会得救，因为他是悖逆的始作俑者。\Person{塔提安}的教义就此为止。\par

% structure: p0033 source=h2 target=subsection reason=relative-heading-rank; base=h2

\subsection{第十章：\Person{赫尔摩革涅}；采纳苏格拉底哲学；关于我们主的诞生和身体的观点。}

但有一个名叫赫尔谟革涅斯的人，自己也以为提出了一些新颖的见解，他说天主用与祂同在的、非受生的原始质料造了万有。因为天主不可能从无中造出受造之物。天主永远是主，永远是创造者，质料永远是服从的（实体），并且是在变化中——但不是全部。因为当质料持续地粗野而混乱地运动时，天主通过以下方式使之有序。祂注视（质料）如同锅内的东西在火下沸腾，遂做了部分分离。祂从整体中取出一部分，加以驯服，另一部分则任其混乱地旋转。他断言被驯服的部分就是世界，另一部分仍保持野性，被称为混沌质料。他认为这就是万有的实体，仿佛为门徒引入了一个新颖的教义。然而，他没有想到这正是苏格拉底式的论述，柏拉图比赫尔谟革涅斯阐述得更为详尽。然而他承认基督是创造万有的天主之子；同时他承认祂由童贞女并（由）圣神而生，正如福音的声音所传。并且（赫尔谟革涅斯主张）基督在受难后以身体复活，向门徒显现，祂升天时将身体留在太阳中，而祂自己则前往父那里。现在（赫尔谟革涅斯）引用证言，试图以经上话语支持自己，即圣咏作者达味所说：\BibleQuote{祂在太阳中安置了祂的帐幕，祂自己如同新郎走出洞房，又如勇士欢然奔路。}这些便是赫尔谟革涅斯也试图建立的意见。\par

% structure: p0035 source=h2 target=subsection reason=relative-heading-rank; base=h2

\subsection{第十一章。十四日派。}

另有一些（异端者），生性好争辩，在知识上完全无知，且行为上（比常人）更为好斗，他们联合（主张）应按照法律的规定，在第一个月的第十四天庆祝复活节，不论那天是星期几。（但这样）他们只注意到法律上记载的，不如此遵守（诫命）者将受诅咒。他们没有留意这一点：法律规定是为犹太人而设的，他们将来要宰杀真正的逾越节羔羊。而这（逾越节祭献的效力）已传到外邦人，凭信德领悟，如今不再仅按文字奉行。他们只注意这一条诫命，却不看宗徒所说的话：\BibleQuote{我再向一切受割礼的人作证：他必须遵守全部法律。}\BibleRef{迦 5:3}除此之外，他们同意宗徒传授给教会的所有传统。\par

% structure: p0037 source=h2 target=subsection reason=relative-heading-rank; base=h2

\subsection{第十二章。孟他努派；普黎史拉和玛克西弥拉——他们的女先知；其中有些是诺厄图派。}

但还有另一些人，他们本性甚至比（上述）更为异端，是弗里吉亚血统。这些人因先前被两个可怜的女人（一个叫普黎史拉，一个叫玛克西弥拉）所迷惑而陷入错误，他们以为这些女人是女先知。他们断言护慰者圣神已降在这些女人身上；而在她们之前，他们同样认为孟他努是先知。由于拥有他们无数的著作，（弗里吉亚人）满受欺骗；他们不按照（理性的）标准判断这些人的任何言论，也不听从有能力判断的人；而是因信赖这些（骗子）而轻率地被冲走。他们声称通过这些女人学到了比法律、先知和福音更多的东西。他们竟将这些可怜的女人高举于宗徒和一切恩宠之上，以至于他们中有些人竟敢断言在这些女人身上有超越基督的事物。这些人与教会一样，承认天父是宇宙之父、万有的创造者，并接受福音为基督所见证的一切。然而他们引入了关于斋戒、节期、干粮餐和萝卜餐的新奇做法，声称是由女人教导的。他们中有些人赞同诺厄图派的异端，认为圣父即是圣子，且这一个（圣子）经历了诞生、苦难和死亡。关于这些人，我将以更细致的方式重新解释；因为这些人的异端已成为许多人的祸根。因此我们以为，当我们简要地向所有人证明他们的大部分著作是愚蠢的，他们的（推理）尝试是软弱的，不值一顾时，关于这些（异端者）的陈述就足够了。但拥有健全心智的人没有必要注意（他们的著作或论证）。\par

% structure: p0039 source=h2 target=subsection reason=relative-heading-rank; base=h2

\subsection{第十三章。克己派的教义。}

然而，有些自稱為\OriginalTerm{禁欲派}{Encratites}的人，在有關天主和基督的一些事上，與教會持相同看法。然而，在生活方式上，他們卻驕傲自大地度日。他們以為，透過禁食肉食、只飲水、禁止婚姻、終身致力於苦修習慣，就能藉著食物抬高自己。但這類人更被視為犬儒派而非基督徒，因為他們不留心宗徒保祿針對他們所說的話。保祿預言將來某些（異端者）會徒然引入新奇事物，於是如此聲明：\Quote{聖神明明地說：在末期，有些人將要離棄真道，聽從欺詐的神和惡魔的教導，以虛偽的謊言行事，良心像被烙鐵烙過，禁止嫁娶，禁戒食物，就是天主所創造，為讓那些信者和認識真理的人以感恩之心領受的；因為天主所造的萬物都是好的，凡以感恩之心領受的，沒有一樣是可棄絕的；因為它藉著天主的話和祈禱得以聖化。}\BibleRef{弟前 4:1}那麼，聖保祿的這番話足以駁斥那些如此生活並自詡為義的人，也足以證明（禁欲派的）這個主張同樣構成異端。但即使還有其他一些被稱為異端的——我指的是\OriginalTerm{該隱派}{Cainites}、\OriginalTerm{蛇派}{Ophites}或\OriginalTerm{諾亞派}{Noachites}，以及諸如此類的其他派別——我也沒有必要解釋他們的言行，以免他們因此自以為重要，或值得關注。然而，既然關於這些的論述似乎已經足夠，我們就轉到萬惡之源——\OriginalTerm{諾威督派}{Noetians}的異端。現在，在揭露了這（異端）的根源，並公開譴責其內藏的毒素後，我們須設法勸阻那些被猛烈的靈——如同奔湧的洪流——驅入此等謬誤的人。\par

\SourceRecord{Translated by J.H. MacMahon.；Ante-Nicene Fathers；Vol. 5.；Edited by Alexander Roberts, James Donaldson, and A. Cleveland Coxe.；Buffalo, NY: Christian Literature Publishing Co.,；1886.；<http://www.newadvent.org/fathers/050108.htm>.}

% structure: p0001 source=h1 target=section reason=page-title

\section{驳斥一切异端（卷九）}

% structure: p0002 source=h2 target=subsection reason=relative-heading-rank; base=h2

\subsection{目录。}

以下是\WorkTitle{驳斥一切异端}第九卷的目录：——\par

诺厄图斯的亵渎愚妄是什么，以及他醉心于晦涩者赫拉克利特的教义，而非基督的教义。\par

以及卡利斯托如何将诺厄图斯的弟子克勒奥美涅斯的异端与特奥多图斯的异端相混合，构造出另一种更新奇的异端，以及这位异端者的生活是怎样的。\par

奇怪的精神埃尔克赛近期（在罗马）的到来，以及他表面上遵守律法如何掩盖了他特有的谬误，而实际上他投身于诺斯替派，甚至占星家的学说，以及巫术的技巧。\par

犹太人的习俗是什么，以及他们中间有多少意见分歧。\par

% structure: p0008 source=h2 target=subsection reason=relative-heading-rank; base=h2

\subsection{第一章 当代异端概述}

那么，既然我们对所有异端进行了旷日持久的争战，我们绝未留下任何未经驳斥者，现在最大的斗争仍然存在，即要陈述并驳斥那些在我们当代兴起的异端，某些无知而狂妄的人借此试图四散扰乱教会，并在全世界所有信友中引起了极大的混乱。因为我们认为，有必要抨击那构成（当代）邪恶首要根源的意见，指出这一意见的起源原则，以便其枝芽成为众所周知，并成为普遍嘲笑的对象。\par

% structure: p0010 source=h2 target=subsection reason=relative-heading-rank; base=h2

\subsection{第二章 诺厄图斯异端的根源；克勒奥美涅斯为其门徒；在泽菲里努斯与卡利斯托任主教期间于罗马出现；希波利图斯在罗马反对诺厄图斯主义}

有一个人，名叫诺厄图斯，出生于斯米纳。此人从赫拉克利特的教义中引入了一种异端。有一个叫厄皮戈努斯的人，成了他的仆从和门徒，此人在罗马逗留期间散布了这无神的见解。但克勒奥美涅斯，成了他的门徒，生活与习惯都与教会格格不入，常为（诺厄图斯）的教义作证。那时，泽菲里努斯自以为管理着教会的事务——一个无知且可耻地腐败的人。他因被给予的好处说服，常纵容那些前来成为克勒奥美涅斯门徒的人。而（泽菲里努斯）自己，久而久之被引诱，急急忙忙陷入了同样的意见；他并有卡利斯托作为顾问，以及这些邪恶教义的战友。但这位（卡利斯托）的生活以及他发明的异端，我稍后将解释。这些异端者的学派，在\Emph{这样的主教}相继期间，不断获得力量和扩张，因为泽菲里努斯和卡利斯托帮助他们得势。然而，我们从未与他们同流合污；但我们常常反对他们，驳斥他们，并迫使他们勉强承认真理。他们被真理羞愧和约束，曾短暂地承认他们的\Emph{错误}，但不久后又再次滚入同样的泥淖。\par

% structure: p0012 source=h2 target=subsection reason=relative-heading-rank; base=h2

\subsection{第三章 诺厄图斯主义是赫拉克利特哲学的一个分支}

但既然我们已经展示了他们传承的谱系，接下来似乎应该也解释他们教义中所包含的败坏教导。\Emph{为此}，我们首先将引述赫拉克利特\Quote{晦涩者}所提出的意见，然后我们揭示这些意见中哪些部分源自赫拉克利特。\Emph{他们体系的这些部分}的当前拥护者并不知道它们属于\Quote{晦涩者}\Emph{哲学家}，却以为它们属于基督。但是，如果他们偶然遇到下面的观察，也许他们会因此感到羞愧，并被诱导停止他们这种无神的亵渎。现在，尽管赫拉克利特的意见先前已在\WorkTitle{哲学论丛}中被我们阐述过，但现在似乎仍应将\Emph{两种体系}并列对比，以便通过这种更切近的驳斥，他们能更清楚地受教。\Emph{我的意思是}，这位（异端者）的追随者，他们自以为基督的门徒，而实际上并非如此，却是\Quote{晦涩者}的门徒。\par

% structure: p0014 source=h2 target=subsection reason=relative-heading-rank; base=h2

\subsection{第四章 赫拉克利特体系概述}

\Person{赫辣克利托斯}随后说，宇宙是\Emph{一}，可分割且不可分割；受造\Emph{而}非受造；可朽\Emph{而}不朽；理性，永恒；圣父、圣子与正义，就是天主。\Quote{因为那些不听从我、却听从道理的人，有智慧承认万有为一，}\Person{赫辣克利托斯}说；因为并非所有人都知道或承认这一点，他以如下言语发出责斥：\Quote{人们不理解，那分歧之物（却）如何与自身合一，正如弓与竖琴的反向和谐。}但那理性永远存在，因为它构成宇宙，并渗透万有，他以这种方式肯定说：\Quote{但关于这永远存在的理性，人们始终缺乏理解，无论在他们听说它之前，还是初次听闻之时。因为尽管万有都按照这理性发生，他们却好似对它毫无经验的人。然而他们仍试图做出如我所述的那类言行，藉着按本性区分，并宣告事物如何。}而圣子即宇宙，且历经无尽世代为万有的永恒君王，他如此断言：\Quote{一个玩耍的孩子，掷着骰子，就是永恒；王国属于一个孩子。}而万有之父是一位非受造的受造物，他是创造者，让我们听\Emph{\Person{赫辣克利托斯}}以这些话肯定：\Quote{矛盾是万有之父和万有之王；它彰显一些为神，另一些为人，使一些为奴，另一些为自由。}而（他同样肯定）存在着一种和谐，如在弓与竖琴中。那隐晦的和谐（更佳），虽不为人类所知且不可见，他以这些话肯定：\Quote{隐晦的和谐优于明显的和谐。}他称赞并钦佩那未知不可见者\Emph{就}其能力而言，胜于已知者。而那为人可见、并非不可发现的\Emph{和谐}更为优越，他以这些话肯定：\Quote{凡视觉、听觉\Emph{和}理智的对象，这些我格外尊崇}，他说；也就是说，\Emph{他偏好}可见之物胜于不可见者。从他的这些表述，容易理解\Emph{其哲学的精髓}。\Quote{人啊，他说，在对明显事物的认识上，他们受欺骗，与比所有希腊人都更智慧的\Person{荷马}相似。因为连杀虱子的孩童也欺骗了他，那时他们说：\InnerQuote{我们所看见并抓住的，我们留下；而我们所没有看见也没有抓住的，我们带走。}}\par

% structure: p0016 source=h2 target=subsection reason=relative-heading-rank; base=h2

\subsection{第五章 \Person{赫辣克利托斯}对\Person{赫西奥德}的评价；\Person{赫辣克利托斯}的悖论；他的末世论；源于\Person{赫辣克利托斯}的\Person{诺厄图斯}异端；\Person{诺厄图斯}关于吾主诞生与受难的看法。}

赫拉克利特以这种方式给予可见者与不可见者同等的地位和荣誉，仿佛可见者与不可见者显然是一回事。因为他说：\Quote{暗合的和谐优于明显的和谐；}以及\Quote{凡是视觉、听觉和理智的对象}，即（有形）器官的对象——\Quote{这些，}他说，\Quote{我特别尊敬，}（这次不是，尽管以前）他曾特别尊敬不可见的事物。因此，赫拉克利特断言，黑暗与光明、恶与善并无不同，而是同一回事。无论如何，他责备赫西俄德，因为他\Emph{不}知道白天和黑夜。因为白天和黑夜，他说，是一回事，他这样表述自己：\Quote{然而，拥有大量知识的老师是赫西俄德，人们认为这位\Emph{诗人}拥有极其丰富的知识，但他却不知道白天和黑夜的本性，因为它们是一回事。}就善与恶而言，（根据赫拉克利特，它们同样）是一回事。\Quote{医生无疑，}赫拉克利特说，\Quote{当他们切开和灼烧时，\Emph{尽管}在各方面邪恶地折磨病人，却抱怨没有从病人那里得到合适的报酬，尽管他们对疾病施行这些有益的治疗。}而直和曲，他说，是相同的。\Quote{梳羊毛者的路是直而曲的；}缩绒工坊中称为\Quote{螺旋}的器具的圆周运动是直而曲的，因为它同时上下旋转。\Quote{因此，}他说，\Quote{\Emph{直和曲是}同一的。}而上和下，他说，是一回事。\Quote{向上的路\Emph{和向下的路}是相同的。}他说，污秽与纯净是一回事，可饮与不可饮是一回事。\Quote{海，}他说，\Quote{是极其纯净又极其污秽的水，对鱼无疑可饮且有益\Emph{于它们}，但不适合人饮用，（对人）有害。}而且，他显然断言不朽的是有死的，有死的是不朽的，用以下表述：\Quote{不朽的是有死的，\Emph{而}有死的是不朽的，\Emph{即}当一方从死亡中取得生命，另一方从生命中取得死亡。}他还肯定我们出生所在的这有形的肉体有复活，并且他知道天主是这复活的原因，这样表述：\Quote{那些在这里的\Emph{，天主将}使他们兴起并成为生者和死者的守护者。}他还肯定世界及其中万物的审判通过火进行，这样表述：\Quote{现在，雷霆驾驭万物，}即，指导\Emph{它们}，以雷霆意指永恒之火。但他也断言这火具有理智，是宇宙管理的原因，并称它为渴望和满足。据他看来，渴望是\Emph{世界的}秩序，而满足是它的毁灭。\Quote{因为，}他说，\Quote{火降临\Emph{大地}，将审判并夺取一切。}\par

但在这一章中，\Person{赫拉克利特}同时解释了其思维方式的全部特点，同时也解释了\Person{诺厄托}异端的特点。我已简要证明\Person{诺厄托}不是基督的门徒，而是赫拉克利特的。因为\Emph{这位哲学家}断言原始世界本身就是\OriginalTerm{造物主}{Demiurge}和自身的创造者，在以下段落中：\Quote{天主是白天，黑夜；冬天，夏天；战争，和平；满足，饥荒。}一切事物都是对立面——这似乎是他的意思——但发生了一种变化，就像香料与其他\Emph{种类}的香料混合，但根据每种\Emph{香料}产生的愉悦感来命名。现在，所有人都清楚，\Person{诺厄托}的愚蠢后继者和其异端的捍卫者，即使他们没有听过赫拉克利特的言论，但无论如何，当他们接受\Person{诺厄托}的观点时，毫不掩饰地承认这些（赫拉克利特式的）主张。因为他们提出这样的说法：同一天主是万物的创造者和圣父；当祂乐意时，祂\Emph{尽管如此}向古代的义人显现，（虽不可见）。因为当祂不被看见时，祂是不可见的；\Emph{祂是}晦涩难懂的，当祂不愿被理解时，但可理解时，当祂被理解时。因此，根据同样的说法，祂是不可战胜的和可战胜的，非受生的和受生的，不朽的和有死的。持有这种观点的人怎能不被证明是赫拉克利特的弟子？难道晦涩的（赫拉克利特）没有用相同的表述方式，在\Person{诺厄托}之前建立了哲学体系吗？\par

现今，无人不知\Person{诺厄提}断言圣子与圣父是同一的。但他如此陈述：\Quote{确实，当圣父尚未被生时，祂仍公义地被称为圣父；而当祂乐意经历受生、被生之后，祂自身便成了祂自己的圣子，而非他人的。}因为他以为这样就能确立\Emph{天主的}主权，声称所谓的圣父与圣子是同一的\OriginalTerm{实体}{substance}，不是由一个不同的者产生的，而是自身出于自身；且按时间的变迁，祂被称为圣父和圣子。但祂是那一位\Emph{在我们中间}显现的，既从童贞女受生，又以人的身份与人交往。因着所发生的诞生，祂向看见祂的人承认自己是圣子，无疑；但祂对能领悟祂的人并不隐瞒祂是圣父。这一位被钉在树上受苦，将祂的灵魂交给自己，\Emph{表面上}死了，而并非（实际）死亡。祂被埋葬在坟墓中，被枪刺穿，被钉子扎透，第三天自己复活了。\Person{克娄美尼}与他的\Emph{追随者}一同断言这一位是宇宙的天主和圣父，\Emph{从而}在许多人心目中引入了一种模糊不清（的思想），如同我们在赫拉克利特的哲学中所见的那样。\par

% structure: p0020 source=h2 target=subsection reason=relative-heading-rank; base=h2

\subsection{第六章 科利奥斯特及则斐琳在诺厄提主义一事上的行为；则斐琳关于耶稣基督的公开意见；希玻律的反对；作为同时期事件，希玻律有能力解释此事。}

\Person{科利奥斯特}试图确认这一异端——此人在邪恶上狡猾，在欺骗上诡诈，被不安的野心驱使要登上主教宝座。\Emph{此人把}则斐琳，一个无知无识、不谙教会定义的人，\Emph{塑造成自己的工具}。又因\Emph{则斐琳}贪财受贿，\Person{科利奥斯特}通过礼物引诱他，并借非法要求，得以诱使他随己意行事。如此，科利奥斯特成功地诱导则斐琳不断在弟兄们中制造纷乱，而他自己随后用诡诈的话语小心地使两派都向他怀有好意。有时，对持正确见解的人，他私下声称他们与自己持相同教义，从而欺骗他们；另一些时候，\Emph{他对持}萨伯里\Emph{观点的人也如法炮制}。但\Person{科利奥斯特}却败坏了\Person{萨伯里}本人，虽然他本有能力纠正\Emph{这异端者的错误}。因为（任何时候）在我们的劝诫中，\Person{萨伯里}并未显明顽固；但当他独自与科利奥斯特在一起时，竟被这同一\Person{科利奥斯特}促使重陷于克娄美尼的体系，而科利奥斯特却声称他与\Emph{克娄美尼}持相似意见。\Person{萨伯里}当时未察觉\Emph{科利奥斯特的}欺诈，但后来他明白了，正如我即将叙述的。\par

现在\Person{科利奥斯特}将\Person{则斐琳}本人推上前台，促使他公开宣认\Emph{以下观点}：\Quote{我知道只有一个天主，耶稣基督；除祂之外，\Emph{我不知道}任何其他受生而可受苦难者。}而在另一场合，当他要说以下话时：\Quote{圣父没有死，而是圣子死了。}\Person{则斐琳}就这样不断在百姓中制造无休止的骚动。我们得知他的观点后，没有让步，而是为真理的缘故责备并抵制他。他因众人都赞同他的伪善而仓促陷入愚妄——我们\Emph{却}没有\Emph{这样做}——他称我们为敬拜两个天主的人，不自觉地吐出心中藏着的毒液。我们认为有必要说明这位\Emph{异端者}的生平，因他与我们几乎同时出生，以便通过揭露这等人物之习性，他所试图确立的异端可轻易被认识，或可能被有理智者视为荒谬。这位\Person{科利奥斯特}在弗斯奇安担任罗马长官期间成了\Quote{殉道者}，他\Quote{殉道}的方式如下。\par

% structure: p0023 source=h2 target=subsection reason=relative-heading-rank; base=h2

\subsection{第七章 科利奥斯特的个人历史；他的银行家职业；对卡尔波弗鲁斯的欺诈；科利奥斯特潜逃；企图自杀；被判磨坊苦役；因长官弗斯奇安的命令再次判刑；放逐到撒丁岛；马尔基昂的干预使科利奥斯特获释；科利奥斯特抵达罗马；维克多教宗将科利奥斯特迁往安提乌姆；维克多死后科利奥斯特返回；则斐琳对他友好；萨伯里指控科利奥斯特；希玻律关于科利奥斯特意见的记述；罗马的科利奥斯特学派及其做法；此教派在希玻律时代仍存在。}

\Emph{\Person{Callistus}}碰巧是某位\Person{Carpophorus}的家仆，此人属于凯撒王室中有信仰的人。因\Emph{\Person{Callistus}}也是有信仰的，\Person{Carpophorus}委托给他相当可观的一笔钱，并指示他通过银行业务赚取利润。他收下\Emph{钱}后，尝试在所谓的\Emph{\Place{公共水池}}（Piscina Publica）开展银行业务。随着时间的推移，许多寡妇和弟兄以\Emph{将钱存放在\Person{Carpophorus}处}的托词，将大量存款托付给他。\Emph{\Person{Callistus}}却挥霍了所有委托给他的钱，陷入财务困境。做出这等事后，有人向\Person{Carpophorus}告发，后者表示要向他算账。\Person{Callistus}察觉此事，预感主人有危险，便偷偷逃走，向海边奔去。他在\Place{Portus}找到一艘准备启航的船，便上了船，打算随船去往任何目的地。但即使如此，他也没能逃脱，因为有人将发生的事告诉了\Person{Carpophorus}。\Emph{\Person{Carpophorus}}得到消息后，立即赶往港口（Portus），试图\Emph{追赶\Person{Callistus}}上船。然而，小船停泊在港口中央；由于船夫动作缓慢，船上的\Person{Callistus}得以远远望见他的主人。他知道自己难免被擒，便不顾性命；认为自己处境绝望，便投身入海。但水手们跳上小船，把他拖了出来，尽管他不愿上来，而岸上的人们大声呼喊。这样，\Emph{\Person{Callistus}}被交给他的主人，被带回罗马，主人将他关进\Emph{\Place{Pistrinum}}（磨坊）。\par

但随着时间过去，正如\Emph{这类事情}常发生的，弟兄们去见\Person{Carpophorus}，恳求他释放这名逃仆免受惩罚，理由是声称\Emph{\Person{Callistus}}承认自己有一些钱\Emph{存放在}某些人那里。但\Person{Carpophorus}作为一个虔诚的人，说他对自己\Emph{财产}并不在意，但他关心那些存款；因为许多人流着泪对他说，他们以\Emph{将钱存放在他那里}的托词，把委托给\Person{Callistus}的东西交给了他。\Person{Carpophorus}被他们的劝说打动，下令释放\Emph{\Person{Callistus}}。然而，\Person{Callistus}无钱可还，又因被监视而无法再次逃跑，便设下一计，希望以此寻死。他假装前往债权人处，在犹太人的安息日匆匆赶到他们聚集的会堂，站在那里，在他们中引起骚乱。犹太人被他搅扰，便侮辱他、殴打他，将他拖到城市长官\Person{Fuscianus}面前。当被问及为何如此对待时，他们这样回答：\Quote{罗马人已准许我们公开阅读我们祖传的法律。但这人进入我们的敬拜场所，阻止我们这样做，在我们中间制造骚乱，并声称自己是基督徒。} \Person{Fuscianus}此时正坐在审判席上；他因犹太人的陈述而对\Person{Callistus}表示愤怒，便有人去告知\Person{Carpophorus}这些事。\Person{Carpophorus}急忙赶到长官的审判席前，喊道：\Quote{我恳求你，我的主\Person{Fuscianus}，不要相信这个\Emph{家伙}；他不是基督徒，而是寻死，他盗走了我大量钱财，我能证明。} 犹太人却以为这是计策，好像\Person{Carpophorus}以此借口要释放\Emph{\Person{Callistus}}，便在长官面前更加敌意地\Emph{对他}叫嚷。\Emph{\Person{Fuscianus}}却被这些\Emph{犹太人}影响，鞭打了\Emph{\Person{Callistus}}，将他\Emph{发配}到\Place{Sardinia}的一个矿场。\par

但过了一段时间，在那个地方有其他殉道者，\Person{玛尔恰}，\Person{科莫都斯}的妾，一位热爱天主的女性，渴望行一些善工，就邀请当时担任教会主教的\Person{真福维克托}到她的面前，询问他\Place{撒丁岛}上有哪些殉道者。他将所有人的名字都交给了她，但没有给出\Person{卡利斯都}的\Emph{名字}，因为他知道\Person{卡利斯都}所冒险做出的那些行为。\Person{玛尔恰}从\Person{科莫都斯}那里获得准许后，将释放信交给一位年纪较大的太监\Person{伊亚钦托}。\Person{伊亚钦托}收到\Emph{信}后，乘船前往\Place{撒丁岛}，将\Emph{信}交给当时该地的总督，成功使那些殉道者获得释放，但\Person{卡利斯都}除外。但\Emph{\Person{卡利斯都}}自己双膝跪下，流着泪恳求也能获得释放。因此，\Person{伊亚钦托}被\Emph{这位囚犯}的央求所打动，就请求总督\Emph{予以释放}，声称\Person{玛尔恰}已授权给他（释放\Person{卡利斯都}），并且他会安排使此事对总督\Emph{没有}风险。于是（总督）被说服，也释放了\Person{卡利斯都}。当\Emph{他}到达\Place{罗马}时，\Person{维克托}对所发生的事非常伤心；但因为他是个富有同情心的人，就没有采取任何行动。然而，为了防范许多人的指责（因为\Emph{这个\Person{卡利斯都}}所做的尝试并非遥远的事情），也因为\Person{卡尔波福}仍然心怀敌意，\Emph{\Person{维克托}}就将\Emph{\Person{卡利斯都}}送到\Place{安提乌姆}居住，并为他规定了每月一定量的膳费。在\Emph{\Person{维克托}死后，\Person{则斐琳}}让\Emph{\Person{卡利斯都}}作为同工协助管理他的神职人员，对他表示尊敬，却损害了自己；并将此人从\Place{安提乌姆}调回，任命他管理坟场。\par

而\Person{卡利斯都}，一向习惯与\Person{则斐琳}交往，并且如我先前所说，对他假意侍奉，则揭示了\Emph{通过对比，\Person{则斐琳}是}一个既不能判断所说的事物、也不能辨别\Person{卡利斯都}意图的人；他习惯于与\Emph{\Person{则斐琳}}谈论令\Emph{\Person{则斐琳}}满意的话题。这样，在\Person{则斐琳}死后，他以为得到了他如此热切追求的地位，就将\Person{撒伯流}逐出教会，认为他抱有不合正统的意见。他这样做是出于对我的畏惧，并以为能以此方式抹去教会中\Emph{对他}的指控，就好像他并不持有怪异的意见。他那时是个骗子和无赖，并逐渐带走了许多人。他甚至心中有剧毒，对任何事物都没有正确的见解，同时又耻于说出真相，\Emph{这个\Person{卡利斯都}}不仅公开指责我们说：\Quote{你们是\OriginalTerm{二位神论者}{Ditheists}，}而且由于他经常被\Person{撒伯流}指控为背弃了他起初的信德，就发明了以下这样的异端。\Emph{\Person{卡利斯都}}声称：圣言自己就是子，而且自己就是父；虽然以\Emph{不同}的名号称呼，但实际上祂是不可分的一位天主。\Emph{他主张}：父不是一个位格而子是另一个，而是同一的；并且万物都充满圣神，无论是上面的还是\Emph{下面的}。\Emph{他又肯定}：在童贞玛利亚内降生成人的圣神与父没有分别，而是同一的。\Emph{他补充说}，这就是\Emph{救主}所宣布的：\BibleQuote{你不信我在父内，父在我内吗？}\BibleRef{若 14:2}因为那看得见的，即人，\Emph{他认为是}子；而在子内的圣神，则是父。\Quote{因为，}\Person{卡利斯都}说\Quote{我不愿相信两位天主，父与子，而只信一位。因为父，祂自己存在于\Emph{子}内，在取了我们肉身后，将肉提升到神性的本质，使之与祂自己结合，并使它成为一体；因此父与子必须被称为一位天主，并且这一位格既然是唯一的，就不能是两位。}就这样，\Emph{\Person{卡利斯都}主张}父与子一同受苦；因为他不想断言父受苦，且是一个位格，小心避免对父的亵渎。（他多么小心啊！）这个无知而无赖的家伙，到处即兴编造亵渎的话，只是为了让自己\Emph{不}显得违背真理，而且不羞于时而陷入\Person{撒伯流}的教义，时而又陷入\Person{特奥多图斯}的学说。\par

那个骗子\Person{卡利斯托}大胆提出这些见解，采纳了前述的教导体系，建立了一个与教会对立的\Emph{神学}学派。他首先发明了纵容人们沉溺于\Emph{肉欲}享乐的手段，说所有人都能藉着他自己获得罪的赦免。因为，凡有习惯参加别人的集会且被称为基督徒的人，若犯了任何过犯，他们就说，只要他赶快\Emph{并加入}卡利斯托的学派，这罪就不算在他身上。许多人因这规定而喜悦，因为他们良心不安，同时又已被许多教派拒绝；而且他们中有些人，根据我们的定罪裁决，已被我们强行逐出教会。\Emph{现在这样的门徒}转而投向这些\Emph{卡利斯托的追随者}，使他的学派拥挤不堪。此人提出主张：如果一个主教犯了任何罪，即使是至于死的罪\BibleRef{若一 5:16}，他也不应被罢免。大约在这人的时代，结过两次婚和三次婚的主教、司铎和执事开始\Emph{被允许}保留他们在圣职人员中的位分。而且，如果任何处于圣秩中的人结了婚，\Person{卡利斯托}允许这样的人继续留在圣秩中，如同没有犯罪一样。作为理由，他声称宗徒所说的这些话是针对此\Emph{那人}的：\Quote{你是谁，竟敢论断别人的仆人？}\BibleRef{罗 14:4}。但他断言同样\Emph{那个}莠子的比喻也是针对此\Emph{一人}的：\Quote{让莠子与麦子一起生长；}\BibleRef{玛 13:30}换句话说，让那些在教会中犯了罪的人\Emph{留在教会内}。他还肯定诺厄的方舟被造是为教会的象征，其中有狗、狼、乌鸦，以及一切洁净和不洁净的；于是他声称教会也应当如此。并且，他收集到的许多\Emph{圣经的}支持这种\Emph{观点}的部分，他都这样解释。\par

并且\Person{卡利斯托}的听众喜悦于他的主张，继续与他同在，\Emph{如此}嘲弄他们自己和许多\Emph{其他人}，成群结队的这些\Emph{受骗者}涌入他的学派。因此\Emph{他的学生}增多，他们以（参加学派的）人群为荣，为了基督所不允许的享乐。但出于对基督的蔑视，他们对任何罪的犯下都不加约束，声称他们赦免那些顺从（卡利斯托观点）的人。因为他甚至允许女性，如果她们未婚且在不适当的年纪欲火中烧，或者她们不愿意通过合法婚姻来损害自己的尊严，那么她们就可以随意选择任何人作为同寝者，无论是奴隶还是自由人；并且\Emph{一个女人}，尽管没有合法结婚，也可以将这样的\Emph{伴侣}视为丈夫。于是，被认为有信德的妇女开始服用避孕药物，并束紧身体，以便排出所怀的胎，因为她们不愿因奴隶或任何下贱人而生孩子，是为了维护其家族和过度的财富。看哪，那个无法无天的\Emph{人}已陷入何等不敬虔，同时灌输奸淫和谋杀！而且，在如此大胆的行为之后，他们不顾羞耻，竟试图自称为天主教会！有些人以为会获得福乐，就附和了他们。在此\Emph{一人}的\Emph{主教任期}内，他们首次妄行第二次圣洗。然后，这些就是那个最令人惊讶的卡利斯托所建立的（做法和主张），他的学派持续存在，保留其习俗和传统，不分辨应与谁共融，\Emph{反而}不加选择地向所有人提供共融。并且他们从他那里得了他们\Emph{名称}的称号；以致，由于卡利斯托是此类做法的主要倡导者，他们应被称为卡利斯托派。\par

% structure: p0030 source=h2 target=subsection reason=relative-heading-rank; base=h2

\subsection{第8章 厄勒克赛派；希波利托斯对其的反对。}

卡利斯托的学说传遍整个世界后，一个狡黠且绝望之徒，名叫阿尔基比亚德，住在叙利亚的阿帕米亚，仔细查考了此事。他自认比卡利斯托更令人畏惧、更善于此类伎俩，便前往罗马；他带来一本书，声称某个义人厄尔克赛从帕提亚的塞拉镇领受了此书，并将其交给了名叫索比亚伊的人。该卷内容据称是由一位天使启示的，其身高为二十四朔伊诺伊，即九十六英里，宽度为四朔伊诺伊，从肩到肩六朔伊诺伊；他足迹的长度延伸至三个半朔伊诺伊，相当于十四英里，宽度为一个半朔伊诺伊，高度为半朔伊诺伊。他还声称有一位女子与他同在，其尺寸据称与前述标准相符。他断言那男子（天使）是天主子，而那女子被称为圣神。通过详述这些奇事，他幻想能迷惑愚人，同时说出以下话语：\Quote{在图拉真在位第三年，有一项新的罪赦被传扬给人。}厄尔克赛还规定了圣洗的性质，这一点我也要说明。他断言，对于那些陷入各种淫荡、污秽和邪恶行为的人，只要其中任何一人是信徒，当他悔改、听从那书并相信其内容时，就应藉圣洗获得罪的赦免。\par

然而，厄尔克赛竟敢继续这些奸计，以卡利斯托所拥护的上述教义为借口。因为他看出许多人对此类应许感到欣喜，便认为自己可以适时进行刚才提到的尝试。尽管我们对此提出抵制，并且不允许许多人长时间成为此迷妄的牺牲品。因为我们使人们确信，这是伪灵的工作，是骄傲膨胀之心所发明的，并且这一位像狼一样起来攻击许多迷途的羊，这些羊是卡利斯托藉其欺骗之术驱散的。但我们既已开始，就不会对此人的意见保持沉默。首先，我们要揭露他的生平，证明他所谓的戒律纯属伪装。其次，我将引述其主张的主要要点，以便读者定睛注视（厄尔克赛的）著作，就能明白此人所大胆尝试的异端是什么性质。\par

% structure: p0033 source=h2 target=subsection reason=relative-heading-rank; base=h2

\subsection{第九章 厄尔克赛的体系源自毕达哥拉斯；行邪术}

厄尔克赛把（经）法律所授权的政体作为诱饵，声称信徒应当受割损并按照法律生活，同时他又强行摘取上述异端中的某些片段。他断言基督是按照与众人相同的方式降生为人，并且基督并非首次由童贞女降生地上，而是从前和以后多次降生。基督将这样从时期到时期显现并存在于我们中间，经历出生的变化，灵魂从身体转移至身体。厄尔克赛采纳了我已经提到的毕达哥拉斯的教义。但厄尔克赛派竟骄傲到如此地步，甚至自称拥有预言未来的能力，其起点显然是上述毕达哥拉斯技艺的尺度和数字。他们也致力于数学家、占星家和巫师的教义，仿佛这些教义是真实的。他们诉诸这些，以迷惑愚人，使（受迷惑的）人以为（这些异端者）参与了某种权能的教义。他们教导为被狗咬、被鬼附及患有其他疾病的人施行某些咒语和程式；我们对于这些（异端者）的这类（行径）也不会保持沉默。既已充分解释了他们的原则和冒昧尝试的原因，我将转而叙述他们的著作，借此读者将认识到厄尔克赛派的这些琐碎而不虔的努力。\par

% structure: p0035 source=h2 target=subsection reason=relative-heading-rank; base=h2

\subsection{第十章 厄尔克赛施行圣洗的方式；咒文}

于是，对于那些接受他口传教导的人，他以如下方式施行圣洗，向受骗者讲说类似这样的话：\Quote{因此，（我的）孩子们，若有人与任何种类的动物交合，或与男性、或与姊妹、或与女儿行淫，或犯奸淫、或行淫乱，且渴望获得罪的赦免，从他听从这书的那一刻起，就让他再次受洗，因至大至高的天主之名，并因祂的圣子、大能君王之名。并且藉着\Emph{圣洗}，让他被净化、被洁净，并让他为自己呼求这书上所记述的那七个见证——天、水、\OriginalTerm{圣神们}{holy spirits}、祈祷的天使、油、盐和地。} 这些便是\Person{厄耳克塞}的令人惊异的奥秘，那些无可言喻且具大能的\Emph{隐秘事}，他传授给配得的门徒。但那位无法无天的\Emph{者}并不满足于此，还在两三个证人面前印证他自己邪恶的\Emph{作为}。他又这样表述：\Quote{再者我说，奸夫淫妇和假先知啊，若你们渴望悔改，使你们的罪得赦免，一旦你们听从这书，并连同你们的衣服再次受洗，平安就归于你们，\Emph{你们的}份与义人同在。} 既然我们说过，这些人为被狗咬的人及其他\Emph{不幸}采用咒语，我们就解释这些。如今\Emph{\Person{厄耳克塞}}使用如下套语：\Quote{若一条因内有毁灭之灵而狂怒的疯狗咬伤任何男人、女人、青年或女孩，或骚扰或触及\Emph{他们}，在那同一时辰，让那人穿着全部衣服奔跑，下到河流或泉水，只要那里有深处。让他（或她）连同全部衣服浸入水中，并怀着心中的信德向至大至高的天主献上祈求，然后让他\Emph{如此}呼求这书上所记述的七个见证：\InnerQuote{看，我呼求天、水、\OriginalTerm{圣神们}{holy spirits}、祈祷的天使、油、盐和地为见证。我凭这七个见证作证：从今以后我不再犯罪，不行奸淫，不偷盗，不行不义，不贪恋，不怀恨，不轻蔑，也不以任何邪恶\Emph{行为}为乐。}说完这些话之后，就让那人穿着全部衣服，因大能至高的天主之名受洗。}\par

% structure: p0037 source=h2 target=subsection reason=relative-heading-rank; base=h2

\subsection{第十一章 \Person{厄耳克塞}的规诫}

但在许多其他方面他胡言乱语，也教导使用这些语句来治疗患痨病的人，并要他们在七天内在冷\Emph{水}中浸入四十次；\Emph{他规定}对附魔者采用类似疗法。啊，无可比拟的智慧和充满权能的咒语！谁不会对如此这般的言辞力量感到惊异呢？但既然我们说过，他们也运用星象学的欺骗，我们将从他们自己的套语中证明\Emph{这一点}；因为\Emph{\Person{厄耳克塞}}如此说：\Quote{存在邪恶的不虔敬之星。这宣告如今已由我们发出，\Emph{哦，你们这些}虔诚\Emph{者}和门徒：要警惕这些\Emph{星}的统治之日\Emph{的}权能，在它们的\Emph{统治}日子不可开始任何工作的开端。在这些\Emph{星}的权能之日，不可为男人或女人施洗，此时月亮从它们中间（出来）运行\Emph{天空}，并与它们同行。当心那天直到\Emph{月亮}从这些\Emph{星}中经过，然后才可施洗并开始你们一切工作的开端。此外，要尊敬安息日，因为那天是这些\Emph{星}掌权之日。但要当心，不可在安息日之后的第三\Emph{天}开始\Emph{你们的工作}，因为当\Person{特罗扬}皇帝统治\Emph{之年}再过三年，从他使帕提亚人臣服于自己权下之时算起——当，\Emph{我说}，三年届满，战争在北部\Emph{星宿}的不虔敬天使之间爆发；\Emph{并且}因此所有不虔敬的王国都处于混乱之中。}\par

% structure: p0039 source=h2 target=subsection reason=relative-heading-rank; base=h2

\subsection{第十二章 \Person{厄耳克塞}派的异端为衍生者}

埃克沙塞（Elchasai）认为，若将这些伟大而不可言喻的奥迹遭受践踏，或交付给许多人，那是对理性的侮辱，故而劝诫应如宝贵的珍珠（\BibleRef{玛 7:6}）加以保存，其言如下：\Quote{不可向众人诵述此记述，要谨慎持守这些诫命，因为并非所有人都是忠信的，也并非所有女子都是正直的。} \Emph{包含这些（主张）的书籍}，埃及的智者们并未藏于圣所，希腊的贤者毕达哥拉斯（Pythagoras）也未\Emph{将其隐于其中}。因为若当时埃克沙塞尚在，又何须毕达哥拉斯、泰勒斯（Thales）、梭伦（Solon）、智慧的柏拉图（Plato）、甚至其余希腊圣贤成为埃及司祭的门徒？他们本可从阿尔喀比亚德（Alcibiades）——那可怜的埃克沙塞最令人惊叹的诠释者——处获得如此智慧。因此，为认识这些人的疯狂而作的陈述，对有健全心智者而言似乎已足够。正因如此，再引述他们更多的格式条文似非明智，因这些条文既繁多又可笑。然而，我们既未省略在我们时代兴起的那些\Emph{做法}，也未对在我们之前\Emph{盛行}的保持沉默，为能通览所有\Emph{体系}而无所遗漏，似乎应当陈述犹太人的（习俗）以及他们中间的纷歧见解，因我想这些仍有待\Emph{我们考察}。现在，我一旦就这些要点打破沉默，便将继续论述真理之道的证明，以便在历经与一切异端的长篇辩争之后，我们虔诚地向着天国冠冕迈进，坚信真理，而不致动摇。\par

% structure: p0041 source=h2 target=subsection reason=relative-heading-rank; base=h2

\subsection{第十三章 犹太人的派别}

起初犹太人中只盛行一种习俗；因为天主赐给他们一位导师，\Emph{即}梅瑟，并由这\Emph{同一位}梅瑟赐下一条法律。并且\Emph{只有}一片旷野和一座西乃山，因为\Emph{只有}一位天主为这些\Emph{犹太人}立法。然而，当他们渡过约但河，并按签分得征服之地后，他们便以各种方式撕裂了天主的法律，各人对天主所作的声明提出\Emph{不同的解释}。他们就这样为自己兴起导师，发明了异端性质的理论，并继续陷入分裂。我目前正打算解释这些犹太人的纷歧。虽然他们长期以来已分裂成极多派别，但我只拟阐明其中主要的，而好学的人自会轻易了解其余。因为他们分为三种：第一种是法利塞人，第二种是撒杜塞人，其余是厄色尼人。后者度更虔诚的生活，充满彼此相爱，并节制有度。他们回避一切过分欲望的行为，甚至不愿\Emph{听说}此类事情。他们拒绝婚姻，但收养别人的男孩，\Emph{从而}为他们生育后裔。他们引导\Emph{这些收养的孩子}遵守自己的特殊习俗，用这种方式教养他们，并促使他们学习科学。然而，他们并不禁止领养的孩子结婚，虽然他们自己避免婚姻。但他们不接纳妇女，即使妇女愿意奉行同样生活方式，因他们对妇女毫无信心。\par

% structure: p0043 source=h2 target=subsection reason=relative-heading-rank; base=h2

\subsection{第十四章 厄色尼人的教理}

他们轻视财富，不拒绝与贫困者分享\Emph{他们的财物}。他们中无人拥有比别人更多的财富。因为他们的规矩是：凡前来\Emph{加入}此派的人必须变卖财产，将\Emph{价款}交给团体。在收到\Emph{钱}后，该\Emph{团体}的会长将其按各人所需分给众人。因此他们中无人困乏。他们不用油，认为傅油是玷污。他们有指定的监督者，负责照料公共财物，且他们常穿白衣。\par

% structure: p0045 source=h2 target=subsection reason=relative-heading-rank; base=h2

\subsection{第十五章 厄色尼人的教理（续）}

但他们并无\Emph{一座}城市，而是多人定居于各城中。若该派成员有来自异\Emph{乡}者，他们认为一切都与他共享，接待素不相识者如同自家亲属。他们遍历本土，每次旅行只携带武器。他们在城中亦有一位会长，负责用为此\Emph{目的}募集的\Emph{款项}为他们购置衣物和食物。他们的衣袍式样朴素。他们不拥有两件外衣或两双鞋；待现有者陈旧后，才另取新的。他们既不买也不卖任何东西；凡有人所有，便给与没有者；凡有人所无，便接受。\par

% structure: p0047 source=h2 target=subsection reason=relative-heading-rank; base=h2

\subsection{第十六章 厄色尼人的教理（续）}

他们有条不紊地持续从清晨开始祈祷，除非先以赞美诗赞美了天主，否则不出一言。然后他们各自出去，从事任何他们喜欢的工作；工作到第五时辰就停止。然后他们又聚集到一处，用亚麻腰带围身，以遮住私处。他们就这样用冷水洗涤；洗净后，他们一同进入一个房间——那时，任何持有不同意见的人都不会进入那房屋——他们开始用早餐。他们静默就座后，按顺序摆放面包，接着摆上某种与\Emph{面包}一同吃的食物，每人从中得到足够的一份。然而，在司铎祝福并祈祷\Emph{食物}之前，无人品尝。早餐后，当司铎再次献上祈祷，如同开始时一样，在用餐结束时他们以赞美诗赞美天主。接着，他们将一起在\Emph{屋内}用餐时所穿的衣服视为神圣而脱下——现在\Emph{这些衣服}是亚麻的——然后穿上\Emph{他们留在门厅的衣服}，急忙从事愉快的活动直到傍晚。他们用晚餐，一切都按类似的方式进行。没有人会大声喊叫，也听不到任何喧闹的声音。但每个人都安静地交谈，并以礼让的方式互相交谈，以至于\Emph{屋内}之人的静默在外人看来像是一种奥秘。他们始终节制，饮食皆有分寸。\par

% structure: p0049 source=h2 target=subsection reason=relative-heading-rank; base=h2

\subsection{第十七章　厄色尼派教义（续）}

所有人都听从会长；他发布的任何命令，他们都当作法律遵守。因为他们渴望向劳累的人施以怜悯和帮助。他们尤其避免恼怒、怒气以及诸如此类的\Emph{情绪}，认为这些对人有害。他们中无人习惯于发誓；但任何人所说的话，都视为比誓言更有约束力。若有人发誓，他将被谴责为不值得信任的人。他们同样关心法律和先知书的诵读；此外，若有任何信徒的著作，\Emph{也同样关心}。他们对植物和石头表现出极大的好奇，更忙于探究这些事物的作用能力，说这些东西并非徒然受造。\par

% structure: p0051 source=h2 target=subsection reason=relative-heading-rank; base=h2

\subsection{第十八章　厄色尼派教义（续）}

但对于那些希望成为该派门徒的人，他们不会立即传授其规则，除非他们已经事先考验了\Emph{他们}。现在，在一年之内，他们向（候选人）提供同样的食物，而后者则继续居住在\Emph{厄色尼派的}聚会处之外的另一所房屋中。他们给（受考验者）一把斧头、亚麻腰带和一件白袍。当一年期满，若有人证明了自己的自制力，他就更接近于该\Emph{派的}生活方式，并比之前更洁净地洗涤。然而，他尚未与\Emph{厄色尼人}一同用餐。因为，在提供证据证明他能获得自制力之后——但这等人的习惯要经过两年的考验——当他显为配得时，他就被算为该\Emph{派的}成员。然而，在他被允许与他们一同用餐之前，他要在可怕的誓言下受约束。首先，他将敬拜天主；其次，他将对人施行公正的\Emph{待人处事}，绝不伤害任何人，并且不恨那伤害他的人或敌视他的人，反而为他们祈祷。他同样发誓，将始终协助正义之人，\Emph{并且}对所有人守信，尤其是对那些统治者。因为，\Emph{他们论证说}，权位不会无缘无故地落在任何人身上，除非出于天主。而如果\Emph{该厄色尼人}本人是统治者，\Emph{他发誓}，他在\Emph{行使}权力时绝不傲慢行事，也不挥霍，或采用任何装饰，或超过习惯\Emph{允许}的更\Emph{豪华状态}。\Emph{他同样发誓}，要热爱真理，斥责说谎者，不偷盗，不为不义之财玷污良心，不向本派成员隐瞒\Emph{任何事}，也不向他人泄露，即使被折磨至死。除了上述\Emph{承诺}之外，他发誓不以不同于他自己领受的方式向任何人传授教义的知识。\par

% structure: p0053 source=h2 target=subsection reason=relative-heading-rank; base=h2

\subsection{第十九章　厄色尼派教义（续）}

于是，他们用这样的誓言约束那些前来的人。然而，若有人因任何罪而被定罪，他就被逐出该团体；但一个如此\Emph{被逐出的人}有时会可怕地死去。因为，既然他受该\Emph{派的}誓言和仪式约束，他不能食用其他人所用的食物。因此，\Emph{那些被逐出的人}有时完全因饥饿而毁坏身体。所以，到了最后，\Emph{厄色尼人}有时会怜悯许多\Emph{他们中}濒临死亡的人，因为他们认为即使死亡这样\Emph{施加的}刑罚在这些\Emph{罪犯}身上，也是足够的\Emph{惩罚}。\par

% structure: p0055 source=h2 target=subsection reason=relative-heading-rank; base=h2

\subsection{第二十章　厄色尼派教义（终）}

但在司法裁决方面，\Emph{厄色尼派}最为精确公正。他们聚齐至少\Emph{一百}人后，便宣判；其判决不可撤销。他们尊崇立法者\Emph{仅次于}天主；若有人亵渎这位\Emph{律法制定者}，必受惩罚。他们被教导要服从统治者和长老；若十人同坐一\Emph{室}，其中\Emph{一人}若不先看是否对那九人有利，便不发言。他们小心不在在场者中间和右边吐唾沫。然而，他们比所有\Emph{其他}犹太人更注重在安息日停工。因为他们不仅提前一天为自己预备食物，以免（在安息日）生火，甚至在那天连一件器皿也不从一个地方挪到另一个地方，也不解手；有些人甚至不从床上起来。但在其他日子，当他们想要解手时，就用镐头挖一个一英尺长的坑——因为这种斧头是起初会长赐给前来报名成为门徒之人的——并用衣服遮盖（这个坑）四周，声称他们并不必要侮辱阳光。然后他们将翻出的土放回坑中；这是他们的习惯，选择较僻静的地方。但做完这件事后，他们立刻沐浴，仿佛粪便玷污了他们。\par

% structure: p0057 source=h2 target=subsection reason=relative-heading-rank; base=h2

\subsection{第二十一章 厄色尼派的不同派别。}

\Emph{厄色尼派}随着时间推移发生了分裂，他们不再以相同方式持守训练体系，因为他们已分裂为四个派别。有些人自我训练超过\Emph{派规}的要求，甚至不肯接触当地流通的钱币，说他们不应携带、看见或制造像；因此他们中无人进城，以免（因此）通过有雕像\Emph{竖立}的城门，认为经过像下是犯法。但另一派的追随者，若偶然听见有人讨论天主及其法律——假设这人是未受割礼者——就会严密监视他，当他们单独在任何地方遇见这种人时，就会威胁要杀死他，除非他接受割礼。如果后者不愿满足此要求，\Emph{厄色尼人}毫不留情，甚至杀害。正是由此事件他们得了名号，被称为（某些人称为）热诚者，另一些人称为匕首党。另一派的追随者除了天主之外不称任何人为主，即使有人折磨甚至杀死他们。但还有\Emph{另有一些}较晚时期的人，他们偏离（派规）训练如此之远，以至于那些持守原始习俗的人甚至不愿触摸这些人。若偶然与他们接触，便立刻沐浴，仿佛摸到了异族人。然而这里也有非常多的人如此长寿，甚至活过一百岁。因此他们断言，这原因在于他们极其虔诚，谴责一切关于（食物）供应的过度，并且他们节制、不易发怒。因此他们蔑视死亡，若能以良心平安结束一生便喜悦。然而，若有人甚至折磨这种人，试图诱使他们中任何人谤法或吃祭偶像之物，他必不能得逞；因为\Emph{这一派的人}宁可受死忍痛，也不愿违背良心。\par

% structure: p0059 source=h2 target=subsection reason=relative-heading-rank; base=h2

\subsection{第二十二章 厄色尼派关于复活的信念；其体系具有启发性。}

如今，复活的教义在\Emph{他们}中也得到了支持；因为他们承认肉身将要复活，且将是不朽的，如同灵魂已经是不灭的。\Emph{他们主张}，灵魂在现世分离后，（进入）一个地方，那里通风良好、光明，\Emph{在那里}，他们说，\Emph{它}安息直到审判。希腊人曾听传闻知道这地方，称之为\Quote{福岛}。他们的其他许多教义被许多希腊人借用，并由此形成自己的观点。因为按照这些（犹太派别），关于神性的训练体系比所有民族的\Emph{都}更古老。因此证据就在眼前：所有那些敢于对天主或现存万物的创造发表言论的（希腊人），其原则都来自犹太立法。其中\Emph{特别可举出}毕达哥拉斯和斯多葛派，他们（在埃及居住期间）成为这些\Emph{犹太人}的门徒而获得（其体系）。如今他们断言将有审判和宇宙的焚烧，恶人将受永罚。在他们中间，预言和预测未来之事也受到培养。\par

% structure: p0061 source=h2 target=subsection reason=relative-heading-rank; base=h2

\subsection{第二十三章 厄色尼派的另一派别：法利塞人。}

于是有厄色尼人的另一支派，他们采用与前几派相同的习俗和规定的生活方式，但在一点上有所不同，即婚姻。他们认为，那些废除婚姻的人犯下了某种可怕的\Emph{罪行，这}是对生命的毁灭，并且不应该断绝子孙的后代；\Emph{因为}如果所有人都持有这种观点，整个人类种族将轻易被灭绝。然而，他们对订婚的女子进行为期三年的考验；当她们经过三次洁净，以证明生育能力后，才结婚。他们不与孕妇同居，表明他们结婚不是出于感官的动机，而是为了子女的利益。妇女们也以类似的方式（与丈夫）进行洗浴，并且也穿着亚麻衣服，如同男子\Emph{都是}束着腰带一般。这些\Emph{事}，就是我所要\Emph{陈述的关于厄色尼人的}。\par

但也有其他一些人，他们自己遵守犹太习俗；这些人在等级和法律方面都被称为法利塞人。他们大部分\Emph{可见于}各地，因为虽然所有人都被称为犹太人，但由于他们提出的意见独特，他们被冠以属于\Emph{各派}的适当称号。他们牢牢持守古代传统，并以辩论精神仔细探究那些按照\Emph{这}法律\Emph{被认为}洁净和不洁净的事物。他们解释法律的\Emph{规条}，并推举教师，\Emph{他们使之合格}\Emph{教导}这些\Emph{事物}。这些\Emph{法利塞人}肯定命运的存在，并认为有些事情在我们的能力之内，而其他事则在命运的控制之下。\Emph{他们主张}某些\Emph{行为}取决于我们自己，而其他的则取决于命运。但他们断言，天主是万有的原因，没有祂的旨意，就没有事物被管理或发生。他们也承认肉身的复活，灵魂不朽，并且将有审判和大焚烧，义人将不朽，而恶人将在不灭的火中承受永罚。\par

% structure: p0064 source=h2 target=subsection reason=relative-heading-rank; base=h2

\subsection{第二十四章 撒杜塞人。}

这些就是法利塞人的主张。然而，撒杜塞人主张废除命运，他们认为天主不做任何邪恶的事，也不对（尘世事物）行使天主眷顾；但\Emph{他们争辩说}善恶的选择在于人的能力。他们否认不仅有肉身的复活，而且他们认为灵魂在死后并不继续存在。\Emph{在死后。他们认为灵魂什么也不是}，仅仅是生命力，正因如此，人才被创造。然而，他们主张复活的观念已完全实现在一个单一情况中，即我们在世上留下子女后结束此生。但他们仍坚持认为，死后不会期待遭受任何痛苦或福乐；因为灵魂和身体都将解体，人归于非存在，与动物的\Emph{物质}相似。至于一个人在生命中可能犯下的任何邪恶，只要他\Emph{向受伤害的一方}和解，他（因过犯）便有所得，因为他逃脱了（否则将由人施加的）惩罚。而一个人可能获得的任何东西，（在任何方面），通过变得富有而获得名声，他便有所得。但他们坚持他们的断言，即天主对个人\Emph{在此世}的事务毫不关怀。法利塞人充满相互关爱，而撒杜塞人则受自爱驱使。这一派别尤其在撒玛黎雅周围的\Emph{区域}有强大据点。他们也遵守法律的习俗，说一个人应当如此生活，好使自己行为端正，并在世上留下子女。然而，他们不关注先知，也不关注其他任何智者，只关注梅瑟法律，\Emph{关于这}法律他们却不作任何解释。这些就是撒杜塞人选择\Emph{教导}的主张。\par

% structure: p0066 source=h2 target=subsection reason=relative-heading-rank; base=h2

\subsection{第二十五章 犹太宗教。}

因此，既然我们已经说明了犹太人之间的差异，似乎也适宜不缄默其宗教体系的内容。关于宗教，全体犹太人的教义分为四类：神学的、自然的、伦理的\Emph{和}礼仪的。他们肯定只有一位天主，祂是宇宙的创造者和主；祂形成了所有这些先前不存在的荣耀工程，并且这也不是出于任何早已存在的同质实体，而是祂的旨意——即有效的原因——是\Emph{要创造}，祂就创造了。他们也主张有天使，这些天使被造是为了服务受造界；但还有一位至高之神灵，永远存在于天主身边，为光荣和赞颂。他们认为受造界中万物都有知觉，没有无生命的东西。他们认真追求严肃的习惯和有节制的生活，这可从他们的法律中得知。\Emph{这些事情}早已由那些在古代领受了天主所定法律的人严格规定，因此读者会惊叹于在有关人的法律习俗中所\Emph{投入}的节制和勤奋。然而，适应于神圣\Emph{敬拜}的礼仪服务，以适合\Emph{宗教}尊严的方式，在他们当中以最高度的精致实行。他们的礼仪优越性，愿意知道的人很容易查明，只要他们阅读提供这些信息的书。\Emph{他们将会看到}犹太\Emph{司铎}如何庄严圣洁地将天主所赐予人类享用和成长的初果奉献给天主；他们\Emph{如何}按天主的诫命，规律而坚定地履行职责。不过，其中有些礼仪用法是撒杜塞人所不接受的，因为他们不愿承认天使或神灵的存在。\par

然而，所有\Emph{派别}都同样期待默西亚，因为法律和先知确实预先宣讲祂将临于世上。但由于犹太人不认识祂来临的时期，他们仍认为关于祂来临的圣经宣告尚未应验。因此直到今天，他们仍在期待\Emph{基督}的未来来临，因为当祂\Emph{在世界}时他们没有认出祂。但看到祂已经在我们当中的时代记号，\Emph{犹太人}感到困扰；他们羞于承认祂已来临，因为他们亲手将祂处死，因他们被祂指证没有遵守法律而恼怒。他们声称天主所派遣的那位不是他们所期待的这位基督；但他们承认另一位\Emph{默西亚}将会来临，目前还不存在；他会显示法律和先知所预示的一些记号，而关于其余的记号，他们认为自己判断有误。因为他们说他的出生将出自达味家族，但不是由童贞女和\Emph{圣神}，而是由一男一女，正如所有人都是从种子受生一样。他们宣称这位\Emph{默西亚}将作他们的君王——一位好战而强有力的人物，在召集全体犹太人民\Emph{并}与所有国家作战之后，将为他们恢复耶路撒冷王城。他将把整个\Emph{希伯来}民族聚集到这城，并再次带回古代的习俗，使他们履行君王和司祭的职责，在足够长的时间内安居。\Emph{在此安息之后，他们认为}战争将随之而来，攻击这些聚集的人；在这冲突中，基督将被刀剑所杀；\Emph{而}不久之后，宇宙的终结和焚烧将随之而来；这样，他们关于复活的意见将得完成，并按各人的行为得到报应。\par

% structure: p0069 source=h2 target=subsection reason=relative-heading-rank; base=h2

\subsection{第二十六章 对本书的结论说明。}

现在我们认为，关于所有希腊人和外邦人的教义，我们已经充分说明，并且无论是哲学所探讨的点，还是异端者提出的主张，都没有遗漏未经反驳。从这些说明本身，异端者的定罪是明显的，因为他们或窃取了教义，或从某些希腊人精心阐述的教义中获取贡献，然后提出这些意见仿佛源于天主。既然我们匆忙地审视了所有这些\Emph{体系}，并以大量劳苦在九卷书中宣告了他们的所有意见，为众人留下了生活中微小的行粮，并为我们同时代的人提供了学习的渴望，带来极大的喜乐和愉悦，我们认为合理，作为整个\Emph{著作}的顶点，引入关于真理的论述，并在一卷书即第十卷中提供我们的描绘。\Emph{我们的目的是}，读者不仅在知悉那些胆敢建立异端者的失败后，蔑视他们的\Emph{空想}，而且在认识真理的力量后，踏上救恩之路，将那配得天主应得的信仰寄托于天主。\par

\SourceRecord{Translated by J.H. MacMahon.；Ante-Nicene Fathers；Vol. 5.；Edited by Alexander Roberts, James Donaldson, and A. Cleveland Coxe.；Buffalo, NY: Christian Literature Publishing Co.,；1886.；<http://www.newadvent.org/fathers/050109.htm>.}

% structure: p0001 source=h1 target=section reason=page-title

\section{\WorkTitle{驳斥一切异端（卷十）}}

% structure: p0002 source=h2 target=subsection reason=relative-heading-rank; base=h2

\subsection{\WorkTitle{目录。}}

以下是第十卷\WorkTitle{驳斥一切异端}的内容：—\par

所有哲学家的概要。\par

所有异端的概要。\par

最后，什么是真理的教义。\par

% structure: p0007 source=h2 target=subsection reason=relative-heading-rank; base=h2

\subsection{第一章 概述。}

我们不是用暴力冲破异端的迷宫，而是仅仅通过反驳，或者换句话说，凭借真理的力量，解开了（其错综复杂之处），然后我们接近真理\Emph{本身}的证明。因为那时，当我们成功地确立真理的定义源自何处时，错误的诡辩将在其所有不一致中被揭露。真理没有从希腊人的智慧中吸取其原则，也没有像秘密奥迹那样从埃及人的教义中借用其道理，埃及人的教义虽然愚蠢，却被他们视为值得信赖而受到宗教性的敬礼。它也不是由加色丁人的好奇心所表达的支离破碎（结论）的谬误形成的。真理的存在也不归功于通过魔鬼的作为而产生的惊异，以及巴比伦人非理性的狂热。但它的定义是按照每个真实定义的方式构成的，即简单而不加修饰。这样的定义，只要彰显出来，就会\Emph{自身}驳斥错误。虽然我们非常频繁地提出了证明，并且充分详尽地向那些愿意（学习）真理之规则的人阐述了；但即使在现在，在讨论了希腊人和异端者提出的所有观点之后，我们仍然决定，无论如何，将这作为\Emph{前面}（九）卷的一个收尾，在第十卷中提出这个证明，并非不合理。\par

% structure: p0009 source=h2 target=subsection reason=relative-heading-rank; base=h2

\subsection{第二章 哲学家意见摘要。}

因此，我们用了四卷涵盖了希腊人中所有智者的教义，五卷涵盖了异端首领提出的道理，现在我们将在一卷中展示关于真理的道理，首先简要介绍所有人\Emph{分别}持有的假设。希腊的独断论者将哲学分为三部分，以这种方式不时设计其思辨体系；有的称其\Emph{体系}为自然（哲学），有的为伦理（哲学），还有的为辩证\Emph{哲学}。那些称其科学为自然哲学的古代思想家，是在\Emph{第一卷}中提到的。他们提供的说明是这样：他们中有些人\Emph{导出}万物从一，而另一些人从多于一个导出万物。在从一导出万物的人中，有些人从无性质者导出，有些人从有性质者导出。在从性质导出万物的人中，有些人从火导出，有些人从气导出，有些人从水导出，有些人从土导出。在从多于一个导出宇宙的人中，有些人从可数的数量导出，而另一些人从无限的数量导出。在从可数数量导出万物的人中，有些人从二导出，有些人从四导出，有些人从五导出，有些人从六导出。在从无限数量导出宇宙的人中，有些人从与生成物相似的东西导出实体，而另一些人从不相似的东西导出实体。在这些中，有些人从不能承受变化的东西导出实体，而另一些人从能承受变化的东西导出实体。因此，斯多亚派从无性质而具有统一性的物体来说明宇宙的产生。因为根据他们，无性质且在各部分都能变化的物质构成了宇宙的起源原则。因为当这个变化发生时，就产生了火、气、水、土。但是，希帕索斯、阿那克西曼德和米利都的泰勒斯追随者\Emph{倾向于认为}万物是从一个有性质的一个（实体）产生的。梅塔蓬图姆的希帕索斯和以弗所的赫拉克利特宣称万物\Emph{的起源}是从火而来，而阿那克西曼德是从气，泰勒斯是从水，色诺芬尼是从土。\Quote{因为从土，}他说道，\Quote{万物都是，万物也终结于土。}\par

% structure: p0011 source=h2 target=subsection reason=relative-heading-rank; base=h2

\subsection{第三章 哲学家意见摘要（续）}

但在那些\Emph{从多于一个且从可数数量导出所有实体}的人中，诗人荷马断言宇宙由两种\Emph{实体}组成，即土和水；有一次他这样表达：—\par

\Quote{众神的源头是海和大地母亲。}\par

在另一个场合\Emph{这样}：—\par

\Quote{但你们确实都可能变成水和土。}\par

科洛丰的色诺芬尼似乎与他一致，因为他说道：—\par

\Quote{我们都源自水和土。}\par

而欧里庇得斯（将宇宙）源自土和气，正如人们可以从他下面的断言中确定的：—\par

\Quote{万物的母亲，气和土，我歌唱。}\par

但恩培多克勒\Emph{将宇宙}从四个原则导出，他这样表达自己：—\par

\Quote{首先听你万物的四个根源：\newline{}灿烂的宙斯，赋予生命的赫拉和哈迪斯，\newline{}以及涅斯蒂斯，她以泪水沾湿凡人之泉。}\par

然而，卢卡尼亚的奥克洛斯和亚里士多德从五个\Emph{原则}导出宇宙；因为，除了四个元素之外，他们还假定\Emph{存在}第五个，并且（这是一个）做圆周运动的物体；他们说，天界事物由此而存在。但恩培多克勒的门徒认为宇宙的产生源自六个\Emph{原则}。因为在他所说的段落中，\Quote{首先听你万物的四个根源，}他从四个\Emph{原则}产生了生成。然而，当他补充说——\par

\Quote{毁灭性的争执与这些并列，在每一点上相等，\newline{}以及与之同等的友谊，在长度和广度上相等，——}\par

他也提出了宇宙的六条原则，其中四条是质料性的——土、水、火与空气；而另外两条是形成性的——友谊与纷争。然而，克拉佐美纳的阿那克萨戈拉、德谟克利特、伊壁鸠鲁及其众多追随者的门徒则认为，宇宙的产生源自无限数量的原子；我们先前已部分提及这些哲学家。但阿那克萨戈拉认为宇宙来自于与所产生的物相似的事物；而德谟克利特和伊壁鸠鲁的追随者则认为宇宙来自于与产物既不同且无感受的事物，即原子。然而，本都的赫拉克利德与阿斯克勒庇阿德的追随者则认为宇宙来自于与产物不同且有感受的事物，如同来自不协调的微粒。但柏拉图的弟子们断言这些实体来自三个原则——天主、质料与范型。然而，他们将质料分为四个原则——火、水、土与空气。并说天主是这质料的创造者，而心智是它的范型。\par

% structure: p0025 source=h2 target=subsection reason=relative-heading-rank; base=h2

\subsection{第四章　哲学家意见综述续}

既深信生理学原则对于所有这些哲学家而言都公开地显得困难重重，我们自己也要无畏地宣告有关真理的典范，即它们如何存在，以及我们确信它们如何存在。但我们之前将以摘要的方式预先解释异端首领的教义，以便通过将所有的教义摆在我们读者面前，使之因这处理方法而广为人知，从而以浅显明了的方式展现真理。\par

% structure: p0027 source=h2 target=subsection reason=relative-heading-rank; base=h2

\subsection{第五章　纳塞尼派}

既然这样做似乎合宜，我们首先从公开崇拜蛇的人开始。纳塞尼派称宇宙的第一原则为一个人，并且这同一原则也是人子；他们将这\Quote{人}分为三部分。因为他们说，他的一部分是理性的，另一部分是精神的，但第三部分是属地的。他们称他为\OriginalTerm{阿达马斯}{Adamas}，并认为关于他的知识是认识天主之能力的根源。纳塞尼派声称，所有这些理性的、精神的和属地的品质都已进入耶稣，并且通过祂，这三个实体同时向宇宙的三个种类说话。他们声称存在三种存在——天使的、精神的和属地的；并且有三个教会——天使的、精神的和属地的；这些的名称是——被拣选的、被召叫的和被俘虏的。这些是他们所提出的教义要点，就简短理解而言。他们声称主的兄弟雅各伯将这些教义传授给玛利亚娜，以此说法来诽谤两者。\par

% structure: p0029 source=h2 target=subsection reason=relative-heading-rank; base=h2

\subsection{第六章　佩拉特派}

佩拉特派，即卡律斯托的阿德墨斯和佩拉特的欧弗拉特，说有一个世界——这是他们所用的名称——并断言它分为三部分。但根据他们，这三重划分有一个原则，如同一个巨大的泉源，能够由理性分割为无限的部分。第一和最近的部分，据他们说，是三元组，被称为完美的善和父性的伟大。但三元组的第二部分是一定数量的，如同无限的能力。第三部分则是形式的。第一是非受生的；因此他们明确声称有三个天主、三个\OriginalTerm{逻各斯}{Logoi}、三个心智、三个人。因为当划分完成后，他们给世界的每一部分都分派了天主、逻各斯、人等其余。但自上，来自非受生性和世界的第一部分，当世界后来达到圆满时，佩拉特派断言，在黑落德时代，有一个具有三重性、三重身体和三重能力的人降临，名为基督，祂从世界的三部分中在自身内拥有世界所有的结合和能力。他们认为这就是所宣告的：\BibleQuote{在祂内住有整个圆满的天主性，具有身体} \BibleRef{哥 2:9}。他们断言，从上方的两个世界——即非受生的和自生的——有各种能力的种子被带入我们所在的这个世界。并且基督从上方的非受生性降临，为要借着祂的下降，所有被分为三部分的事物都能得救。因为，佩拉特派说，从上降下的那些事物将通过祂上升；而那些密谋反对降下之物的事物将被无心拒绝，并送去受罚。佩拉特派说，有两个部分被拯救——即那些位于上方的——通过使之与腐败分离，而第三部分被毁灭，称之为形式世界。这些就是佩拉特派的教义。\par

% structure: p0031 source=h2 target=subsection reason=relative-heading-rank; base=h2

\subsection{第七章　塞特派}

但就塞特派（\OriginalTerm{塞特派}{Sethians}）看来，存在三个被精确界定的本原（\OriginalTerm{本原}{principles}）。每个本原按其本性适于被生成，正如在人的灵魂中，任何可被学习的技艺都能被发展（\Emph{被发展}）出来。其结果（\Emph{其结果}）与一个孩子长期接触乐器而成为音乐家，或接触几何学而成为几何学家，或接触任何其他（\Emph{其他}）技艺而得到类似结果一样。塞特派（\Emph{塞特派}）说，这些本原的实质是光与黑暗。介于两者之间的是纯神（\OriginalTerm{纯神}{pure spirit}）；他们说，这神居于下方的黑暗与上方的光明之间。它不是像一股气流或可感知的微风那样的神，而仿佛是某种由精细混合物制成的香膏或焚香的香气——一种藉着某种不可言喻、超越言语的香气的冲动而扩散自己的能力。因此，既然光明在上，黑暗在下，神介于两者之间，光明也像太阳的光线一样从上照射在下面的黑暗上。神的香气向前飘荡，占据中间位置，向前扩散，如同祭品（放在）火上的焚香的气味扩散一样。既然三分之物的能力是这样，神与光的能力同时存在于下面的黑暗中。然而，他们说，黑暗是恐怖之水（\OriginalTerm{恐怖之水}{horrible water}），光与神被吸入其中，从而被转变到这种本性中。黑暗既然赋有理智，并且知道光一旦离开它，黑暗就会继续荒芜，缺乏光辉（\Emph{和}）光彩、能力（\Emph{和}）效能，且无力，因此这黑暗用尽一切思虑和理性的努力，试图将光辉、光的火花以及神的香气包含在自身之内。他们引入以下比喻来说明这一点，这样表达自己：正如眼珠在下面的体液下看起来是黑暗的，却被神照亮，同样，黑暗切望追求神，并拥有所有想要退去和返回的能力。这些能力是无限无量的，从它们混合中，所有事物都被塑造和生成，如同印记一般。因为正如一个印记接触蜡时产生一个形状，（而印记）本身仍保持它原来的样子，同样，这些能力通过彼此相通，形成了所有无限种类的动物。\Emph{塞特派断言}，因此，从三个本原的初次交会生出了一个大印记（\OriginalTerm{大印记}{great seal}）的形象，即天和地，具有子宫（\OriginalTerm{子宫}{womb}）的形状，中央有一个脐。因此，所有其他事物的形象，如天和地，都被塑造得类似一个子宫。\par

塞特派（\Emph{塞特派}）说，从水中产生了一个首生本原（\OriginalTerm{首生本原}{first-begotten principle}），即一股猛烈狂躁的风（\OriginalTerm{猛烈狂躁的风}{vehement and boisterous wind}），它是所有生成的原因，通过水的运动在世界上产生一种热和运动。他们（\Emph{他们主张}）认为这风（\Emph{风}）被塑造得像蛇的嘶嘶声，成为一个完美的形象。世界注视着它，并匆匆进入生成，如同子宫被点燃；从此他们倾向于（\Emph{倾向于认为}）宇宙的生成已经发生。他们说这风构成一种神，一个完全的天主（\OriginalTerm{完全的天主}{perfect God}）从水的香气、神的香气和灿烂的光中产生。他们（\Emph{他们肯定}）说心智（\OriginalTerm{心智}{mind}）以从女性生成的方式产生——（心智意指）上界的火花（\OriginalTerm{上界的火花}{supernal spark}）——并且它已与身体的复合物混合在下面，急切地想要逃离，好得以离开，不因水的欠缺而消散。因此，按照他们的说法，它习惯从水的混合中呼喊，如圣咏作者所说：\Quote{上界之光的一切焦虑，就是为了把下面的火花从下方的父（\OriginalTerm{下方的父}{Father beneath}）那里解救出来，} 即从风那里。而\Emph{这父}产生热和扰动，并为自身生出一个儿子，即心智，按他们的说法，这心智（\Emph{如}他们声称）并非他自己特有的（\Emph{产物}）。这些异端者（\Emph{这些异端者肯定那子}）看见上界之光的完全的逻各斯（\OriginalTerm{完全的逻各斯}{perfect Logos}），就发生了转变，并以蛇的形状（\OriginalTerm{蛇的形状}{shape of a serpent}）进入一个子宫，以便能够收回那来自光的火花的心智（\Emph{心智}）。这就是所说的：\Quote{祂虽具有天主的形体，却没有把与天主同等视为应紧抓不放的，反而空虚自己，取了奴仆的形体。} \BibleRef{斐 2:6} 可怜而有害的塞特派倾向于（\Emph{倾向于认为}）这就是宗徒所指的奴仆的形体。这些就是这些塞特派（\Emph{塞特派}）也提出的主张。\par

% structure: p0034 source=h2 target=subsection reason=relative-heading-rank; base=h2

\subsection{第八章 西门·马古斯}

然而，那位极其聪明的\Emph{家伙}\Person{西满}如此陈述：存在一个无限的能力，\Emph{并且}这是宇宙的根源。而这无限的能力，他说，就是火，它本身并非如那粗俗的\Emph{空论家们所主张}的那样是单纯的，因为他们断言有四个无混合的元素，并假定火\Emph{作为其中之一}是不复合的。\Emph{\Person{西满}则断言}火的本性是双重的；他将这双重（本性）的一部分称为某种隐秘的，另一部分称为某种显明的。并且\Emph{他断言}，那隐秘的隐藏在火的显明\Emph{部分}中，而火的显明\Emph{部分}是从隐秘中产生的。他说，火的一切部分，无论可见的与不可见的，都被认为具有感知能力。因此，他说，那受生的世界，是从那无生的火产生的。他说，它开始存在的方式如下：那无生者从那火的本质中取了生成本原的六个原始根源。因为\Emph{他坚持认为}，这些根源是从火中成对生成的；并称它们为：理智与智力、声音与名字、推理与反思。并且\Emph{他断言}，在这六个根源中同时存在着那无限的能力，\Emph{这能力}他断言是\OriginalTerm{那曾站立、站立着、将要站立的}{Him that stood, stands, and will stand}。当这\Emph{一位}被塑造成形体时，\Emph{照这异端者说}，祂将实质地\Emph{和}潜在地存在于六种能力中。\Emph{而祂将是}在大小与完美上同那无生而无限的能力完全一样，没有任何方面比那无生、不变、无限的能力更欠缺。然而，如果\Emph{\OriginalTerm{那曾站立、站立着、将要站立的}{Him that stood, stands, and will stand}}只是潜在地存在于六种能力中，而没有采取任何确定的形体，\Emph{\Person{西满}}说，祂就变得完全幻灭而消灭。\Emph{而这发生}的方式与语法或几何能力相同：这种能力虽然被\Emph{植入}人的灵魂，但当它没有得到（帮助）一位\Emph{任一}这类技艺的大师来教导\Emph{那灵魂进入其原理}时，就会\Emph{消灭}。现在\Person{西满}肯定，他本人就是\OriginalTerm{那曾站立、站立着、\Emph{并}将要站立的}{Him that stood, stands, \Emph{and} will stand}，并且祂是超乎万有的能力。以上，便是\Person{西满}的意见。\par

% structure: p0036 source=h2 target=subsection reason=relative-heading-rank; base=h2

\subsection{第九章 \OriginalTerm{瓦伦廷}{Valentinus}。}

\OriginalTerm{瓦伦廷}{Valentinus}以及这一学派的追随者，虽然他们一致断言宇宙的本原就是\Emph{圣父}，但仍被迫采取了\Emph{关于祂}的相反意见。因为他们中的一些人\Emph{主张}（圣父）是独一而能生育的；而另一些人则\Emph{认为}（在祂如在其他情况下）没有女性就不可能生育。因此他们加上了\OriginalTerm{西革}{Sige}作为这位\Emph{圣父}的配偶，并称\Emph{圣父}自己为\OriginalTerm{深渊}{Bythus}。从这位\Emph{圣父}和祂的配偶，一些人\Emph{声称}产生了六个流出体——即\OriginalTerm{理智}{Nous}和\OriginalTerm{真理}{Aletheia}、\OriginalTerm{圣言}{Logos}和\OriginalTerm{生命}{Zoe}、\OriginalTerm{人}{Anthropos}和\OriginalTerm{教会}{Ecclesia}——而这构成了能生育的\OriginalTerm{八元体}{Ogdoad}。而\Emph{瓦伦廷派坚持认为}那些是\Emph{最初的}在界限内发生的流出体，\Emph{并且}被称为\Quote{\Emph{那些}在\OriginalTerm{圆满}{Pleroma}之内者}；第二类是\Quote{在\OriginalTerm{圆满}{Pleroma}之外者}；第三类是\Quote{在界限之外者}。而这些的产生构成了\OriginalTerm{缺欠}{Hysterema} \Emph{\OriginalTerm{阿卡摩斯}{Acamoth}}。他断言，从一个\OriginalTerm{永世}{Aeon}所生、存在于\Emph{\OriginalTerm{缺欠}{Hysterema}}中\Emph{并且}被投射（到界限之外）的，就是创造主。但是\Emph{\Person{瓦伦廷}}并不倾向于\Emph{肯定如此所生者}是原初的神性，而是用贬抑的言辞谈论祂和祂所造之物。（他断言）基督从\OriginalTerm{圆满}{Pleroma}之内降下，为了拯救那已迷失的灵。这灵（照瓦伦廷派说）居于我们内在的人；并且他们说，这\Emph{内在的人}由于这内住的\Emph{灵}而得救恩。然而\Emph{\Person{瓦伦廷}}（为维护教义）断定肉身不得救，并称之为\Quote{一件皮衣}，以及人可朽坏的\Emph{部分}。我（已经）以摘要的方式宣告了这些教理，因为在他们的体系中存在大量的讨论内容以及众说纷纭。因此，以这种方式，对于瓦伦廷学派而言，\Emph{提出他们的意见}似乎是适当的。\par

% structure: p0038 source=h2 target=subsection reason=relative-heading-rank; base=h2

\subsection{第十章 \OriginalTerm{巴西里德}{Basilides}。}

但巴西里德本人也肯定有一个非存有的天主，祂作为非存有者，造了一个非存有的世界，这世界是从不存有的事物中形成的，祂投下一种种子，如同芥菜籽，本身包含茎、叶、枝\Emph{和}果实。或者\Emph{这种子}如同孔雀的蛋，本身包含各种颜色。\Emph{巴西里德派}说，这构成了世界的种子，万物都由它产生。因为\Emph{他们主张}它本身包含一切事物，如同那些\Emph{尚未}存有的事物，并已由非存有的天主预定要使之存有。于是，他说，在种子本身中有一个三重的子性，在各方面都与非存有的天主同体，是由不存有的事物所生的。这子性分为三部分，一部分是精微的，另一部分是粗重的，第三部分需要净化。精微的部分，当非存有的天主最初放下种子时，立即迸发出来，向上上升，走向非存有的\Emph{天主}。因为每一本性都因\Emph{祂}的超越之美而渴望那\Emph{天主}，但不同造物因不同原因渴望祂。但较粗重的部分仍留在种子中；由于它是一种模仿性的\Emph{本性}，不能向上飞升，因为它比精微部分更粗重。然而，\Emph{较粗重的部分}以圣神装备自己，如同翅膀；因为子性\Emph{这样}装扮，向这\Emph{圣神}显示恩惠，并\Emph{反过来}接受恩惠。第三部分子性需要净化，\Emph{因此}它留在所有胚芽的聚集中，并且显示和接受恩惠。巴西里德断言，有一些东西被称为\Quote{世界，}而另一些东西被称为超世界；因为\Emph{存有物}被他分为两个基本类别。介于它们之间的，他称为\Quote{毗连的圣神，}这圣神\Emph{在本身内}具有子性的芬芳。\par

从宇宙种子的一切胚芽的聚集迸发出来，生出了\OriginalTerm{大元首}{Great Archon}，世界的头，\Emph{一位\OriginalTerm{永世}{Aeon}}，其美和伟大不可言喻。这位元首将自己升到穹苍之上，以为没有另一位在祂以上。\Emph{因此}祂变得比下面所有的\Emph{永世}更辉煌、更有能力，除了已被留下的子性，但祂并不知道这子性比祂更有智慧。这位元首将注意力转向创造世界，首先为自己生了一个儿子，比祂自己更优越；\Emph{这儿子}祂让坐在自己右边，这些\Emph{巴西里德派}声称这就是\OriginalTerm{奥格多阿得}{Ogdoad}。大元首自己产生了整个天界的受造物。另一位元首从所有胚芽中升起来，\Emph{他}比所有下面的\Emph{永世}更伟大，除了已被留下的子性，但远逊于前一位。他们称这\Emph{\OriginalTerm{第二位元首}{second Archon}}为\OriginalTerm{赫布多马得}{Hebdomad}。他是管理、创造和控制祂以下的\Emph{一切事物}的那一位，这\Emph{元首}为自己生了一个比\Emph{自己}更有远见和智慧的儿子。现在，他们断言这一切都是根据那位非存有的\Emph{天主}的预定而存在，并且存在无限的诸世界和间隔。\Emph{巴西里德派肯定}，在由玛利亚所生的耶稣身上，福音的能力降下，照亮了奥格多阿得和赫布多马得的儿子。\Emph{这是为了}照明和区分不同的存有秩序，并净化被留下的子性，以便对灵魂施以恩惠，并反过来接受恩惠。他们说他们自己是儿子，在世界上为此目的，即通过教导来净化灵魂，并与子性一同上升到在上的圣父那里，\Emph{从}祂发出最初子性。他们声称世界持续到所有灵魂都随子性一同修复到那里为止。然而，这些就是巴西里德所详述为奇迹的见解，他并不羞于提出。\par

% structure: p0041 source=h2 target=subsection reason=relative-heading-rank; base=h2

\subsection{第十一章 犹斯定}

但是\Person{犹斯定}本人也试图确立与这些相似的见解，并这样表达：宇宙有三个无始的本原，两个男性\Emph{和}一个女性。而在两个男性中，一个本原被称为\Quote{善者}。只有这个本原用这种方式称呼，\Emph{并且}具有对宇宙的预知。而另一个\Emph{是}一切受造之物的父，\Emph{并且}没有预知，不可知，不可见，被称为\OriginalTerm{厄罗音}{Elohim}。这个女性本原没有预知，充满激情，有两个心智，\Emph{以及}两个身体，正如我们在关于这个\Emph{异端者}体系的\Emph{先前}论述中已详细说明的那样。这个女性本原，在她身体上部直至腹股沟处，\Emph{犹斯定派说}，是一个处女，而从腹股沟向下则是一条蛇。她被称为\OriginalTerm{厄登}{Edem}和\Person{以色列}。\Emph{这个异端者}断言这些就是宇宙的本原，万物都从它们产生。而且\Emph{他断言}，厄罗音没有预知，对那半处女产生了过度欲望，与她交合后，生下了十二个天使；而这些天使的名字，他说\Emph{就是先前已经给出过的}。\Emph{在这些天使中}，属于父系的与父相连，属于母系的与母相连。而\Person{犹斯定}主张，这些就是乐园的树木，\Person{梅瑟}以寓意方式论及它们在法律中所写的事。而且\Person{犹斯定}确认，万物都是由厄罗音和厄登造成的。他说，动物以及其他\Emph{这类受造物}来自那类似兽的部分，而人来自腹股沟以上的部分。\Person{犹斯定}认为，厄登将灵魂（这是她自己的力量）赋予了\Emph{人}本身，而厄罗音则赋予了神。而\Person{犹斯定}断言，这个厄罗音在得知\Emph{自己的起源}后，升到了\Emph{善者}那里，并抛弃了厄登。而且\Emph{这个异端者断言，厄登}因受到这样的对待而愤怒，便策划了整个阴谋针对厄罗音安置在人内的神。\Person{犹斯定}告诉我们，为此父派遣了\Person{巴路克}，并向先知们发出指示，以便厄罗音的神能被解救，所有人都能远离厄登的诱惑。但这个异端者断言，就连\Person{赫拉克勒斯}也是一位先知，并且他被\Person{翁法勒}（即\Place{巴贝耳}）击败；而\Emph{犹斯定派}称后者为\Person{维纳斯}。他们又说，后来在\Person{黑落德}的日子，\Person{耶稣}诞生，为\Person{玛利亚}和\Person{若瑟}的儿子，\Person{犹斯定}声称巴路克曾向他说话。\Person{犹斯定}断言，厄登曾谋害这位耶稣，但未能欺骗他；因此，她使他被钉在十字架上。而耶稣的神，\Person{犹斯定}说，升到了\Emph{善者}那里。\Emph{犹斯定派}主张，所有如此服从那些愚蠢而无用的言论的人，他们的神将得救，而厄登的身体和灵魂则被遗弃。但愚蠢的\Person{犹斯定}称这厄登为\Quote{大地}。\par

% structure: p0043 source=h2 target=subsection reason=relative-heading-rank; base=h2

\subsection{第12章。幻象说派。}

幻象说派提出这样的主张：原始神如同一颗无花果树的种子；从这种子生出三个\OriginalTerm{永世}{Aeon}，如同树干、叶子和果实；这些永世又生出三十个永世，每个十个；它们都按十位联合，但仅在位置上不同，有些在其他之前。幻象说派断言，无数的永世被无限地生出，并且所有这些永世都是雌雄同体。他们又说，这些\Emph{永世}策划同时进入一个永世，而从这个中间永世\Emph{和}从\Person{童贞玛利亚}，他们生出了一位万有的救主。\Emph{这位救赎者}在每个方面都与无花果树的第一个种子相似，但在这点上次于它，即祂是被生的；因为无花果树所由来的种子是无始的。\Emph{这}就是永世的伟大光明——完全是光辉——它不接受任何装饰，并在自身中包含\Emph{所有}动物的形式。幻象说派主张，这\Emph{光明}进入下面的混沌时，为被产生的和实际存在的事物提供了存在的原因，并且从上面降下时，它将永恒种属的形式印在了下面的混沌上。因为第三个永世，它曾将自己三重化，当它察觉自己全部特性被强力抽入下方的黑暗时，且既知道黑暗的恐怖又知道光的单纯，便着手创造天；在使中间之物坚固后，它将黑暗与光分开。他说，由于第三个永世的所有形式都被黑暗压制，这个\Emph{永世}的形象甚至成了一团活火，由光产生。从这（源头），他们声称，生成了大掌权者，\Person{梅瑟}论及他说，祂是火神和\OriginalTerm{德穆革}{Demiurge}，祂也不断地将所有（永世）的形式改变为身体。幻象说派声称，这些就是救主为之被生的灵魂，祂指出灵魂逃脱的路，这些灵魂（现在）被（黑暗）压制。幻象说派主张，\Person{耶稣}穿上那独一受生的能力，因此无人能看见祂，因为祂的荣耀极其伟大。他们说，所有发生在祂身上的事正如\WorkTitle{福音}中所记载的那样。\par

% structure: p0045 source=h2 target=subsection reason=relative-heading-rank; base=h2

\subsection{第13章。\Person{莫诺伊摩斯}。}

但阿拉伯人\Person{莫诺伊穆斯}的追随者们断言，宇宙的起源原则是一个原初之人和人子；并且，按照\Person{梅瑟}的说法，被造之物不是由原初之人产生的，而是由那原初之人的\Emph{圣子}产生的，\Emph{然而}不是由整个圣子，而是由祂的一部分。\Person{莫诺伊穆斯}断言，人子是\OriginalTerm{约塔}{\GreekText{ι}}，它代表十，这个首要数字内含有所有数字的基体，并通过它每一个数字得以构成，以及宇宙、火、空气、水和地的生成。但既然这是一个约塔和一个笔画，且完美者出自完美者，\Emph{或者换句话说}，一个笔画从上流下，包含万有于自身；因此，人所拥有的一切，人子之父也\Emph{同样拥有}。所以\Person{梅瑟}说世界在六天内造成，即通过六种权能，世界由那一个笔画造成。因为立方体、八面体、金字塔以及所有类似的具有等面积的面，火、空气、水和地即由它们构成，都是从包含在那约塔的简单笔画中的数字产生出来的，这笔画就是人子。\Person{莫诺伊穆斯}说，当\Person{梅瑟}提到\Emph{为了带来}给埃及降下灾祸而挥舞的棍杖时，他喻指了约塔世界的改变；并且他所造的灾祸没有超过十个。\Emph{然而}他说，如果你想要认识宇宙，要省察你自己内里，是谁在说：\Quote{我的灵魂、我的肉身\Emph{和}我的心智}，是谁将每一样东西归给自己，如同另一个人为自己那样。要明白这是一个从完美者产生的\Emph{完美的一}，并且他将所有所谓的非存有和所有存有视为己有。这些也是\Person{莫诺伊穆斯}的意见。\par

% structure: p0047 source=h2 target=subsection reason=relative-heading-rank; base=h2

\subsection{第十四章　\Person{塔提安}}

\Person{塔提安}却与\Person{瓦伦提努斯}及其他人类似，说存在着某些不可见的\OriginalTerm{永世}{Aeon}，并且下层世界及其中的万物是由其中的某一个造出来的。他使自己习惯于一种非常犬儒式的生活方式，并且在几乎任何方面都与\Person{马尔基昂}无异，无论是关于他的诽谤，还是关于婚姻的法规。\par

% structure: p0049 source=h2 target=subsection reason=relative-heading-rank; base=h2

\subsection{第十五章　\Person{马尔基昂}和\Person{切尔多}}

但\Person{本都}的\Person{马尔基昂}和他的老师\Person{切尔多}，他们自己也设定宇宙的三个本原——善、义和物质。然而他们的一些门徒又添加了\Emph{第四个}，说：善、义、恶\Emph{和}物质。但他们都肯定那善者什么也没有造，尽管有些人也称那义者为恶，另一些人则认为他仅被称为义者。他们声称那义者用底层的物质造了万物，因为他造得不好，而是不合理。因为被造之物应当与造物主相似；因此他们这样使用福音中的比喻，说：\BibleQuote{好树不能结坏果子}\BibleRef{玛 7:18}以及其余段落。\Emph{现在马尔基昂}认为那义者自己所设想得不好的概念构成了这段话的寓意。他说基督是善者的圣子，并被他称为内里的人差来拯救灵魂。他断言祂显现为人，\Emph{尽管}不是人，成为肉身，\Emph{尽管}不是肉身。\Emph{并且他坚持}祂的显现只是幻影，祂并没有经历出生或受难，除了外表。他拒绝承认肉身复活；并且宣告婚姻是败坏，从而引导他的门徒走向非常犬儒式的生活。通过这些\Emph{方式}，他想象自己能激怒造物主，如果他戒绝由祂所造或所定的事物的话。\par

% structure: p0051 source=h2 target=subsection reason=relative-heading-rank; base=h2

\subsection{第十六章　\Person{阿佩莱斯}}

但\Person{阿佩莱斯}，这个\Emph{异端者}的门徒，对他老师提出的论点感到不满，正如我们先前所说，并用另一种理论假设有四位神。他声称第一位是\Quote{善者}，先知们不认识祂，基督是祂的圣子。\Emph{他声称}第二位是宇宙的创造者，并且他不愿称祂为神。\Emph{他说明}第三位神是显现的那位火神，第四位是恶神。\Person{阿佩莱斯}称这些为天使；并通过将基督也加入其中，他将宣称祂是第\Emph{五位}神。但\Emph{这个异端者}习惯研读一本书，他称之为某位\Person{菲卢梅内}的\WorkTitle{启示录}，他视她为女先知。他断言基督不是从童贞玛利亚获得肉身，而是从世界邻近的物质中获得。他以此方式撰写了反对法律和先知的论文，并试图废除它们，仿佛它们说了假话，不认识天主。\Emph{阿佩莱斯}与\Person{马尔基昂}类似，肯定各种肉身都将被消灭。\par

% structure: p0053 source=h2 target=subsection reason=relative-heading-rank; base=h2

\subsection{第十七章　\Person{凯林图斯}}

\Person{凯林图斯}本人曾在埃及受过训练，断定世界不是由第一位神造的，而是由某种天使权能造的。\Emph{这权能}远远分离并远离那超乎整个存在圈之上的主权，并且不认识那超乎万有的神。他说耶稣不是由童贞玛利亚所生，而是由\Person{若瑟}和\Person{玛利亚}所生，\Emph{作为他们的}儿子，与其余人相似；并且他在公义、智慧和理解上超越其余\Emph{所有人类}。\Emph{凯林图斯坚持}，在耶稣领受圣洗后，基督以鸽子形状从超乎整个存在圈之上的主权降临到祂身上，然后祂开始宣讲那未识之父，并施行奇迹。\Emph{他断言}，在受难结束时，基督从耶稣身上飞走了，但耶稣受了苦，而基督不可能受苦，因为祂是主的神。\par

% structure: p0055 source=h2 target=subsection reason=relative-heading-rank; base=h2

\subsection{第十八章　\OriginalTerm{厄彼翁人}{Ebionaeans}}

但厄俾翁派却断言世界是由真天主所造，\newline{}他们\Emph{论及}基督的方式与刻林托类似。然而，他们在一切事上均按梅瑟法律生活，声称如此便得以成义。\par

% structure: p0057 source=h2 target=subsection reason=relative-heading-rank; base=h2

\subsection{第十九章　德奥多托}

拜占庭的德奥多托引入了下述异端，声称万有皆由真天主所造；而基督，他说，以一种与前述诺斯底派相似的方式，按照某种描述显现。\newline{}\Emph{德奥多托肯定}基督是具有与所有人相同本性的人，但他在这一点上超越众人：按照天主的计划，他由童贞女所生，圣神遮蔽了\Emph{他的母亲。这位异端者却坚持，耶稣}并未在童贞女胎中取肉身，而是其后基督在耶稣圣洗时以鸽子形状降临到他身上。由此，\Emph{德奥多托的追随者}肯定，起初奇迹异能并未在\Emph{救主}自身内运作。\newline{}\Emph{德奥多托}却决意否认基督的天主性。德奥多托所倡导的即是此类观点。\par

% structure: p0059 source=h2 target=subsection reason=relative-heading-rank; base=h2

\subsection{第二十章　默基瑟德派}

另有些人也作出与前述完全相同的断言，只引入一点改动，即认为默基瑟德是一种能力。但他们声称默基瑟德自身超越一切能力；并按照他的肖像，他们想要\Emph{主张}基督也是这样\Emph{受生的}。\par

% structure: p0061 source=h2 target=subsection reason=relative-heading-rank; base=h2

\subsection{第二十一章　弗吕吉亚派或蒙丹派}

弗吕吉亚派从其异端原则出自某个蒙丹、普利斯卡和马克西米拉，视这些可怜妇人为女先知，视蒙丹为先知。然而，关于宇宙的起源和受造，\Emph{弗吕吉亚派}被认为表达正确；而在他们关于基督所陈述的信条上，他们的意见并非不当。但他们与\Emph{先前提到的异端者}一同误入歧途，并将其注意力集中于这些人的言论超过福音，从而制定了关于新奇和陌生斋戒的规条。\par

% structure: p0063 source=h2 target=subsection reason=relative-heading-rank; base=h2

\subsection{第二十二章　弗吕吉亚派或蒙丹派（续）}

但他们中另一些人，依附于诺耶图派的异端，持有与\Emph{弗吕吉亚派的}愚蠢妇人和蒙丹相似的意见。然而，关于万有之父的真理，他们犯了亵渎罪，因为他们断言他是子和父、可见和不可见、受生和非受生、可朽和不朽的。这些人是从某个诺耶图那里取得契机来\Emph{提出其异端的}。\par

% structure: p0065 source=h2 target=subsection reason=relative-heading-rank; base=h2

\subsection{第二十三章　诺耶图与卡利斯托}

同样，诺耶图，斯弥纳本地人，一个轻率妄言且诡诈的人，将这一起源于厄丕戈诺的异端引入我们中间。它传至\Emph{罗马}，被克来奥默纳采纳，并一直延续至今在其后继者中。\newline{}\Emph{诺耶图}断言，宇宙只有一位父和天主，他造了万有，且当他不愿时对那些存在者是不可察觉的。\Emph{诺耶图主张，父}在他愿意时显现；他未被看见时是不可见的，但被看见时是可见的。这位\Emph{异端者也声称，父}在他不受生时是非受生的，但在他由童贞女出生时是受生的；同样，当他不受苦、不死时，他是无苦和不死的。然而，当他的苦难临到他时，\Emph{诺耶图承认父}受苦且死。\Emph{诺耶图派}认为这位父本身被称为子（\Emph{反之亦然}），这是针对在各自适当的时候发生在他们各自身上的事件而言。\par

卡利斯托证实了这些\Emph{诺耶图派}的异端，但我们已\Emph{详细}解释了其生平细节。\newline{}\Emph{卡利斯托}本人也产生了一种异端，其出发点取自这些\Emph{诺耶图派}——即他承认有一位父和天主，即宇宙的创造者，并且这位天主被谈论、被称为子的名，但在实体上是一位神。因为神，作为天主，他说，并非任何与\Emph{圣言}不同的\Emph{存有}，圣言也与神无异；因此，这一位（根据卡利斯托）在名分上是分开的，但在实体上并非如此。他认为这唯一的圣言就是天主，并且肯定在\Emph{圣言那里}有降生成人。他倾向于（主张）那位在肉身中被看见并被钉十字架的是子，但父是住在他内的那位。\Emph{卡利斯托因此}时而分支到诺耶图的意见，时而又分支到德奥多托的意见，不持守确定的教义。这些便是卡利斯托的意见。\par

% structure: p0068 source=h2 target=subsection reason=relative-heading-rank; base=h2

\subsection{第二十四章　赫尔默根}

但有一个赫尔默根，他自己也想说些什么，断言天主用与他\Emph{同存}并\Emph{服从他计划}的原始物质造了万有。因为\Emph{赫尔默根}认为，天主除非从已有的东西中，否则不可能造出受造之物。\par

% structure: p0070 source=h2 target=subsection reason=relative-heading-rank; base=h2

\subsection{第二十五章　厄耳卡赛派}

但另有一些人，引入某种新颖的\Emph{教义}，从所有异端中剽窃了\Emph{其体系的组成部分}，并搜集了一册奇异的著作，其扉页上写着\Person{Elchasai}的名字。这些人同样承认宇宙的诸原则是由天主所创立。然而，他们并不告白只有一个基督，而是有一个超乎\Emph{其余者}的基督，且祂经常穿梭于许多身体中，而现在则在耶稣内。同样，\Emph{这些异端者主张}，有时\Emph{基督}由天主所生，有时则成为圣神，有时\Emph{出生于}童贞女，有时则不然。并且\Emph{他们断言}，这耶稣后来也持续不断穿梭于身体中，并在\Emph{不同}时期显现在许多（不同的身体）中。他们又诉诸咒语和在告白元素时的洗礼。他们忙于占星与数学的学问以及巫术技艺。但\Emph{也}自诩拥有预知的能力。\par

% structure: p0072 source=h2 target=subsection reason=relative-heading-rank; base=h2

\subsection{第二十六章 犹太编年史}

……（亚巴郎）奉天主的命令，从美索不达米亚的哈兰城迁居到现今称为巴力斯坦和犹太之地，当时则是客纳罕地区。关于这地域，我们在其他论述中已部分、但并非草率地作了叙述。由此迁移，可追溯犹太地人口增长的起点；犹太地得名于雅各伯的第四子犹大，雅各伯也称为\Emph{以色列}，因将有君王之后裔从他而出。亚巴郎离开美索不达米亚时年七十五岁，百岁时生依撒格。依撒格六十岁时生雅各伯。雅各伯八十六岁时生肋未；肋未四十岁时生子刻哈特；刻哈特下埃及与雅各伯同去时四岁。因此，亚巴郎寄居以及由他通过依撒格所出的整个家族在时称客纳罕地的时期，总共为二百一十五年。这\Emph{亚巴郎}的父亲是特辣黑，\Emph{特辣黑}的父亲是纳曷尔，\Emph{纳曷尔}的父亲是色鲁格，\Emph{色鲁格}的父亲是勒伍，\Emph{勒伍}的父亲是培肋格，\Emph{培肋格}的父亲是厄贝尔（\BibleRef{创 11:16}）。由此，\Emph{犹太人}便以希伯来人的名称被称呼。但在培肋格时代，民族分散了。这些民族共有七十二个，\Emph{与亚巴郎子女的数目相符}。这些民族的名字我们也在其他书中列出，甚至未遗漏此\Emph{点}于其适当之处。\Emph{我们详尽详述的理由}，是渴望向那些有勤勉心志的人表明我们对天主性的爱慕，以及我们在劳作中所获取的关于真理的确切知识。但这厄贝尔的父亲是舍拉；\Emph{舍拉}的父亲是刻南；\Emph{刻南}的父亲是阿帕革沙得，其\Emph{父亲}是闪；\Emph{闪}的父亲是诺厄。在\Emph{诺厄}时代，发生了遍及全世界的洪水，这是埃及人、加色丁人、希腊人都记不得的；因为发生于\Person{厄癸革}与\Person{德乌卡利翁}时代的洪水仅局限于他们居住的地区。因此，在这些族长（即从诺厄到厄贝尔，共五代）中，有四百九十五年。这\Emph{诺厄}是位极虔敬且爱天主的人，独自与妻子儿女以及这三人的三位妻子逃脱了随之而来的洪水。他因一只方舟而得救；这\Emph{方舟}的尺寸与遗物，正如我们所述，至今仍展示在那些位于\Place{阿迪阿贝尼人地区}方向的\Place{亚拉辣特山}上。因此，凡有意勤加探究此事者，都能清楚看出，有一个敬拜\Emph{真}天主的民族，比所有加色丁人、埃及人\Emph{和}希腊人更为古老。然而，目前有何必要详述那些在诺厄之前既虔敬又获准与\Emph{真}天主交谈的人呢？因为就当前所处理的主题而言，这关于\Emph{天主子民}古老性的证据岂不是够了？\par

% structure: p0074 source=h2 target=subsection reason=relative-heading-rank; base=h2

\subsection{第二十七章 犹太编年史（续）}

但既然证明这些全神贯注于\Emph{哲学思辨}的民族，比那些习惯崇拜\Emph{真}天主的民族更为晚近，似乎并非不合理，那么我们有理由说明：\Emph{这些后者的家族源自何处}；以及当他们定居在这些地区时，并非从实际地点得名，而是以最初出生并居住于此的人为自己\Emph{取名字}。诺厄有三个儿子：闪、含\Emph{和}雅斐特。全人类都从他们繁衍，\Emph{大地的}每个角落\Emph{最初都归功于他们}。因为天主的话在他们身上有效，当\Emph{主}说：\BibleQuote{你们要生育繁殖，充满大地。}\Emph{那}唯一的话如此有力，\Emph{以至于}在\Emph{这个}家族中，从\Emph{诺厄的三个儿子}生出了72个子孙：从闪生了25个，从雅斐特生了15个，\Emph{以及}从含生了32个。然而，含的这些32个子孙是按照之前的宣告而生。\Emph{含的子孙中有}：迦南，迦南人从他而来；米兹辣因，埃及人从他而来；雇士，埃塞俄比亚人从他而来；普特，利比亚人从他而来。这些民族按照他们中\Emph{通行}的语言，至今仍以祖先的名号称呼；甚至在希腊语中，他们也以现在所称的名字被称呼。但即使假设这些地方先前无人居住，也无法证明有一族人类从一开始就存在在那里，然而，这些恭敬天主的诺厄的儿子们，\Emph{已足以证明所争论之点。因为显然诺厄}本人必是虔诚之人的门徒，因此他逃脱了那巨大而短暂的洪水威胁。\par

那么，恭敬\Emph{真}天主的人怎么不比所有加色丁人、埃及人\Emph{和}希腊人更古老呢？因为我们必须记住，这些\Emph{外邦人}的先祖是从雅斐特所生，得名雅完，并成了希腊人和伊奥尼亚人的祖先。如果那些致力于哲学问题的民族，被证明比恭敬天主的民族以及\Emph{洪水的时代}都要晚得多，那么那些蛮族民族，以及世上所有已知和未知的\Emph{部落}，不显得比这些更近代吗？因此，你们这些希腊人、埃及人、加色丁人以及全人类，都要精通这门教义，并从我们这些天主的朋友学习：天主的本质是什么，祂井井有条的受造物是什么。我们培养这套\Emph{体系}，并非以浮夸的言辞表达自己，而是以\Emph{证明我们认识真理和实践明智的}术语来撰写我们的论文，我们的目标是证明祂的\Emph{真理}。\par

% structure: p0077 source=h2 target=subsection reason=relative-heading-rank; base=h2

\subsection{第二十八章 真理的教义。}

唯一的天主，既是万有的造物主和主，没有任何事物\Emph{与祂同在}；没有无限的混沌，没有无量的水，没有坚固的地，没有浓密的空气，没有温暖的火焰，没有精微的气息，没有惊人穹苍的蔚蓝天幕。祂是独一的，在祂自身内独自存在。祂凭意志创造了存在的事物，这些事物先前并不存在，除非祂愿意造它们。因为祂完全知晓将要发生的一切，因为预知也临于祂。然而，祂首先制造了将要存在的事物的不同元素，即火、气、水、土，从这些不同的\Emph{元素}祂开始形成自己的创造。有些对象祂用一种本质形成，有些用两种、三种或四种复合而成。那些\Emph{形成}于一种\Emph{本质}的是不朽的，因为\Emph{就它们而言}不会分解，因为单一的事物永远不会被分解。另一方面，由两种、三种或四种\Emph{本质}形成的是可分解的，因此它们被称为可朽的。因为这被称为死亡，即连接着的\Emph{本质}的分解。我现在认为我已经足够回答了那些心智健全的人；如果他们渴望进一步的教导，并愿意准确研究这些事物的本质和整个受造物的原因，那么在阅读我们的作品\WorkTitle{论宇宙的本质}时，他们就会了解这些要点。然而，我认为现在足以阐明那些原因；\Emph{希腊人}不知道这些原因，以浮夸的言辞荣耀受造物的部分，却不认识造物主。异端首领们以此为契机，将那些希腊人先前提出的陈述转化为类似的教义，从而编造了可笑的异端。\par

% structure: p0079 source=h2 target=subsection reason=relative-heading-rank; base=h2

\subsection{第二十九章 真理的教义（续）。}

因此，这独一无上的天主通过反思，首先产生了\OriginalTerm{圣言}{Logos}；不是\Emph{用声音发出的}言语，而是宇宙的推理，在\Emph{神圣心智中}构思并居住。祂独自从存在的事物中产生了圣言；因为\Emph{圣父}自己构成了存在，而从祂所生的圣言是万物的成因。圣言本身就在圣父内，承载着祂的创造者的意志，并且不陌生于圣父的思想。因为当祂从祂的创造者发出时，由于祂是这创造者的\Emph{首生者}，祂在自己内如同声音，拥有在圣父中构思的观念。因此，当圣父命令世界存在时，圣言逐一完成了\Emph{每一个受造物}，从而取悦了天主。有些通过生殖繁衍的事物，祂造为男性和女性；但凡是\Emph{旨在服务和管理}的受造物，祂造为男性，或不需要女性，或既非男性也非女性。因为甚至连这些原初实体，即从无中形成的火、灵、水、土，既非男性也非女性；若没有那一切权柄之源的天主愿意圣言\Emph{协助完成此类创造}，则从这些中的任何一个都不能产生男性或女性。我承认天使是火，并且我坚持女性灵体不与它们同在。而且我认为太阳、月亮和星辰同样\Emph{是由火与灵产生的}，既非男性也非女性。创造者的意愿是，水生和飞行的动物来自水，有男有女。因为天主如此\Emph{命令}，祂愿意存在一种具有生育能力的湿性物质。同样，天主\Emph{命令}从土中生出爬行动物和野兽，以及各种动物的雄性和雌性；因为受造物的本性如此允许。凡祂愿意的，天主都随时创造了。祂通过\Emph{圣言}创造了这些，因为万物只能按它们被创造的方式生成。但当祂按自己的意愿也形成万物时，祂给它们命名，并如此宣告了祂的\Emph{创造之功}。在\Emph{创造}这些时，祂形成了万物的统治者，并用所有复合物质塑造了他。\Emph{创造者}不想把他造为神，也并未失败；也不是天使——不要受骗——而是人。因为如果祂愿意把你造为神，祂本可以做到。你有圣言的例子。但祂的意愿是，你应当是人，\Emph{而}祂把你造为人。但如果你也渴望成为神，就要服从创造你的主，现在不要抗拒，以便你在小事上被证明忠信时，也能受托大事。\par

唯独这\Emph{独一真天主}的\OriginalTerm{圣言}{Logos}是从天主本身而来；因此\Emph{圣言}就是天主，是天主的\Emph{实体}。世界是从无中造的；因此\Emph{它}不是天主；\Emph{也因为}这个\Emph{世界}在创造者愿意时就会消散。但创造\Emph{它}的天主，没有也没有制造恶。祂制造荣耀和卓越的；因为创造\Emph{它}的是善的。人被造时，是有自决能力的受造物，但未拥有至高的理智，未通过反省、权威和能力统管万物，而是\Emph{受制于自己的私欲}，并且自身包含\Emph{各种}对立。但人由于拥有自决能力，就产生了恶，即偶然地；这\Emph{恶}并不完全，除非你实际上犯了某种恶行。因为就在于我们渴望任何邪恶的事，或思想它时，那恶才\Emph{如此}被称为恶。恶从起初并不存在，而是后来才出现的。既然人有自由意志，神圣者就为他制定了一条法律，\Emph{为了引导他}，并非没有良好的目的。因为如果人没有意愿和不愿意的能力，为什么要设立法律？因为法律不会为无理性的动物设立，而是用笼头和鞭子；但给人颁布了诫命和刑罚，以便履行或不执行所命令的事。因为对于如此被造的人，在远古时代由正义之人制定了法律。更近的时代，由梅瑟（前述的那位虔诚且蒙天主喜爱的人）设立了一条严肃而正义的法律。\par

现在天主的\OriginalTerm{圣言}{Logos}掌管这一切；他是圣父的首生之子，是晓光之声，先于晨星。此后，义人出生，成为天主的友人；这些人被称为先知，因为他们预见未来的事。预言的话交托给他们，不只是为\Emph{一个}时代\Emph{而已}；而且历代预言之事的宣告，都极其清晰地赐予了。这事不仅发生在\Emph{异象者}回答当时在场之人的时候；而且历代将要发生的事都已预先显示；因为在谈及已过之事时，\Emph{先知们}使人类重忆往事；在指明当前之事时，他们竭力劝诫人不可懈怠；而预言未来之事时，他们使我们在看到早已预言的事时胆战心惊\Emph{，}并同样期待那些被\Emph{预言为}仍将来临的事。这就是我们的信德，哦，众人啊——\Emph{我说，我们的信德}，我们不是被空洞的言辞说服，也不是被\Emph{那}内心的突然冲动所掳，更未被华丽言谈的动听所欺骗，但我们不拒绝服从由神能所说的话。这些诫命是天主赐给\OriginalTerm{圣言}{Word}的。但圣言通过宣告它们，传布了神圣的诫命，从而使人远离悖逆，并非以必然的强力将人带入奴役，而是通过一种自发的选择召唤人进入自由。\par

这位\OriginalTerm{圣言}{Logos}是圣父在\Emph{末世}所派遣的，不再通过先知说话，也不愿\OriginalTerm{圣言}{Word}被模糊地宣告，成为纯粹猜测的对象，而是愿他被显示出来，好让我们亲眼看见他。这位\Emph{圣言}，我说，是圣父所派遣的，为使\Emph{这}世界看见他时，能尊敬那不以\Emph{那}先知的身份传授诫命、也不以\Emph{那}天使惊吓灵魂的，而是他自己——那曾说话的——\Emph{有形地}临在\Emph{我们中间}。这位\Emph{圣言}我们知道由童贞女取了身体，\Emph{并}经\OriginalTerm{新受造}{new creation}重塑了\OriginalTerm{旧人}{old man}。\Emph{我们相信这位圣言}经历了\Emph{此}生的每个阶段，为使他自身成为每个时代的法律，并藉着临于我们中间，将他的人性展示为所有人的目标。并且藉着他\Emph{亲自}，他证明天主没有创造任何邪恶，而人拥有\OriginalTerm{自决}{self-determination}的能力，因为他能愿意也能不愿意，\Emph{并}赋有能力做两者。\Emph{这}\OriginalTerm{那人}{Man}我们知道是由\Emph{我们人性的}复合体所造的。因为他若不同我们\Emph{本性相同}，他命令我们效法这位师傅就是徒然的。因为那人若碰巧是不同物质的，他为何要对我这生来软弱者施以与他所接受者相似的指令？这怎会是\Emph{那良善与公正者的行为}？然而，为不使他被以为\Emph{与我们不同}，他甚至经受劳苦，甘愿忍受饥饿，不拒绝感到口渴，并沉入安静的睡眠。他并未抗议他的苦难，反而服从至死，并彰显了他的复活。如今在这一切\Emph{行为}中，他献上了他自己的人性作为\OriginalTerm{初果}{first-fruits}，为使你，当你处于患难中时，不致灰心，反而承认自己是一个人（与救赎者本性相同），并盼望也将获得\Emph{圣父}所赐予这\Emph{圣子}的一切。\par

% structure: p0084 source=h2 target=subsection reason=relative-heading-rank; base=h2

\subsection{第三十章 作者的结束语}

这就是关于天主性的真道，哦，你们众人，希腊人和野蛮人，加色丁人和亚述人，埃及人和利比亚人，印度人和埃塞俄比亚人，克尔特人，以及你们统军作战的拉丁人，以及所有居住在欧罗巴、亚细亚和利比亚的人。我已经成为你们的劝诫者，因我是\Emph{那}仁慈的\OriginalTerm{圣言}{Logos}的门徒，\Emph{因而}有人性，为使你们能赶快，并藉着我们被教导谁是真实的天主，以及\Emph{什么}是他井然有序的创造。不要专注于人为言论的谬误，也不要专注于剽窃异端者的虚妄许诺，而要专注于谦逊真理的可敬纯朴。藉着这知识，你们将逃脱\Emph{那}审判之火的迫近威胁，以及阴暗塔尔塔罗斯的无光景象，在那里，\OriginalTerm{圣言}{Word}光照之声的光线永不照耀！\par

\Emph{你将逃脱}地狱永火之湖的滚沸洪流，以及那因罪受罚被锁在\Place{塔尔塔罗斯}的\Emph{堕落}天使们始终以威胁目光注视的眼睛；\Emph{你也将逃脱}那不断在身体周围盘绕觅食的虫子，而身体的污秽生出了\Emph{它}。如今，通过认识真天主的教导，你将避免这样的（痛苦）。你将拥有一个不朽的身体，一个甚至置于腐败可能性之外的身体，如同灵魂一样。你将领受天国，你，在此生旅居时认识天上君王的人。你将与\OriginalTerm{天主}{Deity}为伴，与基督同为承继者，不再受情欲或激情的奴役，也\Emph{再也不会}被疾病消耗。因为你已成为天主：因为你身为人类时所遭受的任何痛苦，祂都赐予了你，因为你是可朽的，但凡是符合天主\Emph{赐予}的，天主都应许赐予你，因为你已被神化，被生入不朽。这就是\Emph{这句谚语的含义}：\Quote{认识你自己；}即在你内发现天主，\Emph{因为}祂按照\Emph{祂自己的肖像}形成了你。因为认识自我与成为天主认识的对象是相连的，因为你被\Emph{天主}亲自召唤。所以，人啊，不要彼此燃起敌意，也不要犹豫，赶快回头。因为基督是至高无上的天主，祂已安排洗去人类身上的罪，使旧人重生。天主从起初就称人为祂的肖像，并以比喻显明了祂对你的爱。只要你服从祂庄重的诫命，成为良善者的忠实追随者，你将相似祂，因为你将获得由祂赐予的尊荣。因为\Emph{天主}（由于俯就）不会减损祂神圣完美的\OriginalTerm{天主}{Deity}丝毫；祂甚至使你成为天主，以光荣祂！\par

\SourceRecord{Translated by J.H. MacMahon.；Ante-Nicene Fathers；Vol. 5.；Edited by Alexander Roberts, James Donaldson, and A. Cleveland Coxe.；Buffalo, NY: Christian Literature Publishing Co.,；1886.；<http://www.newadvent.org/fathers/050110.htm>.}
